\documentclass[french]{book}

\usepackage[utf8]{inputenc}
\usepackage[T1]{fontenc}
\usepackage{lmodern}
\usepackage{geometry}
\usepackage{graphicx}
\usepackage{babel}
\usepackage{amsmath}
\usepackage{amsthm}
\usepackage{amssymb}
\usepackage{hyperref}
\usepackage{stmaryrd}
\usepackage{enumitem}
\usepackage{tcolorbox}
\usepackage{dsfont}
\usepackage{multirow}
\usepackage{etoc}
\usepackage{titlesec}
\usepackage{esint}
\usepackage{cancel}

\setlist[itemize]{label=$\bullet$}
\setcounter{tocdepth}{0}

\renewcommand{\thesection}{\arabic{section}}
\renewcommand{\thesubsection}{\arabic{section}.\alph{subsection}}

\theoremstyle{plain}
\newtheorem{thm}{Théorème}
\newtheorem*{ind}{Indication}

\theoremstyle{definition}
\newtheorem{dfn}{Definition}
\newtheorem{ex}{Exercice}

\theoremstyle{remark}
\newtheorem{sol}{Solution}

\renewcommand{\P}{\mathbb{P}}
\newcommand{\N}{\mathbb{N}}
\newcommand{\Z}{\mathbb{Z}}
\newcommand{\Q}{\mathbb{Q}}
\newcommand{\R}{\mathbb{R}}
\newcommand{\C}{\mathbb{C}}
\newcommand{\K}{\mathbb{K}}
\renewcommand{\L}{\mathbb{L}}
\newcommand{\U}{\mathbb{U}}
\newcommand{\st}{^*}
\newcommand{\tq}{\middle|}
\newcommand{\inv}{^{-1}}
\newcommand{\grad}{\overrightarrow{\text{grad}}}
\newcommand{\rot}{\overrightarrow{\text{rot}}}
\renewcommand{\div}{\text{div}}

\newtcolorbox{rbox}[1]{colback=red!5!white,colframe=red!75!black,fonttitle=\bfseries,title={#1}}
\newtcolorbox{gbox}[1]{colback=lime!5!white,colframe=lime!75!black,fonttitle=\bfseries,title={#1}}
\newtcolorbox{bbox}[1]{colback=cyan!5!white,colframe=cyan!75!black,fonttitle=\bfseries,title={#1}}

\DeclareMathOperator{\card}{Card}
\DeclareMathOperator{\vect}{Vect}
\DeclareMathOperator{\id}{id}



\author{La MP*1 de Niels Vandecasteele (le goat)}
\title{Cours de physique-chimie}

\begin{document}
\maketitle

\tableofcontents

\setcounter{tocdepth}{3}

\part{Mécanique}
\chapter{Mécanique du point matériel}
\localtableofcontents
\section{Cinématique du point}
\subsection{Systèmes usuels de coordonnées}
\subparagraph{Coordonnées cartésiennes}
Vecteur position : $\overrightarrow{OM}=x\overrightarrow{u_x}+y\overrightarrow{u_y}+z\overrightarrow{u_z}$ \\
Vecteur déplacement élémentaire : $d\overrightarrow{l}=d\overrightarrow{OM}=dx\overrightarrow{u_x}+dy\overrightarrow{u_y}+dz\overrightarrow{u_z}$ \\
Vecteur vitesse : $\overrightarrow{v_{M/R}}=\dot{x}\overrightarrow{u_x}+\dot{y}\overrightarrow{u_y}+\dot{z}\overrightarrow{u_z}$ \\
Vecteur accélération : $\overrightarrow{a_{M/R}}=\ddot{x}\overrightarrow{u_x}+\ddot{y}\overrightarrow{u_y}+\ddot{z}\overrightarrow{u_z}$ \\

\subparagraph{Coordonnées cylindriques}
Vecteur position : $\overrightarrow{OM}=r\overrightarrow{u_r}$ \\
Vecteur déplacement élémentaire : $d\overrightarrow{l}=d\overrightarrow{OM}=dr\overrightarrow{u_r}+rd\theta\overrightarrow{u_\theta}+dz\overrightarrow{u_z}$ \\
Vecteur vitesse : $\overrightarrow{v_{M/R}}=\dot{r}\overrightarrow{u_r}+r\dot{\theta}\overrightarrow{u_\theta}+\dot{z}\overrightarrow{u_z}$ \\
Vecteur accélération : $\overrightarrow{a_{M/R}}=(\ddot{r}-r(\dot{\theta}^2))\overrightarrow{u_r}+(2\dot{r}\dot{\theta}+r\ddot{\theta})\overrightarrow{u_\theta}+\ddot{z}\overrightarrow{u_z}$ \\

\subparagraph{Coordonnées sphériques}
Vecteur position : $\overrightarrow{OM}=r\overrightarrow{u_r}$ \\
Vecteur déplacement élémentaire : $d\overrightarrow{l}=d\overrightarrow{OM}=dr\overrightarrow{u_r}+rd\theta\overrightarrow{u_\theta}+r\sin(\theta)d\varphi\overrightarrow{u_\varphi}$ \\
Vecteur vitesse : $\overrightarrow{v_{M/R}}=\dot{r}\overrightarrow{u_r}+r\dot{\theta}\overrightarrow{u_\theta}+r\sin(\theta)\dot{\varphi}\overrightarrow{u_\varphi}$ \\
Vecteur accélération :
\begin{align*}
	\overrightarrow{a_{M/R}}=&\left(\ddot{r}-r(\dot{\theta}^2)-r\sin^2(\theta)(\dot{\varphi})^2\right)\overrightarrow{u_r} \\ &+\left(2\dot{r}\dot{\theta}+r\ddot{\theta}-r\sin(\theta)\cos(\theta)(\dot{\varphi})^2\right)\overrightarrow{u_\theta}\\
	&+\left(2\dot{r}\sin(\theta)\dot{\varphi}+2r\dot{\theta}\cos(\theta)\dot{\varphi}+r\sin(\theta)\ddot{\varphi}\right)\overrightarrow{u_\varphi}
\end{align*} 
\subsection{Cas du mouvement circulaire}
Pour un mouvement circulaire de rayon $R$, en utilisant les coordonnées cylindriques centrées en le centre du cercle et d'axe vertical orthogonal au cercle, on a : $$\overrightarrow{a}=-R(\dot{\theta})^2\overrightarrow{u_r}+R\ddot{\theta}\overrightarrow{u_\theta}$$
En particulier, on a : $\displaystyle a_r=-\frac{v^2}{R}$

\subsection{Base de Frenet}
Pour un mouvement curviligne plan quelconque, on a : $$\overrightarrow{a}=\frac{dv}{dt}\overrightarrow{T}-\frac{v^2}{\rho}\overrightarrow{N}$$ où $\overrightarrow{T}$ est un vecteur unitaire tangent à la trajectoire de même sens que $\overrightarrow{v}$, $\overrightarrow{N}$ est un vecteur unitaire orthogonal à $\overrightarrow{T}$ dans le plan de la trajectoire et dirigé vers l'intérieur de la concavité de la trajectoire, et $\rho$ est le rayon de courbure de la trajectoire à cet endroit (ie le rayon du cercle par lequel on peut localement approximer la trajectoire). Par exemple, si la trajectoire est rectiligne, le rayon de courbure est infini.

\subsection{Précession d'un vecteur}
On dit que le vecteur $\overrightarrow{A}$ précesse à la vitesse angulaire constante $\omega$ autour de l'axe $\overrightarrow{u_z}$ dans le repère $(O;\overrightarrow{u_x},\overrightarrow{u_y},\overrightarrow{u_z})$ lorsqu'il existe des constantes $B$, $C$ et $\varphi$ telles que $$\overrightarrow{A}(t)=B\cos(\omega t+\varphi)\overrightarrow{u_x}+B\sin(\omega t+\varphi)\overrightarrow{u_y}+C\overrightarrow{u_z}$$
On montre que cette condition est réalisée si, et seulement si, $$\frac{d\overrightarrow{A}}{dt}=\omega\overrightarrow{u_z}\wedge\overrightarrow{A}$$

\section{Dynamique du point matériel}
\subsection{Grandeurs cinétiques}
On définit la quantité de mouvement du point matériel $M$ de masse $m$ dans le référentiel  $R$ par : $$\overrightarrow{p_{M/R}}:=m\overrightarrow{v_{M/R}}$$ On définit le moment cinétique vectoriel du point $M$ par rapport au point $A$ dans le référentiel $R$ par : $$\overrightarrow{L_A}(M/R):=\overrightarrow{AM}\wedge\overrightarrow{p_{M/R}}$$
On définit le moment cinétique scalaire du point $M$ par rapport à un axe $\Delta=(A;\overrightarrow{u})$ par : $$L_\Delta:=\overrightarrow{L_A}(M/R)\cdot\overrightarrow{u}$$ Pour calculer un moment cinétique scalaire efficacement, on peut utiliser ce qu'on appelle le bras de levier. Si on note $v_\bot$ la norme de la composante de la vitesse orthogonale à $\overrightarrow{u}$ et $b$ la distance de l'axe $\Delta$ au support de la vitesse $\overrightarrow{v_\bot}$, alors on a : $$\lvert L_\Delta\rvert = bmv_\bot$$ La signe du moment cinétique scalaire par rapport à un axe indique si on tourne ou non dans le sens direct autour de l'axe avec la direction de l'axe imposée par $\overrightarrow{u}$.
\subsection{Forces usuelles}
\paragraph{Poids} Dans le référentiel terrestre, il tient compte de la force de gravitation et du caractère non galiléen du référentiel terrestre (\textit{cf} deuxième chapitre de mécanique).
\paragraph{Loi de la gravitation universelle} $$\overrightarrow{F}_{1\to 2}=-\mathcal{G}\frac{m_1m_2}{(M_1M_2)^3}\overrightarrow{M_1M_2}$$
\paragraph{Loi de Coulomb} (France, XVIIIe-XIXe) Si $q_1$ et $q_2$ sont immobiles, alors : $$\overrightarrow{F}_{1\to 2}=\frac{q_1q_2}{4\pi\varepsilon_0(M_1M_2)^3}\overrightarrow{M_1M_2}$$
\paragraph{Force de Lorentz}$$\overrightarrow{F}=q\left(\overrightarrow{E}+\overrightarrow{v}\wedge\overrightarrow{B}\right)$$
\paragraph{Réaction d'un support} S'il n'y a pas de frottement avec le support, $\overrightarrow{R}$ est normale au plan tangent de contact. Pour plus de détails, voir le chapitre $4$ de mécanique. \textbf{Attention}, la réaction normale peut travailler si on ne se place pas dans le référentiel du support !
\paragraph{Tension d'un fil} Bien lire le sujet. Pour un fil inextensible de masse négligeable tendu, la tension est la même partout. Une poulie parfaite (\textit{ie} de masse négligeable et sans frottement) modifie simplement la direction d'un fil idéal. S'il y a des frottements ou des contraintes variables le long du fil, passer à une étude locale du fil pour trouver une équation différentielle et remonter à l'équation de la tension.
\paragraph{Frottements fluides} Bien lire le sujet ! Il y a différents modèles : $$\begin{cases}
	\overrightarrow{f}=-\lambda\overrightarrow{v} \\
	\overrightarrow{f}=-\mu v\overrightarrow{v} \\
	...
\end{cases}$$
Le premier modèle, dit "linéaire", est utilisé à faibles vitesses, le deuxième, dit "quadratique", est utilisé pour des vitesses plus importantes.
\paragraph{Poussé d'Archimède} SI un corps est entièrement plongé dans un ou des fluide(s) à l'équilibre dans un référentiel galiléen, alors il est soumis à : $$\overrightarrow{\pi}=-m_{\text{fluide déplacé}}\overrightarrow{g}$$ Souvent, la poussée d'Archimède est négligeable devant le poids du dispositif, mais dans certains cas (sous-marins et ballons par exemple), elle est essentielle au bon fonctionnement du système. \\ \par 
Astuce : pour calculer la résultante des forces qui s'appliquent à un objet immergé mais dont un bord touche un solide, ajouter une fine bande de fluide fictive entre le bord et le solide, appliquer Archimède, puis enlever ce qu'on a ajouté en intégrant sur le bord.
\paragraph{Loi de Hooke / Force de rappel d'un ressort} Pour un ressort de raideur $k$ et de longueur à vide $l_0$, la force de rappel exercée par un ressort sur un objet $M$ à son extrémité est : $$\overrightarrow{F}=-k(l-l_0)\overrightarrow{u}$$ où $\overrightarrow{u}$ est un vecteur unitaire dirigé du point d'attache du ressort vers l'objet attaché au ressort.

\subsection{Théorèmes de la dynamique du point}
\begin{thm}[Théorème du moment cinétique vectoriel]
	Si $M$ est un point matériel, $R$ est galiléen et $A$ est fixe dans $R$, alors : $$\frac{d}{dt}\overrightarrow{L_A}(M/R)\bigg\rvert_{R}=\sum_i\overrightarrow{\mathcal{M}_A}\left(\overrightarrow{f_i}\right)$$
\end{thm}
\begin{thm}[Théorème du moment cinétique scalaire]
	Si $M$ est un point matériel, $R$ est galiléen et $\Delta$ est un axe fixe dans $R$, alors : $$\frac{d}{dt}L_\Delta(M/R)=\sum_i\mathcal{M}_\Delta\left(\overrightarrow{f_i}\right)$$
\end{thm}
\subsection{Trajectoire de phase d'un système à un degré de liberté}
\begin{dfn}[Portrait de phase]
	On appelle \textbf{portrait de phase} le graphe de $\dot{f}$ en fonction de $f$. On peut aussi renormaliser l'axe des ordonnées afin d'avoir la même dimension pour les deux axes. On oriente le sens de variation du graphe.
\end{dfn}
Pour un oscillateur harmonique sans second membre, on a un cercle. Avec un second membre, on a plus généralement une ellipse. Si l'oscillateur est amorti, on a une spirale qui s'enroule jusqu'à une position d'équilibre de vitesse nulle.
\section{Énergétique du point matériel}
Pour un point M de masse m dans un référentiel R, son énergie cinétique dans ce référentiel est définie par : $$E_c(M/R)=\frac{1}{2}mv_{M/R}^2$$
Tout comme la vitesse, l'énergie cinétique dépend donc profondément du référentiel d'étude.
\subsection{Théorème de la puissance cinétique, de l'énergie cinétique...}
\begin{thm}[Théorème de la puissance cinétique]
	Si $M$ est un point matériel dans un référentiel R \textbf{galiléen}, alors : $$\frac{d}{dt}E_c(M/R)=\sum_i \overrightarrow{f_i}\cdot\overrightarrow{v}_{M/R}$$
\end{thm}
\begin{thm}[Théorème de l'énergie cinétique]
	Si $M$ est un point matériel et $R$ est galiléen, alors : $$E_c(B/R)-E_c(A/R)=W_A^B$$
\end{thm}
\subsection{Définition d'une force conservative}
Une force $\overrightarrow{F}$ est dite \textbf{conservative} lorsqu'il existe $E_p$ qui ne dépend que de la position et telle que : $$\delta W =\overrightarrow{F}\cdot d\overrightarrow{OM}=-dE_p$$ Pour un point matériel, cela correspond à $\overrightarrow{F}=-\grad(E_p)$.
\begin{thm}[Théorème de la puissance mécanique]
	Si $M$ est un point matériel dans un référentiel R \textbf{galiléen}, alors : $$\frac{d}{dt}E_m(M/R)=\sum \overrightarrow{f}_{\text{non conservatives}}\cdot\overrightarrow{v}_{M/R}$$
\end{thm}
\begin{thm}[Théorème de l'énergie mécanique]
	Si $M$ est un point matériel et $R$ est galiléen, alors : $$E_m(B/R)-E_m(A/R)=W_{A \ \text{non conservatives}}^B$$
\end{thm}
\subsection{Équilibre et stabilité d'un système conservatif à un degré de liberté}
On trace l'énergie potentielle en fonction de la coordonnée qui constitue le degré de liberté. On trace la droite d'équation $y=E_m$. Les zones pour lesquelles $E_p>E_m$ ne sont pas atteignables car $E_c\geqslant0$. Les positions d'équilibre sont repérées par des \textit{extrema} locaux de l'énergie potentielle.
\begin{dfn}[Position d'équilibre stable]
	Une position d'équilibre est \textbf{stable} lorsqu'après une petite perturbation du système, le système tend à revenir vers cette position d'équilibre.
\end{dfn}
\begin{dfn}[Position d'équilibre instable]
	Une position d'équilibre est \textbf{instable} lorsqu'après une petite perturbation du système, le système tend à s'éloigner de cette position d'équilibre et ne pas y revenir sans une autre perturbation.
\end{dfn}
\subsection{Trajectoire de phase d'un système conservatif à un degré de liberté}
Pour un système conservatif à un degré de liberté, deux trajectoires de phase ne peuvent pas se croiser (ou alors il s'agit d'un "faux croisement" et les trajectoires repartent chacune de leur côté). \\ \par En pratique, on trace la courbe de l'énergie potentielle, et on aligne avec le portrait de phase. On trace alors le portrait de phase pour différents seuils d'énergie mécanique.
\section{Oscillateurs}
\subsection{Oscillations au voisinage d'une position d'équilibre stable}
Exemple du pendule simple ou de l'anneau coulissant sans frottement sur un cerceau. On a un système conservatif à un degré de liberté qu'on note $\theta$. On obtient une équation différentielle de la forme : $$\ddot{\theta}+\frac{g}{R}\sin(\theta)=0$$ par une méthode énergétique ou un théorème du moment cinétique. Les positions d'équilibre sont donc $\theta=0$ et $\theta=\pi$. Pour $\theta=0$, en effectuant l'approximation des petits angles, on trouve une pulsation des petites oscillations de la forme $$\omega_0=\sqrt{\frac{g}{R}}$$ En $\theta=\pi$, un DL et la résolution de l'équation différentielle donnent une solution divergente, ce qui confirme bien que cette approximation n'est pas valable très longtemps et que la position d'équilibre est instable.
\subsection{Réponse d'un oscillateur amorti à une sollicitation constante}
\begin{dfn}[Régimes pour une EDL du second ordre]
	On est en régime : 
\begin{itemize}
	\item \textbf{pseudo-périodique} si $Q>\displaystyle\frac{1}{2}$ ;
	\item \textbf{apériodique critique } (ou \textbf{critique}) si $Q=\displaystyle\frac{1}{2}$ ;
	\item \textbf{apériodique} si $Q>\displaystyle\frac{1}{2}$.
\end{itemize}
\end{dfn}
\begin{dfn}[Pseudo-pulsation, pseudo-période]
	En régime pseudo-périodique, on définit la \textbf{pseudo-pulsation} par $$\omega_{pp}=\omega_0\sqrt{1-\frac{1}{4Q^2}}$$ On définit aussi la \textbf{pseudo-période} par $$T_{pp}=\frac{2\pi}{\omega_{pp}}$$
\end{dfn}
\begin{dfn}[Décrément logarithmique]
	En régime pseudo-périodique, on définit le \textbf{décrément logarithmique} par : $$\delta=\ln\left(\frac{x(t)-x(\infty)}{x(t+T_{pp})-x(\infty)}\right)$$ On a : $$\delta=\frac{\pi}{\displaystyle\sqrt{1-\frac{1}{4Q^2}}}$$ Lorsque $Q>>1$, on a $\delta\approx\displaystyle\frac{\pi}{Q}$.
\end{dfn}
\subsection{Réponse d'un oscillateur amorti à une sollicitation sinusoïdale}
\paragraph{Passe-bas} Pour un passe-bas du second ordre, on peut avoir ou non résonance. La fonction de transfert se met sous la forme canonique : $$\underline{H}=\frac{1}{\displaystyle\left(1-\frac{\omega^2}{\omega_0^2}\right)+j\frac{\omega}{\omega_0Q}}$$ \\ \par Condition de résonance : $Q>\displaystyle\frac{1}{2}$. On a lors $\omega_{res}=\omega_0\sqrt{\displaystyle1-\frac{1}{2Q^2}}$. On a toujours $\omega_{res}<\omega_0$.
\paragraph{Passe-bande} La fonction de transfert se met sous la forme canonique : $$\underline{H}=\frac{cte}{1+jQ\displaystyle\left(x-\frac{1}{x}\right)}$$ Si on note $\omega_1<\omega_2$ les deux pulsations de coupure à $-3dB$ -- \textit{ie} telles que $G=\displaystyle\frac{G_{\max}}{\sqrt{2}}$--, alors on a : $$\Delta\omega=\omega_2-\omega_1=\frac{\omega_0}{Q}$$ Il faut savoir refaire la démonstration de ce résultat (résoudre une équation bicarrée et sélectionner les seules solutions "viables" puis en faire la différence).
\paragraph{Rappels sur le mode XY de l'oscilloscope} Pour détecter la résonance, il doit y avoir une droite passant par l'origine de pente positive. En opposition de phase, on obtient une droite passant par l'origine de pente négative. En quadrature de phase, on a une ellipse.
\subsection{Étude d'oscillateurs couplés ; modes propres ; battements}
\begin{dfn}[Mode propre]
	Pour un \textbf{mode propre}, tous les degrés de liberté d'un système sont des fonctions sinusoïdales de la même pulsation, dite pulsation propre.
\end{dfn}
\begin{dfn}[Battement]
	On a un phénomène de \textbf{battement} lorsqu'on a le produit de deux fonctions sinusoïdales dont une pulsation est très grande devant l'autre. On observe alors des oscillations sinusoïdales très rapides dans une enveloppe sinusoïdale beaucoup plus lente.
\end{dfn}
Exemple : les pendules couplés par torsion. \\ \par \textbf{Astuce} : souvent, pour l'étude de système couplés, il peut être intéressant de faire un changement de fonctions inconnues, en posant par exemple la demi-somme et la demi-différence dans le cas des pendules couplés par torsion.
\subsection{Oscillateur isotrope ; trajectoire elliptique}
Exemple : un ressort laissée libre dans l'espace. \\ \par La trajectoire de phase d'un oscillateur isotrope est une ellipse centrée sur l'origine (contrairement au forces newtoniennes dont la trajectoire de phase est une ellipse centrée ailleurs qu'à l'origine).
\subsection{Effets non-linéaires ; non-isochronisme ; apparition d'harmoniques}
Exemple : anneau coulissant sur une barre sans frottement soumis à une force de rappel élastique, avec la barre placée à la distance de la longueur à vide du ressort. \\ \par ==> Faire un DL de l'énergie potentielle au voisinage de la position d'équilibre "au milieu" puis écrire la conservation de l'énergie mécanique. On peut alors écrire $E_m$ en fonction des données avec les conditions à un instant particulier, puis cela nous permet d'obtenir $\dot{x}$ en fonction de $x$ et de constantes. On sépare les variables et on intègre pour obtenir la pulsation des petites oscillations. On montre qu'elle est égale à une constante multipliée par une donnée, si bien que cette pulsation dépend des conditions initiales. Il y a donc \textbf{non-isochronisme des oscillations}. \\ \par $x$ est alors une fonction périodique du temps non-sinusoïdale : son spectre n'est pas réduit à un fondamental.
\section{Mouvements à force centrale}
On rappelle qu'il faut absolument connaître : $$\begin{cases}
	M_{\text{Soleil}}\approx 2.10^{30}\ kg\\
	M_{\text{Terre}}\approx 6.10^{24}\ kg
\end{cases}$$
\subsection{Problème à deux corps}
On suppose le système $\{M_1,M_2\}$ isolé. On note $G$ son centre de masse et $\mu$ la masse réduite du système telle que $\displaystyle\frac{1}{\mu}=\frac{1}{m_1}+\frac{1}{m_2}$. Le référentiel $R\st$ lié à $G$ est en translation rectiligne uniforme par rapport au référentiel de départ, supposé galiléen. Les équations du mouvement sont dans $R\st$ : $$\begin{cases}
	\displaystyle m_1\frac{d^2\overrightarrow{GM_1}}{dt^2}=-\overrightarrow{f}_{1\to2}\\[4mm]
	\displaystyle m_2\frac{d^2\overrightarrow{GM_2}}{dt^2}=\overrightarrow{f}_{1\to2}
\end{cases}$$ On introduit la particule fictive $P$ de masse réduite $\mu$ telle que $\overrightarrow{GP}=\overrightarrow{M_1M_2}$ et on a alors $$\mu\frac{d^2\overrightarrow{GP}}{dt^2}=\overrightarrow{f}_{1\to2}$$ qui peut alors être plus simple résoudre selon les cas. Un cas particulier simple est celui où $m_1=m_2$. \\ \par Attention, ce n'est pas la bonne méthode en général ! Elle est surtout présentée ici à titre culturel, mais on peut toujours tenter de l'utiliser en dernier recours.
\subsection{Conservation du moment cinétique ; conséquences}
Désormais, $M$ est soumis uniquement à $\overrightarrow{f}=f(r)\overrightarrow{u_r}$ en coordonnées sphériques de centre $O$ avec $O$ fixe dans $R$ galiléen. Le théorème du moment cinétique par rapport à $O$ fixe dans $R$ galiléen donne : $$\displaystyle\frac{d}{dt}\overrightarrow{L_O}(M/R)=\overrightarrow{0}$$ (car la force est centrale), si bien que $\overrightarrow{L_O}(M/R)=\overrightarrow{cte}$. Or, $\overrightarrow{OM}\cdot\overrightarrow{L_O}(M/R)=0$ par définition du moment cinétique, donc $M$ reste dans le plan passant par $O$ de vecteur normal $\overrightarrow{L_O}(M/R)$. Si de plus le moment cinétique est nul, le mouvement est rectiligne. \\ \par On utilise alors les coordonnées polaires dans le plan de la trajectoire d'origine $O$. Avec ces coordonnées, on a : $$\overrightarrow{L_O}(M/R)=r\overrightarrow{u_r}\wedge m(\dot{r}\overrightarrow{u_r}+r\dot{\theta}\overrightarrow{u_\theta})=mr^2\dot{\theta}\overrightarrow{u_z}$$ Or, $\overrightarrow{L_O}(M/R)=\overrightarrow{cte}$, donc on obtient la \textbf{loi des aires} : $$r^2\dot{\theta}=cte=\mathcal{C}$$
\subsection{Énergie potentielle effective}
L'énergie mécanique s'écrit $$E_m=\frac{1}{2}mv^2+E_p=\frac{1}{2}m\left(\left(\dot{r}\right)^2+\left(r\dot\theta\right)^2\right)+E_p$$ avec la constante des aires, cela donne : $$E_m=\frac{1}{2}m\left(\dot{r}\right)^2+E_p+\frac{1}{2}m\frac{C^2}{r^2}$$ On définit alors l'\textbf{énergie potentielle effective} par : $$E_{p,eff}=E_p+\frac{1}{2}m\frac{C^2}{r^2}$$ Elle est dite effective car elle cache dans la constante des aires les conditions initiales du problème. L'étude de l'énergie potentielle effective permet de mener l'étude du mouvement radial. \\ \par Par exemple, pour l'oscillateur isotrope, $E_{p,eff}=\displaystyle\frac{1}{2}kr^2+\frac{1}{2}m\frac{C^2}{r^2}$.
\section{Mouvements dans un champ newtonien}
On dit qu'une force $\overrightarrow{F}$ est \textbf{newtonienne} lorsqu'il existe une constante $K$ telle que $$\overrightarrow{F}=\frac{K}{r^2}\overrightarrow{u_r}$$ Par exemple, la force gravitationnelle et la force de Coulomb sont des forces newtoniennes. Notons que toute force newtonienne est conservative, on trouve facilement qu'elle dérive de l'énergie potentielle $K/r$.
\subsection{Première loi de Kepler (admise)}
Dans le système solaire, les planètes décrivent des ellipses dont l'un des foyers est occupé par le Soleil. \\ \par Généralisation : La trajectoire d'un point matériel soumis à une force newtonienne est une conique.
\subsection{États liés : trajectoires circulaires ou elliptiques}
\paragraph{Cas circulaire}On travaille avec $\overrightarrow{F}=-\displaystyle\frac{GMm}{r^2}\overrightarrow{u_r}$. Le PFD donne : $$v=\sqrt{\frac{GM}{r}}$$ donc le mouvement est uniforme. On en déduit que $$E_m=-\frac{GMm}{2r}$$ puis que l'astre vérifie la troisième loi de Képler : $$\frac{T^2}{r^2}=\frac{4\pi^2}{GM}$$
\paragraph{Cas elliptique dans un champ newtonien}
On note $a$ le demi-grand axe, $b$ le demi-petit axe, et $c$ la demi-distance entre les foyers $F$ et $F'$. On rappelle que l'ellipse est le lieu des points M tels que la distance $FM+F'M$ soit constante égale à $2a$. On retrouve donc la formule $a^2=b^2+c^2$. En notant $x$ l'axe passant par $F$ et $F'$ et $y$ l'axe orthogonal à $x$ passant par $O$, on a $$\frac{x^2}{a^2}+\frac{y^2}{b^2}=1$$ En considérant les points d'annulation de $\dot{r}$ (périgée et apogée), et en écrivant l'équation du second degré vérifiée par $r$ en ces points, et en utilisant le fait que $r_P+r_A=2a$ ainsi que la somme des solutions d'une équation de degré 2, on obtient la généralisation : $$E_m=-\frac{GMm}{2a}$$
\subsection{États de diffusion : trajectoires parabolique ou hyperboliques}
\paragraph{Cas hyperbolique} Exemple de la déviation d'une particule. \\ \par \textbf{Astuce} : pour une force newtonienne, écrire le premier principe sous la forme : $$m\frac{d\overrightarrow{v}}{dt}=\frac{K}{r^2}\overrightarrow{u_r}=\frac{K}{r^2}\frac{\dot{\theta}}{\dot{\theta}}\overrightarrow{u_r}=-\frac{K}{\mathcal{C}}\frac{d\overrightarrow{u_\theta}}{dt}$$ permet d'intégrer pour obtenir : $$\overrightarrow{v}-\overrightarrow{v_0}=-\frac{K}{m\mathcal{C}}\left(\overrightarrow{u_\theta}-\overrightarrow{u_{\theta,0}}\right)$$
\paragraph{Cas parabolique} Exemple de la trajectoire parabolique. \\ \par \textbf{Astuce} : se souvenir que pour une trajectoire parabolique en présence d'une force newtonienne, $E_m=0$. On peut alors écrire la conservation du moment cinétique en un point particulier de la trajectoire et utiliser le fait que $E_m=0$ pour déterminer les paramètres manquants qui apparaissent.
\section{Mouvement d'une particule chargée dans un champ statique}
Très souvent, on peut négliger les interactions gravitationnelles devant la force de Lorentz. Le plan de bataille est toujours le suivant : \begin{itemize}
	\item Système, référentiel, bilan des forces
	\item Principe fondamental et projections
	\item La force de Lorentz ne peut coupler que deux équations. Pour résoudre ces équations couplées, on peut en intégrer une et réinjecter dans l'autre, ou alors poser une fonction complexe $u=x+iy$ si $x$ et $y$ sont couplées par exemple, trouver une équation différentielle en $u$, la résoudre, puis passer aux parties réelle et imaginaire.
\end{itemize}
Souvent, si le champ $\overrightarrow{B}$ est constant, on peut poser la pulsation cyclotron $\omega_c=\displaystyle\frac{\lvert q\rvert B}{m}$ et faire les calculs avec $\varepsilon=\pm 1$.

\chapter{Changements de référentiel}
\localtableofcontents
\section{Cinématique des changements de référentiel}
\subsection{Mouvement d'un référentiel par rapport à un autre}
Il faut préciser : 
\begin{itemize}
	\item la position de $O'$
	\item les orientations de $\overrightarrow{e_{x'}}$, $\overrightarrow{e_{y'}}$ et $\overrightarrow{e_{z'}}$ par rapport à $R$
\end{itemize}
\paragraph{Translation}
\begin{dfn}[Translation]
	$R'=(O',x',y',z')$ est en translation par rapport à $R=(O,x,y,z)$ lorsque $\overrightarrow{e_{x'}}$, $\overrightarrow{e_{y'}}$ et $\overrightarrow{e_{z'}}$ sont fixes par rapport à $R$. C'est-à-dire que : $$\begin{cases}
		\displaystyle \frac{d}{dt}\overrightarrow{e_{x'}}\bigg\rvert_R=\overrightarrow{0} \\[4mm]
		\displaystyle \frac{d}{dt}\overrightarrow{e_{y'}}\bigg\rvert_R=\overrightarrow{0} \\[4mm]
		\displaystyle \frac{d}{dt}\overrightarrow{e_{z'}}\bigg\rvert_R=\overrightarrow{0}
	\end{cases}$$
\end{dfn}
Attention, toutes les translations ne sont pas rectilignes. Par exemple, une cabine de grande roue est en translation circulaire.
\paragraph{Rotation uniforme autour d'un axe fixe} On peut prendre par exemple $R'$ en rotation autour de $Oz$ à la vitesse de rotation $\Omega=cte$. On a alors : $$\begin{cases}
	\displaystyle \frac{d}{dt}\overrightarrow{e_{x'}}\bigg\rvert_R=\Omega \overrightarrow{e_z}\wedge\overrightarrow{e_{x'}} \\[4mm]
	\displaystyle \frac{d}{dt}\overrightarrow{e_{y'}}\bigg\rvert_R=\Omega \overrightarrow{e_z}\wedge\overrightarrow{e_{y'}} \\[4mm]
	\displaystyle \frac{d}{dt}\overrightarrow{e_{z'}}\bigg\rvert_R=\Omega \overrightarrow{e_z}\wedge\overrightarrow{e_{z'}}
\end{cases}$$
On remarque qu'on reconnaît un mouvement de précession. Aussi, les relations précédentes restent valables si la vitesse de rotation n'est pas constante (projection dans la base fixe et dérivation).
\subsection{Dérivées temporelles d'un vecteur dans deux référentiels ; relation de la dérivation composée}
\paragraph{Translation} On a la formule de dérivation composée si $R'$ est en translation par rapport à $R$ : $$\frac{d}{dt}\overrightarrow{A}\bigg\rvert_R=\frac{d}{dt}\overrightarrow{A}\bigg\rvert_{R'}$$
\paragraph{Rotation autour d'un axe fixe}On a la formule de dérivation composée si $R'$ est en rotation par rapport à $R$ à $\overrightarrow{\Omega_{R'/R}}$ : $$\frac{d}{dt}\overrightarrow{A}\bigg\rvert_R=\frac{d}{dt}\overrightarrow{A}\bigg\rvert_{R'}+\overrightarrow{\Omega_{R'/R}}\wedge\overrightarrow{A}$$
\subsection{Lois de composition des vitesses et des accélérations}
\paragraph{Translation de $R'$ par rapport à $R$} On a la \textbf{loi de composition des vitesses} : $$\overrightarrow{v}_{M/R}=\overrightarrow{v}_{M/R'}+\overrightarrow{v}_{O'/R}$$En cas de doute sur la présence du point $O'$ dans la formule, prendre $M=O'$ dans la formule pour vérifier. Avec $M=O$, on a $\overrightarrow{v}_{O/R'}=-\overrightarrow{v}_{O'/R}$. \\ \par On a aussi la \textbf{loi de composition des accélérations} : $$\overrightarrow{a}_{M/R}=\overrightarrow{a}_{M/R'}+\overrightarrow{a}_{O'/R}$$
\paragraph{Rotation uniforme de $R'$ autour d'un axe fixe par rapport à $R$ (à $\Omega=\Omega\overrightarrow{e_z}$)} On a la \textbf{loi de composition des vitesses } : $$\overrightarrow{v}_{M/R}=\overrightarrow{v}_{M/R'}+\Omega\wedge\overrightarrow{HM}$$ avec $H$ le projeté orthogonal de $M$ sur $Oz$. \\ \par On a aussi la \textbf{loi de composition des accélérations} : $$\overrightarrow{a}_{M/R}=\overrightarrow{a}_{M/R'}+2\Omega\wedge\overrightarrow{v}_{M/R'}-\Omega^2\overrightarrow{HM}$$
\subsection{Point géométrique coïncident ; résumé}
\begin{dfn}[Point coïncident]
	Le point coïncident à $t_0$, noté $P_{t_0}$ est le point tel que : \begin{itemize}
		\item $M(t_0)=P_{t_0}$
		\item $P_{t_0}$ est fixe dans $R'$
	\end{itemize}
\end{dfn}
\begin{dfn}[Vitesse d'entraînement, accélérations d'entraînement et de Coriolis]
	Dans le cas de la translation ou de la rotation pure, on écrit $\overrightarrow{v}_{M/R}=\overrightarrow{v}_{M/R'}+\overrightarrow{v_e}$ et $\overrightarrow{a}_{M/R}=\overrightarrow{a}_{M/R'}+\overrightarrow{a_e}+\overrightarrow{a_c}$ où $\overrightarrow{v_e}$ est la \textbf{vitesse d'entraînement}, $\overrightarrow{a_e}$ est l'\textbf{accélération d'entraînement} et $\overrightarrow{a_c}$ est l'\textbf{accélération de Coriolis}. La vitesse d'entraînement est la vitesse du point coïncident dans $R$ et l'accélération est l'accélération du point coïncident dans $R$.
\end{dfn}
Voici un résumé de la situation : 
\begin{center}
	\begin{tabular}{|c|c|c|}
		\hline
		&Translation de $R'$ dans $R$&Rotation uniforme de $R'$ dans $R$ à $\overrightarrow{\Omega}$ \\[3mm]
		\hline
		$\overrightarrow{v}$&$\overrightarrow{v_e}=\overline{v}_{O'/R}$&$\overrightarrow{v_e}=\Omega\wedge\overrightarrow{HM}$ \\[3mm]
		\hline
		$\overrightarrow{a}$&$\begin{cases}
			\overrightarrow{a_e}=\overrightarrow{a}_{O'/R} \\
			\overrightarrow{a_c}=\overrightarrow{0}
		\end{cases}$ &$\begin{cases}
		\overrightarrow{a_e}=-\Omega^2\overrightarrow{HM} \\
		\overrightarrow{a_c}=2\overrightarrow{\Omega}\wedge\overrightarrow{v}_{M/R'}
		\end{cases}$ \\
		\hline
	\end{tabular}
\end{center}
Remarquer que dans le cas de la translation, aucun de ces vecteurs ne dépend du point $M$, alors qu'ils dépendent tous \textit{a priori} de $M$ dans le cas de la rotation uniforme autour d'un axe fixe.
\section{Dynamique dans un référentiel non galiléen}
\subsection{Principes de relativité galiléenne}
La relativité galiléenne part de deux postulats fondamentaux :
\begin{itemize}
	\item le "temps absolu" : il y a quelque part une horloge qui donne le même temps pour tous les référentiels qui l'utilisent ;
	\item la "longueur absolue" : il existe un étalon de longueur qui ne dépend pas du référentiel.
\end{itemize}
Ce sont là des approximations qui sont très souvent vérifiées à des échelles ni trop faibles ni trop grandes. En réalité, il y a "dilatation des durées" et "contraction des longueurs". \\ \par De plus, la relativité galiléenne postule qu'on ne peut pas distinguer deux référentiels en translation rectiligne uniforme l'un par rapport à l'autre.
\begin{thm}
	Si $R$ est galiléen et $R'$ est en translation rectiligne uniforme par rapport à $R$, alors $R'$ est galiléen.
\end{thm}

\subsection{Principe fondamental ; pseudo-forces d'inertie}
On définit la pseudo-force d'inertie d'entraînement\textbf{} par $$\overrightarrow{f_{ie}}=-m\overrightarrow{a_e}$$ On définit la \textbf{pseudo-force d'inertie de Coriolis} par : $$\overrightarrow{f_{ic}}=-m\overrightarrow{a_c}$$
On note $\overrightarrow{F}$ la résultante des vraies forces s'appliquant au système dans $R$. Alors on a le :
\begin{thm}[Principe fondamental dans $R'$]
	Si $R$ est galiléen, alors $$m\overrightarrow{a}_{M/R'}=\overrightarrow{F}+\overrightarrow{f_{ie}}+\overrightarrow{f_{ic}}$$
\end{thm}
Attention, les pseudo-forces d'inertie ne sont pas des forces au sens propre car elles ne vérifient pas la troisième loi de Newton. Toutefois, elles modifient le mouvement dans $R'$. \\ \par Dans le cas de la translation rectiligne uniforme de $R'$ par rapport à $R$, le principe fondamental dans $R'$ montre que $R'$ vérifie la deuxième loi de Newton, si bien que $R'$ est galiléen (ouf !).
\subsection{Théorèmes du moment cinétique}
\begin{thm}[Théorème du moment cinétique vectoriel]
	Si \begin{itemize}
		\item $R'$ est non galiléen et on connaît son mouvement par rapport à $R$ galiléen
		\item $A$ est fixe dans $R'$
	\end{itemize}
	alors $$\frac{d}{dt}\overrightarrow{L_A}(M/R')\bigg\rvert_{R'}=\sum_i\overrightarrow{\mathcal{M}_A}\left(\overrightarrow{f_i}\right)+\overrightarrow{\mathcal{M}_A}\left(\overrightarrow{f_{ie}}\right)+\overrightarrow{\mathcal{M}_A}\left(\overrightarrow{f_{ic}}\right)$$
\end{thm}
\begin{thm}[Théorème du moment cinétique scalaire]
	Si \begin{itemize}
		\item $R'$ est non galiléen et on connaît son mouvement par rapport à $R$ galiléen
		\item l'axe $\Delta$ est fixe dans $R'$
	\end{itemize}
	alors $$\frac{d}{dt}L_\Delta(M/R')\bigg\rvert_{R'}=\sum_i\mathcal{M}_\Delta\left(\overrightarrow{f_i}\right)+\mathcal{M}_\Delta\left(\overrightarrow{f_{ie}}\right)+\mathcal{M}_\Delta\left(\overrightarrow{f_{ic}}\right)$$
\end{thm}
\subsection{Théorèmes énergétiques}
\begin{thm}[Théorème de la puissance cinétique]
	Si $M$ est un point matériel et $R'$ est un référentiel non galiléen dont on connaît le mouvement par rapport à $R$ galiléen, alors : $$\frac{d}{dt}\left(\frac{1}{2}mv_{M/R'}^2\right)=\sum_i\overrightarrow{f_i}\cdot\overrightarrow{v}_{M/R'}+\overrightarrow{f_{ie}}\cdot\overrightarrow{v}_{M/R'}+\overrightarrow{f_{ic}}\cdot\overrightarrow{v}_{M/R'}$$ SI $R'$ est en translation ou en rotation pure par rapport à $R$, alors on a même $\overrightarrow{f_{ic}}\cdot\overrightarrow{v}_{M/R'}=0$.
\end{thm}
On dispose des variantes habituelles : théorème de l'énergie cinétique, théorème de la puissance mécanique ou encore théorème de l'énergie mécanique.
\paragraph{Cas de la translation}
Si $R'$ est en translation rectiligne uniformément accélérée par rapport à $R$, alors $\overrightarrow{f_{ie}}$ dérive d'une énergie potentielle : $$E_{p,ie}=m\overrightarrow{a}_{O'/R}\cdot\overrightarrow{OM}+cte$$
\paragraph{Cas de la rotation pure}
Si $R'$ est en rotation uniforme à $\Omega$ autour d'un axe fixe par rapport à $R$, alors $\overrightarrow{f_{ie}}$ dérive d'une énergie potentielle : $$E_{p,ie}=-\frac{1}{2}m\Omega^2HM+cte$$
\section{Applications}
\subsection{Sismographe}
Appliquer le principe fondamental dans le référentiel galiléen du repère ou dans le référentiel non galiléen lié au sismographe et en translation par rapport au premier donne le même résultat (heureusement), en faisant bien attention aux vitesses relatives.
\subsection{Manège}
Des fois, il est préférable d'appliquer le principe dans le référentiel galiléen puis d'effectuer de la cinématique pour passer dans le référentiel dont on connaît le mouvement par rapport au premier référentiel.
\section{Manifestations du caractère non galiléen du référentiel terrestre}
\subsection{Référentiels usuels}
\paragraph{Référentiel terrestre} Il s'agit d'un référentiel lié au sol. Il est en rotation uniforme autour d'un axe fixe dans le référentiel géocentrique.
\paragraph{Référentiel géocentrique} Il s'agit d'un référentiel lié au centre de masse de la Terre et dont les axes sont définis par trois étoiles lointaines considérées comme fixes. Il est en translation elliptique (presque circulaire) dans le référentiel héliocentrique.
\paragraph{Référentiel héliocentrique (ou référentiel de Képler)} Il s'agit d'un référentiel lié au centre de masse du Soleil et dont les axes sont définis par les trois mêmes étoiles lointaines que pour le référentiel géocentrique.
\paragraph{Référentiel de Copernic} Il s'agit d'un référentiel lié au centre de masse du système solaire et dont les axes sont donnés par les trois mêmes étoiles lointaines que pour les référentiels géocentrique et héliocentrique.
\paragraph{Remarque} Tous ces référentiels peuvent ou non être considérés comme galiléens. Cela dépend en réalité de l'échelle du système que l'on a à étudier. Ces distinctions font l'objet des parties 4 et 5 de ce chapitre.
\subsection{Champ de pesanteur}
\begin{dfn}[Poids]
	Le poids d'un objet est l'opposé de la force qu'il faut exercer sur cet objet pour qu'il reste immobile dans le référentiel terrestre.
\end{dfn}
Pour cette définition, on suppose le référentiel géocentrique galiléen. En appliquant le principe fondamental au référentiel terrestre en rotation uniforme autour d'un axe fixe dans le référentiel géocentrique supposé galiléen, l'accélération étant nulle car l'objet est immobile, et la force de Coriolis étant nulle car la vitesse dans le référentiel terrestre est nulle, on obtient : $$\overrightarrow{g}=-\frac{\mathcal{G}M_T}{R_T^2}\overrightarrow{u_r}+\Omega^2\overrightarrow{HM}$$ où $\overrightarrow{u_r}$ est le vecteur unitaire radial en coordonnées sphériques. Il faut connaître l'ordre de grandeur : $$\Omega\approx 7.10^{-5}\ \text{rad}.\text{s}\inv$$ Ainsi, \textbf{le poids rend compte de la force gravitationnelle et du caractère non galiléen du référentiel terrestre}.
\subsection{Déviation vers l'est}
Attention, lorsqu'on étudie un système dans le référentiel terrestre en ne supposant pas le référentiel terrestre galiléen et en supposant le référentiel géocentrique galiléen, la force d'inertie d'entraînement est déjà comptée dans le poids ! \\ \par 
On étudie la chute d'une bille lâchée sans vitesse initiale dans une mine désaffectée plutôt profonde (pour éviter le vent et autres bêtises). Le principe fondamental dans le référentiel terrestre non galiléen s'écrit alors : $$\overrightarrow{a}\rvert_{R_T}=\overrightarrow{g}-2\overrightarrow{\Omega}\wedge\overrightarrow{v}\rvert_{R_T}$$ On résout par une \textbf{méthode perturbative} : on suppose à l'ordre 0 que la vitesse est nulle, et on obtient alors une solution linéaire en $\overrightarrow{v}$. Oà réinjecte à l'ordre 1 en supposant désormais que la vitesse est égale à la solution linéaire trouvée précédemment, on intègre et une composante "vers l'est" apparaît alors. Avec $z$ colinéaire à $\overrightarrow{g}$ et $y$ vers l'est, on trouve à l'ordre 1 : $$\begin{cases}
	x=0 \\
	\displaystyle y=g\Omega\cos(\lambda)\frac{t^3}{3}\\
	\displaystyle z=-g\frac{t^2}{2}+H
\end{cases}$$ où $\lambda$ est la "latitude" repérée en fonction de $\overrightarrow{g}$. Ainsi, à la fin de la chute, il y a bien eu une déviation vers l'Est (prendre la valeur de $t$ pour laquelle $z=0$ et réinjecter dans l'expression de $y$). \\ \par
Cette expérience a été menée par Reich en 1833 : avec une latitude $\lambda=50,5$ degrés, une hauteur $H=158,5\ m$, il trouve une déviation totale vers l'Est de $2,8$ cm. Certes, le résultat est juste, mais trop peu significatif pour mettre clairement en évidence le caractère non galiléen du référentiel terrestre.
\subsection{Cyclones (vernis culturel)}
Les forces de pression sont dirigées des anticyclones vers les dépressions. \\ \par Loin de la dépression, dans l'hémisphère Nord, la force de Coriolis écarte vers l'Est la masse d'air qui se dirige vers l'anticyclone, puis l'anticyclone l'attire à mesure qu'elle s'approche de lui, ce qui fait que l'air s'enroule dans le sens trigonométrique autour de la dépression.
\section{Manifestations du caractère non galiléen du référentiel géocentrique}
Dans cette section, on suppose le référentiel héliocentrique galiléen.
\subsection{Mouvements de la Lune}
Le plan de bataille usuel pour étudier le mouvement de la Lune autour du Soleil consiste à négliger l'action du Soleil sur la Lune et à considérer le référentiel géocentrique comme galiléen. Cependant, quand on calcule l'action du Soleil sur la Lune, on se rend compte que celle-ci est deux fois plus importante que l'action de la Terre sur la Lune (oui, celle qu'on néglige est deux fois plus grande que celle que l'on conserve). Il y a donc un problème ! \\ \par 
En appliquant le principe fondamental à la Lune dans le référentiel géocentrique en translation par rapport au référentiel héliocentrique galiléen, puis en appliquant le principe fondamental à la Terre dans le référentiel héliocentrique galiléen pour trouver l'accélération d'entraînement obtenue dans la première application du principe fondamental, on trouve : $$\overrightarrow{a}\rvert_{R_G}(L)=\overrightarrow{A_T}(L)+\left[\overrightarrow{A_S}(L)-\overrightarrow{A}_{S}(T)\right]-\overrightarrow{A_L}(T)$$ Or, le terme différentiel de marée $\overrightarrow{A_S}(L)-\overrightarrow{A}_{S}(T)$ est très faible devant $\overrightarrow{A_T}(L)$, et $\displaystyle\overrightarrow{A_L}(T)=\frac{M_L}{M_T}\overrightarrow{A_T}(L)$ est négligeable devant $\overrightarrow{A_T}(L)$. Donc on retrouve $\overline{a}\rvert_{R_G}(L)\approx\overrightarrow{A_T}(L)$. \\ \par Ainsi, les deux erreurs commises (négliger l'action du Soleil et négliger le caractère non galiléen du référentiel géocentrique) se compensent.
\subsection{Théorie statique des forces de marée}
On considère une masse $\delta m$ d'eau. En appliquant la même méthode que précédemment, et en négligeant les actions gravitationnelles des astres autres que le Soleil et la Lune, on a : $$\delta m\overrightarrow{a}\rvert_{R_G}(M)=\overrightarrow{f}+\delta m\left(\left[\overrightarrow{A_L}(M)-\overrightarrow{A_L}(T)\right]+\left[\overrightarrow{A_S}(M)-\overrightarrow{A_S}(T)\right]\right)$$ où $\overrightarrow{f}$ est la résultante des autres forces qui ne nous intéressent pas (poids, pression, etc.) \\ \par Considérons le terme différentiel de marée $\overrightarrow{A_L}(M)-\overrightarrow{A_L}(T)$. Le terme $\overrightarrow{A_L}(T)$ s'applique partout sur la Terre de la même façon (\textit{cf} théorème de Gauss dans la chapitre d'électrostatique). Mais le terme $\overrightarrow{A_L}(M)$ dépend de la position de $M$ sur la Terre. On a alors une forme d'océan statique qui ressemble à un ballon de rugby et est invariante par rotation autour de l'axe Terre-Lune. Cette forme de ballon de rugby et la rotation propre de la Terre expliquent la présence de deux marées par jour. \\ \par On pourrait réaliser la même étude pour le terme différentiel de marée dû au Soleil. Toutefois, ce terme est 10 fois moins important que celui lié à la Lune, mais reste très grand devant les termes différentiels dus aux autres planètes. \\ \par On s'intéresse désormais à l'effet combiné de la Lune et du Soleil. Lorsque la Lune et le Soleil sont en phase ou en opposition de phase (pleine Lune ou nouvelle Lune), les effets de la Lune et du Soleil se renforcent : ce sont les \textbf{vives-eaux} (tous les 14 jours). Lorsque la Lune et le Soleil sont en quadrature, leurs effets se compensent : ce sont les \textbf{mortes-eaux}.
\section{Complément : de la transformation de Galilée à celle de Lorentz-Poincaré}
On considère $R'$ en translation rectiligne uniforme par rapport à $R$ à la vitesse $V\overrightarrow{e_x}$. $R'$ et $R$ ont les mêmes axes, et on prend $O'=O$ à $t=0$. Alors, la \textbf{transformation de Galilée} s'écrit : $$\begin{cases}
	x'=x'+Vt \\
	y=y' \\
	z=z' \\
	t=t'
\end{cases}$$ Cette transformation est une \textbf{approximation} de la \textbf{transformation de Lorentz-Poincaré} : $$\begin{cases}
x=\gamma(x'+\beta ct') \\
y=y' \\
z=z' \\
ct=\gamma(ct'+\beta x')
\end{cases}$$ avec $\beta=\displaystyle\frac{v}{c}$ et $\gamma=\displaystyle\frac{1}{\sqrt{1-\beta^2}}$. \\ \par Si on note $\tilde{a}$ l'ordre de grandeur d'une grandeur $a$, alors on retrouve la transformation de Galilée à partir de la transformation de Lorentz-Poincaré si $\tilde{v}\ll c$ et $c\tilde{t}>>\tilde{x}$.

\chapter{Mécanique des systèmes matériels}
\localtableofcontents
\section{Éléments cinétiques des systèmes de points matériels}
\subsection{Deux descriptions}
\paragraph{Description microscopique} On liste les points avec leurs positions et leurs masses.
\paragraph{Description mésoscopique}
On a trois descriptions mésoscopiques : 
\begin{itemize}
	\item \textbf{volumique} : $\delta m=\rho\delta\tau$ avec $\rho$ la masse volumique ;
	\item \textbf{surfacique} : $\delta m=\sigma\delta S$ avec $\sigma$ la masse surfacique ;
	\item \textbf{linéique} : $\delta m=\mu\delta l$ avec $\mu$ la masse linéique.
\end{itemize}
\subsection{Éléments cinétiques}
\paragraph{Masse} Pour un système $\Sigma$ constitué de points matériels $M_i$ de masse $m_i$, on définit la masse du système par $$M_\Sigma=\sum_i m_i$$
\paragraph{Quantité de mouvement (ou résultante cinétique)} Dans le référentiel $R$, elle est définie par $$\overrightarrow{p}(\Sigma/R)=\sum_im_i\overrightarrow{v}(M_i/R)$$ Avec $G$ le centre d'inertie du système $\Sigma$, cela se réécrit : $$\overrightarrow{p}(\Sigma/R)=M_\Sigma\overrightarrow{v}(G/R)$$ par définition du centre d'inertie.
\paragraph{Moment cinétique}Dans le référentiel $R$, par rapport au point $A$, il se définit par $$\overrightarrow{L_A}(\Sigma/R)=\sum_i\overrightarrow{L_A}(M_i/R)=\sum_i m_i\overrightarrow{AM_i}\wedge\overrightarrow{v}(M_i/R)$$ De plus, on a la formule hors-programme du \textbf{premier théorème de König} $$\overrightarrow{L_A}(\Sigma/R)=\overrightarrow{AG}\wedge M_\Sigma\overrightarrow{v}(G/R)+\overrightarrow{L}(\Sigma/R\st)$$ avec $R\st$ le référentiel barycentrique. En effet, comme $\overrightarrow{p}(\Sigma/R\st)=\overrightarrow{0}$, le moment cinétique dans le référentiel barycentrique est un torseur couple. Puis : 
\begin{align*}
	\overrightarrow{L_A}(\Sigma/R)&=\sum_i\overrightarrow{AM_i}\wedge m_i\overrightarrow{v}(M_i/R) \\
	&=\sum_i\overrightarrow{AM_i}\wedge m_i\left(\overrightarrow{v}(M_i/R\st)+\overrightarrow{v}(G/R)\right) \\
	&=\overrightarrow{L_A}(\Sigma/R\st)+\sum_i\overrightarrow{AM_i}\wedge m_i\overrightarrow{v}(G/R) \\
	&=\overrightarrow{L_A}(\Sigma/R\st)+\overrightarrow{AG}\wedge M_\Sigma\overrightarrow{v}(G/R) \\
\end{align*}
\paragraph{Énergie cinétique} L'énergie cinétique du système $\Sigma$ dans $R$ est définie par $$E_c(\Sigma/R)=\sum_i \frac{1}{2}m_i\left(v(M_i/R)\right)^2$$ De plus, on a le \textbf{deuxième théorème de König} (hors-programme) : $$E_c(\Sigma/R)=E_c(\Sigma/R\st)+\frac{1}{2}M_\Sigma v(G/R)^2$$ En effet : 
\begin{align*}
	E_c(\Sigma/R)& =\sum_i \frac{1}{2}m_iv(M_i/R)^2 \\
	& = \sum_i \frac{1}{2}m_i\left\lVert \overrightarrow{v}(M_i/R)\right\rVert^2 \\
	& = \sum_i \frac{1}{2}m_i\left\lVert \overrightarrow{v}(M_i/R\st)+\overrightarrow{v}(G/R)\right\rVert^2 \\
	& = \sum_i \frac{1}{2}m_i\left(v(M_i/R\st)^2+v(G/R)^2+2\overrightarrow{v}(M_i/R\st)\cdot\overrightarrow{v}(G/R)\right) \\
	& = E_c(\Sigma/R\st)+\frac{1}{2}M_\Sigma v(G/R)^2+M_\Sigma\overrightarrow{v}(G/R\st)\cdot\overrightarrow{v}(G/R) \\
	& = E_c(\Sigma/R\st)+\frac{1}{2}M_\Sigma v(G/R)^2
\end{align*}
\section{Dynamique et énergétique des systèmes matériels}
\subsection{Actions subies par un système matériel}
On distingue les forces d'origine intérieure au système qui s'appliquent au système, les forces d'origine extérieure qui s'appliquent au système et les forces hors-sujet qui ne s'appliquent pas au système.
\subsection{Théorèmes dynamiques pour un système matériel}
\begin{dfn}[Système fermé]
	En mécanique, on dira qu'un système de points matériels est fermé lorsqu'il ne "perd pas" de points au cours du temps, et que ses points sont de masse constante.
\end{dfn}
\begin{thm}[Théorème de la résultante]
	Si : \begin{itemize}
		\item $\Sigma$ est un système fermé
		\item $R$ est galiléen
	\end{itemize} alors : $$M_\Sigma\overrightarrow{a}(G/R)=\frac{d}{dt}\left(\sum_i m_i\overrightarrow{v}(M_i/R)\right)=\overrightarrow{F_{ext}}$$ où $\overrightarrow{F_{ext}}$ est la résultante des forces extérieures qui s'appliquent sur le système.
\end{thm}
La troisième loi de Newton montre que la résultante des forces intérieures est nulle. Si $R$ est non galiléen mais on connaît son mouvement par rapport à un référentiel galiléen, il suffit de rajouter la résultante des forces d'inertie d'entraînement et de Coriolis obtenues pour chaque point.
\begin{thm}[Théorème du moment cinétique]
	Si : \begin{itemize}
		\item $\Sigma$ est un système fermé
		\item $R$ est galiléen
		\item $A$ est fixe dans $R$
	\end{itemize} alors : $$\frac{d}{dt}\left(\overrightarrow{L_A}(\Sigma/R)\right)=\overrightarrow{\mathcal{M}_A}_{ext}$$ où $\overrightarrow{\mathcal{M}_A}_{ext}$ est la somme des moments des actions extérieures s'appliquant au système au point $A$.
\end{thm}
Si on a appliqué le théorème de la résultante et le théorème du moment cinétique, appliquer le théorème du moment cinétique en un autre point ne fournit aucune information réellement nouvelle car les deux moments sont alors reliés par la formule BABAR à la résultante. \\ \par La démonstration du théorème (utilisation des deux parties de la troisième loi de Newton) montre que la somme des moments des actions intérieures au point $A$ est nulle. Dans le cas non galiléen, il suffit d'ajouter les moments des forces d'inertie d'entraînement pour chaque point.
\subsection{Théorème énergétiques pour un système matériel}
\textbf{Attention}, la puissance des forces intérieures n'est pas toujours nulle (penser à ses cuisses quand on monte un escalier).
\begin{thm}[Théorème de la puissance cinétique]
	Si \begin{itemize}
		\item $\Sigma$ est un système fermé ;
		\item $R$ est galiléen
	\end{itemize} alors $$\frac{d}{dt}\left(E_c(\Sigma/R)\right)=\mathcal{P}_{ext}+\mathcal{P}_{int}$$ où $\mathcal{P}_{ext}=\displaystyle\sum_i \overrightarrow{f}_{ext\to i}\cdot\overrightarrow{v}(M_i/R)$ et $\mathcal{P}_{int}=\displaystyle\sum_i \overrightarrow{f}_{int\to i}\cdot\overrightarrow{v}(M_i/R)$
\end{thm}
On dispose des variantes habituelles : théorème de l'énergie cinétique, théorème de la puissance mécanique et théorème de l'énergie mécanique. \\ \par Les actions intérieures dérivent d'une énergie potentielle lorsque $$\delta W=-dE_{p,int} \ \ \text{ou} \ \ \mathcal{P}_{int}=-\frac{d}{dt}E_{p,int}$$ où $E_{p,int}$ ne dépend que de la position des points du système.
\section{Statique des fluides}
Quand on parle de statique des fluides, le caractère statique est évidemment macroscopique, il y a en effet toujours un mouvement microscopique. \\ \par Albert Einstein et Jean Perrin mettent en évidence le mouvement brownien en 1905. En 1905, Einstein publie 4 articles majeurs : $E=mc^2$, l'effet photoélectrique, la relativité restreinte et le mouvement brownien.
\subsection{Principe fondamental de la statique des fluides ; équivalent volumique des forces de pression}
On a $$\overrightarrow{\delta F}_{\text{fluide}\to\text{paroi}}=+P\delta S\overrightarrow{n}$$ et $$\overrightarrow{\delta F}_{\text{pression}}=-\grad(P)\delta \tau$$ On dit que $-\grad(P)$ est une densité volumique de force. Elle s'exprime en $N.m^{-3}$. Une pression est une énergie volumique.
\begin{thm}[Principe fondamental de la statique des fluides]
	Si un fluide est macroscopiquement au repos dans un référentiel galiléen, alors : $$\overrightarrow{0}=-\grad(P)\delta\tau+\overrightarrow{f}\delta\tau$$ où $\overrightarrow{f}\delta\tau$ est la résultante des forces autres que la pression s'exerçant sur $\delta\tau$. Si le référentiel n'est pas galiléen : $$\overrightarrow{0}=-\grad(P)\delta\tau+\overrightarrow{f}\delta\tau+\rho\delta\tau\overrightarrow{a_e}+\overrightarrow{0}$$ car $\overrightarrow{f_{ic}}=\overrightarrow{0}$ puisque la vitesse relative est nulle.
\end{thm}
$P$ est continue à la traversée d'une interface si on néglige la tension de surface. En présence d'un ménisque, cela n'est plus vrai sans quoi les masses volumiques de l'air et de l'eau seraient égales.
\subsection{Application : cas d'un liquide incompressible}
\paragraph{Exemple : Expérience de Pascal} On utilise le baromètre de Torricelli rempli de mercure de masse volumique constante. Dans le mercure, en supposant $P=P(z)$, on obtient avec le principe fondamental de la statique des fluides $\displaystyle\frac{dP}{dz}=-\rho g$ donc, puisque la pression de vapeur saturante du mercure est très faible devant la pression atmosphérique $P_0$, on obtient $P_0=\rho g H$ avec $H$ la hauteur de mercure dans le baromètre.
\paragraph{Ordres de grandeur à connaître} Pour l'eau, $\rho =10^3 \ kg.m^{-3}$. En prenant $g\approx 10\ m.s^{-2}$, on trouve que la pression augmente de $1$ bar si on descend de $10$ mètres sous l'eau. \\ Pour le mercure, $\rho_{Hg}\approx 13,6.10^{-3}\ kg.m^{-3}$. Au niveau de la mer (pression atmosphérique valant $1$ bar), on trouve alors $H\approx \ 76\ cm$.
\subsection{Application : cas d'un gaz parfait isotherme, modèle de l'atmosphère}
Les hypothèses de ce modèle sont : 
\begin{itemize}
	\item une atmosphère isotherme à $T_0$ ;
	\item une atmosphère de masse molaire $M$ uniforme ;
	\item l'air est un gaz parfait ;
	\item l'atmosphère est macroscopiquement au repos ;
	\item $P=P(z)$
\end{itemize}
Dans ce cas, $\displaystyle\frac{dP}{dz}=-\rho g$ avec l'équation de la statique des fluides. Or, avec l'équation d'état du gaz parfait appliquée à un volume mésoscopique d'air, $\rho=\displaystyle\frac{PM}{RT_0}$. on résolvant l'équation différentielle, on obtient : $$P(z)=P_0\exp\left(-\frac{gM}{RT_0}z\right)$$ En posant $H=\displaystyle\frac{RT_0}{gM}$, on a en prenant $M\approx 30 \ g.mol\inv$ et $T_0\approx 300\ K$ la valeur approchée $H\approx 8\ km$.
\subsection{Poussée d'Archimède}
Un objet entièrement plongé dans un (des) fluide(s) et immobile dans un référentiel galiléen (ou ayant un mouvement très lent) subit une force $$\overrightarrow{\pi}=-m_{\text{fluide déplacé}}\overrightarrow{g}$$ Le fluide déplacé a un caractère fictif. On rappelle qu'à 1 bar et 20 degrés Celsius, $\rho_{\text{air}}\approx 1,3\ g.L\inv$.
\section{Solide}
\subsection{Définition ; cinématique}
\begin{dfn}[Solide]
	Un objet est un solide lorsqu'il est indéformable, c'est-à-dire que la distance entre deux points quelconques de cet objet est une constante indépendante du temps.
\end{dfn}
En considérant un référentiel lié au solide $S$, on obtient alors : $$\overrightarrow{v}(B\in S/R)=\overrightarrow{v}(A\in S/R)+\overrightarrow{BA}\wedge\overrightarrow{\Omega}(S/R)$$
\subsection{Théorèmes dynamiques}
\begin{thm}[Théorème de la résultante]
	Si $S$ est un solide de masse $M_S$ et $R$ est galiléen, alors : $$\frac{d}{dt}\left(M_S\overrightarrow{v}(G/R)\right)=\overrightarrow{F}_{\text{ext}}$$
\end{thm}
\begin{thm}[Théorème du moment cinétique vectoriel]
	Si $S$ est un solide de masse $M_S$, $R$ est galiléen et $A$ est fixe dans $R$, alors : $$\frac{d}{dt}\left(\overrightarrow{L_A}(S/R)\right)=\overrightarrow{\mathcal{M}_A}_{\text{ext}}$$
\end{thm}
\begin{thm}[Théorème du moment cinétique scalaire]
	Si $S$ est un solide de masse $M_S$, $R$ est galiléen et $\Delta$ est un axe fixe dans $R$, alors : $$\frac{d}{dt}\left(L_\Delta(S/R)\right)=\mathcal{M}_{\Delta,\text{ext}}$$
\end{thm}
Pour l'instant, rien de nouveau par rapport aux systèmes de points matériels.
\subsection{Théorèmes énergétiques}
\begin{thm}
	Si on considère un unique solide, alors $\mathcal{P}_{\text{int}}=0$, et par conséquent $W_{\text{int}}=0$.
\end{thm}
Par conséquent, on obtient une version simplifiée des théorèmes énergétiques : 
\begin{thm}[Théorème de la puissance cinétique pour un solide]
	SI $S$ est un solide et $R$ est galiléen, alors : $$\frac{d}{dt}\left(E_c(S/R)\right)=\mathcal{P}_{\text{ext}}(S/R)$$
\end{thm}
Les variantes usuelles (énergie cinétique, puissance mécanique et énergie mécanique) sont alors vraies avec la simplification.
\subsection{Solide en translation}
Puisque tous les points du solide on la même vitesse que le centre d'inertie du solide si celui-ci est en translation, on a : $$\overrightarrow{L_A}(S/R)=M_S\overrightarrow{AG}\wedge\overrightarrow{v}(G/R)$$ et $$E_c(S/R)=\frac{1}{2}M_S\times v(G/R)^2$$
\section{Solide en rotation autour d'un axe fixe}
\subsection{Moment et énergie cinétiques pour un solide en rotation autour d'un axe fixe}
\begin{dfn}[Moment d'inertie]
	Le moment d'inertie d'un solide $S$ par rapport à un axe fixe $\Delta$ est défini par $$J_\Delta(S)=\sum_{i\in S} m_id(\Delta,M_i)^2$$ Plus généralement, on a en fait $$J_\Delta(S)=\iiint_S r^2\delta m$$ En pratique, on calcule le moment d'inertie élémentaire en remplaçant $\delta$ grâce aux données, puis on intègre sur le paramétrage du solide.
\end{dfn}
Le moment d'inertie est l'analogue de la masse pour les solides en rotation. En effet, on a alors les identités suivantes : 
\begin{thm}
	Si le solide $S$ est en rotation autour de l'axe fixe $\Delta$ à la vitesse angulaire $\Omega$ dans le référentiel $R$ galiléen, alors : $$E_c(S/R)=\frac{1}{2}J_\Delta(S)\Omega^2$$ et $$L_\Delta(S/R)=J_\Delta(S)\Omega$$
\end{thm}
\subsection{Théorème du moment cinétique pour un solide en rotation autour d'un axe fixe}
\begin{thm}
	Si le solide $S$ est en rotation autour d'un axe fixe $\Delta$ dans un référentiel galiléen, alors $$\frac{d}{dt}E_c(S/R)=\mathcal{P}_{\text{ext}}(S/R)$$
\end{thm}
\subsection{Couple}
\begin{dfn}[Couple]
	Une action mécanique est un couple lorsque sa résultante est nulle. Par abus de langage, on appelle souvent couple le moment d'une action mécanique qui est un couple puisque le moment d'un couple ne dépend pas du point où il est calculé.
\end{dfn}
Par exemple, appuyer de façon antisymétrique sur les barres d'un gouvernail est un couple.
\subsection{Liaison pivot}
Le rotor tourne autour du stator. La liaison pivot est dite parfaite (ou sans frottement) lorsque $\mathcal{M}_{\Delta,\text{liaison}}=0$
\subsection{Rotation autour d'un axe mobile (complément)}
Exemple : moto qui fait une roue arrière sans glisser. \\ \par Il faut : \begin{itemize}
	\item appliquer le théorème de la résultante à l'ensemble de la moto dans le référentiel terrestre
	\item appliquer le théorème du moment cinétique au système {châssis, roue avant} par rapport à l'axe de la roue arrière qui est fixe dans le référentiel de la moto, lui-même en translation par rapport au référentiel terrestre supposé galiléen
	\item appliquer le théorème du moment cinétique à la roue arrière par rapport à l'axe de la roue arrière fixe dans le référentiel de la moto
	\item combiner toutes les équations scalaires pour exprimer l'accélération de la moto en fonction de la masse $M$ de la moto, du rayon $R$ de la roue arrière, de $g$, de l'angle $\alpha$ que fait la moto avec le sol, de la distance $d$ entre le centre de la roue arrière et le centre d'inertie de la moto et du moment d'inertie $J$ de la roue arrière par rapport à son axe de rotation.
\end{itemize} On obtient (pour la science) : $$a_0=\frac{Mgd\cos(\alpha)}{\displaystyle\frac{J}{R}+MR+Md\sin(\alpha)}$$
\subsection{Résumé}
Tout est dans les parties précédentes !

\chapter{Lois du frottement solide}
\localtableofcontents
\section{Présentation et définitions}
Il est important de comprendre que, même si nous avons jusqu'à présent très souvent négligé les frottements, ceux-ci sont en réalité nécessaire à presque tout. En effet, sans frottement, on ne pourrait pas avancer lorsqu'on marche (essayez à la patinoire et vous allez rigoler...). Il est donc vital de s'intéresser à ces forces pour mieux comprendre leur fonctionnement. \\ \par
On s'intéresse au contact de deux solides $S_1$ et $S_2$ qui a lieu en un point I. On appelle vitesse de glissement du solide 2 par rapport au solide 1 la vitesse : $$\overrightarrow{v}_{\text{glissement }2/1}:=\overrightarrow{v}_{I\in S_2/R_0}-\overrightarrow{v}_{I\in S_1/R_0}$$
Cette expression est invariante par changement de référentiel puisque les termes d'accélération d'entraînement d'inertie et de Coriolis se compensent alors.
\section{Composante normale de la réaction}
Une liaison unilatérale peut être rompue. Pour une liaison unilatérale, on connaît toujours le sens de la réaction. \\ 
\par Inversement, on ne connaît pas le sens de la réaction d'une liaison bilatérale, mais la liaison ne peut pas être rompue de façon simple.
\section{Composante tangentielle de la réaction}
\subsection{Cas 1/2 : non-glissement ou adhérence}
Si $\overrightarrow{v}_{\text{glissement }2/1}=\overrightarrow{0}$, alors : $$\left\lVert\overrightarrow{R_T}\right\rVert\leqslant f_S\left\lVert\overrightarrow{R_N}\right\rVert$$ avec $f_S$ le coefficient de frottement statique, qui est sans dimension. $f_S$ dépend de la nature des solides, ne dépend pas de l'aire de la surface de contact et ne dépend pas de $R_N$.
\subsection{Cas 2/2 : glissement}
Si $\overrightarrow{v}_{\text{glissement }2/1}\ne\overrightarrow{0}$, alors : $$\left\lVert\overrightarrow{R_T}\right\rVert= f_D\left\lVert\overrightarrow{R_N}\right\rVert$$ avec $f_D$ le coefficient de frottement dynamique (ou cinétique), qui est sans dimension. $f_D$ dépend de la nature des solides, ne dépend pas de l'aire de la surface de contact et ne dépend pas de $R_N$. \\
\par On a toujours $f_D\leqslant f_S$, et, bien souvent, $f_D\approx f_S$.
\subsection{Ordres de grandeur}
Pour un contact métal contre métal, on a : $$f_S\approx f_D\approx 0,2$$
\subsection{Contact sans frottement}
Si $f_S=f_D=0$, alors $\overrightarrow{R_T}=\overrightarrow{0}$ d'après les lois de Coulomb. Dans les exercices, cette condition est traduite par la mention "sans frottement".
\subsection{Interprétation microscopique}
En réalité, une surface de contact telle que nous la modélisons n'est pas réellement pleine de contacts. En fait, les surfaces ont plein de micro-aspérités (de l'ordre de la dizaine de micromètres) et ne se touchent qu'en certains points. \\
\par Pour passer de l'adhérence au glissement, il faut qu'il y ait rupture de ces micro-contacts. \\
\par Ainsi, il ne faut pas confondre la surface apparente de contact (c'est-à-dire macroscopique) avec la surface de contact réelle (c'est-à-dire microscopique et beaucoup moins importante).
\subsection{Exemple 1 : masse initialement au repos sur un plan incliné}
Un solide $S$ de masse $m$ est initialement au repos en étant posé sur un support incliné par rapport au sol d'un angle $\displaystyle\alpha\in\left[0,\frac{\pi}{2}\right]$. Les forces qui s'appliquent sur le solide $S$ sont la réaction normale du support, la réaction tangentielle du support et le poids. On se place dans le repère "naturel" lié à la pente avec $x$ orienté dans le sens de la pente et $z$ orthogonal à la surface dirigé vers l'extérieur de la surface.
\paragraph{Cas 1 : adhérence}
On applique le théorème de la résultante au solide $S$ dans le référentiel galilén lié au support pour obtenir : $$\overrightarrow{0}=m\overrightarrow{g}+\overrightarrow{R_T}+\overrightarrow{R_N}$$ car le solide est alors immobile. En projetant, on obtient :
\begin{itemize}
	\item sur $\overrightarrow{u_x}$ : $\overrightarrow{R_T}\cdot\overrightarrow{u_x}+0+mg\sin(\alpha)=0$
	\item sur $\overrightarrow{u_z}$ : $mg\cos(\alpha)+R_N+0=0$
\end{itemize}
Or d'après la loi de Coulomb sur l'adhérence, on a $R_T\leqslant fR_N$. On obtient donc $mg\sin(\alpha)\leqslant fmg\cos(\alpha)$, soit : $$\tan(\alpha)\leqslant f$$
\paragraph{Cas 2} On suppose $\overrightarrow{v}_{\text{glissement}}$ selon $+\overrightarrow{u_x}$ (solide en translation). On applique le théorème de la résultante au solide $S$ dans le référentiel galiléen lié au support et on projette selon $\overrightarrow{u_x}$ et $\overrightarrow{u_z}$ pour obtenir : $$\begin{cases}
	m\ddot{x}=mg\sin(\alpha)-R_T \\
	0=-mg\cos(\alpha)+R_N
\end{cases}$$ On suppose $\dot{x}>0$. Alors, $\overrightarrow{v}_{\text{glissement}}=\dot{x}\overrightarrow{u_x}\ne\overrightarrow{0}$. Puis, d'après la loi de Coulomb sur le glissement, on a $$\left\lVert\overrightarrow{R_T}\right\rVert=f\left\lVert\overrightarrow{R_N}\right\rVert=fmg\cos(\alpha)$$
\subsection{Cône de frottement ; arc-boutement}
Que ce soit dans le cas du glissement ou de l'adhérence, la réaction du support est contenue dans un cône vertical d'angle au sommet $\alpha$ tel que $\tan(\alpha)=f$. ce cône est appelé cône de frottement. \\ \par Pour passer de l'adhérence au glissement, augmenter la norme de la réaction ne sert à rien car on reste dans le cône de frottement. Il faut plutôt modifier la direction de la réaction. C'est le phénomène d'arc-boutement. --> Exemple du serre-joint 
\section{Aspect énergétique}
\textbf{Attention}, les frottements peuvent être moteurs ! Leur effet global est toujours résistant, mais peut être moteur pour une partie du système --> penser au deuxième principe en thermodynamique : l'entropie d'un système fermé peut diminuer à condition d'augmenter plus vite celle du reste du monde (frigo).
\subsection{Puissance des forces de contact sur un des deux solides}
On suppose que le solide 2 est en translation par rapport au référentiel d'étude. Alors $$\mathcal{P}_{\text{action }1\to 2}=\overrightarrow{R}_{1\to 2}\cdot\overrightarrow{v_{G_2}}$$
\subsection{Puissance totale des forces de contact entre deux solides}
On suppose que les solides $1$ et $2$ sont en translation par rapport au référentiel d'étude. Alors : $$\mathcal{P}_{\text{tot}}=\overrightarrow{R}_{1\to 2}\cdot\left(\overrightarrow{v_{G_2}}-\overrightarrow{v_{G_1}}\right)$$ Ainsi, s'il y a glissement ou si $f=0$, on a $P_{\text{tot}}=0$. Sinon, on a $\mathcal{P}_{tot}<0$. \\ \par Dans tous les cas, $\mathcal{P}_{\text{tot}}\leqslant0$. Il y a conversion d'énergie mécanique en énergie thermique.
\subsection{Exemple 2 : entraînement par frottement}
Exemple d'un cadeau de masse $m$ sur un traineau de masse $M$. Le cadeau frotte sur le traineau, le traineau glisse sans frotter sur le sol et est soumis à une force horizontale $\overrightarrow{F}$ constante qui le tracte. Condition sur $\overrightarrow{F}$ pour que le cadeau ne glisse pas ? \\ \par On distingue les cas où il y a adhérence et où il y a glissement. Dans le cas de l'adhérence, on applique le théorème de la résultante à ${M,m}$ puis le théorème de la résultante à ${m}$. Dans le cas du glissement, on applique le théorème de la résultante à ${m}$ puis on l'applique à ${M}$. La limite est $F=f(m+M)g$ et il y a glissement au-dessus. Tant que le cadeau n'est pas tombé du traineau, les frottements sont moteurs pour lui. 

\part{Électromagnétisme}

\chapter{Analyse vectorielle}
\localtableofcontents
\section{Opérateurs}

\subsection{Gradient}
L'opérateur gradient est un champ de vecteurs : il prend en argument scalaire qui dépend de l'espace et donne un vecteur. On le note $\grad$. Le gradient est linéaire : si $\alpha$ et $\beta$ ne dépendent pas de l'espace et $f$ et $g$ sont des scalaires, alors : $$\grad(\alpha f+\beta g)=\alpha.\grad(f)+\beta.\grad(g)$$
On a : $$\left[\grad(f)\right]=\frac{[f]}{L}$$ 
Le gradient possède une formule intrinsèque qui ne dépend pas du système de coordonnées ni de l'orientation de la base choisie. Cette formule intrinsèque est : $$df=\grad(f)\cdot d\overrightarrow{l}$$ 
En coordonnées cartésiennes : $$\grad(f)=\frac{\partial f}{\partial x}\overrightarrow{u_x}+\frac{\partial f}{\partial y}\overrightarrow{u_y}+\frac{\partial f}{\partial z}\overrightarrow{u_z}$$
En coordonnées cylindriques : $$\grad(f)=\frac{\partial f}{\partial r}\overrightarrow{u_r}+\frac{1}{r}\frac{\partial f}{\partial \theta}\overrightarrow{u_\theta}+\frac{\partial f}{\partial z}\overrightarrow{u_z}$$
En coordonnées sphériques : $$\grad(f)=\frac{\partial f}{\partial \rho}\overrightarrow{u_\rho}+\frac{1}{\rho}\frac{\partial f}{\partial \theta}\overrightarrow{u_\theta}+\frac{1}{\rho\sin(\theta)}\frac{\partial f}{\partial \varphi}\overrightarrow{u_\varphi}$$
On peut introduire le vecteur nabla : $$\overrightarrow{\nabla}:=\begin{pmatrix}
	\displaystyle\frac{\partial}{\partial x} \\[3mm]
	\displaystyle\frac{\partial}{\partial y} \\[3mm]
	\displaystyle\frac{\partial}{\partial z}
\end{pmatrix}$$ qui permet de mémoriser les expressions des différents opérateurs en \textbf{coordonnés cartésiennes uniquement !} Il ne faut \textbf{surtout pas essayer de généraliser ces expressions aux autres systèmes de coordonnées}. Cet avertissement ét	nt fait, on peut retenir le gradient en coordonnées cartésiennes par : $$\grad(f)=\overrightarrow{\nabla}f$$
\subsection{Divergence}
L'opérateur divergence est un champ scalaire : il prend en argument un vecteur et donne un scalaire. On le note $\div$. La divergence est linéaire : si $\alpha$ et $\beta$ ne dépendent pas de l'espace et $\overrightarrow{a}$ et $\overrightarrow{b}$ sont des vecteurs, alors : $$\div(\alpha \overrightarrow{a}+\beta\overrightarrow{b})=\alpha.\div(\overrightarrow{a})+\beta.\div(\overrightarrow{b})$$
On a : $$\left[\div(\overrightarrow{a})\right]=\frac{[a]}{L}$$ 
La divergence possède une formule intrinsèque qui ne dépend pas du système de coordonnées ni de l'orientation de la base choisie. Cette formule intrinsèque est : $$\delta\Phi =\div(\overrightarrow{a})\times\delta\tau$$ avec $\delta\tau$ le volume élémentaire délimité par une surface $\delta S$ et $\delta\Phi$ le flux sortant de $\delta S$. \\ \\
En coordonnées cartésiennes : $$\div(\overrightarrow{a})=\frac{\partial a_x}{\partial x}+\frac{\partial a_y}{\partial y}+\frac{\partial a_z}{\partial z}$$
Pour une expression à symétrie cylindrique ($\overrightarrow{a}=a_r(r)\overrightarrow{u_r}$) : $$\div(\overrightarrow{a})=\frac{1}{r}\frac{\partial(ra_r)}{\partial r}$$
Pour une expression à symétrie sphérique ($\overrightarrow{a}=a_\rho(\rho)\overrightarrow{u_\rho}$) : $$\div(\overrightarrow{a})=\frac{1}{\rho^2}\frac{\partial(\rho^2a_\rho)}{\partial \rho}$$
On a le :
\begin{thm}[Théorème de Green-Ostrogradski]
	Soit $V$ un volume délimité par une surface $S$ fermée. Alors :
	$$\iiint_V\div(\overrightarrow{a}) \ \delta\tau = \oiint_S\overrightarrow{a}\cdot dS\overrightarrow{n_{ext}} $$
\end{thm}
\subsection{Rotationnel}
L'opérateur rotationnel prend un argument vectoriel et donne un résultat vectoriel. On le note $\rot$. Il est linéaire et on a : $$\left[\rot(\overrightarrow{a})\right]=\frac{[\overrightarrow{a}]}{L}$$
En coordonnées cartésiennes, on a : 
$$\rot(\overrightarrow{a})=\left(\frac{\partial a_z}{\partial y}-\frac{\partial a_y}{\partial z}\right)\overrightarrow{u_x}+\left(\frac{\partial a_x}{\partial z}-\frac{\partial a_z}{\partial x}\right)\overrightarrow{u_y}+\left(\frac{\partial a_y}{\partial x}-\frac{\partial a_x}{\partial y}\right)\overrightarrow{u_z}$$
On peut retenir cette formule \textbf{en coordonnées cartésiennes uniquement} avec : $$\rot(\overrightarrow{a})=\overrightarrow{\nabla}\wedge\overrightarrow{a}$$
On a de façon générale le :
\begin{thm}[Théorème de Stokes]
	Soit $\Gamma$ un contour fermé et $S$ une surface s'appuyant sur $\Gamma$. Alors : $$\oint_\Gamma \overrightarrow{a}\cdot d\overrightarrow{l}=\iint_S \rot(\overrightarrow{a})\cdot dS\overrightarrow{n_{ext}}$$ On prendra garde au fait que l'orientation du contour et $\overrightarrow{n_{ext}}$ sont liés par la règle de la main droite.
\end{thm}
On peut aussi écrire cette relation sous forme locale avec $$\delta C=\rot(\overrightarrow{a})\cdot dS\overrightarrow{n}$$ où $\delta C$ est la \textbf{circulation}.
\subsection{Laplacien scalaire}
En coordonnées cartésiennes, le \textbf{laplacien scalaire} de $f$ est défini par : $$\Delta f=\frac{\partial^2f}{\partial x^2}+\frac{\partial^2f}{\partial y^2}+\frac{\partial^2f}{\partial z^2}$$ Le laplacien scalaire est un opérateur linéaire qui prend en argument un scalaire et donne un scalaire. On a : $$\left[\Delta f\right]=\frac{[f]}{L^2}$$
On trouve parfois l'écriture en coordonnées cartésiennes : $$\Delta f=\overrightarrow{\nabla}\cdot\overrightarrow{\nabla}f=\nabla^2f$$
Les expressions du laplacien scalaire en coordonnées cylindriques et sphériques sont bien trop compliquées et hors-programme. Elles seront données si nécessaire. Tenter de les apprendre serait une grave erreur car elles sont beaucoup trop complexes et le risque d'erreur est énorme.
\subsection{Laplacien vectoriel}
En coordonnées cartésiennes, le \textbf{laplacien vectoriel} du champ vectoriel $\overrightarrow{a}$ est défini par : $$\overrightarrow{\Delta}(\overrightarrow{a})=(\Delta a_x)\overrightarrow{u_x}+(\Delta a_y)\overrightarrow{u_y}+(\Delta a_z)\overrightarrow{u_z}$$ où $\Delta a_x$, $\Delta a_y$ et $\Delta a_z$ sont respectivement les laplaciens scalaires en coordonnées cartésiennes des scalaires $a_x$, $a_y$ et $a_z$. \\ \\
Les expressions du laplacien vectoriel en coordonnées cylindriques et sphériques sont bien trop compliquées et hors-programme. Elles seront données si nécessaire. Tenter de les apprendre serait une grave erreur car elles sont beaucoup trop complexes et le risque d'erreur est énorme.
\section{Relations d'analyse vectorielle}
$$\div\left(\rot(\overrightarrow{a})\right)=0$$
$$\rot\left(\grad (f)\right)=\overrightarrow{0}$$
$$\div\left(\grad(f)\right)=\Delta f$$
$$\rot\left(\rot(\overrightarrow{a})\right)=\grad\left(\div(\overrightarrow{a})\right)-\overrightarrow{\Delta}(\overrightarrow{a})$$
$$\div\left(f\overrightarrow{a}\right)=f\div(\overrightarrow{a})+\grad(f)\cdot\overrightarrow{a}$$
$$\rot\left(f\overrightarrow{a}\right)=f\rot(\overrightarrow{a})+\grad(f)\wedge\overrightarrow{a}$$
$$\grad(fg)=f\grad(g)+g\grad(f)$$
$$\div\left(\overrightarrow{a}\wedge\overrightarrow{b}\right)=\rot(\overrightarrow{a})\cdot\overrightarrow{b}-\rot(\overrightarrow{b})\cdot\overrightarrow{a}$$
\section{Deux cas particuliers}
\subsection{Définitions}
\begin{dfn}[Ligne de champ]
	En tout point $M$ d'une ligne de champ $\mathcal{L}$, $\overrightarrow{E}(M)$ est tangent à $\mathcal{L}$ en $M$.
\end{dfn}
Par exemple, la trajectoire d'un point matériel peut être considérée comme la ligne de champ de $\overrightarrow{v}$, à ceci près que $\overrightarrow{v}$ ne dépend pas de l'espace mais du temps. L'image est tout de même celle qu'il faut retenir.
\begin{dfn}[Tube de champ]
	Un tube de champ est un ensemble de lignes de champ qui s'appuient sur un contour fermé et l'intérieur de la zone délimitée par ces lignes de champ.
\end{dfn}
\subsection{Champ à circulation conservative}
Un champ $\overrightarrow{f}$ est dit \textbf{à circulation conservative} quand il vérifie les propriétés équivalentes suivantes :
\begin{itemize}
	\item Formulation intégrale 1 : Pour tous contours $\Gamma_1$ et $\Gamma_2$ reliant deux points $A$ et $B$ en allant de $A$ à $B$, on a $\displaystyle\oint_{\Gamma_1}\overrightarrow{f}\cdot\overrightarrow{dl}=\oint_{\Gamma_2}\overrightarrow{f}\cdot\overrightarrow{dl}$
	\item Formulation intégrale 2 : Pour tout contour fermé $\Gamma$, $\displaystyle\oint_\Gamma\overrightarrow{f}\cdot\overrightarrow{dl}=0$
	\item Formulation locale 1 : En tout point, $\rot(\overrightarrow{f})=\overrightarrow{0}$
	\item Formulation locale 2 : Il existe un champ scalaire $u$ tel que $\overrightarrow{f}=\grad(u)$
\end{itemize}
Une force conservative est un cas particulier de champ vectoriel à circulation conservative.
\subsection{Champ à flux conservatif}
Un champ $\overrightarrow{f}$ est dit \textbf{à flux conservatif} quand il vérifie les propriétés équivalentes suivantes :
\begin{itemize}
	\item Formulation intégrale 1 : Pour toute surface fermée $S$, on a $\displaystyle\oiint_S\overrightarrow{f}\cdot\overrightarrow{dS}=0$
	\item Formulation intégrale 2 : Si $S_1$ est une surface s'appuyant sur le contour fermé $\Gamma_1$ et $S_2$ s'appuie sur $\Gamma_2$ où $\Gamma_1$ et $\Gamma_2$ délimitent un tube de champ, alors $\displaystyle\iint_{S_1}\overrightarrow{f}\cdot\overrightarrow{dS}=\iint_{S_2}\overrightarrow{f}\cdot\overrightarrow{dS}$
	\item Formulation intégrale 3 : Pour tout contour fermé $\Gamma$ et toutes surfaces $S_1$ et $S_2$ s'appuyant sur $\Gamma$, alors $\displaystyle\iint_{S_1}\overrightarrow{f}\cdot\overrightarrow{dS}=\iint_{S_2}\overrightarrow{f}\cdot\overrightarrow{dS}$
	\item Formulation locale 1 : En tout point, $\div(\overrightarrow{f})=\overrightarrow{0}$
	\item Formulation locale 2 : Il existe un champ vectoriel $\overrightarrow{A}$ tel que $\overrightarrow{f}=\rot(\overrightarrow{A})$
\end{itemize}
Remarquer que la dernière formulation intégrale prouve la validité du théorème de Stokes peu importe le choix de la surface.

\chapter{Électrostatique}
\localtableofcontents
\section{Des charges au champ électrique}
\subsection{Distributions de charges}
\paragraph{Description microscopique}
On liste les charges et leurs positions.
\paragraph{Description mésoscopique}
On dispose de trois modélisation microscopiques :
\begin{itemize}
	\item linéique : $\delta Q=\lambda \delta l$ avec $\lambda$ la charge linéique en $C.m\inv$
	\item surfacique : $\delta Q=\sigma \delta S$ avec $\sigma$ la charge surfacique en $C.m^{-2}$
	\item volumique : $\delta Q=\rho \delta t$ avec $\rho$ la charge volumique en $C.m^{-3}$
\end{itemize}
Attention la modélisation surfacique est une simplification car l'épaisseur de la distribution de charges n'est pas nulle en réalité.
\paragraph{Passage d'une description à l'autre} Pour passer d'une description à l'autre, on se place à grande distance du système et on regarde les charges dans un volume mésoscopique.
\subsection{Loi de Coulomb et principe de superposition}
Pour deux charges $(P_1,q_1)$ et $(P_2,q_2)$, on a la loi de Coulomb électrostatique : $$\overrightarrow{F_{1\to 2}}=\frac{q_1q_2\overrightarrow{P_1P_2}}{4\pi\varepsilon_0(P_1P_2)^3}$$ On peut l'écrire $\overrightarrow{F_{1\to 2}}=q_2\overrightarrow{E}(P_2)$ avec $$\overrightarrow{E}(M)=\frac{q_1\overrightarrow{P_1M}}{4\pi\varepsilon_0(P_1M)^3}$$ \par Plus généralement, si on prend un nombre quelconque de charges, le champ électrique s'écrit selon le \textbf{principe de superposition} : $$\overrightarrow{E}(M)=\sum_i \frac{q_i\overrightarrow{P_iM}}{4\pi\varepsilon_0(P_iM)^3}$$
\subsection{Propriétés de symétrie ; principe de Curie}
\paragraph{Plans de symétrie des charges}
Pour un plan de symétrie des charges, un point et son symétrique ont des champs électriques symétriques par rapport à ce plan. Par conséquent, pour tout point appartenant à un plan de symétrie des charges, le champ électrique en ce point est contenu dans le pan de symétrie.
\paragraph{Plans d'antisymétrie des charges}
Pour un plan d'antisymétrie des charges, tout point contenu dans le plan a un champ électrique orthogonal à ce plan d'antisymétrie.
\paragraph{Invariances}Lorsqu'on a une invariance par rotation ou par translation de la distribution des charges, toutes les composantes du champ électrique ne dépendent pas de la coordonnée qui correspond à l'invariance.
\paragraph{Principe de Curie} "Les effets ont au moins les symétries des causes." Dans notre cas, cela signifie que $\overrightarrow{E}$ possède globalement les mêmes symétries que la distribution de charges.
\subsection{Exemple (complément)}
Attention, le calcul direct n'est pas la bonne méthode la plupart du temps. il convient toujours d'étudier les symétries et les invariances puis de tenter d'appliquer le théorème de Gauss si l'on dispose de suffisamment de symétries et d'invariances.
\section{Flux du champ électrique ; théorème de Gauss}
\subsection{Formulation intégrale ; théorème de Gauss}
\begin{thm}[Théorème de Gauss électrostatique]
	Pour toute surface \textbf{fermée} $S$ et tout champ électrique $\overrightarrow{E}$, on a : $$\oiint_S \overrightarrow{E}\cdot\overrightarrow{dS}=\frac{Q_{int}}{\varepsilon_0}$$ avec $Q_{int}$ la charge totale contenue à l'intérieur de la surface $S$.
\end{thm}
\subsection{Notion d'angle solide (complément)}
\subsection{Formulation locale ; équation de Maxwell-Gauss}
\begin{thm}[Équation de Maxwell-Gauss]
	On a $$\div(\overrightarrow{E})=\frac{\rho}{\varepsilon_0}$$
\end{thm}
Le théorème de Green-Ostrogradski permet de montrer que le théorème de Gauss électrostatique et l'équation de Maxwell-Gauss sont équivalents.
\subsection{Exemples}
\paragraph{Fil infini}On considère un fil infini de charge linéique $\lambda$ uniforme. En coordonnées cylindriques, les symétries montrent que le champ n'a qu'une composante radiale. Les invariances montrent que cette composante radiale ne dépend que de la distance radiale. On choisit comme surface de Gauss un cylindre de rayon $r$ et de hauteur $H$ quelconque. L'application du théorème de Gauss donne : $$\overrightarrow{E}=\frac{\lambda}{2\pi\varepsilon_0 r}\overrightarrow{u_r}$$
\paragraph{Sphère chargée en volume}On considère une boule uniformément chargée en volume de charge volumique $\rho$ de rayon $R$. On se place en coordonnées sphériques de centre le centre de la sphère chargée. Les symétries montrent que le champ n'a qu'une composante selon $\overrightarrow{u_r}$, et l'étude des symétries montre que cette composante ne dépend que de $r$. On applique le théorème de Gauss à une sphère de même centre que la sphère chargée et on obtient : $$\overrightarrow{E}=\begin{cases}
	\displaystyle\frac{\rho r}{3\varepsilon_0}\overrightarrow{u_r} \ \ \text{si} \ \ r\leqslant R \\[4mm]
	\displaystyle \frac{\rho R^3}{3\varepsilon_0 r^2}\overrightarrow{u_r} \ \  \text{si} \ \ r\geqslant R 
\end{cases}$$
\paragraph{Sphère chargée en surface}On considère une sphère de rayon $R$ uniformément chargée en surface à $\sigma$. De même que pour la boule, on a $\overrightarrow{E}=E_r(r)\overrightarrow{u_r}$. En appliquant le théorème de Gauss à une sphère de même centre que la sphère chargée, on obtient : $$\overrightarrow{E}=\begin{cases}
	\overrightarrow{0} \ \ \text{si} \ \ r<R \\[1mm]
	\displaystyle\frac{\sigma R^2}{\varepsilon_0 r^2}\overrightarrow{u_r} \ \ \text{si} \ \ r>R
\end{cases}$$ Attention, si $r=R$, on ne peut rien dire. Noter la discontinuité de $\displaystyle\frac{\sigma}{\varepsilon_0}$ en $r=R$.
\paragraph{Généralisation en symétrie sphérique}
A l'extérieur d'une distribution de charges à symétrie sphérique, le champ électrique est le même que si toute la charge était concentrée au centre de la sphère (on le retrouve en appliquant le théorème de Gauss à une sphère de rayon $r>R$ et en constatant que le champ est le même que celui créé par une charge ponctuelle placée au centre de la sphère). \\ \par Attention, cela n'est valable qu'à l'extérieur de la sphère.
\paragraph{Plan infini chargé en surface}On considère un plan infini chargé uniformément en surface à $\sigma$. On repère ce plan par l'équation $z=0$ en coordonnées cartésiennes. L'étude des symétries montre que $\overrightarrow{E}=E_z(z)\overrightarrow{u_z}$ et que la composante $E_z$ est impaire. On applique le théorème de Gauss à un cylindre dont les plans sont $z$ et $-z$ pour $z>0$. On obtient : $$\overrightarrow{E}=\begin{cases}
	\displaystyle\frac{\sigma}{2\varepsilon_0}\overrightarrow{u_z} \ \ \text{si} \ \ z>0 \\[4mm]
	\displaystyle\frac{-\sigma}{2\varepsilon_0}\overrightarrow{u_z} \ \ \text{si} \ \ z<0
\end{cases}$$ Attention, si $z=0$, on ne peut rien dire, il y a un problème de modélisation. Noter la discontinuité de $\displaystyle\frac{\sigma}{\varepsilon_0}$ en $z=0$.
\section{Circulation du champ électrique ; potentiel électrique}
\subsection{Cas d'une charge ponctuelle ; principe de superposition}
\paragraph{Deux charges ponctuelles en interaction}
$$\overrightarrow{E}(M)=\frac{q\overrightarrow{PM}}{4\pi\varepsilon_0(PM)^3}=-\grad\left(\frac{q}{4\pi\varepsilon_0 PM}\right)$$ si bien que si $\Gamma$ est fermé, alors $\displaystyle\oint_\Gamma\overrightarrow{E}\cdot\overrightarrow{dl}=0$. On définit alors le potentiel créé par la charge $q$ placée en $P$ au point $M$ par : $$V(M)=\frac{q}{4\pi\varepsilon_0 PM}+cte$$
\paragraph{Nombre quelconque de charges ; principe de superposition}
Si les charges $q_i$ sont placées aux points $P_i$ alors le potentiel en $M$ est donné par le principe de superposition : $$V(M)=\sum_i\frac{q_i}{4\pi\varepsilon_0 P_iM}+cte$$
\paragraph{Choix de la constante}
On a l'apparition d'une constante à cause du gradient. Lorsque les distributions sont d'extension finie, on peut facilement prendre une constante nulle à l'infini par exemple. Sinon, on peut s'arranger pour prendre la constante nulle en un point particulier du système. Toutefois, lorsqu'on a une distribution de charges à l'infini, on ne peut pas toujours prendre une constante nulle à l'infini.
\paragraph{Remarques}
En \textbf{électrostatique} :
\begin{itemize}
	\item $\overrightarrow{E}$ est un gradient
	\item $\overrightarrow{E}$ est à circulation conservative
\end{itemize} 
\subsection{Relation champ-potentiel ; équation de Maxwell-Faraday statique}
\begin{thm}[Équation de Maxwell-Faraday statique]
	En \textbf{statique}, on a : $$\rot(\overrightarrow{E})=\overrightarrow{0}$$
\end{thm}
Attention, en régime variable on a en réalité : $$\rot(\overrightarrow{E})=-\frac{\partial\overrightarrow{B}}{\partial t}$$ et $$\overrightarrow{E}=-\grad(V)-\frac{\partial \overrightarrow{A}}{\partial t}$$
\subsection{Exemples}
Reprendre les exemples de calcul de champ électrostatique et intégrer. Pour se débarrasser des constantes d'intégration, on peut utiliser la continuité du potentiel aux points de jonction des domaines. Sinon, on peut aussi effectuer un calcul direct, ce qui est beaucoup plus pénible.
\subsection{Continuité de potentiel}
En électrostatique, le potentiel est continu, même en présence d'une charge surfacique.
\begin{center}
	\begin{tabular}{|c|c|c|}
		\hline
		&$\overrightarrow{E}$&$V$ \\
		\hline
		$\rho$&continu&continu \\
		\hline
		$\sigma$&discontinu&continu \\
		\hline
		$\lambda$&diverge ($1/r$)&diverge ($\ln(r)$) \\
		\hline
		$q$&diverge ($1/r^2$)&diverge ($1/r$)\\
		\hline
	\end{tabular}
\end{center}
\section{Conséquence des deux propriétés précédentes}
\subsection{Relation de passage}
A la traversée d'une surface chargée à $\sigma$, en notant $\overrightarrow{E_2}$ le champ à l'extérieur de la surface, $\overrightarrow{E_1}$ le champ à l'intérieur de la surface, et $\overrightarrow{n_{1\to 2}}$ le vecteur normal à la surface dirigé de $1$ vers $2$, on a une discontinuité du champ électrique donnée par la \textbf{relation de passage} : $$\overrightarrow{E_2}-\overrightarrow{E_1}=\frac{\sigma}{\varepsilon_0}\overrightarrow{n_{1\to 2}}$$ Pour le prouver, appliquer un théorème de Gauss à un cylindre dont on fait tendre les dimensions vers $0$. La composante normale est discontinue de $\displaystyle\frac{\sigma}{\varepsilon_0}$ (\textit{cf} l'exemple sur le plan chargé). Pour les autres composantes, on utilise le fait que $\overrightarrow{E}$ est à circulation conservative et que certains flux tendent vers $0$.
\subsection{Équations de Poisson et de Laplace ; théorème d'unicité}
\begin{thm}[Équation de Poisson]
	En \textbf{statique}, on a : $$\Delta V+\frac{\rho}{\varepsilon_0}=0$$
\end{thm}
\begin{thm}[Équation de Laplace]
	En \textbf{statique}, dans une zone vide de charges : $$\Delta V=0$$
\end{thm}
En pratique, l'équation de Poisson et l'équation de Laplace permettent de passer immédiatement des charges au potentiel sans passer par le champ, et donc sans utiliser l'équation de Maxwell-Gauss ou le théorème de Gauss. Néanmoins, on ne peut les utiliser que si on est en coordonnées cartésiennes ou si on dispose d'un formulaire donnant le laplacien scalaire en coordonnées cylindriques ou sphériques.
\begin{thm}[Théorème d'unicité]
	Soit $V$ un volume délimité par une surface $S$. Si
	\begin{itemize}
		\item $V=0$ en tout point de $S$
		\item $\Delta V=0$ en tout point intérieur de $V$
	\end{itemize} alors $V=0$ en tout point de $V$.
\end{thm}
\subsection{Théorème de l'extremum}
Le potentiel ne peut pas atteindre un maximum en un point vide de charges. Il est donc impossible de confiner des charges avec un potentiel électrostatique seul. \\ \par Pour des exemples clairs, \textit{cf} transferts thermiques et bulles de savon.
\section{Topographie du champ électrique}
\subsection{Définitions}
\paragraph{Ligne de champ}Le champ électrique est tangent en tout point à la ligne de champ.
\paragraph{Équipotentielle}L'équipotentielle à $V_0$ est l'ensemble des points $M$ tels que $V(M)=V_0$.
\paragraph{Exemple de la charge ponctuelle chargée positivement} Les lignes de champ sont des demi-droites qui partent de la charge ponctuelle. Les équipotentielles sont des sphères concentriques centrées sur la charge ponctuelle.
\subsection{Champ à circulation conservative ; conséquence}
Les lignes de champ pour un champ $\overrightarrow{E}$ statique ne sont pas fermées : le long d'une ligne de champ, le potentiel décroît strictement, si bien qu'une ligne de champ ne peut pas revenir au potentiel initial.
\subsection{Champ à flux conservatif dans une zone vide de charges ; conséquence}
Dans une zone vide de charges, le champ $\overrightarrow{E}$ est à flux conservatif. Ainsi, dans une zone vide de charges, la norme du champ $\overrightarrow{E}$ diminue lorsque les lignes de champ s'écartent.
\section{Analogie avec la gravitation}
\paragraph{Généralités et passage d'un domaine à l'autre} \textbf{Très important} : il suffit de se rappeler quelles constantes doit-on remplacer par lesquelles, et quelle lettre par laquelle, et les théorèmes gravitationnels sont alors les mêmes que les théorèmes électrostatiques modulo les changements de lettres et de constantes.
\begin{center}
	\begin{tabular}{|c|c|}
		\hline
		$\overrightarrow{E}$&$\overrightarrow{A}$ (champ de gravitation)\\[4mm]
		\hline
		$\displaystyle\overrightarrow{E}(M)=\sum_i\frac{q_i\overrightarrow{P_iM}}{4\pi\varepsilon_0(P_iM)^3}$&$\displaystyle\overrightarrow{A}(M)=\sum_i-\frac{\mathcal{G}m_i\overrightarrow{P_iM}}{(P_iM)^3}$ \\[4mm]
		\hline
		$1/4\pi\varepsilon_0$&$-\mathcal{G}$ \\[4mm]
		\hline
		$q_i$&$m_i$ \\[4mm]
		\hline
	\end{tabular}
\end{center}
\paragraph{Circulation}
\begin{center}
	\begin{tabular}{|c|c|}
		\hline
		$\displaystyle\oint_\Gamma\overrightarrow{E}\cdot\overrightarrow{dl}=0$&$\displaystyle\oint_\Gamma\overrightarrow{A}\cdot\overrightarrow{dl}=0$ \\[4mm]
		\hline
		$\overrightarrow{E}=-\grad(V)$&$\overrightarrow{A}=-\grad(\Phi)$ \\[4mm]
		\hline
		$E_p=q_{test}V$&$E_p=m_{test}\Phi$ \\[4mm]
		\hline
		$\rot(\overrightarrow{E})=\overrightarrow{0}$&$\rot(\overrightarrow{A})=\overrightarrow{0}$ \\[4mm]
		\hline
	\end{tabular}
\end{center}
\paragraph{Flux} Probablement la partie la plus importante : retenir le théorème de Gauss gravitationnel, l'équation de Maxwell-Gauss gravitationnelle et l'équation de Poisson gravitationnelle
\begin{center}
	\begin{tabular}{|c|c|}
		\hline
		$\displaystyle\oiint_S\overrightarrow{E}\cdot\overrightarrow{dS}=\frac{Q_{int}}{\varepsilon_0}$&$\displaystyle\oiint_S\overrightarrow{A}\cdot\overrightarrow{dS}=-4\pi \mathcal{G}M_{int}$\\[4mm]
		\hline
		$\div(\overrightarrow{E})=\displaystyle\frac{\rho}{\varepsilon_0}$&$\div(\overrightarrow{A})=-4\pi \mathcal{G}\mu$\\[4mm]
		\hline
		$\Delta V+\displaystyle\frac{\rho}{\varepsilon_0}=0$&$\displaystyle\Delta \Phi-4\pi \mathcal{G}\mu=0$\\
		\hline
	\end{tabular}
\end{center}
\section{Énergie potentielle électrostatique}
\subsection{Charge ponctuelle dans un champ extérieur}
On s'intéresse à l'énergie potentielle de la charge $q_{test}$ placée en $M$, dans le champ électrostatique créé par les charges $q_i$. On a :$$\overrightarrow{F_{q_{test}}}=q_{test}\overrightarrow{E}(M)=-\grad(q_{test}V)$$ La force électrostatique dérive donc d'une énergie potentielle $E_p$ telle que $$E_p=qV+cte$$
\subsection{Énergie d'un système de charges ponctuelles en interaction}
Pour  un système de charges en interaction, on a : $$E_{p,int}=\sum_{i<j}\frac{q_iq_j}{4\pi\varepsilon_0 P_iP_j}=\frac{1}{2}\sum_{i,j}\frac{q_iq_j}{4\pi\varepsilon_0P_iP_j}$$
\section{Modélisation des condensateurs}
Pour nos condensateurs, en prenant la convention de prendre la tension $u$ de la charge négative vers la charge positive (la flèche sur le schéma électrique dans ce sens-là), on $$Q=Cu$$ avec $C$ la capacité du condensateur en Farad ($F$).
\subsection{Condensateur plan}
On se place dans le cas où deux armatures planes de surface $S$ chargées uniformément en surface à $\sigma$ et $-\sigma$ sont espacées parallèlement d'une épaisseur $e$. On néglige les effets de bord, si bien que $\overrightarrow{E}=E_z(z)\overrightarrow{u_z}$ où $\overrightarrow{u_z}$ est unitaire dirigé de l'armature négative vers l'armature positive. On place le plan positif à $\displaystyle\frac{e}{2}$ et le négatif à $-\displaystyle\frac{e}{2}$.\\ \par On applique le principe de superposition avec les deux armatures pour obtenir un champ nul dès que $\lvert z\rvert >\displaystyle\frac{e}{2}$ et valant $\displaystyle\frac{-\sigma}{\varepsilon_0}$ si $\lvert z\rvert >\displaystyle\frac{e}{2}$. Ainsi, pour $\lvert z\rvert <\displaystyle\frac{e}{2}$, $V(z)=\displaystyle\frac{\sigma}{\varepsilon_0}z+cte$. Donc $$u=V\left(\frac{e}{2}\right)-V\left(-\frac{e}{2}\right)=\frac{\sigma}{\varepsilon_0}e$$Comme $Q=\sigma S$, on en déduit $$C=\frac{\varepsilon_0 S}{e}$$
\begin{thm}[Condensateur plan]
	Pour un condensateur plan, en négligeant les effets de bord ($e^2\ll S$), on a : $$C=\frac{\varepsilon_0 S}{e}$$ où $\varepsilon_0$ est la permittivité diélectrique du vide qui vaut $8,89.10^{-12}\ F.m\inv$. \\ SI on a un matériau isolant entre les deux armatures qui n'est pas du vide, on remplace $\varepsilon_0$ par $\varepsilon_r\varepsilon_0$ qui est la permittivité relative du matériau. Si c'est le cas, cela est donné par le sujet (\textit{cf} chimie).
\end{thm}
\subsection{Condensateur cylindrique}
On prend un cylindre de rayon $R_1$ chargé $Q$ et un cylindre concentrique chargé $-Q$ de rayon $R_2>R_1$. Les deux cylindres sont de hauteur $H$. On néglige les effets de bord en supposant les armatures infinies pour calculer le champ électrique. \\ \par En appliquant le théorème de Gauss, on trouve que $E_r(r)=\displaystyle\frac{Q}{2\pi\varepsilon_0 rH}$ lorsque $R_1<r<R_2$. En intégrant, $V(r)=-\displaystyle\frac{Q}{2\pi\varepsilon_0 H}\ln\left(\frac{r}{a}\right)$ pour $R_1<r<R_2$ où $a$ est une longueur pour la prise de potentiel nul. On obtient alors : $$u=\frac{Q}{2\pi\varepsilon_0 H}\ln\left(\frac{R_2}{R_1}\right)=\frac{Q}{C}$$ D'où : $$C=\frac{2\pi\varepsilon_0 H}{\ln\left(\frac{R_2}{R_1}\right)}$$
\paragraph{Remarques} Un condensateur de hauteur $H''=H+H'$ s'obtient en mettant en parallèle deux condensateurs de hauteurs $H$ et $H'$. L'expression de la capacité obtenue dans les deux cas précédents est alors compatible avec $C''=C'+C$ en parallèle. \\ \par Si on place des condensateurs en série, alors les expressions des capacités sont compatibles avec la relation : $$\frac{1}{C_{1,3}}=\frac{1}{C_{1,2}}+\frac{1}{C_{2,3}}$$ \\ \par Ainsi, les relations obtenues sont compatibles avec les relations sur les impédances complexes.\\ \par Dans le cas limite où $R_1=R$ et $R_2=R+e$ avec $e\ll R$, on retrouve l'expression de la capacité d'un condensateur plan en prenant un équivalent de $\ln(1+u)$ pour $u\ll 1$.
\subsection{Condensateur sphérique}
\textit{cf} TD. L'idée est la même que pour le condensateur cylindrique.
\subsection{Énergie électrique emmagasinée}
\paragraph{Énergie emmagasinée} On rappelle que $$W=\frac12Cu^2=\frac12\frac{Q^2}{C}$$ On admet que $$W=\iiint_{\R^3}\frac{\varepsilon_0E^2}{2}\delta \tau$$ où $\displaystyle\frac{\varepsilon_0E^2}{2}$ est en $J.m^{-3}$
\paragraph{Cas du condensateur plan}
Avec la formule admise, puisque le champ "à l'intérieur" vaut $\displaystyle\frac{\sigma}{\varepsilon_0}$ et celui "à l'extérieur est nul, on a $\displaystyle W=\frac{\varepsilon_0}{2}\left(\frac{\sigma}{\varepsilon_0}\right)^2eS$. Or, $\displaystyle W=\frac{1}{2}\frac{(\sigma S)^2}{C}$. On retrouve alors la formule $$C=\frac{\varepsilon_0 S}{e}$$
\paragraph{Condensateur cylindrique}
Le champ "à l'intérieur" vaut $\displaystyle\frac{Q}{2\pi\varepsilon_0 rH}$ "à l'intérieur" et est nul "à l'extérieur". Avec la même idée que pour le condensateur plan, on retrouve en intégrant et en égalant la formule : $$C=\frac{2\pi\varepsilon_0 H}{\ln\left(\frac{R_2}{R_1}\right)}$$
\section{Conducteur à l'équilibre (complément)}
Dans un conducteur, si $\overrightarrow{E}\ne\overrightarrow{0}$, alors il y a un mouvement macroscopique des charges.
\paragraph{Conducteur en équilibre} On a alors $\overrightarrow{E}=\overrightarrow{0}$. On en déduit que $V$ est uniforme à l'intérieur du conducteur. Aussi, $\rho=0$ à l'intérieur du conducteur grâce à l'équation de Maxwell-Gauss. Ainsi, les charges sont donc présente en surface. On a donc finalement $$Q=\oiint_S\sigma \ dS$$ D'après la relation de passage, juste à l'extérieur du conducteur proche de celui-ci, on a : $$\overrightarrow{E}=\frac{\sigma}{\varepsilon_0}\overrightarrow{n}$$ où $\overrightarrow{n}$ est unitaire dirigé de l'intérieur du conducteur vers l'extérieur.
\paragraph{Influence} Charger un conducteur par influence consiste à charger ce conducteur sans le toucher, en approchant une charge par exemple. \\ \par Par exemple, si on considère un conducteur en équilibre relié à une masse et qu'on en approche une charge positive, alors le conducteur se charge négativement en surface.
\paragraph{Effet de point} Si on place une pointe, les équipotentielles doivent nécessairement se resserrer au voisinage de cette pointe. Cela fait que le champ électrique est beaucoup plus intense au voisinage de la pointe : c'est l'effet de pointe. \\ \par Cet effet explique que la foudre tombe le plus souvent sur des objets pointus.
\chapter{Dipôle électrostatique}
\localtableofcontents
\section{Champ et potentiel crées par un dipôle}
\subsection{Présentation}
Exemples : les molécules polaires comme $H_20$, $HCl$ ou encore $CO_2$. On rappelle que la charge formelle consiste à couper toutes les liaisons en deux puis à faire un bilan des charges et enlever la charge de valence. Alors, la charge totale de la molécule ou de l'ion est égale à la somme des charges partielles. Attention, il ne faut pas confondre la charge formelle avec la règle de l'octet ou la règle du duet, ou encore avec les nombres d'oxydation. \\ \\
On rappelle que la géométrie $AXE$ correspond pour le $X$ au nombre d'atomes auquel l'atome central est lié, et le $E$ au nombre de doublets non liants autour de l'atome central. \\ \\
Une distribution dipolaire de charges est une distribution d'extension finie, globalement neutre et dont le barycentre des charges négatives est distinct du barycentre des charges positives. \\ \\
On se place dans le cas d'un dipôle doublet $\left\{P,N\right\}$ où $P$ est chargé $+q$ et $N$ est chargé $-q$ avec $q>0$. On note $a$ la longueur $PN$ et on note $O$ le milieu du segment $[PN]$. On note $z$ l'axe vertical allant de $N$ vers $P$. On se place en coordonnées sphériques $(r,\theta,\varphi)$ centrées en $O$ et avec la colatitude définie par rapport à $Oz$. On a $E_\varphi=0$ car pour tout point $M$, $(M,\overrightarrow{u_r},\overrightarrow{u_\theta})$ est un plan de symétrie des charges. On a une invariance par rotation autour de l'axe $Oz$ donc $E_r$ et $E_\theta$ ne dépendent pas de $\varphi$.
\subsection{Potentiel créé par le dipôle dans l'approximation dipolaire}
Approximation dipolaire : $r\gg a$ \\ \\
Dans l'approximation dipolaire, en utilisant le principe de superposition et le potentiel créé par une charge ponctuelle, puis en négligeant les termes d'ordre $2$ et en faisant un DL, on obtient \textbf{dans l'approximation dipolaire} : $$V(M)=\frac{qa}{4\pi\varepsilon_0}\times\frac{\cos(\theta)}{r^2}$$
\subsection{Moment dipolaire}
On définit le \textbf{moment dipolaire} $\overrightarrow{p}$ par $\overrightarrow{p}:=q\overrightarrow{NP}$. On pose $p=qa$. On appelle indifféremment $\overrightarrow{p}$ et $p$ le moment dipolaire. L'unité SI du moment dipolaire est le $C.m$. \\ \\
Plus généralement, pour une distribution dipolaire quelconque, son moment dipolaire est défini par : $$\overrightarrow{p}=\sum_i q_i\overrightarrow{AP_i}$$ et cette définition est indépendante du point $A$ de référence choisi car $\displaystyle\sum_iq_i=0$. \\ \\
Ordre de grandeur : pour une molécule
\begin{itemize}
	\item $a\approx 10^{-10}m$
	\item $q\approx e = 1,60.10^{-19}C$
	\item donc $p\approx 1,60.10^{-29}C.m$
\end{itemize}
On introduit alors une unité, le Debye (noté $D$), telle que $1D\approx 3,33.10^{-30}C.m$. \\ \\
On dispose d'une formule intrinsèque pour le potentiel créé par le dipôle dans l'approximation dipolaire. Avec $r=OM$ et $\overrightarrow{u}=\displaystyle\frac{\overrightarrow{OM}}{OM}$, on a : $$V(M)=\frac{\overrightarrow{p}\cdot\overrightarrow{u}}{4\pi\varepsilon_0 r^2}$$
\subsection{Champ créé par le dipôle dans l'approximation dipolaire}
Pour le dipôle doublet ${P,N}$, on a : $$\overrightarrow{E}=\frac{p}{4\pi\varepsilon_0 r^3}(2\cos(\theta)\overrightarrow{u_r}+\sin(\theta)\overrightarrow{u_\theta})$$
Plus généralement, on dispose de la formule intrinsèque suivante, avec les mêmes conventions qu'au paragraphe précédent : $$\overrightarrow{E}=\frac{1}{4\pi\varepsilon_0 r^3}\left(3(\overrightarrow{p}\cdot\overrightarrow{u})\overrightarrow{u}-\overrightarrow{p}\right)$$
\subsection{Carte de champ et équipotentielles}
Attention, il y a une zone proche du dipôle où l'on ne peut rien dire du fait de l'approximation dipolaire !
\section{Actions mécaniques d'un champ extérieur sur un dipôle rigide}
\subsection{Attention}
Attention, on change dans cette partie de point de vue.
\subsection{Cas d'un champ extérieur uniforme}
Pour créer un champ extérieur, rappelons que l'on peut par exemple utiliser un condensateur plan. \\ \\
Dans un champ extérieur uniforme :
\begin{itemize}
	\item $\overrightarrow{R}=\overrightarrow{0}$
	\item $\overrightarrow{\Gamma}=\overrightarrow{p}\wedge\overrightarrow{E_{ext}}$
\end{itemize}
On a donc un couple.
\subsection{Cas d'un champ non uniforme dans l'approximation dipolaire}
On suppose donc $a\ll d$ avec $d$ la distance caractéristique de variation de $\overrightarrow{E_{ext}}$. On a : $$\overrightarrow{R}=(\overrightarrow{p}\cdot\grad)\overrightarrow{E}=(p_x\partial_x+p_y\partial_y+p_z\partial_z)\overrightarrow{E}$$
et $$\overrightarrow{\Gamma}=\overrightarrow{p}\wedge\overrightarrow{E_{ext}}$$
\subsection{Énergie potentielle extérieure ; positions d'équilibre et stabilité}
Dans l'approximation dipolaire : $$E_p=-\overrightarrow{p}\cdot\overrightarrow{E}$$ On a donc affaire à un système conservatif à un degré de liberté : 
\begin{itemize}
	\item les actions de liaison ne travaillent pas ;
	\item l'action de $\overrightarrow{E}$ dérive de $E_p$ ;
	\item $W_{int}=0$ car on considère le dipôle rigide (\textit{ie} c'est un solide).
\end{itemize}
Les positions d'équilibre correspondent à $\displaystyle\frac{dE_p}{d\theta}=0$. Or, $E_p=-pE\cos(\theta)$ (puis visualiser le graphe de $E_p$). Alors, $\theta=0$ est une position d'équilibre stable, et $\theta=\pi$ est une position d'équilibre instable. Si de plus on connaît $J_\Delta$, on peut exprimer la pulsation des petites oscillations autour de $\theta=0$.

\chapter{Magnétostatique}
\localtableofcontents
Le champ magnétique est caractérisé par son action mécanique au travers de la force de Lorentz : $$\overrightarrow{F}=q\left(\overrightarrow{E}+\overrightarrow{v}\wedge\overrightarrow{B}\right)$$ Son unité est le Tesla (T). Le champ à la surface de la terre est de l'ordre de $10^{-5}$ T. Pour avoir un ordre d'idée, les champs de l'ordre de 1 T sont ceux que l'on retrouve dans un appareil à IRM, au CERN ou encore au LHC.
\section{Des courants aux champs magnétiques}
\subsection{Distribution de courants}
\paragraph{Différentes modélisations}
On dispose des différents types de modélisations : 
\begin{itemize}
	\item 1D / linéique : on s'intéresse à l'intensité I en Ampère (A)
	\item 2D / surfacique : le vecteur densité de courant électrique surfacique $\overrightarrow{j_s}$ en A.m$\inv$
	\item 3D / volumique : le vecteur densité de courant électrique volumique $\overrightarrow{j}$ en A.m$^{-2}$
\end{itemize}
\paragraph{Intensité à travers une surface orientée} Pour une surface $S$ orientée par un vecteur unitaire $\overrightarrow{n}$, alors l'intensité au travers de cette surface est définie par :
$$I=\iint_S\overrightarrow{j_{el}}\cdot dS\overrightarrow{n}$$
\paragraph{Modèle simpliste} Si tous les porteurs de charge ont la même charge et la même vitesse, on a : $$\overrightarrow{j}=n_0q\overrightarrow{v} \ \text{\textit{ie}}\ \ \overrightarrow{j}=\rho\overrightarrow{v}$$ avec $n_0$ la densité particulaire de charges en m$^{-3}$.  \\ \par \textbf{Attention}, si tous les porteurs de charge n'ont pas la même vitesse, on peut très bien avoir $\overrightarrow{j}\ne\overrightarrow{0}$ et pourtant avoir $\rho=0$. \\ \par Si on peut regrouper les porteurs de charges selon des paquets $i$ où tous les porteurs de charges d'un même paquet ont même charge et même vitesse, on a : $$\overrightarrow{j}=\sum_i n_iq_i\overrightarrow{v_i}$$
\paragraph{Passage d'une modélisation à une autre} On prend une modélisation où $$\overrightarrow{j}=\begin{cases}
	\displaystyle j_0\frac{r}{a} \ \ \text{si }r<a \\
	0 \ \ \text{sinon}
\end{cases}$$ On peut modéliser cette distribution "de loin" par une distribution linéique avec $$I=\int_{0}^{2\pi}\int_{0}^{a}\overrightarrow{j}\cdot\overrightarrow{dS}=\int_{0}^{2\pi}\int_{0}^{a}j_0\frac{r}{a}rdrd\theta=\frac{2\pi}{3}j_0a^2$$
\paragraph{Passage d'une modélisation à une autre (bis)}
Si un grand nombre de spires sont parcourues par un courant $I$, on peut modéliser le tout par une densité de courant électrique surfacique donnée par $$\overrightarrow{j_s}=nI\overrightarrow{u}$$
\subsection{Un mot sur les aimants}
cf. chapitre EM2' Dipôle magnétostatique
\paragraph{Hypothèse d'Ampère} Le champ magnétique des aimants serait dû aux "courants microscopiques".
\subsection{Propriétés de symétrie du champ magnétique}
\paragraph{Pseudo-vecteurs} Le champ magnétique $\overrightarrow{B}$ est un \textbf{pseudo-vecteur} (ou plus anciennement un \textbf{vecteur axial}) : lorsqu'on symétrise le problème par rapport à un axe, le champ $\overrightarrow{B}$ est transformé en son opposé.
\begin{center}
	\begin{tabular}{|c|c|}
		\hline
		\textbf{Vrais vecteurs}&\textbf{Pseudo-vecteurs} \\
		\hline
		Champ électrique $\overrightarrow{E}$ & Champ magnétique $\overrightarrow{B}$ \\
		\hline
		Vitesse $\overrightarrow{v}$ & Rotation $\overrightarrow{\Omega}$\\
		\hline
		Position, accélération, etc.& Moments cinétiques \\
		\hline
		Force $\overrightarrow{F}$ & Moments de forces \\
		\hline
		Moment dipolaire électrique&Moment dipolaire magnétique \\
		\hline 
		Vecteur de Poynting $\displaystyle\frac{\overrightarrow{E}\wedge\overrightarrow{B}}{\mu_0}$ & \\
		\hline
	\end{tabular}
\end{center}
\paragraph{Pseudo-scalaires} De même, on appelle pseudo-scalaire un scalaire qui est changé en son opposé par symétrisation du problème par rapport à un axe. Par exemple, $\overrightarrow{E}\cdot\overrightarrow{B}$ est un pseudo-scalaire.
\paragraph{Nouvelle condition d'homogénéité} Un vecteur non-nul ne peut pas être à la fois un vrai vecteur et un pseudo-vecteur. Cela permet de détecter rapidement de potentielles erreurs dans des résultats vectoriels.
\paragraph{Plans de symétrie et d'antisymétrie des courants} Si un plan est tel que deux courants $i_1$ et $i_2$ soient symétriques par rapport à ce plan, on dit que ce plan est un plan :
\begin{itemize}
	\item de \textbf{symétrie des courants} si $i_1=i_2$ ;
	\item d'\textbf{antisymétrie des courants} si $i_1=-i_2$
\end{itemize}
\paragraph{Propriétés des plans de symétrie des courants} Si $\Pi$ est un plan de symétrie des courants et $M'$ est le symétrique de $M$ par rapport à $\Pi$, alors $\overrightarrow{B}(M')$ est l'opposé du symétrique de $\overrightarrow{B}(M)$ par rapport à $\Pi$. En particulier, si $N\in\Pi$, alors $\overrightarrow{B}(N)$ est orthogonal à $\Pi$.
\paragraph{Propriétés des plans d'antisymétrie des courants} Si $\Pi\st$ est un plan d'antisymétrie des courants et $M'$ est le symétrique de $M$ par rapport à $\Pi$, alors $\overrightarrow{B}(M')$ est le symétrique de $\overrightarrow{B}(M)$ par rapport à $\Pi\st$. En particulier, si $N\in\Pi\st$, alors $\overrightarrow{B}(N)$ est inclus dans $\Pi\st$.
\paragraph{Remarque} Il faut noter l'inversion des propriétés par rapport au champ $\overrightarrow{E}$.
\paragraph{Invariance par translation} Si la distribution de courants est invariante par translation selon la direction $z$, alors on a : $$\overrightarrow{B}=\overrightarrow{B}(x,y,\xcancel{z})$$
\paragraph{Invariance par rotation} Si la distribution des courants est invariante par rotation autour de l'axe $(Oz)$, alors on a $$B_r=B_r(r,\xcancel{\theta},z)\ \ ; \ \ B_\theta=B_\theta(r,\xcancel{\theta},z) \ \ ; \ \ B_z=B_z(r,\xcancel{\theta},z)$$ \textbf{Attention}, en revanche, on a généralement $\overrightarrow{B}(r,\theta_1,z)\ne\overrightarrow{B}(r,\theta_2,z)$
\subsection{Loi de Biot et Savart (HP)}
\paragraph{Loi de Biot et Savart} Il s'agit de la relation : $$\overrightarrow{B}(M)=\frac{\mu_0}{4\pi}\oint i\overrightarrow{dl}\wedge\frac{\overrightarrow{PM}}{PM^3}$$ Noter l'analogie avec  $$\overrightarrow{E}(M)=\frac{1}{4\pi\varepsilon_0}\iiint_{P\in V} \rho\delta\tau\frac{\overrightarrow{PM}}{PM^3}$$ Grâce à la loi de Biot et Savart, on voit aisément que $\overrightarrow{B}$ est un pseudo-vecteur.
\paragraph{Exemple d'application : Calcul du champ sur l'axe d'une spire} On considère une spire de centre O et de rayon $a$ parcourue par une intensité $i$. Elle est orthogonale à l'axe $(Oz)$ et on souhaite déterminer le champ magnétique en $M\in(Oz)$. Tout plan contenant $(Oz)$ est un plan d'antisymétrie des courants, donc $\overrightarrow{B}(M)=B_z\overrightarrow{u_z}$. Or, la loi de Biot et Savart donne : $$\overrightarrow{u_z}\cdot\overrightarrow{B}(M)=\frac{\mu_0}{4\pi}\oint\frac{i}{PM^3}\left(\overrightarrow{dl}\wedge\overrightarrow{PM}\right)\cdot\overrightarrow{u_z}$$ Or, un découpage en $\overrightarrow{PM}=\overrightarrow{PO}+\overrightarrow{OM}$ donne rapidement $$\left(\overrightarrow{dl}\wedge\overrightarrow{PM}\right)\cdot\overrightarrow{u_z}=a\times dl$$ donc $$B_z=\frac{\mu_0 i}{2}\frac{a^2}{(a^2+z^2)^{3/2}}$$
\subsection{Cartes de champ}
\paragraph{Fil infini} Les lignes de champ sont des cercles concentriques autour de l'axe et orientés par le sens du courant.
\paragraph{Spire} Proche de la spire, les lignes de champ sont presque des cercles. Sinon, elles sont coincées entre la spire et l'axe de la spire, puis s'écartent de plus en plus lorsqu'elles s'éloignent de l'axe.
\section{Flux du champ magnétostatique}
\subsection{Formulations intégrales}
Le champ magnétique est à flux conservatif. Cela se traduit des trois façons suivantes :
\begin{itemize}
	\item si $S$ est une surface fermée, on a : $$\oiint_S\overrightarrow{B}\cdot\overrightarrow{dS}=0$$
	\item si $S_1$ est la face d'entrée d'un tube de champ orientée par le vecteur entrant dans le tube de champ et $S_2$ est la face de sortie de ce tube de champ orientée par le vecteur normal sortant du tube de champ, on a : $$\iint_{S_1}\overrightarrow{B}\cdot\overrightarrow{dS}=\iint_{S_2}\overrightarrow{B}\cdot\overrightarrow{dS}$$
	\item Si $\Gamma$ est un contour orienté fermé sur lequel s'appuient les surfaces $S_1$ et $S_2$, alors $$\iint_{S_1}\overrightarrow{B}\cdot\overrightarrow{dS}=\iint_{S_2}\overrightarrow{B}\cdot\overrightarrow{dS}$$
\end{itemize}
\paragraph{Remarque} La troisième formulation assure donc que la loi de Faraday a du sens puisqu'elle est vraie pour toute surface.
\paragraph{Remarque} $\overrightarrow{B}$ est \textbf{toujours} à flux conservatif, même en régime variable.
\subsection{Application}
On se donne une spire de rayon $a$ autour d'un axe de direction $z$ et un tube de champ crée par la spire. On note $r_0$ le rayon du tube de champ à l'intérieur de la spire. On suppose le rayon $r(z)$ du tube de champ négligeable devant le rayon $a$. Calculer $r(z)$. \\ \par 
$\overrightarrow{B}$ est à flux conservatif donc on sait que $\pi r_0^2B_z(z=0)=\pi r(z)^2B_z(z)$ grâce à notre approximation, car on peut considérer que le champ est égal à celui sur l'axe pour les surfaces sur lesquelles on calcule le flux. Cela se retraduit, en vertu de l'expression du champ crée par une spire sur son axe : $$r(z)=r_0\left(\frac{a^2+z^2}{a^2}\right)^{3/4}=r_0\left(1+\frac{z^2}{a^2}\right)^{3/4}$$
\subsection{Formulation locale}
\paragraph{Équation de Maxwell-Thomson} Anciennement appelée équation de Maxwell-Flux, elle est la quatrième des équations de Maxwell et stipule qu'on a en tout point, que ce soit en régime permanent ou en régime variable, l'égalité : $$\div(\overrightarrow{B})=0$$
\section{Circulation du champ magnétique ; théorème d'Ampère}
\textbf{Attention}, les résultats qui vont être énoncés dans cette section sont faux en régime variable.
\subsection{Formulation intégrale}
\paragraph{Théorème d'Ampère} En \textbf{régime permanent}, si $\Gamma$ est un contour fermé orienté, on a : $$\oint_\Gamma\overrightarrow{B}\cdot\overrightarrow{dl}=\mu_0I_\text{enlacée}$$ où $I_\text{enlacée}$ est le flux de $\overrightarrow{j_\text{el}}$ à travers une surface $S$ s'appuyant sur $\Gamma$.
\paragraph{Remarque} En statique, $\overrightarrow{j_\text{el}}$ est à flux conservatif (c'est la loi des nœuds) donc $I_\text{enlacée}$ ne dépend pas du choix de la surface $S$.
\paragraph{Remarque} L'orientation de $S$ est imposée par l'orientation de $\Gamma$.
\subsection{Formulation locale}
\paragraph{Équation de Maxwell-Ampère stationnaire} En régime stationnaire, on a la relation : $$\rot\left(\overrightarrow{B}\right)=\mu_0\overrightarrow{j}$$
\paragraph{Conséquence} Cette équation implique que $$\div\left(\overrightarrow{j}\right)=0$$ C'est la loi des nœuds.
\paragraph{Équivalence entre les formes intégrale et locale} Elle s'obtient aisément avec le théorème de Stokes.
\paragraph{Conservation de la charge} Soit $V$ un volume délimité par une surface $S$. On a le \textbf{bilan de charges} suivant : $$\frac{Q(t+dt)-Q(t)}{dt}=-\oiint_S\overrightarrow{j}\cdot\overrightarrow{dS}$$ et on a $\displaystyle\frac{Q(t+dt)-Q(t)}{dt}=0$ en régime permanent. Il faut noter le parallèle avec les transferts thermiques, mais remarquer qu'il n'y a jamais de terme source pour la charge.
\section{Applications}
\subsection{La recette}
\begin{itemize}
	\item Étude des symétries et des invariances
	\item Recherche d'un bon contour d'Ampère : il faut de préférence avoir un contour dont les composantes sont orthogonales au champ magnétique ou parallèles au champ magnétique
	\item Conclure avec le théorème d'Ampère
\end{itemize}
\subsection{Fil rectiligne "infini"}
On considère une fil rectiligne infini d'axe $(Oz)$ parcouru par une intensité $I$ constante dirigée selon $\overrightarrow{u_z}$. \\ \par Pour tout point M, le plan $(MOz)$ est un plan de symétrie des courants. Et puisque la distribution des courants est invariante par translation selon l'axe $(Oz)$ et est invariante par rotation autour de l'axe $(Oz)$, on a : $$\overrightarrow{B}=B_\theta(r,\xcancel{\theta},\xcancel{z})\overrightarrow{u_\theta}$$ En appliquant le théorème de Gauss sur un cercle de rayon $r$, puisque $I_\text{enlacée}=I$, on trouve : $$\overrightarrow{B}=\frac{\mu_0 I}{2\pi r}\overrightarrow{u_\theta}$$
\subsection{Cylindre "infini" parcouru par une courant inhomogène}
On considère un cylindre infini d'axe de révolution $(Oz)$ et de rayon $a$ tel que $$\overrightarrow{j}=\begin{cases}
	\displaystyle\frac{j_0 r^2}{a^2}\overrightarrow{u_z} \ \ \text{si }r<a \\[3mm]
	\overrightarrow{0} \ \ \text{sinon}
\end{cases}$$ Les symétries et les invariances sont les mêmes que dans l'exemple précédent. On choisit encore comme contour d'Ampère un cercle de rayon $r$ et d'axe $(Oz)$. On applique le théorème d'Ampère : $$\oint\overrightarrow{B}\cdot\overrightarrow{dl}=\mu_0I_\text{enlacée}$$ soit $2\pi rB_\theta(r)=\mu_0I_\text{enlacée}$. Or, si $r\leqslant a$, on a $\displaystyle I_\text{enlacée}=\frac{\pi j_0r^4}{2a^2}$ et si $r\geqslant a$, on a : $I_\text{enlacée}=\frac{\pi j_0 a^2}{2}$. On en déduit : $$\overrightarrow{B}=\begin{cases}
	\displaystyle \mu_0\frac{j_0}{4}\frac{r^3}{a^2}\overrightarrow{u_\theta} \ \ \text{si }r<a \\[3mm]
	\displaystyle \mu_0\frac{j_0}{4}\frac{a^2}{r}\overrightarrow{u_\theta} \ \ \text{si }r\geqslant a
\end{cases}$$
\subsection{Solénoïde "infini"}
On prend l'exemple d'un solénoïde infini avec $n$ spires par unité de longueur. En admettant que $\overrightarrow{B_\text{ext}}=\overrightarrow{0}$, on a par symétrie et invariance $\overrightarrow{B_\text{int}}=B_z(r)\overrightarrow{u_z}$ puis d’après le théorème d’Ampère : $$B_z=\mu_0nI$$
On peut montrer que $\overrightarrow{B_\text{ext}}=\overrightarrow{0}$ en trouvant le champ sur l’axe pour une spire puis en intégrant (principe de superposition).
\subsection{Nappe de courant ; relation de passage}
\paragraph{Nappe de courant} On considère une nappe de courant située dans le plan $(Oxy)$ avec $\varepsilon$ la distance entre deux fils et les intensités $I$ toutes constantes orientées selon $-\overrightarrow{u_y}$ et dans le plan $z=0$. \\ \par $(M,z,y)$ est un plan de symétrie et on a invariance par translation dans les directions $x$ et $y$ donc on a : $$\overrightarrow{B}=B_x(\xcancel{x},\xcancel{y},z)\overrightarrow{u_x}$$ De plus, $z=0$ est un plan de symétrie, donc $B_x(-z)=-B_x(z)$. On applique le théorème d'Ampère sur un rectangle de longueur arbitraire et on fait bien attention au signe de $I_\text{enlacée}$. On obtient alors : $$\overrightarrow{B}=\begin{cases}
	\displaystyle -\mu_0\frac{I}{2\varepsilon}\overrightarrow{u_x} \ \ \text{si }z>0 \\[3mm]
	\displaystyle \mu_0\frac{I}{2\varepsilon}\overrightarrow{u_x} \ \ \text{si }z<0
\end{cases}$$
\paragraph{Généralisation : relation de passage}
Pour une densité surfacique de courant $\overrightarrow{j_\text{s}}$, on a la \textbf{relation de passage} admise $$\overrightarrow{B_2}-\overrightarrow{B_1}=\mu_0\overrightarrow{j_\text{s}}\wedge\overrightarrow{n_{1\to 2}}$$ Il faut penser à l’homogénéité vecteurs/pseudo-vecteurs et au champ qui s’enroule pour s'en souvenir.
\section{Comparaison magnétostatique - électrostatique}
\begin{center}
	\begin{tabular}{|c|c|}
		\hline
		Électrostatique&Magnétostatique \\
		\hline
		Plan de symétrie : $\overrightarrow{E}$ dans ce plan.&Plan de symétrie : $\overrightarrow{B}$ orthogonal à ce plan. \\
		\hline 
		Plan d'antisymétrie : $\overrightarrow{E}$ orthogonal à ce plan.& Plan d'antisymétrie : $\overrightarrow{B}$ dans ce plan \\
		\hline 
		Loi de Coulomb&Loi de Biot et Savart \\
		\hline
		Théorème de Gauss&$\overrightarrow{B}$ est à flux conservatif\\
		\hline
		Maxwell-Gauss (valable en régime variable)&Maxwell-Thomson (valable en régime variable) \\
		\hline
		$\overrightarrow{E}$ est à circulation conservative (statique seulement)&Théorème d'Ampère (statique seulement)\\
		\hline
		Maxwell-Faraday \textbf{statique}& Maxwell-Ampère \textbf{statique} \\
		\hline
		$\overrightarrow{E}=-\grad(V)$&$\overrightarrow{B}=\rot(\overrightarrow{A})$ avec $\overrightarrow{A}$ le potentiel vecteur \\
		\hline
		$\overrightarrow{E_2}-\overrightarrow{E_1}=\frac{\sigma}{\varepsilon_0}\overrightarrow{n}_{1\to 2}$&$\overrightarrow{B_2}-\overrightarrow{B_1}=\mu_0 \overrightarrow{j}\wedge\overrightarrow{n}_{1\to 2}$ \\
		\hline
	\end{tabular}
\end{center}
\chapter{Dipôle magnétostatique}
\localtableofcontents
\section{Champ créé par un dipôle magnétostatique}
\subsection{Définition ; moment magnétique}
\paragraph{Définition} Pour un circuit plan orienté parcouru par un courant $I$ entourant une surface $S$ orientée par un vecteur normal $\overrightarrow{n}$, le \textbf{moment magnétique} est donné par $$\overrightarrow{M}=IS\overrightarrow{n}$$ Il s'exprime en A.m$^2$.
\paragraph{Remarque} Le sens de $I$ et de $\overrightarrow{n}$ sont liés par la règle de la main droite.
\paragraph{Cas de la bobine plate} Pour une bobine plate (les spires sont "confondues") de $N$ spires, on a : $$\overrightarrow{M}=NIS\overrightarrow{n}$$
\subsection{Moment magnétique à l'échelle atomique}
\paragraph{Moment magnétique orbital} Dans le modèle de Bohr, un électron sur une orbite circulaire de moment cinétique quantifié $\theta_0=n\hbar$ n'émet pas. Son moment magnétique vérifie alors : $$\overrightarrow{M}=\frac{-e}{2m_e}\overrightarrow{\theta_0}$$ $\displaystyle \frac{-e}{2m_e}$ est le \textbf{facteur gyromagnétique} de l'électron. Le moment magnétique est alors de l'ordre du \textbf{magnéton de Bohr} : $$\mu_b=\frac{e\hbar}{2m_e}\approx 9,274.10^{-24}\ \text{J.T}\inv$$
\paragraph{"Rappel"} On rappelle que $$L=\sqrt{l(l+1)}\hbar \ \ \text{et} \ \ L_z=n_l\hbar$$
\paragraph{Spin} Le \textbf{spin} est un autre moment défini par : $$\overrightarrow{M}=\frac{-eg}{2m_e}\overrightarrow{S}$$ où $g$ est la facteur de Landé de l'électron (vaut environ 2). On retrouve l'idée de configuration électronique (n, l, $n_e$ et $n_s$).
\paragraph{Ordres de grandeur} A l'échelle atomique, $M$ est de l'ordre du magnéton de Bohr. Pour un aimant de quelques grammes, $M$ est de l'ordre de 1 A.m$^2$.
\subsection{Champ créé par un dipôle dans l'approximation dipolaire (admis)}
\paragraph{Champ créé} On admet que, dans l'approximation dipolaire (le dipôle est "vu de loin"), le champ créé est le même que pour un dipôle électrostatique. On a donc : $$\overrightarrow{B}=\frac{\mu_0 M}{4\pi r^3}\bigg(2\cos(\theta)\overrightarrow{u_\text{r}}+\sin(\theta)\overrightarrow{u_\theta}\bigg)$$
\paragraph{Remarque} En fait, tout est relié par la relation $\mu_0\varepsilon_0c^2=1$.
\subsection{Carte de champ}
On a une ligne de champ alignée avec le moment magnétique. Les autres lignes de champ sortent du "haut" du dipôle, font comme une patate qui grandit et rentrent par le "bas" du dipôle. Il y a une zone proche du dipôle où on ne peut rien dire car on travaille dans l'approximation dipolaire. Noter la ressemblance avec la carte de champ du champ créé par une spire.
\subsection{Analogies et différences avec le dipôle électrostatique}
En dehors des charges et des courants : 
\begin{center}
	\begin{tabular}{cc}
		$\div\overrightarrow{E}=0$&$\div\overrightarrow{B}=0$\\
		$\rot\overrightarrow{E}=\overrightarrow{0}$&$\rot\overrightarrow{B}=\overrightarrow{0}$
	\end{tabular}
\end{center}
$\overrightarrow{M}$ est un pseudo-vecteur puisque symétriser le problème revient à prendre l'opposé de l'intensité pour qu'elle soit comptée comme strictement positive et donc à inverser le vecteur normal, et donc $\overrightarrow{M}$. \\ \par Tout comme le champ électrique, la composante sur $\overrightarrow{u_\varphi}$ de $\overrightarrow{B}$ est nulle car le plan $(M,\overrightarrow{u_r},\overrightarrow{u_\theta})$ est un plan d'antisymétrie des courants. On a aussi, comme pour le dipôle électrostatique, une invariance par rotation autour de l'axe $(Oz)$.
\section{Actions subies par un dipôle dans un champ magnétique extérieur}
\subsection{Attention}
Dans cette section, on considère un champ magnétique extérieur créé par d'autres courants.
\subsection{Résultante et moment}
\paragraph{Dipôle dans un champ extérieur uniforme} Si le champ $\overrightarrow{B_\text{ext}}$ est uniforme, la résultante est nulle, c'est la force de Laplace $$\overrightarrow{R}=\oint_\Gamma i\overrightarrow{dl}\wedge\overrightarrow{B_\text{ext}}$$ En revanche, le moment n'a aucune raison d'être nul et vaut $$\overrightarrow{\Gamma}=\overrightarrow{M}\wedge\overrightarrow{B_\text{ext}}$$ Dans ce cas, l'action mécanique est un couple de forces.
\paragraph{Remarque} Il y a des cas où la formule du moment est sans intérêt, par exemple celui où il y a une liaison pivot dans le circuit.
\paragraph{Dipôle dans un champ extérieur non uniforme}Si le champ $\overrightarrow{B_\text{ext}}$ est non uniforme mais qu'on se place dans l'approximation dipolaire, (le circuit est petit devant les variations de $\overrightarrow{B_\text{ext}}$), le moment vaut toujours \textbf{à l'ordre 1} : $$\overrightarrow{\Gamma}=\overrightarrow{M}\wedge\overrightarrow{B_\text{ext}}$$
\subsection{Énergie potentielle magnétique (admis)}
\paragraph{Énergie potentielle magnétique}L'énergie potentielle s'écrit alors $$E_p=-\overrightarrow{M}\cdot\overrightarrow{B_\text{ext}}$$ (penser à une boussole), et on a alors : $$\overrightarrow{R}=-\grad\left(\overrightarrow{M}\cdot\overrightarrow{B_\text{ext}}\right)$$
\paragraph{Remarque} $\overrightarrow{M}$ et $\overrightarrow{B_\text{ext}}$ sont des pseudo-vecteurs donc $E_p$ est un "vrai" scalaire.
\paragraph{Exemple de la boussole} Le moment magnétique pivote sur on axe grâce à une liaison pivot parfaite. L'angle entre $\overrightarrow{B_\text{ext}}$ et $\overrightarrow{M}$ est noté $\theta$. On a alors grâce à la formule : $$E_p=-MB_\text{ext}\cos(\theta)$$ $\theta=0$ est donc une position d'équilibre stable et $\theta=\pi$ est une position d'équilibre instable. Avec une information supplémentaire comme le moment d'inertie $J$ de la boussole autour de son axe de rotation, on peut déterminer la pulsation des petites oscillations autour d'une position d'équilibre stable.
\paragraph{Remarque} Dans l'exemple de la boussole, ce sont en réalité les frottements qui permettent à la boussole de s'aligner parfaitement avec le Nord.
\subsection{Généralisation : énergie d'interaction entre deux dipôles}
\paragraph{Énergie d'interaction} On admet que $$E_p=-\overrightarrow{M_1}\cdot\overrightarrow{B_2}(P_1)=-\overrightarrow{M_2}\cdot\overrightarrow{B_1}(P_2)$$
\paragraph{Remarque} Penser à utiliser les formules intrinsèques dans l'étude de l'interaction entre plusieurs dipôles, car les formules générales du cours avec les systèmes de coordonnées ne sont valables que pour le système de coordonnées du dipôle considéré : les systèmes de coordonnées de différents dipôles ne sont donc pas compatibles.
\section{Expérience de Stern et Gerlash}
\paragraph{Expérience de Stern et Gerlash} La projection de $\overrightarrow{S}$ est quantifiée : $$\overrightarrow{S}\cdot\overrightarrow{u_z}=\pm\frac{1}{2}\hbar$$ L'expérience consiste à faire sublimer des atomes d'argent, garder ceux ayant uniquement une vitesse selon $\overrightarrow{u_x}$ et les faire passer dans un champ magnétique inhomogène orienté selon $\overrightarrow{u_z}$. On s'attendrait alors à observer une distribution continue (déviation selon $\overrightarrow{M}\cdot\overrightarrow{u_z}$) mais on n'observe dans les faits que 2 raies.
\chapter{Équations de Maxwell}
\localtableofcontents
\section{Équations de Maxwell ; conséquences}
\subsection{Postulats de l'électromagnétisme}
\paragraph{Force de Lorentz} $$\overrightarrow{F}=q\bigg(\overrightarrow{E}+\overrightarrow{v}\wedge\overrightarrow{B}\bigg)$$
\paragraph{Équation de Maxwell-Thomson ou équation de Maxwell-flux} $$\div\left(\overrightarrow{B}\right)=0$$
\paragraph{Équation de Maxwell-Gauss} $$\div\left(\overrightarrow{E}\right)=\frac{\rho}{\varepsilon_0}$$
\paragraph{Équation de Maxwell-Faraday} $$\rot\left(\overrightarrow{E}\right)=-\frac{\partial\overrightarrow{B}}{\partial t}$$
\paragraph{Équation de Maxwell-Ampère} $$\rot\left(\overrightarrow{B}\right)=\mu_0\overrightarrow{j}+\mu_0\varepsilon_0\frac{\partial\overrightarrow{E}}{\partial t}$$
\subsection{Cas statique}
Si $\overrightarrow{E}$ et $\overrightarrow{B}$ sont indépendants du temps, les équations sont découplées et on retrouve les équations données dans les chapitres d'électrostatique et de magnétostatique :
\begin{center}
	\begin{tabular}{|c|c|}
		\hline
		Électrostatique&Magnétostatique \\[3mm]
		\hline
		$\div\overrightarrow{E}=\frac{\rho}{\varepsilon_0}$&$\div\overrightarrow{B}=0$ \\[3mm]
		\hline
		$\rot\overrightarrow{E}=\overrightarrow{0}$&$\rot\overrightarrow{B}=\mu_0\overrightarrow{j}$ \\[3mm]
		\hline
	\end{tabular}
\end{center}
\subsection{Transformations non relativistes des champs}
Pour plus de précisions, \textit{cf}. la page Wikipédia "Transformation de Lorentz du champ magnétique".
\paragraph{Attention} On s'apprête ici à effectuer un raisonnement fallacieux qui donne néanmoins une formule juste.
\paragraph{Situation} On se place dans le cas d'un référentiel $R'$ en translation rectiligne uniforme par rapport à un référentiel $R$ avec une vitesse d'entraînement $\overrightarrow{v_e}=\overrightarrow{v}(O'/R)$. On pose $\displaystyle\beta=\frac{v_e}{c}$ et $\displaystyle\gamma=\frac{1}{\sqrt{1-\beta^2}}$.
\paragraph{Résultat} On cherche à obtenir la formule de Lorentz avec $v_e\ll c$ et $cB\st\gg E\st$, c'est-à-dire la limite non relativiste "magnétique" : $$\overrightarrow{B'}=\overrightarrow{B} \ ; \ E_x'=E_x \ ; \ E_y'=E_y-\beta cB_z \ ; \ E_z'=E_z+\beta cB_y$$
\paragraph{"Démonstration"} On utilise l'invariance de la force de Lorentz par changement de référentiel et la loi de composition des vitesses classique. On a $\overrightarrow{F'}=\overrightarrow{F}$, soit $$q\left(\overrightarrow{E}+\overrightarrow{v}\wedge\overrightarrow{B}\right)=q\left(\overrightarrow{E'}+\overrightarrow{v'}\wedge\overrightarrow{B'}\right)$$ Or, d'après la loi de composition des vitesses, $\overrightarrow{v}=\overrightarrow{v'}+\overrightarrow{v_e}$ donc on a pour tout $\overrightarrow{v'}$ : $$\overrightarrow{E}+\left(\overrightarrow{v'}+\overrightarrow{v_e}\right)\wedge\overrightarrow{B}=\overrightarrow{E'}+\overrightarrow{v'}\wedge\overrightarrow{B'}$$ En particulier, pour $\overrightarrow{v'}=\overrightarrow{0}$, on obtient $\overrightarrow{E}+\overrightarrow{v_e}\wedge\overrightarrow{B}=\overrightarrow{E'}$ d'où on déduit que pour tout $\overrightarrow{v'}$, $\overrightarrow{v'}\wedge\overrightarrow{B}=\overrightarrow{v'}\wedge\overrightarrow{B'}$ et donc que $\overrightarrow{B}=\overrightarrow{B'}$.
\paragraph{Limite non relativiste "magnétique"} Elle est très utile pour l'induction : $$\overrightarrow{E'}=\overrightarrow{E}+\overrightarrow{v_e}\wedge\overrightarrow{B} \ \ \text{et} \ \ \overrightarrow{B'}=\overrightarrow{B}$$
\paragraph{Exemple pour souligner l'importance de l'hypothèse "dominante magnétique"}
On s'intéresse à un cylindre infini d'axe (Oz) et de rayon $a$ dans un référentiel $R$. La densité volumique de charge est de $\rho_0$ dans le cylindre et $0$ en-dehors. On regarde ensuite ce cylindre dans un référentiel en translation à $\overrightarrow{v_e}$ avec $\overrightarrow{v_e}=-v_e\overrightarrow{u_z}$, les charges ont donc un mouvement à $-\overrightarrow{v_e}$. On cherche à exprimer $\overrightarrow{E}$, $\overrightarrow{B}$, $\overrightarrow{E'}$ et $\overrightarrow{B'}$ pour $r<a$. \\ \par Les symétries et les invariances dans le cas de $R$ donnent rapidement $\overrightarrow{E}=E_r(r,\xcancel{\theta},\xcancel{z})\overrightarrow{u_r}$ puis le théorème de Gauss fournit $\displaystyle\overrightarrow{E}=\frac{\rho r}{2\pi\varepsilon_0}\overrightarrow{u_r}$ si $r<a$. Et on a $\overrightarrow{B}=\overrightarrow{0}$ car il n'y a pas de courant. \\ \par Dans $R'$, $\overrightarrow{E}$ ne change pas. On a $\overrightarrow{j}=-nq\overrightarrow{v_e}=-\rho_0\overrightarrow{v_e}$ car il n'y a qu'un seul type de porteurs de charges. Les symétries et les invariances donnent rapidement $\overrightarrow{B'}=B_\theta(r,\xcancel{\theta},\xcancel{z})\overrightarrow{u_\theta}$ L'application du théorème d'Ampère sur un cercle d'axe $(Oz)$ et de rayon $r<a$ fournit rapidement $\overrightarrow{B'}=-\mu_0\rho_0 v_e\frac{r}{2}\overrightarrow{u_\theta}$
\paragraph{Limite non relativiste "électrique"} De façon symétrique, il existe une limite non relativiste à dominante "électrique" qui stipule que $$\overrightarrow{E'}=\overrightarrow{E} \ \ \text{et} \ \ \overrightarrow{B'}\ne\overrightarrow{B}$$ 
\subsection{Conservation de la charge}
\paragraph{Conservation de la charge} On l'obtient en 1D en faisant un bilan $$Q(t+dt)-Q(t)=I(x)-I(x+dx)$$ ou en utilisant $$\frac{Q(t+dt)-Q(t)}{dt}=\Phi$$ et le théorème de Green-Ostrogradski. Elle s'écrit $$\frac{\partial\rho}{\partial t}+\div\left(\overrightarrow{j}\right)=0$$ Les équations de Maxwell permettent de la retrouver : il suffit d'appliquer l'opérateur divergence à l'équation de Maxwell-Ampère.

\subsection{Formes intégrales des équations de Maxwell}
\paragraph{Maxwell-Thomson intégrale} $\overrightarrow{B}$ est à flux conservatif (théorème de Green-Ostrogradski)
\paragraph{Maxwell-Gauss intégrale} C'est le théorème de Gauss : $$\oiint_S\overrightarrow{E}\cdot dS\overrightarrow{n}=\frac{Q_\text{int}}{\varepsilon_0}$$ (théorème de Green-Ostrogradski)
\paragraph{Maxwell-Faraday intégrale} Avec le théorème de Stokes, on montre que pour un contour fermé $\Gamma$ : $$\oint_\Gamma\overrightarrow{E}\cdot\overrightarrow{dl}=\frac{d\Phi_{\overrightarrow{B}}}{dt}=-\frac{\partial}{\partial t}\iint_S\overrightarrow{B}\cdot dS\overrightarrow{n}$$
\paragraph{Loi de Faraday (induction de Neumann)} Si le circuit / contour $\Gamma$ est fixe, on a $$e=-\frac{\partial \Phi_{\overrightarrow{B}}}{\partial t}$$ où $\Phi_{\overrightarrow{B}}$ est le flux de $\overrightarrow{B}$ à travers une surface de bord $\gamma$ et $e$ est la force électromotrice, \textit{ie} l'énergie fournie par unité de charge. \textbf{Attention}, en réalité, le "générateur" est réparti sur l'ensemble du circuit avec la loi de Faraday.
\paragraph{Maxwell-Ampère intégrale / Théorème d'Ampère généralisé} Avec le théorème de Stokes, on obtient : $$\oint_\Gamma\overrightarrow{B}\cdot\overrightarrow{dl}=\mu_0I_\text{enlacée}+\mu_0\varepsilon_0\frac{\partial \Phi_{\overrightarrow{E}}}{\partial t}$$ Attention, $I_\text{enlacée}$ et $\Phi_{\overrightarrow{E}}$ dépendent de $S$, mais leur somme ne dépend que de $\Gamma$.
\paragraph{Exemple du condensateur} On considère un condensateur plan de surface $S$ avec une intensité $i$ et une normale $\overrightarrow{u_z}$ en régime lentement variable (ARQS) et on nélgie les effets de bord, si bien que l'on a $$\overrightarrow{E}=\begin{cases}
	\overrightarrow{0} \ \ \text{hors du condensateur}\\[1mm]
	\displaystyle\frac{Q(t)}{\varepsilon_0 S}\overrightarrow{u_z} \ \ \text{dans le condensateur}
\end{cases}$$
On considère un contour circulaire $\Gamma$ autour de l'axe $z$. On choisit comme surfaces $S_1$ la surface du disque et $S_2$ une surface qui enveloppe une armature du condensateur. On a donc avec le théorème d'Ampère généralisé : 
\begin{itemize}
	\item pour $S_1$ : $$\oint_\Gamma\overrightarrow{B}\cdot\overrightarrow{dl}=\mu_0 i +0$$
	\item pour $S_2$ : $$\oint_\Gamma\overrightarrow{B}\cdot\overrightarrow{dl}=0+\mu_0\varepsilon_0\frac{d}{dt}\left(\frac{Q(t)S}{\varepsilon_0 S}\right)=\mu_0\frac{dQ}{dt}$$
\end{itemize} Si bien que l'on retrouve $i=\displaystyle\frac{dQ}{dt}$
\subsection{Potentiels associés au champ}
\paragraph{Potentiels (HP)} $\div\left(\overrightarrow{B}\right)=0$ donc il existe $\overrightarrow{A}$ tel que $\overrightarrow{B}=\rot\left(\overrightarrow{A}\right)$. On en déduit d'après l'équation de Maxwell-Faraday que $$\rot\left(E+\frac{\partial\overrightarrow{A}}{\partial t}\right)=\overrightarrow{0}$$ On en déduit que $$\overrightarrow{E}=-\grad(V)-\frac{\partial \overrightarrow{A}}{\partial t}$$
\subsection{Généralisation des relations de passage}
\paragraph{Relations de passage} Ce sont les mêmes qu'en électrostatique et en magnétostatique : $$\overrightarrow{E_2}-\overrightarrow{E_1}=\frac{\sigma}{\varepsilon_0}\overrightarrow{n_{1\to 2}}$$
$$\overrightarrow{B_2}-\overrightarrow{B_1}=\mu_0\overrightarrow{j_\text{s}}\wedge\overrightarrow{n_{1\to 2}}$$
\section{Propagation du champ électromagnétique}
\subsection{Équation de propagation dans le vide}
\paragraph{Onde dans le vide et équation de d'Alembert} Pour trouver l'équation de d'Alembert, on part de la relation : $$\rot\left(\rot\right)=\grad\left(\div\right)-\overrightarrow{\Delta}$$ et des équations de Maxwell dans le vide ($\rho=0$ et $\overrightarrow{j}=\overrightarrow{0}$) pour arriver à \textbf{l'équation de d'Alembert} : $$\overrightarrow{\Delta}\overrightarrow{E}=\mu_0\varepsilon_0\frac{\partial^2\overrightarrow{E}}{\partial t^2}$$ et de même pour $\overrightarrow{B}$. Les ondes progressives se déplaçant à une célérité $c$ telle que $c^2=\displaystyle\frac{1}{\mu_0\varepsilon_0}$ sont solution de l'équation de d'Alembert.
\paragraph{Cas 1D} On se place dans le cas où $\overrightarrow{E}=E_x(x,t)\overrightarrow{u_x}$. Alors, l'équation de d'Alembert devient $$\frac{\partial^2 E_x}{\partial x^2}=\mu_0\varepsilon_0\frac{\partial^2 E_x}{\partial t^2}$$ Toute fonction de la forme $(x,t)\mapsto f(x-ct)$ avec $\mu_0\varepsilon_0c^2=1$ est alors solution. 
\paragraph{Remarque} L'équation de d'Alembert (qu'on appelle aussi parfois équation d'onde) est invariante par renversement du temps, contrairement à l'équation de la diffusion thermique.
\subsection{Propagation par couplage des champs}
On prend l'exemple d'une suite de masses repérées par leurs abscisses $x_i$ et toutes de masse $m$. Elles sont reliées à la précédente et à la suivante par des ressorts de mêmes longueur à vide $l$ et de même raideur $k$. On note $l_i=x_{i+1}-x_i$ et $\displaystyle v_i=\frac{dx_i}{dt}$. On a alors $\displaystyle\frac{dl_i}{dt}=v_{i+1}-v_i$, cette équation couplant une dérivée temporelle (membre de gauche) et une dérivée spatiale (membre de droite). \\ \par On applique le principe fondamental à la masse i, ce qui donne : $$m\frac{dv_i}{dt}=k(l_i-l)-k(l_{i-1}-l)$$ et on a donc un couplage entre une dérivée temporelle et une dérivée spatiale : $$\frac{dv_i}{dt}=\frac{k}{m}\left(l_i-l_{i-1}\right)$$

\subsection{La lumière, une onde électromagnétique}

\paragraph{Hertz} La première vérification expérimentale de l'existence d'onde électromagnétiques a été effectuée par Hertz, à qui on doit les ondes hertziennes et l'unité Hz.

\paragraph{Postulat de Maxwell} On trouve que $c=\displaystyle\frac{1}{\sqrt{\mu_0\varepsilon_0}}$ pour la lumière, donc Maxwell postule que la lumière est une onde électromagnétique.

\section{Énergie du champ électromagnétique}

\subsection{Interaction volumique entre le champ et les porteurs ; exemples}

\paragraph{Force volumique} La force de Lorentz sur un petit volume s'écrit : $$\delta\overrightarrow{F}=\sum_{q_i\in\delta\tau}q_i\bigg(\overrightarrow{E}+\overrightarrow{v_i}\wedge\overrightarrow{B}\bigg)=\rho\delta\tau\overrightarrow{E}+\overrightarrow{j}\delta\tau\wedge\overrightarrow{B}$$ La densité volumique de la force de Lorentz s'écrit donc $$\overrightarrow{f_v}=\rho\overrightarrow{E}+\overrightarrow{j}\wedge\overrightarrow{B}$$

\paragraph{Puissance volumique cédée aux porteurs} On la calcule à partir du travail de la force de Lorentz puisque la partie magnétique de la force de Lorentz ne travaille pas : $$P_v=\overrightarrow{j}\cdot\overrightarrow{E}$$

\paragraph{Exemple}On considère un fil de rayon $a$ orienté selon un axe $\overrightarrow{u_z}$. Il est constitué d'un plasma neutre avec des ions fixes et des électrons homocinétiques se déplaçant à une vitesse $+v\overrightarrow{u_z}$ avec une densité $n$. Le fil est parcouru par une intensité $I$ compté positivement (donc orientée selon $-\overrightarrow{u_z}$). On suppose que la pression à l'extérieur du fil est négligeable. Quelle est la pression en $r=0$ ? \\ \par Stratégie : on détermine $\overrightarrow{B}$ avec les symétries et le théorème d'Ampère, puis on calcule $\overrightarrow{j}\wedge\overrightarrow{B}$ et on utilise l'équation de la statique des fluides. \\ \par On calcule $\overrightarrow{B}$ pour $r<a$. Le plan $(M,\overrightarrow{u_r},\overrightarrow{u_z})$ est un plan de symétrie des courants et on a des invariances pour rotation autour de l'axe $\overrightarrow{u_z}$ et par translation selon $\overrightarrow{u_z}$, si bien que $$\overrightarrow{B}=B_\theta(r,\xcancel{\theta},\xcancel{z})\overrightarrow{u_\theta}$$ Le théorème d'Ampère sur un cercle d'axe (Oz) et de rayon $r$ fournit alors : $$\overrightarrow{B}=-\frac{\mu_0 I r}{2\pi a^2}\overrightarrow{u_\theta}$$ Or, puisqu'on a un seul type de porteurs de charges mobiles, on a $$\overrightarrow{j}=-nev\overrightarrow{u_z}=-\frac{I}{\pi a^2}\overrightarrow{u_z}$$ La densité volumique de force vaut donc $$\overrightarrow{f_v}=\overrightarrow{j}\wedge\overrightarrow{B}=-\frac{\mu_0 I^2 r}{2\pi^2 a^4}\overrightarrow{u_r}$$ Or, l'équation de la statique des fluides s'écrit $\overrightarrow{0}=\overrightarrow{f_v}\delta\tau-\grad(P)\delta\tau$ donc après projection sur $\overrightarrow{u_r}$ : $$0=-\frac{\mu_0 I^2}{2\pi^2 a^4}r-\frac{dP}{dr}$$ d'où après intégration entre 0 et a avec l'hypothèse de négligeabilité de $P(a)$ : $$P(0)=\frac{\mu_0 I^2}{4\pi^2 a^2}$$ Cette expression est bien homogène (utiliser $\overrightarrow{B}$ et $\overrightarrow{F}_\text{Laplace}$)
\subsection{A la recherche d'un bilan sur l'énergie des champs}
\paragraph{Bilan de puissance} On introduit l'énergie magnétique $U$ contenue dans un volume $V$ et $\omega$ l'énergie volumique correspondante : $$U=\iiint_V\omega\delta\tau$$ On introduit le vecteur de Poynting $\overrightarrow{\pi}$ tel que le flux de $\overrightarrow{\pi}$ à travers $S$ soit la puissance transverse. Le bilan s'écrit alors : $$\frac{\partial U}{\partial t}=-\oiint_S\overrightarrow{\pi}\cdot dS\overrightarrow{n}-\iiint_V\overrightarrow{j}\cdot\overrightarrow{E}$$ le premier terme étant la puissance transmise à l'extérieur et le second terme étant la puissance transmise à l'intérieur. Grâce au théorème de Green-Ostrogradski, on en déduit sa version locale : $$\frac{\partial\omega}{\partial t}+\div\left(\overrightarrow{\pi}\right)=-\overrightarrow{j}\cdot\overrightarrow{E}$$ Il faut noter l'analogie avec la conservation de la charge et les bilans en transferts thermiques. On peut retrouver cette expression à partir des équations de Maxwell en effectuant $\overrightarrow{E}$ scalaire Maxwell-Ampère moins $\overrightarrow{B}$ scalaire Maxwell-Faraday et en utilisant l'identité $$\div\left(\overrightarrow{E}\wedge\overrightarrow{B}\right)=\overrightarrow{B}\cdot\rot\left(\overrightarrow{E}\right)-\overrightarrow{E}\cdot\rot\left(\overrightarrow{B}\right)$$ Cela donne les expressions : $$\overrightarrow{\pi}=\frac{\overrightarrow{E}\wedge\overrightarrow{B}}{\mu_0} \ \ \text{et} \ \ \omega=\frac{\varepsilon_0E^2}{2}+\frac{B^2}{2\mu_0}$$
On retrouve ainsi les expressions des énergies stockées dans un condensateur et une bobine.
\subsection{Identité de Poynting}
On a les différentes formules : $$\frac{\partial\omega}{\partial t}+\div\left(\overrightarrow{\pi}\right)=-\overrightarrow{j}\cdot\overrightarrow{E}$$ $$\omega=\frac{\varepsilon_0E^2}{2}+\frac{B^2}{2\mu_0}$$ $$\overrightarrow{\pi}=\frac{\overrightarrow{E}\wedge\overrightarrow{B}}{\mu_0}$$
\subsection{Deux exemples}
\paragraph{Exemple 1 : Condensateur dans l'ARQS électrique}
Dans l'ARQS électrique, le champ dans la condensateur de section $S$, de charge $Q(t)$ et d'épaisseur $e$ vaut $$\overrightarrow{E_1}=\frac{Q(t)}{\varepsilon_0 S}\overrightarrow{u_z}$$ Ce champ électrique est la source d'un champ magnétique $\overrightarrow{B_2}$. D'après l'équation de Maxwell-Ampère, $\displaystyle \rot\overrightarrow{B_2}=\mu_0\overrightarrow{j}+\mu_0\varepsilon_0\frac{\partial \overrightarrow{E_1}}{\partial t}$ avec $\overrightarrow{j}=\overrightarrow{0}$. Or, le plan $(M,\overrightarrow{u_r},\overrightarrow{u_\theta})$ est un plan de symétrie de $\overrightarrow{E_1}$ donc $\overrightarrow{B_2}=B_\theta(r,\theta,z)\overrightarrow{u\theta}$. L' théorème d'Ampère généralisé sur un cercle d'axe $(Oz)$ de rayon $r$ fournit : $$\oint\overrightarrow{B_2}\cdot\overrightarrow{dl}=0+\mu_0\varepsilon_0\frac{d\Phi_{\overrightarrow{E}}}{dt}$$ soit $$2\pi rB_\theta=\varepsilon_0\mu_0\pi r^2\frac{1}{\varepsilon_0 S}\frac{dQ}{dt}$$ donc $$\overrightarrow{B_2}=\frac{\mu_0 r}{2S}\frac{dQ}{dt}\overrightarrow{u_\theta}$$ A son tour, $\overrightarrow{B_2}$ crée $\overrightarrow{E_3}$, qui crée $\overrightarrow{B_4}$, etc.
\paragraph{Exemple 2 : Solénoïde dans l'ARQS magnétique}
Dans l'ARQS magnétique, $\overrightarrow{B_0}=\mu_0 ni(t)\overrightarrow{u_z}$. $\overrightarrow{B_0}(t)$ est la source de $\overrightarrow{E_1}(t)$. La plan $(M,\overrightarrow{u_r},\overrightarrow{u_z})$ est un plan d'antisymétrie des courants (causes indirectes de $\overrightarrow{E_1}$) et par invariances par rotation et translation, on obtient : $$\overrightarrow{E_1}=E_\theta(r,\xcancel{\theta},\xcancel{z})\overrightarrow{u_\theta}$$
Remarquons que les lignes de champ de $\overrightarrow{E_1}$ sont des cercles et sont donc fermées, \textbf{ce qui ne pose pas de problème car on est en régime variable}. On applique la loi de Faraday avec comme contour $\Gamma$ un cercle d'axe $(Oz)$ et de rayon $r$ et comme surface associée $S$ la surface du disque associé. On a alors : $$\oint_\Gamma\overrightarrow{E}\cdot\overrightarrow{dl}=-\frac{d\Phi_{\overrightarrow{B}}}{dt}$$ soit $\displaystyle 2\pi rE_1=-\mu_0 n\pi r^2\frac{di}{dt}$ donc on obtient, dans le solénoïde : $$\overrightarrow{E_1}=-\frac{\mu_0 nr}{2}\frac{di}{dt}\overrightarrow{u_\theta}$$ Mais $\overrightarrow{E_1}$ est la cause de $\overrightarrow{B_2}$. Puisque $(M,\overrightarrow{u_r},\overrightarrow{u_\theta})$ est un plan de symétrie et par invariances par rotation et translation, on a : $$\overrightarrow{B_2}=B_z(r,\xcancel{\theta},\xcancel{z})\overrightarrow{u_z}$$ On admet que $B_z(r=0)=0$. En appliquant le théorème d'Ampère sur un cadre de hauteur $H$ et de longueur $r$ dont l'un des côtés est selon $(Oz)$, on obtient : $$0+0+0-B_2H=-\mu_0\frac{n}{c^2}\frac{d^2i}{d^2t}H\frac{r^2}{4}$$ et par conséquent : $$\overrightarrow{B_2}=\frac{\mu_0 n}{c^2}\frac{d^2i}{d^2t}\frac{r^2}{4}\overrightarrow{u_z}$$ Si le courant est sinusoïdal de pulsation $\omega$, on a : $$\overrightarrow{B_2}=-\frac{r^2\omega^2}{4c^2}\overrightarrow{B_0}$$ En notant $a$ le rayon du solénoïde, la condition de validité de l'ARQS magnétique s'écrit alors : $$a\omega \ll c$$
\paragraph{Approximation des régimes quasi-stationnaires (ARQS)} La plus fréquente est la magnétique $E\ll cB$ : on remplace $I(t)$ par I, la propagation est instantanée et on néglige $\rho$ ainsi que le courant de déplacement $\displaystyle\frac{\partial \overrightarrow{E}}{\partial t}$. On considère alors que $\overrightarrow{B}$ est la seule source de $\overrightarrow{E}$. Parfois, on considère l'ARQS électrique $E\gg cB$ : on néglige alors $\displaystyle\frac{\partial\overrightarrow{B}}{\partial t}$ et $\overrightarrow{j}$.
\chapter{Conducteur soumis à un champ électromagnétique}
\localtableofcontents
\section{Conduction électrique dans les métaux}
\subsection{Modèle de Drude en régime stationnaire}
\paragraph{Modèle de Drude} Pour un métal, les porteurs sont les électrons de valence. 
\begin{itemize}
	\item Un seul électron par atome métaux alcalins
	\item Les électrons se déplacent librement, le noyau et les électrons de cœur (cations) sont fixes
\end{itemize}
On en déduit la densité moyenne d'électrons $n_e=\displaystyle\frac{\rho N_A}{M}=10^{29}$ m$^{-3}$. Le mouvement des électrons est décrit par deux modélisations : 
\begin{itemize}
	\item Lorentz plus frottements : $\displaystyle\frac{m}{\tau}\overrightarrow{v}$
	\item les chocs se produisent tous les $\tau$ et renvoient l'atome avec une vitesse nulle
\end{itemize}
Dans les deux cas, on a : $$\langle v\rangle=\frac{\tau q}{m}\overrightarrow{E}$$
\paragraph{Ordres de grandeurs} Pour le cuivre, $\tau\approx 10^{-14}$ s.
\subsection{Loi d'Ohm statique : formulations locale et intégrale}
\paragraph{Loi d'Ohm} Si $T\gg\tau$, on a $\overrightarrow{j}=\sigma\overrightarrow{E}$ avec $\sigma>0$ la conductance. Dans le modèle de Drude, on a : $$\sigma=\frac{n_0q^2\tau}{m}$$ Connaissant $\sigma$, on peut estimer $\tau$, ce qui est valable si $f\ll 10^{14}$ Hz (donc valable pour les ondes radio mais pas pour les ondes visibles). La puissance dissipée par effet Joule s'écrit alors : $$P_v=\overrightarrow{j}\cdot\overrightarrow{E}=\sigma E^2=\frac{j^2}{\sigma}$$
\paragraph{Loi d'Ohm dans un conducteur} On justifie l'ARQS magnétique en établissant l'équation différentielle sur $\rho$ à l'aide de la loi d'Ohm $\overrightarrow{j}=\sigma\overrightarrow{E}$ et des équations de Maxwell. On trouve un temps caractéristique $\displaystyle\tau=\frac{\varepsilon_0}{\sigma}\approx 10^{-18}$s : l'électroneutralité du conducteur est donc très vite atteinte. On montre également que $$\varepsilon\omega E=\varepsilon_0\frac{\partial \overrightarrow{E}}{\partial t}\ll\overrightarrow{j}=\sigma\overrightarrow{E}$$
\paragraph{Loi d'Ohm intégrale} Comme pour les résistances thermiques et les capacités, on calcule la résistance d'un conducteur de conductivité $\gamma$, de longueur $L$ et de section $S$ par : $$R=\frac{L}{\gamma S}$$
\paragraph{Exemple biscornu} \textbf{A COMPLÉTER}
\paragraph{Remarque : solution particulière de l'équation de Laplace en cylindriques} On vient d'observer que $V=\displaystyle\frac{\theta}{r}$ est une solution en coordonnées cylindriques de l'équation $\Delta V=0$.
\paragraph{Remarque : densité de charges en statique dans un conducteur} En statique, dans un conducteur ($\gamma$ uniforme), on a $\rho=0$. En effet, $\rho=\varepsilon_0 \div\overrightarrow{E}$ d'après l'équation de Maxwell-Gauss, mais $\overrightarrow{E}=\displaystyle\frac{\overrightarrow{j}}{\gamma}$ d'après la loi d'Ohm donc $\rho=\displaystyle\frac{\varepsilon_0}{\gamma}\div\overrightarrow{j}$. Or, $\div\overrightarrow{j}=0$ d'après l'équation de Maxwell-Ampère en statique, donc $\rho=0$.
\subsection{Puissance Joule}
\paragraph{Puissance dissipée par unité de volume}
La puissance dissipée par effet Joule par unité de volume vaut $$p_v=\overrightarrow{j}\cdot\overrightarrow{E}=\gamma E^2=\frac{j^2}{\gamma}$$
\paragraph{Puissance dissipée} La puissance dissipée par effet Joule dans un conducteur métallique vaut $$\mathcal{P}=Ri^2=\frac{u^2}{R}$$
\paragraph{Preuve dans un cas simple} On se place dans le cas d'une résistance cylindrique de longueur $L$, de section $S$, de conductivité $\gamma$ parcourue par une intensité $I$ et avec une tension $u$. On a $E=\displaystyle\frac{\lvert u\rvert}{L}$ et alors : $$\mathcal{P}=\iiint_V\overrightarrow{j}\cdot\overrightarrow{E}\delta\tau=SL\times\gamma\frac{u^2}{L^2}=\frac{u^2}{R}$$
\subsection{Effet Hall} On applique le principe fondamental à un électron : $$m\frac{d\overrightarrow{v}}{dt}=-e\overrightarrow{E_0}-e\overrightarrow{E_\text{Hall}}-\frac{m}{\tau}\overrightarrow{v}-e\overrightarrow{v}\wedge\overrightarrow{B}$$ Or, en régime permanent, $m\displaystyle\frac{d\overrightarrow{v}}{dt}=\overrightarrow{0}$ donc en projection sur $\overrightarrow{u_y}$ en régime permanent, on obtient : $$\overrightarrow{E_\text{Hall}}\cdot\overrightarrow{u_y}=+v_xB_z$$

\subsection{Généralisation aux régimes variables ; domaine de validité de la loi d'Ohm}

\paragraph{Modèle de Drude en régime sinusoïdal forcé} En régime sinusoïdal avec $e^{+i\omega t}$, cela devient $$m\frac{d\overrightarrow{v}}{dt}=-e\overrightarrow{E}-\cancel{e\overrightarrow{v}\wedge\overrightarrow{B}}-\frac{m}{\tau}\overrightarrow{v}$$ On néglige le terme en $e\overrightarrow{v}\wedge\overrightarrow{B}$ car il est d'ordre $\displaystyle\frac{E}{c}$.

\paragraph{Généralisation de la loi d'Ohm} On passe en notation complexe et le modèle de Drude en régime sinusoïdal forcé devient $im\omega\underline{\overrightarrow{v}}=-\underline{\overrightarrow{E}}-\displaystyle\frac{m}{\tau}\underline{\overrightarrow{v}}$. On en déduit : $$\underline{\overrightarrow{v}}=\frac{-e}{\frac{m}{\tau}+im\omega}\underline{\overrightarrow{E}}$$ Or, $\underline{\overrightarrow{j}}=\overrightarrow{0}-ne\underline{\overrightarrow{v}}$, le terme $\overrightarrow{0}$ correspondant aux ions fixes. Or, $\gamma_0=\displaystyle\frac{ne^2\tau}{m}$ donc on obtient une \textbf{loi d'Ohm généralisée} $$\underline{\overrightarrow{j}}=\frac{\gamma_0}{1+i\omega\tau}\underline{\overrightarrow{E}}$$ Ainsi, $\underline{\overrightarrow{j}}$ et $\underline{\overrightarrow{E}}$ sont diphasés. A basses fréquences, il y a de l'effet Joule ; à hautes fréquences, $\underline{\overrightarrow{j}}$ et $\underline{\overrightarrow{E}}$ sont en quadrature et il n'y a pas d'effet Joule.

\paragraph{Remarque} Si $\underline{\gamma}\in\R_+\st$, il y a de l'effet Joule. Si $\underline{\gamma}\in i\R$, il n'y a pas d'effet Joule.

\subsection{ARQS dans un conducteur ohmique}

\paragraph{Électroneutralité} La conservation de la charge s'écrit $\displaystyle\frac{\partial\rho}{\partial t}+\div\overrightarrow{j}=0$. Or, $\overrightarrow{j}=\gamma_0\overrightarrow{E}$ d'après la loi d'Ohm. En passant en notation complexe et en utilisant l'équation de Maxwell-Gauss, on obtient : $$\underline{\rho}\left(\frac{\gamma_0}{\varepsilon_0}+i\omega\right)=0$$ donc $\underline{\rho}=0$.
\paragraph{Condition 1} Loi d'Ohm statique, $\omega\tau\ll 1$. Pour le cuivre, cela correspond à $\omega \ll 10^{14}$ rad/s, ce qui est peu restrictif.
\paragraph{Condition 2} Négliger le courant de déplacement : $\omega\ll\displaystyle\frac{\gamma_0}{\varepsilon_0}$ Pour le cuivre, $\gamma_0\approx 10^7$ S/m et $\varepsilon_0=8,89.10^{-12}$ F/m, donc cela correspond à $\omega\ll 10^{18}$ rad/s, ce qui est très peu restrictif.
\paragraph{Résumé} $\overrightarrow{j}=\gamma_0\overrightarrow{E}$ ; $\rot\overrightarrow{B}=\mu_0\overrightarrow{j}$ et $\div\overrightarrow{E}=0$
\section{Actions de Laplace sur un conducteur}
\subsection{Densité volumique de force de Laplace}
\paragraph{Force de Laplace} Un circuit $\Gamma$ dans un champ $\overrightarrow{B}$ subit une force locale appelée \textbf{force de Laplace} $$\delta\overrightarrow{F}=i\overrightarrow{dl}\wedge\overrightarrow{B}$$ 
\paragraph{Force de Laplace volumique}La force volumique associée est donnée par $$\overrightarrow{f_\text{v}}=\overrightarrow{j}\wedge\overrightarrow{B}$$
\subsection{Application : ancienne définition de l'Ampère}
La force de Laplace servait de définition à l'Ampère : 2 fils infinis parallèles espacés d'un mètre parcourus par 1 Ampère subissent une force linéique de $2.10^{-7}$ N.m$\inv$.
\subsection{Actions de Laplace en champ uniforme}
\paragraph{Résultante et moment pour l'intégralité du circuit} La résultante et le moment des forces de Laplace sont données par : $$\overrightarrow{R}=\oint_Ci\overrightarrow{dl}\wedge\overrightarrow{B_\text{ext}} \ \ \text{et} \ \ \overrightarrow{\Gamma}=\oint_C \overrightarrow{OP}\wedge\bigg(i\overrightarrow{dl}\wedge\overrightarrow{B_\text{ext}}\bigg)$$
La partie due au champ propre est nulle (elle peut déformer le circuit, mais est très faible). Ces expressions se simplifient :
\begin{itemize}
	\item si $\overrightarrow{B_\text{ext}}$ est uniforme, $\overrightarrow{R}=\overrightarrow{0}$ et $\overrightarrow{\Gamma}=\overrightarrow{M}\wedge\overrightarrow{B_\text{ext}}$
	\item si $\overrightarrow{B_\text{ext}}$ varie peu, on est dans l'approximation dipolaire et alors $\overrightarrow{R}=-\grad(E_p)=\grad\left(\overrightarrow{M}\cdot\overrightarrow{B}\right)$
\end{itemize}
Cela est faux si on ne considère qu'une partie du circuit.
\paragraph{Résultante et moment pour une portion de circuit} Si on considère uniquement une portion du circuit, il faut dessiner les forces élémentaires de Laplace sur cette portion et intégrer uniquement sur la portion en question.
\paragraph{Exemple pour une portion de circuit} On considère la partie du circuit qui est une barre en rotation autour d'un axe fixe et qui a un contact électrique sans frottement à l'autre bout. On a : $$\delta\overrightarrow{F}=I\overrightarrow{dl}\wedge\overrightarrow{B}=-Idr\overrightarrow{u_\theta}$$ On en déduit alors que $$\overrightarrow{R}=\int_\text{barre}\delta\overrightarrow{F}=-ILB\overrightarrow{u_\theta}$$ Puis, avec un bras de levier, on a : $$\delta\mathcal{M}_{Oz}=\left(\overrightarrow{OP}\wedge\delta\overrightarrow{F}\right)\cdot\overrightarrow{u_z}=-IrBdr$$ donc on obtient finalement \textbf{pour la barre} : $$\mathcal{M}_{Oz}=-\frac12IL^2B$$
\subsection{Exemples}
\paragraph{Exercice 12 TD7} On cherche à calculer le champ au voisinage de l'axe connaissant le champ sur l'axe. On utilise le fait que $\overrightarrow{B}$ est à flux conservatif sur un cylindre de côtes $z$ et $z+dz$ et de rayon r : $$0=\pi r^2\left(B_0(z+dz)-B_0(z)\right)+B_r^{(1)}(r,z)2\pi r dz$$ donc on en déduit qu'à l'ordre 1 : $$B_r=-\frac{r}{2}\frac{dB_0}{dz}$$ Au voisinage de l'axe, il n'y a pas de courant donc $\overrightarrow{B}$ y est à circulation conservative. On calcule alors la circulation sur un cadre rectangulaire de côtes $z$ et $z+dz$ qui passe par l'axe Oz et dont la "hauteur" vaut $r$ : $$0=B_0(z)dz+\int_{0}^{r}B_r(z+dz,r')dr'-B_z(z,r)dz-\int_{0}^{r}B_r(z,r')dr'$$ donc on en déduit que $$B_z=B_0+\frac{d}{dz}\int_{0}^{r}B_(z,r')dr'=B_0-\frac{r^2}{4}\frac{d^2B_0}{dz^2}+O(r^4)$$
\paragraph{Exemple : Action d'une spire sur une petite spire}On considère une grande spire de rayon $a$ parcourue par un courant $i$ et une petite spire (assimilée à un moment magnétique) de rayon $a'$ parcourue par un courant $i'$. Les spires sont coaxiales et à une distance $d$ l'une de l'autre. On cherche à calculer la résultante des actions de la spire sur le dipôle. On admet que sur l'axe : $$\overrightarrow{B_{\text{spire}}}=\frac{\mu_0 ia^2}{2(a^2+z^2)^{3/2}}\overrightarrow{u_z}$$
\subparagraph{Plan A} On utilise $E_p=-\overrightarrow{M}\cdot\overrightarrow{B_\text{spire}}$ puis $\overrightarrow{F}=-\grad(E_p)$
\subparagraph{Plan B} On utilise la troisième loi de Newton et le champ crée par le dipôle (c'est la seule approche possible si on ne nous donne pas $\overrightarrow{B_\text{spire}}$)
\subparagraph{Plan C} Calculer $\overrightarrow{B_\text{spire}}$ au voisinage de l'axe (cf. exercice 12 TD 7, l'exemple précédent) puis intégrer les forces élémentaires de Laplace.
\subparagraph{Résultat} On trouve un résultat qui vaut $$\overrightarrow{F}=-\frac{3\pi}{2}\mu_0 ii'a^2a'^2\frac{z}{(z^2+a^2)^{3/2}}\overrightarrow{u_z}$$
\section{Induction de Neumann}
On considère ici $\overrightarrow{B}$ variable et un circuit fixe.
\subsection{Loi de Faraday ; loi de Lenz}
\paragraph{Loi de Faraday} On dispose alors des lois de Faraday locale ou intégrale $$\rot\left(\overrightarrow{E}\right)=-\frac{\partial \overrightarrow{B}}{\partial t} \ \ \text{ou} \ \ e=\oint\overrightarrow{E}\cdot\overrightarrow{dl}=-\frac{d\Phi_{\overrightarrow{B}}}{dt}$$
Les orientations de $\overrightarrow{n}$ et de $e$ (pour $\Phi_{\overrightarrow{B}}$) sont liées.
\paragraph{Loi de Lenz} L'induction, par ses effets, s'oppose aux causes qui lui ont donné naissance.
\paragraph{Remarque} Lors du calcul du flux $\Phi_{\overrightarrow{B}}$, on a par principe de superposition $$\Phi_{\overrightarrow{B}}=\Phi_\text{propre}+\Phi_\text{ext}$$
\subsection{Auto-induction}
\paragraph{Flux propre ; inductance propre} Le flux propre est proportionnel à $i$, on note $L$ l'inductance, le coefficient de proportionnalité associé : $$\Phi_\text{propre}=Li$$ $L$ est positif et s'exprime en Henry (H). L ne dépend pas de la convention d'orientation choisie pour $i$. L peut être calculé en exprimant le flux ou en utilisant une approche énergétique avec la formule : $$\frac{1}{2}Li^2=\frac{B^2}{2\mu_0}$$
\paragraph{Exemple : un calcul d'inductance propre} On considère un solénoïde quasi-infini de $N$ spires et de longueur $H$ parcouru par une intensité $i$. Le champ magnétique crée vaut alors : $$\overrightarrow{B}=\begin{cases}
	\displaystyle\mu_0 \frac{N}{H} i \overrightarrow{u_z} \ \text{dans le solénoïde} \\[2mm]
	\overrightarrow{0} \ \text{en-dehors}
\end{cases}$$ Le flux propre vaut alors : $$\Phi_\text{propre}=N\Phi_\text{à travers une spire}=N\pi a^2\mu_0\frac{N}{H}i$$ On en déduit que $$L=\frac{\pi a^2\mu_0 N^2}{H}$$
On vérifie bien que $L\geqslant0$, que $L$ s'exprime en $H$ et que la dépendance est en $N^2$.

\subsection{Couplage de deux circuits par induction mutuelle}
\paragraph{Exemples de systèmes fonctionnant par induction mutuelle} Les transformateurs, les chargeurs sans contact, les plaques à induction, etc.
\paragraph{Couplage de bobines} En l’absence de flux extérieur, on a avec le théorème de superposition : $$\Phi_\text{à travers 1}=\Phi_\text{propre}+\Phi_{2\to 1}=L_1i_1+M_{2\to 1}i_2$$
\paragraph{Théorème de Neumann (admis)} Le théorème de Neumann stipule que $$M_{1\to 2}=M_{2\to 1}$$ qu'on appelle alors l'\textbf{inductance mutuelle}, et qu'on note alors $M$. $M$ s'exprime en Henry (H). L'inductance mutuelle dépend de la position relative des deux circuits et dépend des orientations de $i_1$ et $i_2$ : elle peut donc être négative.
\paragraph{Exemple : un calcul de M} On considère deux solénoïdes quasi-infinis de même longueur $H$. Le solénoïde 1 a $N_1$ spires, est de rayon $a_1$ et est parcouru par une courant $i_1$. Il contient le solénoïde 2 qui a $N_1$ spires, est de rayon $a_2$ et est parcouru par une courant $i_2$. On a au niveau du circuit 2 : $$\overrightarrow{B_1}=\mu_0i_1\frac{N_1}{H}\overrightarrow{u_z}$$ donc on en déduit que $$\Phi_{1\to 2}=N_2\Phi_{1\to \text{une spire de }2}=\frac{1}{H}\mu_0N_1N_2\pi a_2^2i_1$$ et donc on obtient la valeur : $$M=\pi\mu_0N_1N_2\frac{a_2^2}{H}$$ si $a_2<a_1$. $M$ est bien en H et est positive. M a une dépendance en $N_1N_2$ ce qui est rassurant si on augmente le nombre de spires de $1$ ou de $2$.
\paragraph{Énergie stockée} L'énergie stockée dans un couplage par inductance mutuelle est $$E_\text{mag}=\frac{1}{2}L_1i_1^2+\frac{1}{2}L_2i_2^2+Mi_1i_2$$ (faire un bilan électrique sur les deux circuits, puis sommer les équations électriques multipliées par la bonne intensité).
\paragraph{Inégalité sur $M$} Or, on sait que $E_\text{mag}$ est positive comme l'intégrale d'une énergie positive. Si on fixe $i_2\ne0$, on a alors pour tout $i_1$ : $$\frac{1}{2}L_1i_1^2+\frac{1}{2}L_2i_2^2+Mi_1i_2\geqslant 0$$ En termes de discriminant, cela se traduit par $\Delta=(M^2-L_1l_2)i_2\leqslant0$ et par conséquent on a l'inégalité : $$M^2\leqslant L_1L_2$$
\paragraph{Exemple du transformateur} On utilise deux solénoïdes coaxiaux de rayons très proches $a_1\approx a_2\approx a$ et les expressions précédentes des inductances et de la mutuelle. On cherche à exprimer $u_1$ en secondaire ouvert. La loi de Faraday fournit : $$e=-\frac{d}{dt}\left(\Phi_\text{propre 1}+\Phi_{2\to 1}\right)=-\frac{d}{dt}\left(L_1i_1+Mi_2\right)=-L_1\frac{di_1}{dt}-M\frac{di_2}{dt}$$La loi des mailles sur la maille $1$ donne $u_1+e_1=R_1i_1$ donc $$u_1=r_1i_1+L_1\frac{d_i1}{dt}+\underbrace{M\frac{di_2}{dt}}_{=0}$$ De même, $$u_2=\underbrace{L_2\frac{di_2}{dt}}_{=0}+M\frac{di_1}{dt}+\underbrace{R_2i_2}_{=0}$$ En passant en régime sinusoïdal forcé, on a les expressions : $$\begin{cases}
	\underline{u_1}=(R_1+jL_1\omega)\underline{i_1}\\
	\underline{u_2}=jM\omega\underline{i_1}
\end{cases}$$ D'où on en déduit la formule : $$\underline{u_2}=\frac{jM\omega}{R_1+jL_1\omega}\underline{u_1}$$ Si $R_1\ll L\omega$, on a $$\underline{u_2}=\frac{M}{L_1}\underline{u_1}=\frac{N_2}{N_1}\underline{u_1}$$ Dans ce cas, si $N_2>N_1$, alors $\underline{u_2}>\underline{u_1}$, et on parle de transformateur élévateur de tension. Dans le cas contraire, on parle de transformateur abaisseur de tension.
\subsection{Courants de Foucault}
\paragraph{Courants de Foucault} Les courants de Foucault sont dûs à l'induction dans un conducteur "massif". Il peuvent notamment servir au freinage magnétique et à la cuisson par induction.
\paragraph{Méthode} On part de $\overrightarrow{B}$ variable dans le temps. On en déduit $\overrightarrow{E}$ avec l'équation de Maxwell-Faraday, puis on déduit $\overrightarrow{j}$ de $\overrightarrow{E}$ grâce à la loi d'Ohm statique. On peut ensuite obtenir la puissance volumique avec la formule $$p_v=\overrightarrow{j}\cdot\overrightarrow{E}=\gamma E^2=\frac{j^2}{\gamma}$$
\paragraph{Exemple} Voir l'exercice 13 du TD8, corrigé en séance de TD par VdC
\section{Induction de Lorentz}
\subsection{Loi de Faraday ; loi de Lenz}
\paragraph{Loi de Lenz} L'induction, par ses effet, s'oppose aux causes qui lui ont donné naissance.
\paragraph{Loi de Faraday} Pour un circuit filiforme, on a : $$e=-\frac{d\Phi}{dt}$$ le signe $-$ traduit la loi de Lenz.
\paragraph{Roue de Barlow} Dans le cas de la roue de Barlow (le système est complété par une roue qui tourne sans frottement), il faut faire autrement car le circuit n'est plus filiforme. 
\paragraph{Champ électro-moteur} Pour se sortir du pétrin de la roue de Barlow, on peut appliquer la loi d'Ohm locale dans le référentiel lié au conducteur $\overrightarrow{j'}=\gamma\overrightarrow{E'}$ puis effectuer une transformation avec la force de Lorentz : $$\overrightarrow{j'}=\gamma\left(\overrightarrow{E}+\underbrace{\overrightarrow{v}_\text{conducteur}\wedge\overrightarrow{B}}_{\text{champ électro-moteur}}\right)$$
\subsection{Exemple/recette}
\paragraph{Analyse qualitative} La force électromotrice $E$ imposée par le générateur donne l'intensité $i$ par la loi des mailles. Cette intensité permet d'exprimer les forces de Laplace élémentaires et/ou leur moment, puis on peut appliquer le théorème mécanique adapté (théorème de la résultante ou du moment cinétique) pour déterminer la grandeur géométrique qui varie (angle ou longueur). Cette variation géométrique du système fait changer le flux de $\overrightarrow{B}$, ce qui crée une force électromotrice induite d'après la loi de Faraday.
\paragraph{Recette}Il faut expliquer tout cela à l'oral et mettre en forme les équations caractérisant les différentes étapes de ce processus pour répondre à la question posée. Notamment, on établira les équations électrique et mécanique, et on pourra ajouter l'équation électrique multipliée par $i$ à l'équation mécanique multipliée par $v$ ou $\dot{\theta}$ pour obtenir le bilan de puissance.
\paragraph{Exemple}
On applique le théorème de la résultante à la tige dans le référentiel galiléen lié aux rails, en projection sur $\overrightarrow{u_x}$ : $$m\ddot{l}=-k(l-l_0)+\underbrace{\overrightarrow{N_1}\cdot\overrightarrow{u_x}+\overrightarrow{N_2}\cdot\overrightarrow{u_x}}_{=0\text{ (pas de frottement)}}+\overrightarrow{F_\text{Lap}}\cdot\overrightarrow{u_x}$$ Or, $\delta\overrightarrow{F_\text{Lap}}=i\overrightarrow{dl}\wedge\overrightarrow{B}=-idlB\overrightarrow{u_x}$ donc $\overrightarrow{F_\text{Lap}}$ On obtient donc l'équation mécanique (M) : $$m\ddot{l}=-k(l-l_0)-iBL$$
En appliquant la loi des mailles au circuit électrique équivalent, on obtient $E=Ri-e$. Or, d'après la loi de Faraday, on a $$e=-\frac{d\Phi}{dt}$$ avec $\Phi=-LB(a+l)$, donc on obtient l'équation électrique (E) : $$E=Ri-BL\dot{l}$$ En faisant (E)$\times i-$(M)$\times\dot{l}$, on obtient finalement, après élimination des termes nuls et écriture sous forme primitivée : $$Ei-Ri^2=\frac{d}{dt}\left(\frac{1}{2}k(l-l_0)^2+\frac{1}{2}m\dot{l}^2\right)$$
\subsection{Puissance des actions de Laplace}
\paragraph{Induction de Lorentz} Couplage mécanique-électrique (alternateur, moteur) : $\overrightarrow{B}$ est fixe et le circuit est mobile. La loi de Faraday reste vraie si l’on ne modifie par le circuit (interrupteurs, boucles...). Deux approches classiques :
\begin{itemize}
	\item Calculer $\Phi$ puis $\displaystyle e=-\frac{d\Phi}{dt}$, établir l’équation électrique (E) puis l’équation mécanique (M) et effectuer le bilan de puissance $(E)\times i+(M)\times\dot{x}$
	\item Calculer la puissance des forces de Laplace et en déduire la force électromotrice induite (même puissance), si l’on ne peut pas calculer le flux.
\end{itemize}
\chapter{Propagation d'ondes électromagnétiques dans le vide}
\localtableofcontents
\section{Rappels}
\subsection{Équation de la propagation dans le vide}
\paragraph{Définition du vide} Le vide est défini dans ce cadre comme l'absence de charge ($\rho=0$) et l'absence de courant ($\overrightarrow{j}=\overrightarrow{0}$).
\paragraph{Propagation du champ électrique dans le vide} Avec $\rot(\rot)=\grad(\div)-\overrightarrow{\Delta}$, on obtient : $$\overrightarrow{\Delta}\overrightarrow{E}=\mu_0\varepsilon_0\frac{\partial^2 \overrightarrow{E}}{\partial t^2}$$
\paragraph{Propagation du champ magnétique dans le vide} Avec $\rot(\rot)=\grad(\div)-\overrightarrow{\Delta}$, on obtient : $$\overrightarrow{\Delta}\overrightarrow{B}=\mu_0\varepsilon_0\frac{\partial^2 \overrightarrow{B}}{\partial t^2}$$
\subsection{Relation de dispersion dans le vide}
\paragraph{Relation de dispersion} C'est la relation qui lie $\omega$ et $\overrightarrow{k}$. Elle est obtenue en cherchant des solutions de la forme $e^{i(\omega t-\overrightarrow{k}\cdot\overrightarrow{r})}$ dans l’équation de propagation. Dans l’équation de d’Alembert, on trouve, $\displaystyle k=\frac{\omega}{c}$ donc il n’y a pas de dispersion (même vitesse de propagation pour toutes les ondes).
\paragraph{Méthode} Il faut injecter une OPPH dans l'équation de propagation pour trouver la relation de dispersion.
\section{Ondes électromagnétiques planes}
\subsection{Ondes sphériques et ondes planes}
\paragraph{Ondes planes} Il s'agit d'ondes uniformes dans un plan perpendiculaire à la direction de propagation $\overrightarrow{u}$, ne dépend que de la coordonnée sur $\overrightarrow{u}$ : $$\overrightarrow{E}=\overrightarrow{E}\left(\overrightarrow{r}\cdot\overrightarrow{u},t\right) \ \text{et} \ \overrightarrow{B}=\overrightarrow{B}\left(\overrightarrow{r}\cdot\overrightarrow{u},t\right)$$
\paragraph{Limite du modèle} Pas physiquement réalisable (infini), mais bonne approximation locale
\subsection{Solutions de l'équation de d'Alembert en dimension 1}
On prend $E_x=f(x,t)$. L'équation de d'Alembert devient : $$\frac{\partial^2f}{\partial x^2}=\frac{1}{c^2}\frac{\partial^2 f}{\partial t^2}$$ On vérifie alors que toute fonction de la forme $E_x=f(x-ct)$ ou $E_x=g(x+ct)$ est solution de l'équation de d'Alembert en dimension 1. Réciproquement, on \textbf{admet} que toute solution de l'équation de d'Alembert en dimension est somme de deux fonctions de la forme $f(x-ct)$ et $f(x+ct)$.
\subsection{Structure d'une OPP}
\paragraph{Ondes planes progressives (OPP)} Ce sont des ondes progressives donc $$E,B=f\left(ct-\overrightarrow{r}\cdot\overrightarrow{u}\right)$$ Dans le vide, les équations de Maxwell donnent pour ce type de solutions la relation de structure : $$\overrightarrow{u}\wedge\overrightarrow{E}=c\overrightarrow{B}$$ L'onde est donc transverse et $\left(\overrightarrow{u},\overrightarrow{E},\overrightarrow{B}\right)$ forme un trièdre direct.
\section{Ondes planes progressives harmoniques}
\subsection{Modèle des OPPH}
\paragraph{OPP harmoniques/sinusoïdales/monochromatiques (OPPH)} Toutes les composantes de $\overrightarrow{E}$ et de $\overrightarrow{B}$ sont des fonctions sinusoïdales de $\omega t-\overrightarrow{k}\cdot\overrightarrow{r}$. On appelle \textbf{vecteur d'onde} le vecteur $\displaystyle\overrightarrow{k}=\frac{2\pi}{\lambda}\overrightarrow{u}$ et \textbf{pulsation} le scalaire $\omega=\displaystyle\frac{2\pi}{T}$. En notation complexe, $\overrightarrow{E}=Re\left(\overrightarrow{\underline{E}}\right)$. On a : $$\overrightarrow{\underline{E}}=e^{i(\omega t-\overrightarrow{k}\cdot\overrightarrow{r})}\bigg(E_{0,x}e^{i\varphi_x}\overrightarrow{u_x}+E_{0,y}e^{i\varphi_y}\overrightarrow{u_y}+E_{0,z}e^{i\varphi_z}\overrightarrow{u_z}\bigg)$$
\subsection{Utilisation de la notation complexe ; relations de structure d'une OPPH}
\paragraph{Astuces de calcul en notation complexe} les relations suivantes simplifient les calculs en notation complexe :
\begin{itemize}
	\item $\displaystyle \frac{d}{dt}\left(\overrightarrow{\underline{E}}\right)=i\omega\overrightarrow{\underline{E}}$
	\item $\displaystyle \div\left(\overrightarrow{\underline{E}}\right)=-i\overrightarrow{k}\cdot\overrightarrow{\underline{E}}$
	\item $\displaystyle \rot\left(\overrightarrow{\underline{E}}\right)=-i\overrightarrow{k}\wedge\overrightarrow{\underline{E}}$
	\item $\displaystyle\overrightarrow{\Delta}\overrightarrow{\underline{E}}=-k^2\overrightarrow{\underline{E}}$
\end{itemize}
\paragraph{Caractère transverse} On considère $\overrightarrow{E}=\text{Re}\left(\overrightarrow{\underline{E}}\right)$ et $\overrightarrow{B}=\text{Re}\left(\overrightarrow{\underline{B}}\right)$ dans le vide où  $$\overrightarrow{\underline{E}}=\overrightarrow{\underline{E_0}}e^{i(\omega t-\overrightarrow{k}\cdot\overrightarrow{r})} \ \text{et} \ \overrightarrow{\underline{B}}=\overrightarrow{\underline{B_0}}e^{i(\omega t-\overrightarrow{k}\cdot\overrightarrow{r})}$$ Avec les équation de Maxwell-Gauss et 
de Maxwell-Thomson en notation complexe, on obtient $\overrightarrow{k}\cdot\overrightarrow{\underline{E}}=0$ et $\overrightarrow{k}\cdot\overrightarrow{\underline{B}}=0$ donc l'onde est transverse électrique et transverse magnétique.
\paragraph{Relation de structure} Toujours avec les mêmes notations et dans le vide, en notant $\overrightarrow{u}$ le renormalisé de $\overrightarrow{k}$, on obtient grâce aux équation de Maxwell-Ampère et de Maxwell-Faraday : $$-i\overrightarrow{k}\wedge\overrightarrow{\underline{B}}=\frac{1}{c^2}i\omega\overrightarrow{\underline{E}} \ \ \text{et} \ \ -i\overrightarrow{k}\wedge\overrightarrow{\underline{E}}=-i\omega\overrightarrow{\underline{B}}$$ Ce qui se réécrit, avec $\omega=kc$ (relation de dispersion dans le vide) et le vecteur $\overrightarrow{u}$ : $$\overrightarrow{u}\wedge\overrightarrow{\underline{E}}=c\overrightarrow{\underline{B}} \ \ \text{et} \ \ c\overrightarrow{\underline{B}}\wedge\overrightarrow{u}=\overrightarrow{\underline{E}}$$
\paragraph{Résumé} Pour une OPPH dans le vide, on a : 
$$\begin{cases}
	\overrightarrow{u}\cdot\overrightarrow{\underline{E}}=0 \\
	\overrightarrow{u}\cdot\overrightarrow{\underline{B}}=0 \\
	\overrightarrow{u}\wedge\overrightarrow{\underline{E}}=c\overrightarrow{\underline{B}}\\
	c\overrightarrow{\underline{B}}\wedge\overrightarrow{u}=\overrightarrow{\underline{E}}\\
	E=cB
\end{cases}$$ Pour ces relations, \textbf{toujours} faire un petit schéma et bien se souvenir que le trièdre $\left(\overrightarrow{\underline{E}},c\overrightarrow{\underline{B}},\overrightarrow{u}\right)$ est \textbf{direct}.
\subsection{Intérêt et limites du modèle de l'OPPH}
\paragraph{Limite} Une onde plane en général n'est pas réalisable expérimentalement car elle devrait être uniforme dans tout plan perpendiculaire à $\overrightarrow{u}$ et possèderait alors une énergie infinie.
\paragraph{Intérêt} En revanche, ces ondes présentent un intérêt car toute onde est une superposition d'OPPH (\textit{cf}. décomposition de Fourier, admis).
\section{États de polarisation}
On rappelle ici que dans le vide, une onde électromagnétique est transverse électrique et transverse magnétique.
\paragraph{Polarisation de la lumière} Cela correspond aux signaux sinusoïdaux en mode XY dans le plan perpendiculaire à la propagation à $z$ fixé : 
\begin{itemize}
	\item la polarisation est rectiligne si les signaux dessinent un segment (si $E_x=0$ ou $E_y=0$ ou $x$ et $y$ en phase ou en opposition de phase)
	\item la polarisation est circulaire si les signaux selon $x$ et $y$ ont même norme et sont en quadrature
	\item sinon la polarisation est elliptique
\end{itemize}
\paragraph{Attention} Ne pas confondre la direction de propagation et la direction de polarisation.
\subsection{Polarisation rectiligne}
\paragraph{Onde polarisée rectilignement} Une onde est polarisée rectilignement lorsque pour tout point $M$, $\overrightarrow{E}(M)$ garde une direction fixe.
\paragraph{Exemple} $\overrightarrow{\underline{E}}=E_0\overrightarrow{u_x}e^{i(\omega t-kz)}$ se propage selon $z$ et est polarisée selon $\overrightarrow{u_x}$.
\paragraph{Remarque} Toute OPPH $\overrightarrow{\underline{E}}=(E_{0x}e^{i\varphi_x}\overrightarrow{u_x}+E_{0y}e^{i\varphi_y}\overrightarrow{u_y})e^{i(\omega t-kz)}$ est la somme de deux ondes polarisées rectilignement. 
\paragraph{Caractérisation des ondes polarisées rectilignement} L'onde $$(E_{0x}e^{i\varphi_x}\overrightarrow{u_x}+E_{0y}e^{i\varphi_y}\overrightarrow{u_y})e^{i(\omega t-kz)}$$ est polarisée rectilignement si, et seulement si, au moins une des conditions suivantes est respectée :
\begin{itemize}
	\item $\varphi_x \equiv \varphi_y \ [2\pi]$
	\item $\varphi_x \equiv \varphi_y +\pi\ [2\pi]$
	\item $E_{0x}=0$
	\item $E_{0y}=0$
\end{itemize}
Pour les deux premières conditions, penser au mode XY de l'oscilloscope.
\subsection{Polarisation circulaire}
\paragraph{Onde polarisée circulairement}
Une onde est polarisée circulairement lorsque pour tout point $M$, l'extrémité de $\overrightarrow{E}(M,t)$ décrit un cercle. Lorsque le cercle est décrit dans le sens direct, l'onde est dite polarisée circulairement gauche, et lorsque le cercle est décrit dans le sens indirect, l'onde est dite polarisée circulairement droite (convention volant de voiture).

\paragraph{Caractérisation des ondes polarisées circulairement} L'onde $$(E_{0x}e^{i\varphi_x}\overrightarrow{u_x}+E_{0y}e^{i\varphi_y}\overrightarrow{u_y})e^{i(\omega t-kz)}$$ est polarisée circulairement si, et seulement si, les deux conditions suivantes sont respectées :
\begin{itemize}
	\item $\varphi_x\equiv\varphi_y\pm\frac{\pi}{2} \ [2\pi]$
	\item et $E_{0x}=E_{0y}$
\end{itemize}
\paragraph{Autre formulation de la caractérisation} L'onde précédente est polarisée circulairement droite si, et seulement si $\underline{E_{0x}}=i\underline{E_{0y}}$ et elle est polarisée circulairement gauche si, et seulement si, $\underline{E_{0x}}=-i\underline{E_{0y}}$. ET rappelons que $\underline{E_{0x}}=E_{0x}e^{i\varphi_x}$ et symétriquement pour $y$.
\paragraph{Remarque} La somme de deux ondes polarisées circulairement droites est une onde polarisée circulairement droite.
\paragraph{Remarque} Toute polarisation est la somme d'une onde polarisée circulairement gauche et d'une onde polarisée circulairement droite.
\paragraph{Exemple d'application de la remarque} Par exemple, l'onde $\overrightarrow{\underline{E}}=E_0e^{i(\omega t-kz)}\overrightarrow{u_x}$ polarisée rectilignement est la somme de $$\overrightarrow{\underline{E_1}}=\frac{E_0}{2}\begin{pmatrix}
	1 \\
	i \\
	0
\end{pmatrix}e^{i(\omega t-kz)}$$ (onde polarisée circulairement droite, passer en notation réelle pour le voir) et de l'onde $$\overrightarrow{\underline{E_2}}=\frac{E_0}{2}\begin{pmatrix}
	1 \\
	-i \\
	0
\end{pmatrix}e^{i(\omega t-kz)}$$ (onde polarisée circulairement gauche, passer en notation réelle pour le voir). Cela résulte simplement des formules d'Euler.
\paragraph{Moyen mnémotechnique personnel} Pour retenir ces histoires de droite et gauche, retenir que cela fonctionne à l'envers. Avec du $+$, c'est dans le sens négatif \textit{ie} droite, et avec du $-$, c'est dans le sens positif, \textit{ie} gauche (cf. la deuxième formulation de la caractérisation et l'exemple)
\subsection{Autres états de polarisation (culture)}
Dans le cas général, l'onde $$(E_{0x}e^{i\varphi_x}\overrightarrow{u_x}+E_{0y}e^{i\varphi_y}\overrightarrow{u_y})e^{i(\omega t-kz)}$$ est polarisée elliptiquement.
\subsection{Ondes électromagnétiques émises par des sources réelles}
\paragraph{Ondes radio/Wifi} Les ondes monochromatiques sont un "bon" modèle pour ce type d'ondes. Ces ondes sont polarisées.
\paragraph{Ondes lumineuses} \textit{cf} les chapitres O2 et O2' et le modèle des trains d'onde : il est difficile d'obtenir des interférences en optique. La lumière émise par le Soleil ou une ampoule est non polarisée.
\subsection{Polariseur ; loi de Malus}
\paragraph{Définition d'un polariseur} Un polariseur/polaroïde idéal laisse passer tous les signaux dans la direction de polarisation et aucun perpendiculaire à cette direction. Il renvoie donc une lumière polarisée.
\paragraph{Cas des ondes centimétriques} Comme leur nom l'indique, $\lambda$ est de l'ordre du cm, et par conséquent $f$ est de l'ordre de 10 GHz. Pour mettre en œuvre un polariseur, on place une sorte de peigne fin dont les dents sont espacées de quelques cm et orientés selon $y$ et parallèles à $x$. La partie de l'onde polarisée selon $x$ est transmise tandis que celle sur $y$ est réfléchie. En sortie, l'onde est polarisée selon $\overrightarrow{u_x}$. 
\paragraph{Cas des ondes lumineuses} Le montage est constitué de deux polariseurs successifs : le premier est fixe, le deuxième, appelé analyseur, à une direction de polarisation formant un angle $\alpha$ avec le premier. On ne connaît pas $\overrightarrow{E}_avant$. On a $$\overrightarrow{E}_\text{après polariseur}=E_0e^{i(\omega t-kz)}\overrightarrow{u_1}$$ et $$\overrightarrow{E}_{\text{après analyseur}}=E_0\cos(\alpha)e^{i(\omega t-kz)}\overrightarrow{u_2}$$ En effet, $\overrightarrow{\underline{E}}_\text{après polariseur}=\overrightarrow{\underline{E}}_\text{après analyseur}$ donc en projetant avec on obtient l'expression.
\paragraph{Exemple} Des molécules conductrices espacées de quelques centaines de nm forment un polariseur pour des ondes lumineuses.
\paragraph{Remarque} Un polariseur permet de savoir si une onde est polarisée rectilignement : il suffit de tourner le polariseur et de voir si on observe une extinction (absence d'onde lumineuse transmise). S'il y a une extinction, c'est que l'onde initiale est polarisée rectilignement perpendiculairement à la direction sur laquelle on s'est arrêté avec le polariseur.
\paragraph{Loi de Malus} Si l'analyseur forme un angle $\alpha$ avec le polariseur, alors : $$I_\text{après analyseur}=\cos(\alpha)^2I_\text{avant analyseur}$$ avec $I=\langle \overrightarrow{pi}\cdot\overrightarrow{u_z}\rangle$. En effet, $I$ est proportionnelle à $\langle\lVert\overrightarrow{E}\rVert^2\rangle$, par conséquent, on a une constant multiplicative près: $$I_\text{après analyseur}=\langle\lVert E_0\cos(\alpha)\cos(\omega t-kz)\overrightarrow{u_x}\rVert^2\rangle=E_0^2\cos^2(\alpha)\frac{1}{2}=\cos^2(\alpha)I_0$$
\paragraph{Polariseur-analyseur croisés} Pour $\alpha=\displaystyle\pm\frac{\pi}{2}$, on parle de \textbf{polariseur-analyseur croisés}.
\paragraph{Repérer un polariseur en TP} En TP, on peut facilement repérer les polariseurs car ils sont toujours avec l'analyseur et se ressemblent. Il est donc facile de les distinguer sur le banc optique.
\paragraph{Montage alternatif} Parfois, on considèrera un polariseur-analyseur croisé et on intercalera un troisième polariseur au milieu.
\section{Énergie associée à une onde électromagnétique}
\subsection{Compatibilité avec le théorème de Poynting}
\paragraph{Aspects énergétiques} On s'intéresse à une OPPH se propageant selon $\overrightarrow{u_z}$ pour une onde polarisée selon $\overrightarrow{u_x}$ : on a $$\overrightarrow{E}=E_0\cos\left(\omega t-kz\right)$$ On peut alors déterminer $\overrightarrow{B}$ avec la relation de structure dans le vide. On obtient : $$c\overrightarrow{\underline{B}}=\overrightarrow{u_z}\wedge\overrightarrow{\underline{E}}=E_0e^{i(\omega t-kz)}\overrightarrow{u_y}$$ puis \textbf{en notation réelle} $$\overrightarrow{\pi}=\frac{\overrightarrow{E}\wedge\overrightarrow{B}}{\mu_0}=\frac{E_0^2}{\mu_0 c}\cos^2(\omega t-kz)\overrightarrow{u_z}$$ et alors $$u_{em}=\frac{\varepsilon_0E^2}{2}+\frac{B^2}{2\mu_0}=\varepsilon_0 E_0^2\cos^2(\omega t-kz)$$  
\paragraph{Conclusions} Le théorème de Poynting $$\frac{\partial u_{em}}{\partial t}+\div(\overrightarrow{\pi})=-\overrightarrow{j}\cdot\overrightarrow{E}$$est vérifié (calculer $\div(\overrightarrow{\pi})$ et la dérivée partielle de l'énergie électromagnétique volumique par rapport au temps). Il y a \textbf{équiproportion de l'énergie pour une OPPH} : $$\frac{\varepsilon_0 E^2}{2}=\frac{B^2}{2\mu_0}$$
\subsection{Puissance instantanée d'une OPP à travers une surface ; vitesse d'énergie}
\paragraph{Vitesse de propagation de l’énergie} On applique la relation de structure pour les OPPH dans le vide qui fournit en notation réelle : $$\overrightarrow{\pi}=\frac{\overrightarrow{E}\wedge\overrightarrow{B}}{\mu_0}=\frac{E^2}{\mu_0c}\overrightarrow{u}=\varepsilon_0cE^2\overrightarrow{u}$$ Si on considère désormais l'énergie qui traverse une surface $S$ de vecteur normal extérieur $\overrightarrow{n}$ entre $t$ et $t+dt$ en arrivant à $\overrightarrow{v_e}$, l'énergie vaut : $$\overrightarrow{\pi}\cdot S\overrightarrow{n}dt=u_{em}Sv_edt\cos(\overrightarrow{v_e}\angle\overrightarrow{n})=u_{em}S\overrightarrow{v_e}\cdot\overrightarrow{n}dt$$ Par conséquent, ceci valant pour tout vecteur $\overrightarrow{n}$, on en déduit que $\overrightarrow{\pi}=u_{em}\overrightarrow{v_e}$, si bien que $$\overrightarrow{v_e}=c\overrightarrow{u}$$
\subsection{Notation complexe et moyennes de grandeurs quadratiques}
\paragraph{Notation complexe et moyenne} Pour $f$ et $g$ sinusoïdales synchrones, on a : $$\langle fg\rangle=\frac12Re\left(\underline{fg\st}\right)$$ On peut ainsi calculer les moyennes des grandeurs quadratiques en notation complexe : $$\langle \overrightarrow{\pi}\rangle=\frac{1}{2\mu_0}\text{Re}\left( \overrightarrow{\underline{E}}\wedge\overrightarrow{\underline{B}^\star}\right)$$
\paragraph{Démonstration} Il suffit d'effectuer le calcul avec la formule de trigonométrie $\cos(a)\cos(b)=\frac{1}{2}\left(\cos(a-b)+\cos(a+b)\right)$ et de remarquer que cela coïncide effectivement avec la partie réelle recherchée.
\paragraph{Remarque} Cette formule est bien homogène, les deux membres sont $\R$-linéaires en $f$ et elle est invariante par échange de $f$ et $g$.
\paragraph{Valeur efficace} On rencontre parfois la \textbf{valeur efficace}, aussi appelée \textit{Root Mean Square} en anglais (souvent alors abrégée en RMS), qui est définie par : $$f_\text{eff}=\sqrt{\langle f^2\rangle}$$ On a alors pour une fonction sinusoïdale d'amplitude $A$ : $$f_\text{eff}=\frac{A}{\sqrt{2}}$$
\chapter{Propagation d'ondes électromagnétiques dans un milieu dispersif}
\localtableofcontents
\section{Modélisation plus réaliste des signaux}
\subsection{Insuffisance du modèle de l'OPPH}
Ces ondes ne sont pas réalisables expérimentalement puisqu'elles devraient posséder une énergie infinie.
\subsection{Exemple simple : superposition de deux ondes OPPH de pulsations voisines}
\paragraph{Vitesse de propagation des OPPH seules} On considère deux OPPH $f=1\cos(\omega_1t-k_1x)$ et $g=A\cos(\omega_2t-k_2x)$. Alors, $f$ se propage à $\displaystyle\frac{\omega_1}{k_1}$ et $g$ se propage à $\displaystyle\frac{\omega_2}{k_2}$. 
\paragraph{Somme des deux OPPH} On a par formule de trigonométrie : $$f+g=\underbrace{2A\cos\left(\frac{\omega_1-\omega_2}{2}-\frac{k_1-k_2}{2}\right)}_\text{enveloppe}\underbrace{\cos\left(\frac{\omega_1+\omega_2}{2}-\frac{k_1+k_2}{2}\right)}_\text{porteuse}$$ On observe donc des battements.
\paragraph{Vitesses de phase et de groupe dans ce cas particulier} On définit la vitesse de phase par $$v_\varphi=\frac{\omega_1+\omega_2}{k_1+k_2}$$ ainsi que la vitesse de groupe par $$v_g=\frac{\omega_1-\omega_2}{k_1-k_2}$$ \textbf{La porteuse se déplace à la vitesse de phase, l'enveloppe se déplace à la vitesse de groupe.}
\subsection{Forme générale d'un paquet d'ondes ; spectre en amplitude}
\paragraph{Paquet d'ondes} Pour un obtenir un signale borné dans le temps, on somme une infinité de sinusoïdes : $$s=\sum_{k}\underline{a}(k)e^{i(\omega_k t-kx)} \ \ \text{ou}\ \ s=\int \underline{a}(k)e^{i(\omega_k t-kx)} dk$$ on admet qu'une distribution de $\underline{a}(k)$ en gaussienne autour d'une valeur $k_0$ donne pour enveloppe une bulle dont la longueur $\Delta x$ est d'autant plus grande que $\Delta k$ est faible.
\subsection{Évolution d'un paquet d'ondes ; vitesses de phase et de groupe}
$$s=e^{i(\omega_0t-k_0x)}\sum_{k}\underline{a}(k)e^{i\left((\omega-\omega_0)t-(k-k_0)x\right)}$$ où le premier facteur correspond à la sinusoïde et la somme correspond à l'enveloppe. Si $\underline{a}(k)$ ne prend des valeurs significatives que sur $\Delta k\ll k_0$ autour de $k_0$, alors $\omega_k-\omega_0\approx(k-k_0)\frac{d\omega}{dk}(k_0)$. Ainsi, l'enveloppe se déplace à la vitesse de groupe $$v_g=\frac{d\omega}{dk}(k_0)$$ et la sinusoïde à la vitesse de phase $$v_\varphi=\frac{\omega_0}{k_0}$$ Cette approximation n'est valable que sur un temps très court car on la multiplie par $t$. On a bien une enveloppe car la longueur d'onde $\displaystyle\frac{1}{k-k_0}$ est très grande devant celle de la sinusoïde $\displaystyle\frac{1}{k}$.
\section{Propagation d'un signal dans l'ionosphère}
\subsection{Modélisation de l'ionosphère}
\paragraph{Plasma} Il s’agit d’un milieu ionisé où les électrons et les ions sont libres. Pour simplifier, on suppose que chaque ion est un proton. On a ainsi une unique densité $n$ de charges pour les protons et électrons.
\paragraph{Ionosphère}
L'ionosphère est située entre 80 kilomètres et 800 kilomètres au-dessus de la Terre. C'est un milieu ionisé, c'est-à-dire un plasma, qui comporte notamment ds électrons, des ions, du dioxygène, du néon, de l'hélium, ...
\paragraph{Nombre d'électrons par unité de volume} Dans l'ionosphère et dans un plasma en général, la densité volumique d'électrons est $n=10^{12}$m$^{-3}$.
\paragraph{Rappel} On rappelle que pour un métal conducteur, la densité volumique d'électrons est de l'ordre de $n=10^{30}$m$^{-3}$.
\subsection{Réaction du plasma au passage d'une OPPH transverse}
\paragraph{Attention} On prend ici une hypothèse supplémentaire \textit{a priori} non nécessaire, celle du fait que l'onde est transverse.
\paragraph{Courant et charges} On applique le principe fondamental de la dynamique à un électron pour obtenir : $$m_e\overrightarrow{a}=-e\left(\overrightarrow{E}+\overrightarrow{v}\wedge\overrightarrow{B}\right)\approx-e\overrightarrow{E_\text{onde}}$$ car pour un électron non relativiste, $vB\ll cB\approx E$, et puisque le plsama est dilué, $E\approx E_\text{onde}$. On en déduit $\overrightarrow{\underline{v}}$ puis $\overrightarrow{\underline{j}}$ : $$\overrightarrow{\underline{j}}=\underline{\gamma}\overrightarrow{\underline{E}} \ \ \text{avec} \ \ \underline{\gamma}=-i\frac{e^2n}{\omega m_e}$$ où $\underline{\gamma}$ est la conductivité complexe. Il n'y a pas de loi d'effet Joule ici, la puissance moyenne vaut $$\langle \overrightarrow{j}\cdot\overrightarrow{E}\rangle=\frac12Re\left(\overrightarrow{\underline{j}}\cdot\overrightarrow{\underline{E}\st}\right)=0$$ pour les charges, on montre grâce à la conservation de $\rho$ que $\displaystyle \left(i\omega-\frac{\underline{\gamma}}{\varepsilon_0}\right)\underline{\rho}=0$ donc $\rho=0$ sauf à la \textbf{pulsation plasma} : $$\omega_p=\sqrt{\frac{e^2n_0}{m_e\varepsilon_0}}$$
\subsection{Relation de dispersion pour des ondes transverses ; pulsation plasma}
\paragraph{Équation de propagation} On part de $\rot(\rot)=\grad(\div)-\overrightarrow{\Delta}$ pour montrer $$c^2\overrightarrow{\Delta}\overrightarrow{E}=\frac{d^2\overrightarrow{E}}{dt^2}+\omega_p^2\overrightarrow{E}$$ 
\paragraph{Relation de dispersion}On en déduit, en injectant une OPPH dans l'équation de propagation, la relation dispersion : $$k^2c^2=\omega^2-\omega_p^2$$ On remarque que cette relation est bien homogène et qu'on retrouve la relation de dispersion dans le vide lorsque $\omega_p=0$, c'est-à-dire lorsque $n=0$.
\paragraph{Analyse de la relation de dispersion}Ainsi
\begin{itemize}
	\item si $\omega<\omega_p$, il n’y a pas de solution avec $k$ réel, on a une solution de la forme $e^{i\omega t}e^{-\frac{z}{\delta}}$ qui est une onde évanescente (stationnaire) ne se propageant pas, elle correspond à une réflexion d’énergie (TSF, radio basses fréquences réfléchies par l’ionosphère)
	\item si $\omega>\omega_p$, on a une solution correspondant à une onde se déplaçant à $\displaystyle\frac{\omega}{\sqrt{\omega^2-\omega_p^2}}c>c$. Cette onde n'est pas réalisable car elle est infinie dans le temps et dans l'espace.
\end{itemize}
\subsection{Cas d'un paquet d'ondes $\omega_0>\omega_p$}
Dans cette sous-section, $\omega_0>\omega_p$ est à comprendre au sens que toutes les pulsations du paquet d'ondes sont supérieures à $\omega_p$.
\paragraph{Détermination des vitesses de phase et de groupe} Dans un plasma, si $\omega>\omega_p$, on a $\omega^2-\omega_p^2=k^2c^2$ donc $$v_\varphi=\frac{c}{\sqrt{1-\left(\frac{\omega_0}{\omega}\right)^2}}>c$$ mais en différenciant on trouve $$v_g=\frac{c^2}{v_\varphi}=c\sqrt{1-\left(\frac{\omega_0}{\omega}\right)^2}<c$$
\paragraph{Remarque} On est rassuré par le fait que $v_g<c$ car c'est la vitesse d'un paquet d'ondes réalisable expérimentalement. D'autre part, $v_\varphi>c$, ce qui n'est pas réellement un problème puisqu'une OPPH n'est pas réalisable expérimentalement et si la porteuse va plus vite que l'enveloppe, son amplitude devient rapidement négligeable.
\paragraph{Indice du plasma} L'indice du plasma, défini par $n=\displaystyle\frac{c}{v_\varphi}$, est donc strictement inférieur à $1$ dans un plasma avec $\omega_0>\omega_p$, contrairement à ce dont on a l'habitude en optique !
\paragraph{Relation de structure} On a les relations de structure suivantes dans ce cas (toujours faire le petit dessin du trièdre direct) :
\begin{itemize}
	\item $\overrightarrow{k}\cdot\underline{\overrightarrow{B}}=0$ (équation de Maxwell-Thomson) ;
	\item $\overrightarrow{k}\cdot\underline{\overrightarrow{E}}=0$ (équation de Maxwell-Gauss) ;
	\item $\overrightarrow{u}\wedge\underline{\overrightarrow{E}}=v_\varphi\underline{\overrightarrow{B}}$ avec $\overrightarrow{u}=\displaystyle\frac{\overrightarrow{k}}{\lVert\overrightarrow{k\rVert}}$ (équation de Maxwell-Faraday).
\end{itemize}
\paragraph{Remarque} On ne peut plus généraliser aux OPP (sans le H) puisque $v_\varphi=v_\varphi(k)$.
\paragraph{Méthode générale de détermination des vitesses de phase et de groupe} En général, on peut trouver facilement la vitesse de phase à partir de la relation de dispersion. Pour la vitesse de groupe, il faut le plus souvent dériver par rapport à $k$ la relation de dispersion et utiliser l'expression de la vitesse de phase trouvée précédemment.
\subsection{Cas d'un paquet d'ondes $\omega_0<\omega_p$}
Dans cette sous-section, $\omega_0>\omega_p$ est à comprendre au sens que toutes les pulsations du paquet d'ondes sont inférieures à $\omega_p$.
\paragraph{Calcul de $\underline{k}$} On considère désormais une OPPH de la forme : $$\overrightarrow{E}=\text{Re}(\underline{\overrightarrow{E}}) \ \ \text{avec} \ \ \underline{\overrightarrow{E}}=\underline{\overrightarrow{E_0}}e^{i(\omega t-\underline{k}z)}$$ La relation de dispersion devient alors $c^2\underline{k}^2=\omega^2-\omega_p^2$, ce qui donne $$\underline{k}=\pm i \frac{\sqrt{\omega_p^2-\omega^2}}{c}$$
\paragraph{Calcul de l'onde résultante} En notation réelle, on obtient $$\overrightarrow{E}\cdot\overrightarrow{u_x}=E_{0,x}e^{\pm z\frac{\sqrt{\omega_p^2-\omega^2}}{c}}\cos(\omega t+\varphi_x)$$ et un terme similaire sur $\overrightarrow{u_y}$. C'est une onde stationnaire puisque $E_x=f(t)g(z)$ \textbf{en notation réelle}.
\paragraph{Onde stationnaire}Les ondes stationnaires sont celles qui présentent un découplage spatio-temporel \textbf{en notation réelle}. En effet, l'onde $$e^{i(\omega t-kz)}=e^{i\omega t}e^{-ikz}$$ est bien découplée au niveau spatio-temporel, mais ce n'est pas une onde stationnaire.
\paragraph{Cas d'un plasma occupant le demi-espace $z>0$, onde évanescente} Pour un plasma occupant le demi-espace $z>0$, le terme en $e^{+ z\frac{\sqrt{\omega_p^2-\omega^2}}{c}}$ divergerait lorsque $z\to+\infty$, donc on ne garde que le terme en $e^{-z\frac{\sqrt{\omega_p^2-\omega^2}}{c}}$. On parle d'\textbf{onde évanescente} car son amplitude tend vers 0 lorsque $z$ tend vers $+\infty$.
\paragraph{Relation de structure} Dans ce cas, l'équation de Maxwell-Faraday fournit : $$\underline{k}\overrightarrow{u}\wedge\underline{\overrightarrow{E}}=\omega\underline{\overrightarrow{B}}$$ Ainsi, $\underline{E_x}$ et $\underline{B_y}$ sont en quadrature.
\subsection{Résumé : signal arrivant sur l'ionosphère}
\paragraph{Antenne radio} Il y a réflexion totale sur l'ionosphère pour les "ondes larges".
\paragraph{Cas où $\omega_0 >\omega_p$} Dans ce cas, une partie de l'onde est transmise dans le plasma et une autre partie de l'onde est réfléchie. La partie transmise possède alors une vitesse plus faible que la partie réfléchie puisque $c>v_g$ (logique, le plasma "ralentit" l'onde).
\paragraph{Cas où $\omega_0<\omega_p$} Dans ce cas, l'onde est totalement réfléchie par le plasma, il n'y a pas de transmission.
\paragraph{Ordres de grandeur} Avec $n=10^{12}$m$^{-3}$, on a $f_p\approx10$ MHz. En guise de comparaison, la bande radio FM a une fréquence de l'ordre de 100 MHz, et la (regrettée par VdC) bande AM a une fréquence de l'ordre de 1 MHz.
\subsection{Comparaison avec les ondes électromagnétiques dans le vide}
\paragraph{EM5 ou EM6 avec $\omega\gg\omega_p$} On a $$c^2\overrightarrow{\Delta}\overrightarrow{E}=\frac{\partial^2\overrightarrow{E}}{\partial t^2}$$ Si $\underline{\overrightarrow{E}}$ est de la forme $e^{i(\omega t-kz)}$, alors on a le relation de dispersion $$\omega^2=k^2c^2$$ et on a $c\underline{\overrightarrow{B}}=\overrightarrow{u}\wedge\underline{\overrightarrow{E}}$. On a de plus $c\overrightarrow{B}=\overrightarrow{u}\wedge\overrightarrow{E}$ de façon générale pour les OPP.
\paragraph{EM6 avec $\omega>\omega_p$} On a $$c^2\overrightarrow{\Delta}\overrightarrow{E}=\frac{\partial^2\overrightarrow{E}}{\partial t^2}+\omega_p^2\overrightarrow{E}$$ Si $\underline{\overrightarrow{E}}$ est de la forme $e^{i(\omega t-kz)}$, alors on a le relation de dispersion $$\omega^2-\omega_p^2=k^2c^2$$ On dispose de la relation de structure suivante, valable uniquement pour des OPPH (on ne peut plus généraliser aux OPP) :  $\overrightarrow{u}\wedge\underline{\overrightarrow{E}}=v_\varphi\underline{\overrightarrow{B}}$
\paragraph{EM6 avec $\omega<\omega_p$} On a $$c^2\overrightarrow{\Delta}\overrightarrow{E}=\frac{\partial^2\overrightarrow{E}}{\partial t^2}+\omega_p^2\overrightarrow{E}$$ Si $\underline{\overrightarrow{E}}$ est de la forme $e^{i(\omega t-kz)}$, alors on a le relation de dispersion $$\omega^2-\omega_p^2=\underline{k}^2c^2$$ On dispose de la relation de structure suivante $\underline{k}\overrightarrow{u}\wedge\underline{\overrightarrow{E}}=\omega\underline{\overrightarrow{B}}$ avec $\underline{k}$ un imaginaire pur. $\underline{E_x}$ et $\underline{B_y}$ sont en quadrature.

\chapter{Propagation d'ondes électromagnétiques dans un conducteur ohmique ; réflexion, guides d'ondes}
\localtableofcontents
\section{Propagation d'ondes électromagnétiques dans un milieu ohmique}
\subsection{Position du problème ; discussion des approximations}
\paragraph{Trois approximations} Dans le conducteur, on suppose :
\begin{itemize}
	\item l'électroneutralité : $\rho=0$. Cela est une bonne approximation si $\tau\omega\ll 1$ ;
	\item la loi d'Ohm : $\overrightarrow{j}=\gamma\overrightarrow{E}$. Cela est une bonne approximation si $\tau\omega\ll 1$ ;
	\item le courant de déplacement négligeable : $\displaystyle\overrightarrow{j}\gg\varepsilon_0\frac{\partial\overrightarrow{E}}{\partial t}$
\end{itemize}
\paragraph{Exemple}Pour les ondes radio ou wifi, on peut utiliser la loi d'Ohm statique.
\paragraph{Résumé} Ainsi dans cette section, on a : 
\begin{itemize}
	\item $\div\overrightarrow{E}=0$
	\item $\displaystyle\rot\underline{\overrightarrow{E}}=-\frac{\partial\underline{\overrightarrow{B}}}{\partial t}$
	\item $\div\underline{\overrightarrow{B}}=0$
	\item $\rot\underline{\overrightarrow{B}}=\mu_0\underline{\overrightarrow{j}}$
	\item $\underline{\overrightarrow{j}}=\gamma_0\underline{\overrightarrow{E}}$
\end{itemize}
\subsection{Équation de propagation ; relation de dispersion}
\paragraph{Équation de propagation} Avec ces trois approximations, on obtient l'équation de propagation en calculant $\rot\rot\overrightarrow{E}$ de deux façons différentes avec la formule ainsi que les équations de Maxwell-Faraday et de Maxwell-Thomson : $$\overrightarrow{\Delta}\overrightarrow{E}=\mu_0\gamma_0\frac{\partial\overrightarrow{E}}{\partial t}$$ Cette équation est linéaire, mais contrairement à l'équation de d'Alembert, elle n'est pas stable par renversement du temps : elle est donc irréversible, ce qui est bien en accord avec l'effet Joule dans le milieu.
\paragraph{Relation de dispersion}En injectant une OPPH dans l'équation de propagation, on en tire l'équation de dispersion : $$\underline{k}^2=-i\gamma\mu_0\omega$$
\subsection{Forme des champs ; effet de peau}
\paragraph{Épaisseur de peau}En posant $\displaystyle\delta=\sqrt{\frac{2}{\gamma\mu_0\omega}}$, on obtient $$\displaystyle\underline{k}=\pm\frac{1-i}{\delta}$$$\delta$ est l'\textbf{épaisseur de peau} : elle correspond à la distance sur laquelle le signal est non nul dans le conducteur.
\paragraph{Calcul de $\overrightarrow{E}$}Si on se place dans le demi-espace $z>0$ (on enlève la solution divergente en $+z/\delta$), on a pour une onde polarisée selon $\overrightarrow{u_x}$ : $$\underline{\overrightarrow{E}}=\underline{E_0}\overset{\text{atténuation}}{\overbrace{e^{-\frac{z}{\delta}}}}\underset{\text{propagation}}{\underbrace{e^{i\left(\omega t-\frac{z}{\delta}\right)}}}\overrightarrow{u_x}$$
\paragraph{Calcul de $\overrightarrow{B}$} On peut calculer $\overrightarrow{B}$ à l'aide de l'équation de Maxwell-Faraday en complexes. On trouve que $\overrightarrow{B}$ et $\overrightarrow{E}$ sont déphasés de $\displaystyle\frac{\pi}{2}$ : les relations de structure dans le vide ne sont pas valables.
\paragraph{Effet de peau} Au-delà de quelques $\delta$, les champs $\overrightarrow{E}$ et $\overrightarrow{B}$ sont nuls dans le conducteur : c'est l'\textbf{effet de peau}.
\paragraph{Ordres de grandeur} Pour le cuivre :
\begin{itemize}
	\item pour $f=1$ GHz, on obtient $\delta=2\ \mu$m et $\lambda_0=30$ cm ;
	\item pour $f=10$ MHz, on obtient $\delta=20\ \mu$m et $\lambda_0=30$ m ;
	\item pour $f=10$ Hz, on obtient $\delta=2$ cm et $\lambda_0=3.10^7$ m
\end{itemize}
\subsection{Bilan énergétique}
\paragraph{Bilan énergétique} On se limite aux grandeurs moyennes dans le temps et on a : $$\underset{=0}{\underbrace{\left\langle\frac{\partial u_{em}}{\partial t}\right\rangle}}+\div\langle\overrightarrow{\pi}\rangle=-\left\langle\overrightarrow{j}\cdot\overrightarrow{E}\right\rangle$$ En effet, $\overrightarrow{E}$ et $\overrightarrow{B}$ son périodiques, donc $u_{em}$ aussi. On peut vérifier cette relation en calculant $\langle\overrightarrow{\pi}\rangle$ et $\left\langle\overrightarrow{j}\cdot\overrightarrow{E}\right\rangle$ grâce à la formule $$\left\langle f\times g\right\rangle=\frac12\text{Re}\left(\underline{f}\times\underline{g}\st\right)$$

\section{Réflexion d'une onde sur un conducteur parfait}
\paragraph{Méthode} L’onde $\overrightarrow{E_i}$ arrivant sur un conducteur donne lieu à une onde transmise $\overrightarrow{E_t}$ et une onde réfléchie $\overrightarrow{E_r}$ (loi de Descartes). On se limite à une OPPM linéairement polarisée (les relations sont linéaires). On admet que les ondes transmises et réfléchies sont planes progressives et monochromatiques de même pulsation (l’onde agite les électrons à $\omega$ qui émettent alors une onde à $\omega$).
\paragraph{Objectifs} On souhaite avoir : 
\begin{itemize}
	\item $\delta\ll\lambda_0$ ;
	\item $\delta$ négligeable devant l'épaisseur du conducteur ;
	\item $\delta$ négligeable devant le rayon de courbure du conducteur.
\end{itemize}
\subsection{Cas limite du conducteur parfait}
\paragraph{Modèle du conducteur parfait} On se limite à une onde arrivant normalement sur un conducteur parfait ($\gamma\to\infty$). Les symétries donnent que les ondes transmises et réfléchies sont polarisées linéairement selon $x$ et invariantes par translation selon $x$ et $y$. On les écrit alors : $$\underline{\overrightarrow{E_r}}=\underline{r}\underline{E_0}e^{i(\omega t+k_1z)}\overrightarrow{u_x}$$ $$\underline{\overrightarrow{E_t}}=\underline{t}\underline{E_0}e^{i(\omega t-k_2z)}\overrightarrow{u_x}$$ et $k_2$ est donné par la relation de dispersion dans la milieu 2.
\paragraph{Calcul de $\overrightarrow{E}$ et $\overline{B}$ dans un conducteur parfait} On utilise un argument énergétique. En effet, la puissance dissipée par effet Joule par unité de volume dans le conducteur est $p_v=\overrightarrow{j}\cdot\overrightarrow{E}=\gamma E^2$. Or, $\gamma\to +\infty$ et comme la puissance dissipée par effet Joule est bornée, on a nécessairement $E\to0$. L'équation de Maxwell-Faraday nous permet d'en conclure que $\overrightarrow{B}$ est un indépendant du temps. Ainsi, dans un conducteur parfait :
\begin{itemize}
	\item $\overrightarrow{E}=\overrightarrow{0}$ ;
	\item $\overrightarrow{B}$ est stationnaire (au sens d'indépendant du temps) ;
	\item $\overrightarrow{B}_\text{onde}=\overrightarrow{0}$
\end{itemize}
\subsection{Existence d'une onde réfléchie} On considère une onde incidente $$\underline{\overrightarrow{E_i}}=\underline{\overrightarrow{E_{0,i}}}e^{i(\omega t-kz)}$$ dans le vide qui arrive à incidence normale sur un plan $z=0$ d'un conducteur parfait. 
\paragraph{Première présentation}La relation de passage en $z=0$ fournit alors $$\underline{\overrightarrow{E_\text{tot}}}(0^+)-\underline{\overrightarrow{E_\text{tot}}}(0^-)=\frac{\underline{\sigma}}{\varepsilon_0}\overrightarrow{u_z}$$ Or, puisqu'on a un conducteur parfait, $\underline{\overrightarrow{E_\text{tot}}}(0^+)=\overrightarrow{0}$. Et s'il n'y a pas d'onde réfléchie, le terme $\underline{\overrightarrow{E_\text{tot}}}(0^-)$ vaut $\underline{\overrightarrow{E_\text{i}}}(0^-)$ qui est orthogonal à $\overrightarrow{u_z}$, donc tout est nul, ce qui est absurde. Ainsi, \textbf{il existe une onde réfléchie}.
\paragraph{Autre présentation} Le champ électrique de l'onde incidente met les électrons dans le métal en mouvement. Ce mouvement d'électrons crée alors une onde électromagnétique.
\subsection{Structure de l'onde résultante}
\paragraph{Calcul de $\underline{r}$ et $\underline{t}$ dans le conducteur parfait} On cherche $\underline{r}$ et $\underline{t}$ dans le modèle du conducteur parfait. On a : 
\begin{enumerate}
	\item $\overrightarrow{E}=\overrightarrow{0}$ car $\overrightarrow{j}\cdot\overrightarrow{E}=\gamma E^2$ est borné alors que $\gamma\to\infty$ ;
	\item $\overrightarrow{B}=\overrightarrow{0}$ car $\overrightarrow{B}$ ne dépend pas du temps d'après l'équation de Maxwell-Faraday et le premier point, et on ne s'intéresse qu'à la partie variable ;
	\item $\rho=0$ d'après l'équation de Maxwell-Gauss ;
	\item $\overrightarrow{j}=\overrightarrow{0}$ d'après les premier et deuxième points.
\end{enumerate} 
Toutefois, il peut y avoir des densités surfaciques de charges et de courants. On utilise donc la relation de passage : $$\overrightarrow{E}(0^+)-\overrightarrow{E}(0^-)=\frac{\sigma}{\varepsilon_0}\overrightarrow{u_z}$$ Dans le conducteur, il n'y a pas d'onde donc $\underline{t}=0$. En projetant la relation de passage sur $\overrightarrow{u_x}$, on en déduit que $\underline{r}=-1$.
\paragraph{Calcul du champ électrique total} On obtient le champ total en sommant les deux. On trouve une onde stationnaire pour $\overrightarrow{E}$ lorsque $z<0$ : $$\underline{\overrightarrow{E_\text{tot}}}=-2i\underline{E_0}e^{i\omega t}\sin(kz)\overrightarrow{u_x}$$
\paragraph{Remarque} $\underline{\overrightarrow{E_\text{tot}}}(z=0)=\overrightarrow{0}$ donc on est content d'avoir un sinus.
\paragraph{Attention} Cette onde n'est pas prgressive. On ne peut pas utiliser la relation de structure pour $\underline{\overrightarrow{E_\text{tot}}}$
\paragraph{Calcul du champ magnétique total} On utilise les relations de structure sur l'onde incidente et l'onde réfléchie, puis on somme les champ $\overrightarrow{B}$ pour obtenir : $$c\underline{\overrightarrow{B_\text{tot}}}=2\underline{E_0}e^{i\omega t}\cos(kz)\overrightarrow{u_y}$$
\paragraph{Expressions en notation réelle} En prenant $\underline{E_0}=E_0e^{i\varphi}$, on obtient $$\overrightarrow{E_\text{tot}}=2E_0\sin(kz)\sin(\omega t+\varphi)\overrightarrow{u_x}$$
$$\overrightarrow{N_\text{tot}}= \frac{2E_0}{c}\cos(kz)\cos(\omega t+\varphi)\overrightarrow{u_y}$$ Ce sont des ondes \textbf{stationnaires}. Les nœuds de $\overrightarrow{E_\text{tot}}$ sont les ventres de $\overrightarrow{B_\text{tot}}$ et réciproquement.
\paragraph{Remarque} Pour une onde stationnaire, on a $\left\{\overrightarrow{\pi}\right\}=\overrightarrow{0}$
\subsection{Charges et courants surfaciques ; pression de radiation}
\paragraph{Charge surfacique} Pour une onde à incidence normale, la relation de passage pour $\underline{\overrightarrow{E}}$ donne $\sigma=0$ (voir partie précédente).
\paragraph{Courant surfacique} Déterminons désormais $\underline{j_s}$. On applique la relation de passage pour $\underline{\overrightarrow{B}}$ en $z=0$ : $$\underbrace{\underline{\overrightarrow{B_\text{tot}}}(z=0^+)}_{=\overrightarrow{0}}-\underline{\overrightarrow{B_\text{tot}}}(z=0^-)=\mu_0\underline{\overrightarrow{j_s}}\wedge\overrightarrow{u_z}$$
Par conséquent, comme le courant surfacique ne présente pas de composante selon $\overrightarrow{u_z}$ (car il est tangent à la surface), on a : $$\underline{\overrightarrow{j_s}}=\frac{2E_0}{\mu_0 c}e^{i\omega t}\overrightarrow{u_x}$$
\paragraph{Pression de radiation} La pression de radiation est la force exercée sur le conducteur par les photons réfléchis par	unité de surface. On a envie d'écrire $\delta\overrightarrow{F}=\overrightarrow{j_s}dS\wedge\overrightarrow{B}$. Mais quel sens donner à $\overrightarrow{B}$ ? En effet, $\overrightarrow{B_\text{tot}}(0)$ n'a pas de sens.
\paragraph{Première présentation : vision particulaire} L'onde émise par les électrons de l'interface n'exerce aucune force sur l'interface puisque l'émission est symétrique entre les $z<0$ et les $z>0$. Ainsi, on a : $$\delta\overrightarrow{F}=\overrightarrow{j_s}dS\wedge\underbrace{\overrightarrow{B_i}(0)}_{\text{a un sens}}$$ Alors, en prenant la moyenne temporelle avec la valeur moyenne d'un $\cos^2$ qui vaut $1/2$, on obtient $$\langle\delta\overrightarrow{F}\rangle=\frac{2}{\mu_0 c^2}E_0^2\frac{1}{2}dS\overrightarrow{u_z}$$ Par conséquent, on obtient : $$\langle\delta\overrightarrow{F}_\text{sur le conducteur}\rangle=dS\varepsilon_0E_0^2\overrightarrow{u_z}$$ On définit alors la pression de radiation par $$P_\text{radiation}=\varepsilon_0E_0^2$$
\paragraph{Deuxième présentation : vision corpusculaire} On prend pour système les photons qui se réfléchissent sur $dS$ entre $t$ et $t + dt$. On calcule la variation de leur quantité de mouvement et on en déduit la force qu’ils ont subie, \textit{ie}	l’opposé de la force qu’ils ont exercé. (troisième loi de Newton). Pour un photon dans le vide, les relations de De Broglie et de Planck-Einstein fournissent $p=\displaystyle\frac{E}{c}$. Or, on a $$\overrightarrow{p}_{t+dt}-\overrightarrow{p}_t=\langle\delta\overrightarrow{F}_{\text{métal}\to\text{lumière}}dt$$ Les quantités de mouvement sont $$\begin{cases}\overrightarrow{p}_t=ndScdt\frac{E}{c}\overrightarrow{u_z}\\\overrightarrow{p}_{t+dt}=-ndScdt\frac{E}{c}\overrightarrow{u_z}\end{cases}$$ Ainsi, on obtient $\langle\delta\overrightarrow{F}_{\text{métal}\to\text{lumière}}=-2nEdS\overrightarrow{u_x}$ Par conséquent, la troisième loi de Newton fournit alors $$P_\text{radiation}=2nE$$ Or, $nE=\langle u_{em,i}\rangle$ et pour l'onde incidente on a $\langle u_{em,i}\rangle=\frac{1}{2}\varepsilon_0E_0^2$. Finalement, on retrouve bien : $$P_\text{radiation}=\varepsilon_0E_0^2$$
\subsection{Réflexion sous incidence oblique sur un conducteur parfait}
On considère une onde incidente qui arrive sur un conducteur parfait avec un angle d'incidence $i$ : $$\underline{\overrightarrow{E_\text{i}}}=\underline{E_0}e^{i(\omega t-k(x\sin(i)+z\cos(i)))}\overrightarrow{u_y}$$ On cherche l'onde réfléchie sous la forme $$\underline{\overrightarrow{E_r}}=\underline{\overrightarrow{E_{0,r}}}e^{i(\omega't-\overrightarrow{k'}\cdot\overrightarrow{r})}$$ où $\underline{\overrightarrow{E_{0,r}}}$ est selon $\overrightarrow{u_y}$ par symétrie. La relation de passage à l'interface s'écrit $$\underline{\overrightarrow{E_\text{tot}}}(z=0^+)-\underline{\overrightarrow{E_\text{tot}}}(z=0^-)=\frac{\underline{\sigma}}{\varepsilon_0}\overrightarrow{u_z}$$ soit ici avec nos expressions $$-\left(\underline{\overrightarrow{E_i}}(0)+\underline{\overrightarrow{E_r}}(0)\right)=\frac{\underline{\sigma}}{\varepsilon_0}\overrightarrow{u_z}$$ Après projection sur $\overrightarrow{u_y}$, on obtient : $$\underline{E_0}e^{i(\omega t-k(x\sin(i)+z\cos(i)))}=-\underline{E_{0,r}}e^{i(\omega't-(k_x'x+k_y'y))}$$
\paragraph{Calcul du coefficient de réflexion} En $x=y=0$ et en $t=0$, la formule précédente donne $\underline{E_0}=-\underline{E_{0,r}}$, soit $\underline{r}=-1$.
\paragraph{Calcul de $\omega'$} En dérivant l'expression précédente par rapport au temps puis en évaluant de même pour $x=y=0$ et $t=0$, on obtient : $$\omega'=\omega$$
\paragraph{Calcul de $k_x'$ et $k_y'$} De même, en dérivant par rapport à $x$ et $y$ puis en évaluant en $x=y=0$ et $t=0$, on obtient : $$\begin{cases}
	k_x'=k\sin(i)\\
	k_y'=0
\end{cases}$$
\paragraph{Calcul de $k_z'$} Les relations de dispersion dans le vide pour l'onde incidente et l'onde réfléchie ainsi que l'égalité $\omega'=\omega$ fournissent : $$k^2=k'^2$$ On obtient alors $k_z'=-k\cos(i)$. En effet, la solution en $+k\cos(i)$ correspond à $\underline{E_\text{tot}}=\overrightarrow{0}$, ce qui est assez peu passionnant (pour citer le goat).
\paragraph{Calcul de l'onde totale} En sommant et avec les formules d'Euler, on obtient l'onde totale : $$\underline{\overrightarrow{E_\text{tot}}}=-2iE_0sin(kz\cos(i))\underline{E_0}e^{i(\omega t-kxsin(i))}\overrightarrow{u_y}$$ D'où l'expression finale de l'onde totale réelle : $$\overrightarrow{E_\text{tot}}=2E_0\sin(kz\cos(i))\sin\left(\omega t-kx\sin(i)+\varphi\right)$$ Cette onde est :
\begin{itemize}
	\item harmonique
	\item progressive de vitesse $\displaystyle\frac{\omega}{k\sin(i)}\overrightarrow{u_x}$ ;
	\item stationnaire selon $z$ à $x$ fixé ;
	\item polarisée rectilignement selon $\overrightarrow{u_y}$.
\end{itemize}
\subsection{Méthode générale : coefficients de réflexion et de transmission en amplitude et en puissance}
Voici les étapes à suivre pour un exercice qui traite des coefficients de réflexion et de transmission : 
\begin{enumerate}
	\item Lire le sujet pour avoir la forme des champs.
	\item Utiliser la relation de passage pour $\underline{\overrightarrow{E}}$.
	\item Si cela ne suffit pas, calculer $\underline{\overrightarrow{B_\text{inc}}}$, $\underline{\overrightarrow{B_\text{trans}}}$ et $\underline{\overrightarrow{B_\text{tot}}}$ à l'aide des relations de structure dans les différents milieux (au besoin utiliser Maxwell-Faraday).
	\item Utiliser la relation de passage pour $\underline{\overrightarrow{B}}$.
\end{enumerate}
On rencontrera parfois aussi les coefficients de réflexion et de transmission en puissance $R$ et $T$, voire le coefficient d'absorption en puissance $A$. On doit toujours avoir $R+A+T=1$ à la fin.
\section{Cavité résonante ; propagation guidée}
\subsection{Cavité 1D, modes propres, forme des champs, aspects énergétique}
\paragraph{Cavité} On appelle cavité une portion vide de l’espace délimitée par des conducteurs parfaits. 
\paragraph{Cas de la cavité à 1D}On cherche une solution sous la forme $\overrightarrow{E}=f(z)g(t)\overrightarrow{u_x}$ et on en déduit que dans la cavité $$\frac{f''(z)}{f(z)}=\frac{g''(t)}{c^2g(t)}=A= \text{cte}$$ L’onde respectant les conditions aux limites $f(0)=f(L)=0$ (relations de passage en $z=0$ et $z=L$ projetée sur $\overrightarrow{u_x}$), elle est non-nulle si $A<0$. Dans ce cas, c’est alors une sinusoïde quantifiée : $$\overrightarrow{E}=\overrightarrow{E_0}\sin\left(\frac{n\pi}{L}x\right)\cos\left(\frac{n\pi c}{L}t+\varphi\right)$$ Il existe des solutions non stationnaires, mais ce sont alors des sommes des modes propres précédents.
\paragraph{Calcul de $\overrightarrow{B}$} Dans ce cas, d'après l'équation de Maxwell-Faraday, on a : $$\rot\overrightarrow{E}=-\frac{\partial\overrightarrow{B}}{\partial t}$$ En projection sur $\overrightarrow{u_y}$, cela donne $$-kE_0\cos(\omega t+\varphi)\cos(kz)=-\frac{\partial B_y}{\partial t}$$ Après primitivation et avec la relation de dispersion $\omega=kc$, on obtient finalement : $$\overrightarrow{B}=-\frac{E_0}{c}\sin(\omega t+\varphi)\cos(kz)\overrightarrow{u_y}$$
\subsection{Propagation guidée ; hyperfréquence entre deux plans conducteurs}
\paragraph{Propagation guidée} Cette fois, on s'intéresse à la propagation d'une onde selon $\overrightarrow{u_z}$ entre deux conducteurs parfaits entre $x<0$ et $x>L$ : $$\overrightarrow{E}=f(x)e^{i(\omega t-k_z z)}\overrightarrow{u_y}$$ En effet, il peut y avoir une discontinuité du champ selon $\overrightarrow{u_x}$ donc il n'y a pas de conditions aux limites sur cette polarité ; et on choisit une invariance selon $\overrightarrow{u_y}$. L'équation de d'Alembert fournit alors une condition sur $f$ : $$f''(x)=-\underbrace{\left(\frac{\omega^2}{c^2}-k_z^2\right)}_{=A}f(x)$$ Les conditions aux limites sont $f(0)=f(L)=0$ (avec la relation de passage en $x=0$ et $x=L$ projetée sur $\overrightarrow{u_y}$). Si $A>0$, $f$ est non nulle et $A$ est quantifiée : $$\omega^2-\omega_{c,n}^2=c^2k_z^2 \ \text{avec} \ \omega_{c,n}=\frac{n\pi c}{L} \ \text{où} \ n\in\N\st$$ La relation de dispersion est semblable à celle du plasma. On fixe $n$ expérimentalement : le système est alors un passe-bas de pulsation de coupure $\omega_{c,n}$. Le calcul de $\overrightarrow{B}$ montre que $\overrightarrow{B}$ n'est pas transverse car $\overrightarrow{B}\cdot\overrightarrow{u_z}\ne0$. Enfin, on peut calculer les vitesses de propagation $v_g$ et $v_\varphi$.
\paragraph{Remarque : "vrai" guide d'onde} Un vrai guide d'onde est plutôt un parallélépipède rectangle. Si on le prend de longueur $b$ selon $y$, de hauteur $a$ selon $x$ et de profondeur selon $z$, on a une expression générale de $\overrightarrow{E}$ de la forme $$\overrightarrow{E}=E_0\sin\left(\frac{n\pi}{a}x\right)\sin\left(\frac{n\pi}{b}y\right)\cos\left(\omega t-k_zz+\varphi\right)\overrightarrow{u_x}$$ ou avec un $\overrightarrow{u_y}$ à la place du $\overrightarrow{u_x}$.
\paragraph{Remarque : base des modes propres sur l'exemple de la cavité} Il n'existe pas d'autres modes propres que ceux que l'on a trouvés. Une solution du problème (EDP et CL) n'est pas nécessairement un mode propre, nous les avons simplement cherchés ici sous la forme d'un mode propre. On admet que toute solution est une superposition des modes propres trouvés précédemment (il s'agit d'une décomposition de Fourier, admise).
\chapter{Rayonnement dipolaire électrique}
\localtableofcontents
\section{Les trois échelles de longueur pertinentes}
\subsection{Modèle du dipôle oscillant}
\paragraph{Trois modèles} Pour créer un moment dipolaire oscillant, on dispose de trois modèles :
\begin{itemize}
	\item Deux charges identiques constantes mobiles : $-q$ en $x=0$ et $+q$ en $x=d\cos(\omega t)$. On obtient un moment dipolaire oscillant $\overrightarrow{p}=qd\cos(\omega t)\overrightarrow{u_x}$ ;
	\item Deux charges variables fixes : $Q(t)=q\cos(\omega t)$ en $x=d$ et $-Q(t)$ en $x=0$. On obtient un moment dipolaire oscillant $\overrightarrow{p}=qd\cos(\omega t)\overrightarrow{u_x}$ ;
	\item Un petit élément de courant avec $i=\displaystyle\frac{dQ}{dt}$ : on a alors $\displaystyle\frac{d}{dt}\overrightarrow{p}=Idz\overrightarrow{u_z}$
\end{itemize}
\paragraph{Applications des modèles} Les atomes et les antennes.
\paragraph{Trois échelles de longueur} On note $l$ la taille du dipôle, $r$ la distance entre l'observateur et le dipôle, et $\lambda$ la longueur d'onde correspondant à l'oscillation.
\subsection{Approximation dipolaire}
Dans la suite du chapitre, on se placera toujours dans l'approximation dipolaire.
\paragraph{Approximations dipolaire} On parle de l'approximation \textbf{dipolaire} lorsque $d\ll r$.
\paragraph{Exemple de validité de l'approximation dipolaire}Elle est valable pour les atomes et les antennes.
\subsection{Approximation non relativiste}
Dans la suite du chapitre, on se placera toujours dans l'approximation non relativiste.
\paragraph{Approximation non relativiste} On parle de l'approximation \textbf{non relativiste} lorsque $l\ll\lambda$.
\paragraph{Exemple de validité (ou non) de l'approximation non relativiste}Elle est valable pour les atomes (car $10^{-10}\ll 10^{-7}$) mais n'est pas valable pour les antennes : il faut alors considérer des petits éléments d'antenne et intégrer.
\paragraph{Explication de l'origine du nom "non relativiste"}Cette approximation provient de l'approximation $v\ll c$ dans laquelle on prend $c=lf$ et $c=\lambda f$.
\paragraph{Remarque} L'approximation non relativiste est toujours respectée dans le cas statique.
\paragraph{Ordres de grandeur} Pour les ondes lumineuses et les atomes, $l\approx10^{-10}$ m et $\lambda\approx100$ nm donc on respecte bien l'approximation non relativiste. Pour les antennes radio/wifi, on a $f\approx1$ GHz et $\lambda=30$ cm. $l$ est de l'ordre du cm donc l'approximation non relativiste n'est pas valable.
\paragraph{Cas des antennes} Pour traiter le cas des antennes, il faut considérer une portion infinitésimale d'antenne dz, établir des relations sur cette portion puis intégrer ces relations.
\subsection{Zones d'observation}
\paragraph{Trois zones d'observation} On distingue trois zones d'observation :
\begin{itemize}
	\item \textbf{quasi statique} lorsque $l\ll r\ll \lambda$. On retrouve alors les résultats des chapitres précédents.
	\item cas intermédiaire lorsque $l\ll r\approx\lambda$
	\item \textbf{rayonnement} lorsque $l\ll\lambda\ll r$
\end{itemize}
\section{Champ rayonné}
\paragraph{Champs rayonnés} En coordonnées sphériques pour un dipôle en $O$ orienté selon $\overrightarrow{u_z}$, on a : $$\overrightarrow{E}(M,t)=\frac{1}{4\pi\varepsilon_0 r^3}\left[\left(p(t')+\frac{r}{c}\dot{p}(t')\right)2\cos(\theta)\overrightarrow{u_r}+\left(p(t')+\frac{r}{c}\dot{p}(t')+\frac{r^2}{c^2}\ddot{p}(t')\right)\sin(\theta)\overrightarrow{u_\theta}\right]$$ où $t'=\displaystyle t-\frac{r}{c}$ est le temps de retard ou temps de propagation. Dans la limite statique, les dérivées s'annulent et on retrouve la formule connue pour le dipôle électrostatique. Dans la limite $\lambda\ll r$, cela devient : $$\overrightarrow{E}(M,t)=\frac{1}{4\pi\varepsilon_0 r^3}\frac{r^2}{c^2}\ddot{p}(t')\sin(\theta)\overrightarrow{u_\theta}$$
\subsection{Expression dans la zone de rayonnement}
\paragraph{Champs rayonnés dans la limite $\lambda\ll r$} Dans la limite $\lambda\ll r$, les expressions des champs sont alors : $$\overrightarrow{E}(M,t)=\frac{\mu_0\ddot{p}(t')}{4\pi r}\sin(\theta)\overrightarrow{u_\theta}=c\overrightarrow{B}\wedge\overrightarrow{u_r}$$ $$\overrightarrow{B}(M,t)=\frac{\mu_0\ddot{p}(t')}{4\pi rc}\sin(\theta)\overrightarrow{u_\varphi}=\frac{\mu_0}{4\pi rc}\overrightarrow{\ddot{p}}(t')\wedge\overrightarrow{u_r}$$ où $\displaystyle t'=t-\frac{r}{c}$ est le temps de retard, ou temps de propagation (voire temps retardé).
\subsection{Invariances et symétries}
\paragraph{Symétries} Quel que soit le point M, le plan $(M,\overrightarrow{u_r},\overrightarrow{u_\theta})$ est un plan de symétrie du dipôle (donc des courants et des charges) et par conséquent $\overrightarrow{B}$ est selon $\overrightarrow{u_\varphi}$ et que $\overrightarrow{E}$ a une composante nulle selon $\overrightarrow{u_\varphi}$. Cela ne permet pas de prouver que $\overrightarrow{E}$ est selon $\overrightarrow{u_\theta}$, on ne peut pas le prouver et on l'admet.
\paragraph{Invariances} La distribution de charges et de courants est invariante par rotation autour de l'axe $(Oz)$ donc : $$\begin{cases}
	E_\theta=E_\theta(r,\theta,\xcancel{\varphi}) \\
	B_\varphi=B_\varphi(r,\theta,\xcancel{\varphi})
\end{cases}$$
\paragraph{Dépendance en $\theta$}La dépendance en $\sin(\theta)$ se retient avec les symétries : par exemple, les champs doivent être nuls sur l'axe.
\subsection{Temps de propagation}
\paragraph{Temps de propagation} Le temps de propagation $$t'=t-\frac{r}{c}$$ correspond à la date d'émission. Le terme $\displaystyle\frac{r}{c}$ correspond au temps de propagation de l'onde.
\paragraph{Cas d'un moment dipolaire sinusoïdal} Pour un moment dipolaire $p=p_0\cos(\omega t)$, on a : $$\ddot{p}(t')=-\omega^2\cos\left(\omega t-\omega\frac{r}{c}\right)$$ On retrouve donc la relation de dispersion des OPPH dans le vide $\omega=kc$, ce qui n'était pas évident au premier abord.
\subsection{Structure locale d'onde plane progressive}
\paragraph{Justification qualitative} Loin du dipôle, les ondes sphériques doivent paraître planes ($\overrightarrow{u_r}\approx\overrightarrow{\text{cte}}$ et $1/r\approx\text{cte}$).
\paragraph{Relations de structure} On a les relations de structure : $$\begin{cases}
	\overrightarrow{u_r}\wedge\overrightarrow{E}=c\overrightarrow{B} \\
	\overrightarrow{u_r}\cdot\overrightarrow{E}=0 \\
	c\overrightarrow{B}\wedge\overrightarrow{u_r}=\overrightarrow{E}
\end{cases}$$ Ce sont les mêmes relations de structure que pour les OPPH de direction de propagation $\overrightarrow{u_r}$. Cependant, cette relation n'est que locale puisque $\overrightarrow{u_r}$ peut changer en fonction de la position.
\paragraph{Application} Ces relations de structure permettent de calculer $\overrightarrow{B}$ à partir de $\overrightarrow{E}$ et réciproquement, sans avoir à passer par la formule admise.
\subsection{Polarisation}
\paragraph{Polarisations} L'onde est polarisée rectilignement. Précisément, $\overrightarrow{E}$ est polarisé rectilignement selon $\overrightarrow{u_\theta}$
et $\overrightarrow{B}$ est polarisé rectilignement selon $\overrightarrow{u_\varphi}$.
\paragraph{Attention} Cela signifie bien que, pour une \textbf{position fixée}, $\overrightarrow{E}$ et $\overrightarrow{B}$ gardent des directions constantes au cours du temps.
\section{Puissance rayonnée}
\subsection{Décroissance de l'amplitude}
\paragraph{Attention} Cela n'est pas de l'absorption.
\paragraph{Calcul de la puissance moyenne} On calcule tout d'abord le vecteur de Poynting avec les expressions admises, ce qui donne : $$\overrightarrow{\pi}=\frac{\mu_0}{c}\left(\frac{\sin(\theta)}{4\pi r}\ddot{p}(t')\right)^2\overrightarrow{u_r}$$ La puissance moyenne sortant de la sphère de rayon $r$ et de centre $O$ est alors : $$P_\text{moy}=\frac{\mu_0}{c(4\pi)^2}\langle(\ddot{p})^2\rangle\iint\frac{\sin^2(\theta)}{r^2}rd\theta r\sin(\theta)d\varphi$$ soit en simplifiant : $$P_\text{moy}=\frac{\mu_0}{8\pi c}\langle(\ddot{p})^2\rangle\int_{0}^{2\pi}\sin^3(\theta)d\theta$$
Or, en linéarisant le sinus, on trouve $\sin^3(\theta)=-\frac{1}{4}\sin(3\theta)+\frac{3}{4}\sin(\theta)$ et l'intégrale cherchée vaut alors $\frac{4}{3}$. On a donc $$P_\text{moy}=\frac{\mu_0}{6\pi c}\langle(\ddot{p})^2\rangle$$
\paragraph{Dépendance en $r$} La dépendance en $r$ s'explique par le fait que l'onde ne perd pas d'énergie entre deux sphères puisqu'il n'y a pas d'absorption. Par conséquent, le flux du vecteur de Poynting vérifie $\phi_r=\phi_{r+dr}$ en moyenne. En intégrant cette relation, on trouve : $$\phi_r=\frac{\omega^4\mu_0 p^2}{12\pi c}$$ $\omega^4$ explique la diffusion de Rayleigh : il y a 16 fois plus de diffusion dans le bleu-violet que dans le rouge.
\subsection{Dépendance en $1/\lambda^4$}
Comprendre ici une dépendance en $\omega^4$.
\paragraph{Diffusion de Rayleigh}
On a : $$\phi_r=\frac{\omega^4\mu_0 p^2}{12\pi c}$$ $\omega^4$ explique la diffusion de Rayleigh : il y a 16 fois plus de diffusion dans le bleu-violet que dans le rouge.
\paragraph{Conséquence : le ciel est bleu}
Les centres diffuseurs de l'atmosphère sont O$_2$, N$_2$, etc. sont de taille $l\approx1$ nm donc $l\ll\lambda\ll r$ avec $\lambda\approx600$ np et $r\approx1$ km. Par conséquent, les courtes longueurs d'onde ($\lambda_V=400$ nm) sont $16=2^4$ fois mieux déviées que les grandes ($\lambda_R=800$ nm). Ainsi le spectre du Soleil s'équilibre autour du bleu. On pourrait se demander si $p_0$ est indépendant de $\omega$. On aura pour cela besoin d'introduire le modèle de l'électron élastiquement lié. Tout cela est détaillé dans la dernière partie.
\subsection{Anisotropie ; dépendance en $\theta$ ; indicatrice de rayonnement}
\paragraph{Indicatrice de rayonnement} Elle permet de caractériser quelles directions reçoivent bien le rayonnement du dipôle. Un dipôle oscillant \textbf{de direction fixe} n'émet pas selon son axe (\textit{cf} polarisation de la lumière diffusée par l'atome).
\subsection{Puissance totale}
On a la formule : $$P_\text{moyenne}=\frac{1}{12\pi\varepsilon_0}\frac{p_0^2\omega^4}{c^3}$$
\subsection{Formule de Larmor ; rayonnement synchrotron ; rayonnement de freinage}
\paragraph{Rayonnement synchrotron} Une particule chargée accélérée rayonne comme un dipôle (enlever la particule fixe revient à soustraire un champ en $-1/r^2$). Par conséquent, pour une particule de charge $q$ et d’accélération $a$, on a $p_0\omega^2=dq\omega^2=qa$ car $a = d\omega^2$ pour un mouvement sinusoïdal. On obtient donc la \textbf{formule de Larmor} qui donne la puissance émise : $$P=\frac{q^2}{6\pi c^3}\langle a^2\rangle$$
\paragraph{Remarque} Le rayonnement synchrotron est utilisé dans le Soleil du plateau de Saclay.
\paragraph{Si le dipôle n'a pas de direction fixe} On a alors la relation de structure suivante, où $\overrightarrow{u_r}$ est le vecteur position : $$\overrightarrow{E}=c\overrightarrow{B}\wedge\overrightarrow{u_r}$$ On dispose alors aussi de l'expression suivante pour $\overrightarrow{B}$ : $$\overrightarrow{B}=\frac{\mu_0}{4\pi cr}\frac{d^2\overrightarrow{p}}{dt^2}(t')\wedge\overrightarrow{u_r}$$ 
\section{Approche documentaire : diffusion Rayleigh}
\subsection{Modèle de l'électron élastiquement lié}
\paragraph{Description et validité du modèle}Le modèle de l'électron élastiquement lié provient initialement de la modélisation de Thomson de l'atome, proposée en 1904 suite à sa découvert de l'électron en 1897. Celui-ci proposa que les électrons étaient plongés dans une boule de rayon $R$ (de l'ordre de $10^{-10}$ m), uniformément chargée de charges positives qui équilibrent la charge négative des électrons. Ce modèle fut invalidé en 1911 lorsque Rutherford découvrit, grâce à ses expériences de diffusion de particules $\alpha$ à travers une très fine feuille d'or, l'existence du noyau atomique, dont la taille est très petite devant le rayon des orbites électroniques.
\paragraph{Inventaire des forces}Selon ce modèle, un électron atomique est rappelé vers le centre de l'atome par une force élastique $-k\overrightarrow{r}$. Les pertes d'énergie dues au rayonnement sont modélisées par une force de frottement fluide $-\beta\overrightarrow{v}$. On pose alors $\tau=\displaystyle\frac{m_e}{\beta}$ le temps caractéristique d'amortissement et $\omega_0=\displaystyle\sqrt{\frac{k}{m_e}}$.
\paragraph{Ordre de grandeur} Les électrons des molécules de dioxygène, de diazote ou d'ozone présentes dans la haute atmosphère ont des fréquences propres d'oscillation de l'ordre de $10^{17}$ Hz, et un temps caractéristique d'amortissement de l'ordre de $10^{-8}$ s.
\subsection{Diffusion d'un rayonnement par un électron atomique}
Considérons une onde électromagnétique plane progressive monochromatique, polarisée rectilignement, qui arrive sur un atome ou une molécule initialement non polaire, mais polarisable. Sous l'action du champ électrique incident $\overrightarrow{E_\text{inc}}$ de cette onde, l'atome ou la molécule acquiert en régime sinusoïdal forcé un moment dipolaire électrique induit $\overrightarrow{p}$. Si on note $\underline{\overrightarrow{E_\text{inc}}}=E_me^{j\omega t}\overrightarrow{u_x}$, alors on a : $$\underline{\overrightarrow{p}}=\underline{p_0}e^{j\omega t}\overrightarrow{u_x} \ \ \text{avec} \ \ \underline{p_0}=\underline{\alpha}(\omega)E_m$$ où $\underline{\alpha}(\omega)$ est un coefficient complexe de polarisabilité. Ce dipôle électrique oscillant induit rayonne dans tout l'espace un champ magnétique qui transporte une puissance dont la moyenne temporelle obéit à la formule de Larmor : $$\langle P\rangle=\frac{\mu_0\omega^4\lvert\underline{p_0}\rvert^2}{12\pi c}$$ On note en général, pour la diffusion par un électron atomique oscillant : $\langle P\rangle=\sigma(\omega)I_\text{inc}$ où $$I_\text{inc}=\left\langle\left\lVert\overrightarrow{\pi_{\text{inc}}}\right\rVert \right\rangle=\frac{\varepsilon_0cE_m^2}{2}$$ est la puissance surfacique moyenne de l'onde incidente et $$\sigma(\omega)=\frac{8\pi r_e^2}{3}\frac{\omega^4}{(\omega_0^2-\omega^2)^2+\omega^2/\tau^2}$$ la section efficace de diffusion, où $\displaystyle r_e=\frac{e^2}{4\pi\varepsilon_0 m_ec^2}$. Cette puissance $\langle P\rangle$ est dite "diffusée" car elle provient de l'onde incidente qui excite le moment dipolaire.
\paragraph{Comment faire les calculs pour le prouver ?} Déterminer le champ électrique selon le modèle de Thomson. Appliquer le PFD à un électron puis passer en notation complexe. Utiliser le fait qu'on est en optique donc $\omega\ll\omega_0$ pour simplifier grossièrement $\sigma(\omega)$, ce qui servira dans la suite.
\subsection{Le bleu du ciel, la couleur rouge du Soleil au couchant}
\paragraph{Le bleu du ciel}Le jour et sans nuage, le ciel apparaît bleu pour un observateur terrestre qui ne regarde pas dans la direction du Soleil. En effet, la lumière reçue par l'œil a été diffusée par les molécules de l'atmosphère éclairées par le Soleil, et le spectre de la lumière diffusée est celui du spectre solaire pondéré par un facteur $\displaystyle\frac{1}{\lambda^4}$, ce qui fait que le bleu y est beaucoup plus représenté que le rouge.
\paragraph{Le rouge du Soleil levant et du Soleil couchant}Lorsque l'observateur regarde le soleil couchant (ou levant), la lumière qui lui parvient arrive avec une faible inclinaison par rapport au sol. Elle traverse une grande longueur d'atmosphère et subit donc une absorption considérable à cause de la diffusion (la puissance diffusée est prélevée à celle de l'onde incidente) : si l'on note $I(\omega,z)$ l'intensité lumineuse d'une OPPH de pulsation $\omega$ se propageant vers les $z$ croissants, on a : $$I(\omega,z)=I_0e^{-z/\delta(\omega)}$$ où $I_0$ est l'intensité de cette onde avant de rentre dans l'atmosphère et $\delta(\omega)$ est la distance caractéristique d'absorption, inversement proportionnelle à $\sigma(\omega)$. Cette absorption est d'autant plus force que la distance traversée est grande et que la longueur d'onde est faible : la lumière transmise devient donc de plus en plus rouge.
\paragraph{Comment faire les calculs pour le prouver ?} Effectuer un bilan de puissance à un cylindre de section $S$ et de longueur $dz$. On note $n$ le nombre de centre diffuseurs par unité de volume. La puissance diffusée vaut alors $Sndz\sigma(\omega)I$ et entre $SI(z)$ et sort $SI(z+dz)$. On obtient une équation différentielle qui se résout en une exponentielle.
\paragraph{Remarque : et les nuages ?} Les nuages diffusent al lumière selon un autre processus appelé "diffusion de Mie". La diffusion Rayleigh n'est plus valable pour les particules diffusantes des nuages, car ce sont des gouttes d'eau dont la taille n'est pas petite devant les longueurs d'onde de la lumière visible. Donc \textbf{les nuages ne sont pas bleus}.
\paragraph{Remarque : de sages paroles}"Si vous arrivez à m'expliquer grâce au cours que les nuages sont bleus, soit vous avez un problème de vue, soit vous avez un problème d'honnêteté." Niels VDC
\subsection{Polarisation de la lumière du ciel}
La lumière émise par le Soleil n'est pas polarisée. On peut décomposer chaque composante monochromatique de l'onde arrivant du Soleil sur une molécule en deux ondes polarisées rectilignement dans deux directions orthogonales arbitraires d'un plan transverse, sans relation de phase entre elles (incohérentes). Notons $O$ la position d'une molécule et $M$ la position de l'observateur. On choisit de décomposer le champ électrique d'une telle onde incidente $\overrightarrow{E_\text{inc}}$ en une composante $\overrightarrow{E_{\text{inc}||}}$ contenue dans le plan qui contient la direction de propagation de l'onde incidente et la droite $(OM)$, et une composante $\overrightarrow{E_{\text{inc}\bot}}$ orthogonale à ce plan. Les puissances surfaciques moyennes associées à ces deux ondes incidentes polarisées rectilignement sont égales : $$\left\langle\left\lVert\overrightarrow{\pi_{\text{inc}||}}\right\rVert \right\rangle=\left\langle\left\lVert\overrightarrow{\pi_{\text{inc}\bot}}\right\rVert \right\rangle$$ Les deux champs électriques ainsi considérés excitent des composantes orthogonales $\overrightarrow{p_{||}}(t)$ et $\overrightarrow{p_{\bot}}(t)$ du moment dipolaire de la molécule, également incohérentes entre elles, de sorte que l'on peut sommer les puissances diffusées par ces deux composantes. On note $\alpha$ l'angle que fait $(OM)$ avec la direction de propagation de l'onde incidente. Comme la puissance surfacique moyenne rayonnée par un dipôle de moment dipolaire $\overrightarrow{p}$ est proportionnelle à $\sin^2(\theta)$ où $\theta=(\overrightarrow{p},\overrightarrow{OM})$, les puissances surfaciques moyennes rayonnées par $\overrightarrow{p_{||}}(t)$ et $\overrightarrow{p_{\bot}}(t)$ ne seront généralement pas les mêmes en $M$ : $$\displaystyle \frac{\left\langle\left\lVert\overrightarrow{\pi_{\text{diff}||}}\right\rVert \right\rangle}{\left\langle\left\lVert\overrightarrow{\pi_{\text{diff}\bot}}\right\rVert \right\rangle}=\cos^2(\alpha)\leqslant 1 \ \ \text{dans la direction de M}$$ Ce rapport est indépendant de la longueur d'onde. Il en résulte que la lumière diffusée par l'atmosphère et reçue par l'observateur M est partiellement polarisée, sauf théoriquement si $\alpha=0$ (inobservable alors car l'observateur regarde dès lors le Soleil), ou si $\alpha=\pi$ (tout autant inobservable car l'observateur ne peut regarder le ciel dans la direction opposée au Soleil). Selon cette analyse, si $\alpha=\displaystyle\frac{\pi}{2}$, on devrait observer une lumière polarisée rectilignement dans la direction de $\overrightarrow{p_\bot}$ : en fait, il existe d'autres mécanismes de diffusion de la lumière solaire par les particules de l'atmosphère et la polarisation observée n'est pas rigoureusement rectiligne.


\part{Thermodynamique}
\chapter{Principes de la thermodynamique}
\localtableofcontents
\section{Modélisation des systèmes thermodynamiques}

\subsection{Systèmes thermodynamiques}
On utilise trois échelles différentes :
\begin{itemize}
	\item macroscopique ($mm$ et plus) ;
	\item mésoscopique ($\mu m$) ;
	\item microscopique ($nm$ et moins).
\end{itemize}
Ordre de grandeur à connaître : dans 1 $\mu m^3$ d'air, en supposant l'air un gaz un gaz parfait et en se plaçant à $P=1\ bar$ et $T=300 \ K$, on a dans ce volume d'air environ $N=25.10^6$ molécules. \\ \par
On distingue les systèmes : 
\begin{itemize}
	\item ouverts ;
	\item fermés : pas d'échange de matière avec l'extérieur ;
	\item isolés : pas d'échange de matière ni d'énergie avec l'extérieur.
\end{itemize} \par
Ne pas confondre :
\begin{itemize}
	\item équilibre : pas d'évolution des paramètres intensifs et pas d'échange aux frontières ; 
	\item état stationnaire : pas d'évolution temporelle.
\end{itemize} \par 
Hypothèse de l'équilibre thermodynamique local : $T$, $P$ sont définies localement mais ne sont pas nécessairement uniformes sur l'ensemble du système
\subsection{Modèle du gaz parfait}
Point de vue microscopique :
\begin{itemize}
	\item On néglige la taille des particules devant la distance entre les particules.
	\item On néglige les interactions à distance entre les particules.
\end{itemize} \par 
Point de vue macroscopique : 
\begin{itemize}
	\item Équation d'état : $PV=nRT$
	\item Vérifie la première loi de Joule : $U_m=U_m(T)$
	\item Vérifie la deuxième loi de Joule : $H_m=H_m(T)$
\end{itemize}
On note $u=\sqrt{<v^2>}$ la vitesse quadratique moyenne en $m.s\inv$. On définit alors la température cinétique par : $$\frac{1}{2}mu^2=\frac32 k_BT$$ et la pression cinétique par $$P=\frac13n\st mu^2$$ avec $n\st$ la densité particulaire en $m^{-3}$. On en déduit alors $PV=nRT$ en posant $R=N_A\times k_B$.
\subsection{Modèle du gaz de Van der Waals (complément)}
$$\left(P+a\frac{n^2}{V^2}\right)\left(V-nb\right)=nRT$$ où $a$ et $b$ sont des constantes strictement positives qui dépendent du gaz mais ne dépendent pas de $n$, $P$, $T$, etc. \\ \par 
Attention : penser à vérifier "l'homogénéité intensive/extensive" des grandeurs. \\ \par 
Le but de ce modèle est de :
\begin{itemize}
	\item prendre en compte le volume des particules ;
	\item prendre en compte les interactions à distance.
\end{itemize}
\subsection{Modèle de la phase condensée idéale}
Il s'agit d'un modèle où la phase condensée est :
\begin{itemize}
	\item \textbf{incompressible} : $V$ ne varie pas lorsque $P$ varie ; 
	\item \textbf{indilatable} : $V$ ne varie pas lorsque $T$ varie.
\end{itemize}  \par  On définit le coefficient de compressibilité isotherme par : $$\chi_T=-\frac{1}{V}\frac{\partial V}{\partial P}\bigg\rvert_T=\frac{1}{\rho}\frac{\partial \rho}{\partial P}\bigg\rvert_T$$ On définit le coefficient de compressibilité isentropique par $$\chi_S=-\frac{1}{V}\frac{\partial V}{\partial P}\bigg\rvert_S$$ Pour une phase condensée idéale, $U_m=U_m(T)$ car on ne peut pas faire varier $V$.
\section{Premier principe ; bilans d'énergie}
\subsection{Énergie interne ; enthalpie ; capacités thermiques}
L'\textbf{énergie interne} $U$ est une fonction d'état extensive qui s'exprime en Joule. Son existence et les propriétés sus-citées forment la première partie du premier principe. \\ \par L'\textbf{enthalpie} $H=U+PV$ est une fonction d'état extensive. Elle s'exprime en Joule. \\ \par Pour un système fermé, on définit la \textbf{capacité thermique à volume constant} (ou \textbf{capacité calorifique}) par $$C_p=\frac{\partial U}{\partial T}\bigg\rvert_V$$ \par Pour un système fermé, on définit la \textbf{capacité thermique à pression constante} par : $$C_P=\frac{\partial H}{\partial T}\bigg\rvert_P$$ \par On rappelle que si $X$ est une grandeur extensive, alors : 
\begin{itemize}
	\item $x$ est la grandeur massique associée ;
	\item $X_m$ est la grandeur molaire associée.
\end{itemize}
Les seules exceptions sont $n$, $m$ et $M$. \\ \par Pour une phase condensée idéale, on a $C_V\approx C_P$. \\ \par Pour l'eau, on a $c=4,18 \ kJ.kg\inv.K\inv$. \\ \par Pour un \textbf{gaz parfait monoatomique} : $$\begin{cases}
	U=\displaystyle\frac32nRT \\[2mm] U_m=\displaystyle\frac32RT \\[2mm] u=\displaystyle\frac32\frac{RT}{M}
\end{cases} \text{et} \ \ \begin{cases}
C_V=\displaystyle\frac32nR \\[2mm] C_{V,m}=\displaystyle\frac32R \\[2mm] c_m=\displaystyle\frac32\frac{R}{M}
\end{cases}$$ \par Pour un \textbf{gaz parfait diatomique} lorsque $T_{rot}\ll T\ll T_{vib}$ (avec en pratique $T_{rot}=10 \ K$ et $T_{vib}=1000\ K$), on a : $$C_V\approx\frac{5}{2}nR$$ \par
Pour un gaz parfait, on définit le \textbf{coefficient de Laplace}, ou \textbf{exposant adiabatique}, ou encore \textbf{rapport des capacités thermiques} $\gamma$, par : $$\gamma=\frac{C_P}{C_V}$$ \par 
Pour un gaz parfait, on a la \textbf{relation de Mayer} : $$C_p=C_V+nR$$ \par 
Pour un \textbf{gaz parfait} : $$C_V=\frac{nR}{\gamma-1} \ \ \text{et} \ \ C_P=\frac{nR\gamma}{\gamma-1}$$
\subsection{Échanges d'énergie}
Travail des forces de pression (qui doivent être uniformes sur l'ensemble du système pour que la formule soit valable) : $$\delta W=-P_{ext}\times dV$$ \par Travail électrique en convention récepteur : $$\delta W=ui\times dt$$ \par Transfert thermique : cf. premier principe
\subsection{Modèles de transformation}
Une transformation est dite :
\begin{itemize}
	\item \textbf{adiabatique} si elle a lieu sans transfert thermique ; 
	\item \textbf{monotherme} si tous les transferts thermiques se font avec un même thermostat porté à une température constante ;
	\item \textbf{isotherme} si la température du système est définie uniformément sur le système et est indépendante du temps ;
	\item \textbf{monobare} si la pression extérieure est constante ;
	\item \textbf{isobare} si la pression du système est définie uniformément et indépendante du temps ;
	\item \textbf{isochore} si le volume du système est constant ;
	\item \textbf{mécaniquement réversible} si à tout instant $P_{ext}=P$ ;
	\item \textbf{quasistatique} si le système est à tout instant infiniment proche d'un état d'équilibre ;
	\item \textbf{réversible} si elle est quasistatique, mécaniquement réversible, et que tout chemin emprunté entre deux états d'équilibre infiniment proches peut être emprunté dans l'autre sens.
\end{itemize}
\subsection{Énoncés du premier principe}
\begin{thm}[Premier principe]
	Lors d'une transformation $I\to F$ d'un système fermé : $$E_F-E_I+U_F-U_I=W+Q$$ 
\end{thm}Très souvent, le terme $E_F-E_I$ est négligeable. \\ \par \begin{thm}[Premier principe infinitésimal]
	Lors d'une transformation d'un système fermé : $$dE+dU=\delta W+\delta Q$$
\end{thm} \par Ordre de grandeur : pour une masse $m$ d'eau, on a, lorsque $\Delta T=1\ K$ et $\Delta z=400\ m$, $\Delta E_{pp}=\Delta U$.
\begin{thm}[Premier principe monobare avec équilibre mécanique initial et final]
	Pour une transformation monobare avec équilibre mécanique initial et final d'un système fermé : $$E_F-E_I+H_F-H_I=W_{utile}+Q$$ où $W_{utile}=W-W_{pression}$
\end{thm}
\subsection{Lois de Laplace}
\begin{thm}[Lois de Laplace]
	Lors d'une transformation $I\to F$ adiabatique, mécaniquement réversible d'un gaz parfait avec $\gamma$ constant, immobilité marcoscopique et sans travail utile, on a : $$ P_IV_I^\gamma=P_FV_F^\gamma \ \ \text{et} \ \ T_I^\gamma P^{1-\gamma}=T_F^\gamma P_F^{1-\gamma} \ \ \text{et} \ \ T_IV_I^{\gamma-1}=T_FV_F^{\gamma-1}$$
\end{thm}
On rappelle que pour un gaz parfait monoatomique $\gamma=\displaystyle\frac53$ et que pour un gaz parfait diatomique avec $T_{rot}\ll T\ll T_{vib}$ $\gamma=\displaystyle\frac75$.
\section{Deuxième principe ; bilans d'entropie}
Pourquoi la nécessité d'un deuxième principe ? Le premier principe est un principe de conservation : pour un système isolé, $U_F=U_I$. Le premier principe ne permet alors pas de savoir si la transformation se déroule dans le sens $I\to F$ ou dans le sens $F\to I$. \\ \par On a alors la nécessité d'un principe d'évolution, le second principe, pour savoir dans quel sens se déroule la transformation.
\subsection{Réversibilité - irréversibilité}
Idée intuitive de l'irréversibilité : le film est complètement incohérent si on le projette à l'envers. \\ \par Par exemple, une solution de l'équation $m\ddot{x}=kx$ est réversible car si $t\mapsto x(t)$ est solution, alors $t\mapsto x(-t)$ est aussi solution. En revanche, ce n'est plus le cas pour $m\ddot{x}=\lambda\dot{x}$, sauf si on prend la solution nulle. \\ \par Le caractère réversible ou non est une question d'échelle des temps : si on prend l'exemple d'un oscillateur très faiblement amorti que l'on observe sur une courte période de temps, alors on ne voit pas l'irréversibilité car l'amortissement est trop faible. En revanche, si on prend une échelle de temps beaucoup plus importante, on peut commencer à apercevoir l'irréversibilité.
\subsection{Énoncés du deuxième principe}
\begin{thm}[Deuxième principe]
	Il existe une fonction d'état extensive appelée \textbf{entropie} et notée $S$ qui s'exprime en $J.K\inv$ telle que pour toute transformation $I\to F$ d'un système fermé : $$S_F-S_I=S_{ech}+S_{créée}$$ avec $S_{créée}\geqslant0$ et $S_{ech}$ l'entropie d'échange qui vaut $0$ si la transformation est adiabatique et $\displaystyle\sum_i \frac{Q_i}{T_i}$ si le système reçoit $Q_i$ du thermostat $T_i$
\end{thm}
\begin{thm}[Deuxième principe infinitésimal]
	Pour une transformation élémentaire d'un système fermé : $$dS=\delta S_{ech}+\delta S_{créée}$$ avec $\displaystyle\delta S_{ech}=\frac{\delta Q}{T}$ si la transformation est thermiquement réversible.
\end{thm}
\subsection{Interprétation de l'entropie (culture)}
Attention, l'entropie n'est pas une mesure du désordre. Elle est plutôt une \textbf{mesure du manque d'information} (exemple simple du nombre de complexions lors d'une détente de Joule Gay-Lussac).
\subsection{Calculs de variation d'entropie}
\paragraph{Cas du gaz parfait}En empruntant un chemin réversible et appliquant les deux premiers principes et l'équation d'état, on trouve $$dS=\frac{nR}{\gamma-1}\frac{dT}{T}+nR\frac{dV}{V}$$ et donc $$S-S_0=\frac{nR}{\gamma-1}\ln\left(\frac{T}{T_0}\right)+nR\ln\left(\frac{V}{V_0}\right)$$ \paragraph{Cas de la phase condensée idéale}On a : $$dS=mc\frac{dT}{T}$$ et donc en supposant $c$ indépendante de $T$ $$S-S_0=mc\ln\left(\frac{T}{T_0}\right)$$ Si $c$ n'est plus supposée indépendante de $T$, il faut alors intégrer la relation infinitésimale.
\section{Applications}
\subsection{Comparaison de deux détentes}
\paragraph{Détente de Joule et Gay-Lussac}
Compartiment calorifugé rigide de volume $2V_0$. Initialement, un gaz parfait est contenu dans le volume $V_0$ avec $n$ moles à la température $T_I$. \\ \par Premier principe : $U_F-U_I=Q+W$ avec $Q=0$ car les parois sont calorifugées. $W=0$ car la détente s'effectue dans du vide. Donc $\Delta U=0$. On en déduit $T_F=T_I$ car $\Delta U=C_V\Delta T$. \\ \par Or, $S_{ech}=0$ car la transformation est adiabatique, donc $\Delta S=S_{créée}=nR\ln(2)$ car $T_F=T_I$ et $V_F=2V_0$ alors que $V_I=V_0$.
\paragraph{Détente isotherme mécaniquement réversible d'un gaz parfait}Un piston détend lentement un gaz parfait d'un état initial $(V_0,T_0,n)$ vers un état final $(2V_0,T_0,n)$ avec une transformation monotherme avec l'atmosphère à $T_0$. \\ \par On a $\Delta U=C_V\Delta T=0$ donc $Q=-W$ d'après le premier principe. Or, $\delta W=-nRT\displaystyle\frac{dV}{V}$ car on a un gaz parfait et la transformation est mécaniquement réversible, donc $W=-nRT\ln(2)$ et $Q=nRT\ln(2)$. \\ \par $\displaystyle\Delta S=\frac{nR}{\gamma-1}\ln\left(\frac{T_0}{T_0}\right)+nR\ln\left(\frac{2V_0}{V_0}\right)=nR\ln(2)$. Or, $\delta S_{ech}=\displaystyle\frac{\delta Q}{T_0}$ donc $S_{ech}=nR\ln(2)$ d'où $S_{créée}=0$ si bien que la transformation est réversible.
\paragraph{Comparaison}
\begin{center}
	\begin{tabular}{|c|c|c|}
		\hline
		& Joule et Gay-Lussac & Isotherme mécaniquement réversible \tabularnewline
		\hline 
		$\Delta U$ & 0 & 0 \\
		\hline
		W&0&$-nRT_0\ln(2)$\\
		\hline
		Q&0&$nRT_0\ln(2)$\\
		\hline
		$S_{ech}$&0&$nR\ln(2)$\\
		\hline
		$\Delta S$&$nR\ln(2)$&$nR\ln(2)$\\
		\hline
		$S_{créée}$&$nR\ln(2)$&0\\
		\hline		
	\end{tabular}
\end{center}
\subsection{Transfert thermique spontané}
On considère deux compartiments $(C_1,T_1)$ et $(C_2,T_2)$ qui sont "en contact thermique et rien d'autre". On suppose $T_1>T_2$ initialement et on veut montrer $dT_1<0$. \\ \par Premier principe à ${1,2}$ : $dU=C_1dT_1+C_2dT_2=\delta W+\delta Q$ avec $\delta W=0$ et $\delta Q=0$ car on considère que le système est isolé de tout. \\ \par Deuxième principe à ${1,2}$ : $dS=\delta S_{ech}+\delta S_{créée}$ avec $\delta S_{ech}$ car la transformation de ${1,2}$ est adiabatique. On a donc $dS\geqslant 0$. Or, $dS=dS_1+dS_2=C_1\displaystyle\frac{dT_1}{T_1}+C_2\frac{dT_2}{T_2}$ \\ \par Si $C_1=C_2=C$, alors $dS=C\left(\displaystyle\frac{1}{T_1}-\frac{1}{T_2}\right)dT_1$ et comme $dS\geqslant0$, $\displaystyle \frac{1}{T_1}-\frac{1}{T_2}\leqslant0$ et $C\geqslant0$, alors $dT_1\leqslant0$.
\section{Troisième principe de la thermodynamique (complément)}
Le troisième principe de la thermodynamique, aussi appelé principe de Nernst, énonce que 
\begin{center}
	\textbf{La valeur de l'entropie de tout corps pur dans l'état de cristal parfait est nulle à la température de 0 Kelvin.}
\end{center}
Ce principe permet d'avoir une valeur déterminée de l'entropie, et non pas une valeur à constante additive près. Ce principe a été généralisé aux solides quelconques, aux liquides et aux gaz. En revanche, il ne s'applique pas aux systèmes quantiques : il est lié à l'indiscernabilité quantique des particules identiques.

\chapter{Changements d'état d'un corps pur}
\localtableofcontents
\section{Variance}
La variance est le nombre maximum de paramètres intensifs qu'un opérateur peut fixer sans rompre l'équilibre du système.
\section{Diagramme $P-T$}
\subsection{Rappels}
Diagramme $P$ en fonction de $T$. La zone solide se trouve à gauche, celle liquide au milieu en haut et celle du gaz en bas à droite. Le point d'intersection des trois courbes est le \textbf{point triple}. Le point au bout de la courbe d'équilibre liquide-vapeur est le \textbf{point critique}. La variance est nulle au point triple. Pour un équilibre entre deux états, la variance vaut 1 : usuellement, on prend alors la pression en fonction de la température. \\ \\ Pour l'eau (et aussi pour le bismuth Bi, le gallium Ga, etc.), la pente de l'équilibre solide-liquide est négative, contrairement à celle pour la plupart des autres corps purs.
\subsection{Formule de Clapeyron}
Pour une transition de phase $1\to 2$, on a la \textbf{formule de Clapeyron} : $$\frac{dP_{1\to 2}}{dT}=\frac{L_{1\to 2}}{T(v_{2}-v_{1})}$$ où $v_1$ et $v_2$ sont pris à $T$ et $P_{1\to 2}(T)$, et où $L_{1\to2}=h_2-h_1$ est la chaleur latente de changement d'état. \\ \\ En particulier, on sait que $\rho_S<\rho_L$ pour l'eau, donc $v_S>v_L$ ce qui explique la pente négative annoncée précédemment.
\section{Diagramme $P-v$}
\subsection{Rappels}
Le diagramme représente la pression P en fonction du volume massique v. La courbe s'appelle \textbf{courbe de saturation}. Son point culminant est le \textbf{point critique}. La partie de la courbe à gauche du point critique est la \textbf{courbe d'ébullition}. La partie de la courbe à droite du point critique est la \textbf{courbe de rosée}. Pour la température critique, l'isotherme est une courbe qui frôle le diagramme au point critique en y faisant un point d'inflexion. Pour une température inférieur à la température critique, elle est quasiment verticale dans le domaine de liquide, horizontale dans le domaine liquide-vapeur (à température constante, le changement d'état d'un corps pur se fait à pression constante) et redevient courbe décroissante dans le domaine vapeur. Pour une température plus grande que la température critique, l'isotherme est au-dessus de la courbe $P-v$, dans la zone du fluide supercritique.
\subsection{Loi des moments}
On appelle titre massique la quantité : $$x_L=\frac{m_L}{m_L+m_V}$$
La \textbf{loi des moments} stipule que : $$x_L=\frac{v_M-v_B}{v_A-v_B}$$ sur une isotherme dans un diagramme $P-v$. Le résultat est identique pour un diagramme $T-s$ sur une isobare.
\section{Enthalpie et entropie de changement d'état}
En réalité, on parle ici d'enthalpies et d'entropies \textbf{massiques} de changement d'état.
\subsection{Enthalpie de changement d'état}
Il s'agit de : $$\Delta_{1\to2}h=h_2(T,P_{1\to2}(T))-h_1(T,P_{1\to2}(T))$$ Ordres de grandeur pour l'eau : \begin{itemize}
	\item $\Delta_{\text{fus}}h\approx 300\ kJ/kg$
	\item $\Delta_{\text{vap}}h\approx 2000\ kJ/kg$
\end{itemize}
\subsection{Entropie de changement d'état}
Il s'agit de : $$\Delta_{1\to2}s=s_2(T,P_{1\to2}(T))-s_1(T,P_{1\to2}(T))$$ Elle est reliée à l'enthalpie massique de changement d'état par : $$\Delta_{1\to2}s(T)=\frac{\Delta_{1\to2}h(T)}{T}$$ (imaginer un chemin réversible, appliquer le premier principe monobare et utiliser le fait qu'on a un thermostat).
\section{Exemples}
\subsection{Cessation d'un état métastable}
TD4, exercice 13 : exemple de la surfusion. \\ Il faut faire des hypothèses sur la constitution finale du système (équilibre diphasé ou non), puis, en remarquant que cette transformation étant rapide, on peut la considérer comme adiabatique et alors puisqu'elle est monobare, la variation d'entropie est nulle. \\ Il faut alors imaginer une suite de transformation fictives simples (échauffements ou refroidissements sans changement d'état, ou changement d'état à pression et température constante) pour calculer la variation d'enthalpie d'une autre manière en utilisant que l'enthalpie est une fonction d'état, ce qui permet alors d'obtenir le titre en liquide ou en masse. \\ Il convient alors de vérifier la validité de l'hypothèse effectuée au départ. \\Ensuite, pour calculer la variation d'entropie, on suit les mêmes transformation fictives en disant que l'entropie est une fonction d'état.
\subsection{Exemple de changement d'état irréversible}
Exemple : Fusion d'un glaçon en contact avec un thermostat à $T_{ext}$. Le glaçon est initialement à $T_{\text{fus}}$. On applique le deuxième principe qui donne $$S_{\text{créée}}=\Delta S-S_{ech}$$ Or, le premier principe monobare donne $\Delta H =Q=m\Delta_{\text{fus}}h$ et on a $\Delta S=m\Delta_\text{fus}s=\displaystyle m\frac{\Delta_\text{fus}h}{T_{\text{fus}}}$. On en tire $$S_{\text{créée}}=m\Delta_\text{fus}h\left(\frac{1}{T_\text{fus}}-\frac{1}{T_\text{ext}}\right)$$ comme $T_\text{fus}<T_\text{ext}$, on en déduit que la transformation est irréversible.
\subsection{Refroidissement isochore d'un mélange liquide - vapeur}
\textit{cf} TP du vendredi 08/11/2024 : exemple du fluide à très haute pression dans un tube et de sa transformation isochore (ce sont des droites verticales sur le diagramme $P-v$). Notamment, on observe un brouillard étrange à la température critique et au volume correspondant au point critique.

\chapter{Systèmes ouverts ; machines thermiques}
\localtableofcontents
\section{Bilans pour un système ouvert en régime stationnaire}
\subsection{Choix d'un système fermé}
On prend $\Sigma\st$ le fluide qui est présent dans la machine à $t$ et entre dans la machine à $t$, et qui est présent dans la machine à $t+dt$ et sort de la machine à $t+dt$. Ce système est fermé. \\ \par Puisque $\Sigma\st$ est un système fermé, on a : $$m_{\Sigma\st}(t)=m_{\Sigma\st}(t+dt)$$ Par extensivité de la masse, cela donne : $$\delta m_e+m_{\Sigma}(t)=m_{\Sigma}(t+dt)+\delta m_s$$ Or, en régime permanent, on a : $m_{\Sigma}(t)=m_{\Sigma}(t+dt)$. On en déduit : $$\delta m_e=\delta m_s$$ que l'on notera $\delta m$ dans la suite. On définit alors le \textbf{débit massique} par : $$D_m=\frac{\delta m}{dt}$$
\subsection{Bilan d'une grandeur extensive quelconque} Le but est de calculer $X_{\Sigma\st}(t+dt)-X_{\Sigma\st}(t)$ où $X$ est une grandeur extensive. On a : \begin{align*}
	X_{\Sigma\st}(t+dt)-X_{\Sigma\st}(t)&=\left(X_{\Sigma}(t+dt)+\delta m_s x_s\right)-\left(X_{\Sigma}(t)+\delta m_e x_e\right)\\
	&= \delta m\left(x_s-x_e\right)
\end{align*} car en régime permanent $X_{\Sigma}(t+dt)=X_{\Sigma}(t)$.
\subsection{Bilan énergétique}
On applique le premier principe au système \textbf{fermé} $\Sigma\st$. On obtient alors : $$\delta m\left[e_m+u\right]_e^s=\delta W+\delta Q$$
\paragraph{Calcul du travail d'admission} En notant $\sigma_e$ la section du tuyau d'entrée, on a : $$\delta W_{\text{admission}}=P_e\sigma_e\overrightarrow{n}\cdot\overrightarrow{dl_e}$$ soit $$\delta W_{\text{admission}}=P_ev_e\delta m$$
\paragraph{Calcul du travail de refoulement} En notant $\sigma_s$ la section du tuyau de sortie, on a : $$\delta W_{\text{refoulement}}=P_s\sigma_s\overrightarrow{n}\cdot\overrightarrow{dl_s}$$ soit $$\delta W_{\text{refoulement}}=-P_sv_s\delta m$$
\paragraph{Premier principe pour les écoulement permanents} En réinjectant dans l'expression trouvée précédemment, on obtient ce qu' l'on appelle aussi parfois le "premier principe industriel" : en régime permanent, on a :  $$\delta m\left[e_m+h\right]_e^s=\delta W_{\text{utile}}+\delta Q$$ En notant $w_{\text{utile}}$ le travail utile massique et $q$ le transfert thermique massique, cela peut se réécrire : $$\left[e_m+h\right]_e^s=w_{\text{utile}}+q$$ Encore, en notant $P_{\text{utile}}$ la puissance utile et en notant $\Phi$ le flux thermique (puissance thermique), cela peut finalement s'écrire : $$D_m\left[e_m+h\right]_e^s=P_{\text{utile}}+\Phi$$
\paragraph{Exemple} Si la seule énergie "macroscopique" est l'énergie potentielle de pesanteur, on a : $$e_m=\frac{\delta m\frac{c^2}{2}+gz\delta m}{\delta m}=\frac{c^2}{2}+gz$$ avec $c$ la vitesse d'ensemble du fluide et $z$ l'altitude en prenant l'axe $\overrightarrow{u_z}$ vertical orienté vers le haut.
\subsection{Bilan entropique}
On applique le deuxième principe au système \textbf{fermé} $\Sigma\st$ pour obtenir la formule : $$\delta m\left[s\right]_e^s=\delta S_\text{ech}+\delta S_\text{créée}$$
\subsection{Exemples}
\paragraph{Détente de Joule Kelvin}
Un fluide traverse en régime permanent un compartiment actif ou un matériau poreux en circulant dans un tuyau calorifugé. On note $P_e$ et $P_s$ les pressions d'entrée et de sortie, et $T_e$ et $T_s$ les températures d'entrée et de sortie. On applique le premier principe pour un écoulement permanent : $$\delta\left[h\right]_e^s=\delta W_\text{utile}+\delta Q$$ Or, $\delta W_\text{utile}=0$ car il n'y a a pas de paroi mobile, et $\delta Q=0$ car les parois du tuyau sont calorifugées. On a donc $$h_e=h_s$$ On dit que la détente de Joule-Kelvin est une détente \textbf{isenthalpe}. Si on rajoute l'hypothèse "gaz parfait", on obtient $$T_e=T_s$$ Montrons que la détente de Joule-Kelvin est irréversible. Puisque la transformation est adiabatique, le deuxième principe donne : $$\delta m\left[s\right]_e^s=S_\text{créée}$$ Or, pour $n$ moles d'un gaz parfait, on a : $$S-S_0=\frac{nR\gamma}{\gamma-1}\ln\left(\frac{T}{T_0}\right)-nR\ln\left(\frac{P}{P_0}\right)$$si bien que $$S_\text{créée}=\frac{\delta m R}{M}\ln\left(\frac{P_e}{P_s}\right)$$ Or, l'expérience avec ce sens d'écoulement n'est possible que si $P_e>P_s$, donc $S_\text{créée}>0$ et la détente est irréversible.
\paragraph{Tuyère}
Le fluide traverse une tuyère calorifugée. Les paramètres à l'entrée de la tuyère sont $T_e$, $P_e$ et $c_e$. Les paramètres à la sortie de la tuyère sont $T_s$, $P_s$ et $c_s$. $c_s$ n'est alors plus négligeable, et c'est même précisément ce qui nous intéresse. \\ \par On applique le premier principe pour les écoulements permanents en négligeant les variations d'énergie potentielle devant les variations d'énergie cinétique, qui donne $$\delta m\left[e_c+h\right]_e^s=\delta W_\text{utile}+\delta Q$$ Or, $\delta W_\text{utile}=0$ car il n'y a pas de paroi mobile, et $\delta Q=0$ car la tuyère est calorifugée. On obtient donc, en supposant que le fluide est un gaz parfait d'exposant adiabatique $\gamma$ et de masse molaire $M$ : $$c_s=\sqrt{\frac{2R\gamma}{(\gamma-1)M}\left(T_e-T_s\right)}$$
\section{Machines thermiques}
\subsection{Rappels}
On a le \begin{thm}[2ème principe de Kelvin]
	Il n'existe pas de moteur cyclique monotherme.
\end{thm}
\begin{proof}
	On applique le premier principe sur un cycle au fluide caloporteur : $$\Delta U=W+Q$$Or, $U$ est une fonction d'état, donc sur un cycle, $\Delta U=0$ et on en déduit que $W=-Q$. On applique le deuxième principe sur un cycle au fluide caloporteur : $$\Delta S=S_\text{ech}+S_\text{créée}$$Or, on a de même $\Delta S=0$, et on sait que $S_\text{ech}=\displaystyle\frac{Q}{T}$. On en déduit que $W=TS_\text{créée}\geqslant0$. Ainsi, on peut avoir un radiateur, mais pas un moteur cyclique monotherme.
\end{proof}
Rappelons ici le résumé des formules à connaître pour les moteurs, pompes à chaleur et machines frigorifiques : 
\begin{center}
	\begin{tabular}{|c|c|c|c|c|}
	\hline
	&Intérêt&Coût&Rendement ou efficacité&Rendement ou efficacité de Carnot\\[3mm]
	\hline
	Moteur&$-W$&$Q_C$&$\displaystyle\eta=\frac{-W}{Q_C}$&$\displaystyle\eta_\text{Carnot}=1-\frac{T_F}{T_C}$ \\[5mm]
	\hline
	PAC&$-Q_C$&$W$&$\displaystyle e_\text{PAC}=\frac{-Q_C}{W}$&$\displaystyle e_\text{PAC, Carnot}=\frac{T_C}{T_C-T_F}$ \\[5mm]
	\hline
	Frigo&$Q_F$&$W$&$\displaystyle e_\text{Frigo}=\frac{Q_F}{W}$&$\displaystyle e_\text{Frigo, Carnot}=\frac{T_F}{T_C-T_F}$ \\[5mm]
	\hline
\end{tabular}
\end{center}
Remarque : Si $T_F=T_C$, alors : \begin{itemize}
	\item $\eta_\text{Carnot}=0$ (en accord avec le deuxième principe de Kelvin)
	\item $e_\text{PAC, Carnot}\longrightarrow +\infty$
	\item $e_\text{Frigo, Carnot}\longrightarrow +\infty$
\end{itemize}
\paragraph{Cycle de Carnot} C'est un cycle réversible avec 2 adiabatiques et 2 isothermes (une à $T_F$ et une à $T_C$). Remarque : sa puissance est nulle...
\subsection{Écoulements dans les principaux éléments d'une machine thermique}
\paragraph{Détendeur} Le premier principe industriel donne en négligeant les variations d'énergie mécanique macroscopique $\delta m\left[h\right]_e^s=\delta Q+\delta W_\text{utile}$. Or, on a souvent $\delta Q=0$ et $\delta W_\text{utile}=0$, donc on a souvent $$h_e=h_s$$
\paragraph{Échangeurs thermiques} Par exemple, ce sont les évaporateurs, les condenseurs, etc. Le premier principe industriel donne $h_s-h_e=q+w_\text{utile}$ avec très souvent $w_\text{utile}=0$.
\paragraph{Tuyère} On néglige plus les variations d'énergie cinétique. On a alors : $$\left[e_c+h\right]_e^s=q+w_\text{utile}$$ avec souvent $q=0$, $w_\text{utile}=0$ et $c_e\ll c_s$.
\paragraph{Pompe/compresseur/turbine} On a $h_s-h_e=q+w_\text{utile}$ avec $q=0$. Pour une turbine, $w_\text{utile}<0$.

\subsection{Utilisation du diagramme de Clapeyron}
Pour une machine thermique, la machine est un moteur si elle décrit son cycle dans le sens horaire dans le diagramme de Clapeyron (diagramme $P-v$). Et on a les machines frigorifiques (et donc les PAC) qui correspondent au sens trigonométrique. Retenir \textbf{frigo/trigo}. La démonstration est la même que dans le paragraphe suivant.
\subsection{Utilisation du diagramme entropique}
Pour une machine thermique, la machine est un moteur si elle décrit son cycle dans le sens horaire dans le diagramme entropique (diagramme $T-s$). Et on a les machines frigorifiques (et donc les PAC) qui correspondent au sens trigonométrique. Retenir \textbf{frigo/trigo}. \\ \par Preuve : On réalise la preuve pour un cycle de Carnot. On a alors : $Q_C+Q_F=-W$ d'après le premier principe appliqué sur un cycle. Or, on a $$\oint_{\text{Cycle}}T\ dS=Q_C+Q_F$$ En parcourant le cycle dans le sens horaire, on trouve $Q_F+Q_C>0$ donc $W<0$. Dans le sens trigonométrique, on trouve $W>0$.
\subsection{Utilisation du diagramme enthalpique (ou des "frigoristes")}
\textit{cf} DM sur le diagramme des frigoristes et le cours de sup \\ \par Prendre au garde que parfois le diagramme représente $P$ en fonction de $h$, mais le plus souvent, bien qu'on le note encore $P-h$, il s'agit en réalité de $\ln(P)$ en fonction de $h$. Il faut savoir quelles sont les différentes courbes remarquables de ce diagramme. Souvent, placer les points remarquables de la transformation grâce aux courbes remarquables et appliquer le premier principe industriel permettra par lecture graphique d'arriver très vite au rendement de la machine. \\ \par Les \textbf{isothermes} sont des droites verticales dans le domaine liquide (le modèle de la phase condensée idéale est applicable et l'enthalpie est alors constante si $T$ est constante). Dans le domaine liquide-vapeur, ce sont des droites horizontales car à température donnée, un changement d'état s'effectue à pression fixée. Dans le domaine gazeux, pour un gaz parfait, ce doit être une droite verticale. Ainsi, c'est une courbe concave qui tend à devenir une droite verticale à mesure que la pression diminue, le gaz tendant alors à se comporter comme un gaz parfait (\textit{cf} le diagramme d'Amagat $PV-P$) \\ \par Les \textbf{isotitres} sont de courbes qui partent du point critique et descendent dans la zone liquide-vapeur jusqu'à l'axe des abscisses sans se croiser. \\ \par Dans le modèle du gaz parfait, les \textbf{isentropiques} ont dans la zone du gaz une équation de la forme $$\ln(P)=\frac{c_PM}{R}\ln\left(\frac{h}{c_P}+A\right)+B$$ où $A$ et $B$ sont des constantes.
\section{Parenthèse : identités thermodynamiques}
\paragraph{Première identité thermodynamique} Pour un système fermé sans changement de constitution, on a : $$dU=TdS-PdV$$ (appliquer le premier et le deuxième principe pour une réversible fictive)
\paragraph{Deuxième identité thermodynamique} Pour système fermé sans changement de constitution, on a : $$dH=TdS+VdP$$
($dH=dU+PdV+VdP$ et utiliser la première identité thermodynamique)

\chapter{Transferts thermiques}
\localtableofcontents
\section{Modes de transferts thermique : diffusion (\textit{ie} conduction), convection et rayonnement}
\paragraph{Modes de transfert} Il y en a trois 
\begin{itemize}
	\item la \textbf{conduction}, de proche en proche sans mouvement d’ensemble
	\item la \textbf{convection}, de proche en proche avec mouvement macroscopique
	\item le \textbf{rayonnement}, possible sans support matériel
\end{itemize}
\paragraph{Exemple} Le vase Dewar contient une partie de vide pour éliminer la conduction et la convection qui nécessitent un support matériel, et est recouvert par une paroi métallique pour éliminer le rayonnement.
\section{Conduction thermique}
\subsection{Vecteur densité de courant thermique ; flux thermique à travers une surface orientée}
On note $\phi$ le flux thermique, c'est une puissance en W. Le flux thermique à travers une surface orientée $S$ vérifie : $$\phi=\iint_S\overrightarrow{j_{\text{th}}}\cdot\overrightarrow{dS}$$ où $\overrightarrow{j_\text{th}}$ est le vecteur densité de courant thermique en W.m$^{-2}$.
\paragraph{Analogie} On a une analogie avec le courant électrique : $\phi$ s'apparente à $I$ et $\overrightarrow{j_\text{th}}$ s'apparente à $\overrightarrow{j_\text{el}}$. En effet, pour des électrons circulant à la vitesse $\overrightarrow{v}$ à travers une surface $S$, on a : $$\delta Q=-en_\text{el}vS\cos(\theta)dt$$ où $n_\text{el}$ est la densité électronique et $\theta$ est l'angle entre la normale à $S$ et $\overrightarrow{v}$. Ici, le terme $vS\cos(\theta)dt$ représente donc un volume. Or, $I=\overrightarrow{j_\text{el}}\cdot S\overrightarrow{n}$ car $\overrightarrow{j_\text{el}}$ est constant ici. Donc $\overrightarrow{j_\text{el}}=-en_\text{el}\overrightarrow{v}$.
\subsection{Loi de Fourier ; conductivité thermique} 
\paragraph{Loi de Fourier}
Dans le cas de la conduction, on a : $$\overrightarrow{j_\text{th}}=-\lambda\grad(T)$$ où $\lambda$ est la conductivité thermique du milieu en W.m$\inv$.K$\inv$. $\overrightarrow{j_\text{th}}$ est donc perpendiculaire aux isothermes et orienté vers les zones froides (transfert spontané). Cette loi nécessite un équilibre thermodynamique local pour que $T$ soit définie localement. Cette loi est \textbf{phénoménologique} : elle n'est valable que si la température ne varie pas sur des échelles trop petites et à basses fréquences de variations de $T$.
\paragraph{Ordres de grandeurs de conductivités thermiques}
Sans convection, on a : 
\begin{itemize}
	\item $\lambda_\text{isolant parfait}=0$ W.m$\inv$.K$\inv$ ;
	\item $\lambda_\text{air}=10^{-2}$ W.m$\inv$.K$\inv$ (valable plus généralement pour un gaz parfait dans les CNTP sans la convection);
	\item $\lambda_\text{bois}=1$ W.m$\inv$.K$\inv$ ;
	\item $\lambda_\text{cuivre}=4.10^{2}$ W.m$\inv$.K$\inv$
\end{itemize}
\section{Bilans d'énergie ou de puissance}

\subsection{Cas sans terme source ni puits dans les trois géométries}
\paragraph{Bilan thermique} S'il n'y a que de la conduction, on a : $$\mu c\frac{\partial T}{\partial t}=-\div\left(\overrightarrow{j_\text{th}}\right)$$ avec $\mu$ la masse volumique du solide et $c$ sa capacité thermique massique.
\paragraph{Méthode} On a applique le premier principe industriel qui fournit $$dH=\phi dt$$ puis le théorème de Green-Ostrogradski fournit : $$\iiint_V\mu_Mc_M\frac{\partial T}{\partial t}d\tau=-\iiint_V\div\left(\overrightarrow{j_\text{th}}\right)d\tau=\oiint_S\overrightarrow{j_\text{th}}\cdot \overrightarrow{dS}$$ En pratique, on choisit un petit volume pour que $T$ soit définie et on choisit ce volume adapté aux éventuelles symétries du problème. On exprime alors $\phi$ en fonction de $\overrightarrow{j_\text{th}}$ pour se ramener à l'expression de la divergence.
\subsection{Cas avec production ou disparition d'énergie thermique}
Avec des termes source, on rajoute le travail ou transfert thermique additionnel dans le bilan enthalpique. \paragraph{Attention} L'énergie totale est conservée. On convertit une autre forme d'énergie en énergie thermique via : 
\begin{itemize}
	\item des réactions chimiques ;
	\item des changements d'états ;
	\item des réactions nucléaires ;
	\item l'effet Joule ; 
	\item des frottements.	
\end{itemize}
\paragraph{Exemple : Crayon d'uranium pour une réaction nucléaire} On considère un crayon d'uranium, \textit{ie} un cylindre de section $S$ d'uranium, orienté selon une direction $x$. On note $\rho$ sa masse volumique et $c$ sa capacité thermique massique. Un volume élémentaire $\delta\tau$ reçoit la puissance thermique $p_v\delta\tau$ avec $p_v$ en W.m$^{-3}$. Quelle est l'équation vérifiée par $\overrightarrow{j}=j_x(x)\overrightarrow{u_x}$ en régime permanent ? \\ \par On effectue un bilan de puissance à la portion de crayon nucléaire entre $x$ et $x+dx$ : $$\frac{U(t+dt)-U(t)}{dt}=p_vSdx+\Phi_x-\Phi_{x+dx}$$ ce qui se réécrit en régime permanent et avec les expressions des flux en fonction de la section et des densités de courant thermique : $$0=p_vSdx+Sj_x(x)-Sj_x(x+dx)$$ On en tire alors la relation : $$\frac{\partial j_x}{\partial x}=p_v$$
\section{Équation de la diffusion thermique ; diffusivité thermique}
\subsection{Sans sources ni puits}
\paragraph{Équation de la chaleur} En combinant la loi de Fourier et les bilans thermiques, sans source ni puits, on obtient $$\frac{\partial T}{\partial t}=\frac{\lambda}{\mu c}\Delta T$$ et on pose alors $\displaystyle D:=\frac{\lambda}{\mu c}$ la diffusivité thermique en m$^2$.s$\inv$. Plus $D$ est grand, plus la température s'homogénéise rapidement.
\paragraph{Remarque} Cette équation est \textbf{linéaire} : on peut donc utiliser le principe de superposition et les conditions aux limites.
\subsection{Avec sources ou puits}
On reprend l'exemple du crayon d'uranium et on a dès lors : $$\rho c\frac{\partial T}{\partial t}=\lambda\Delta T+p_v$$
\subsection{Irréversibilité}
Cette équation est irréversible : si $t\mapsto T(t)$ est solution, alors $t\mapsto T(-t)$ ne l'est pas (sauf si la solution est constante par rapport à $T$). Cela est dû à la "dérivée seconde" qu'est le Laplacien d'un côté, par rapport à l'unique dérivée partielle de l'autre côté.
\subsection{Analyse en ordre de grandeur, temps caractéristique de diffusion}
Si on multiplie la longueur caractéristique par 2, le temps caractéristique est multiplié par 4 : en effet, si $T(x,t)$ est solution, alors $T(2x,4t)$ est aussi solution. \\ \par On note $\tau$ la durée caractéristique et $L$ la longueur caractéristique. "En ordre de grandeur", l'équation de la chaleur donne $\displaystyle\frac{1}{\tau}=\frac{D}{L^2}$, ou de façon plus claire : $$\tau=\frac{L^2}{D}$$
\section{Conditions aux limites}
On analyse ce qui se passe aux bords du système, là où l'équation de la chaleur n'est plus valable.
\subsection{Sur le vecteur densité de courant}
\paragraph{Paroi calorifugée} Pour une paroi calorifugée, on a $\overrightarrow{j_\text{th}}\cdot\overrightarrow{n}=0$.
\paragraph{Continuité du flux}
A travers toute surface, on a : $$\overrightarrow{j_\text{th}}(0^-)=\overrightarrow{j_\text{th}}(0^+)$$
\subsection{Contact parfait entre deux solides}
Pour un contact parfait entre deux solides, on a $T(z=0,t)=T_0$.
\subsection{Transferts conducto-convectifs à l'interface avec un fluide ; loi de Newton}
On considère un solide et un fluide. Une couche limite d’épaisseur
$e$ du fluide à la surface du solide est immobile (seulement de la conduction). Au-delà, la convection est parfaite (température homogène). La \textbf{loi de Newton} s’écrit alors : $$\overrightarrow{j_\text{th}}=h\left(T_\text{solide}-T_\text{fluide}\right)\overrightarrow{n}_{\text{solide}\to\text{fluide}}$$ avec $\displaystyle h=\frac{\lambda}{e}$. En effet, un bilan puis une loi de Fourier sur la couche limite donne une loi de température affine.
\section{Résolution en régime stationnaire : exemple de solutions}
Pour trouver des solutions en régime variable, penser au principe de superposition et à découpler les variables temporelle et spatiale. Il faut essayer de transformer $T(x,t)$ en $f(x)g(t)$.
\section{Régime stationnaire ; résistance thermique}
En régime permanent, sans terme source, l’énergie thermique est conservée. On a par analogie avec la loi d’Ohm : $V_1-V_2=RI$ la relation $T_1-T_2=R_\text{th}\phi$ (en convention récepteur). Dans le cas d’un cylindre d’axe z, de section S, longueur L et conductivité $\lambda$ avec $T=f(z)$, on a $$R_\text{th}=\frac{L}{\lambda S}$$
Partir du bilan $\Delta H=0$ donc $\phi(z)=\phi(0)$ d’où $S_zj_z(z)=S(0)j_z(0)$...
Raisonner sur les cylindre-résistances en série/parallèle pour trouver la proportionnalité en $L$ et $\displaystyle\frac{1}{S}$ ainsi que la proportionnalité entre $R_\text{th}$ et $\displaystyle\frac{1}{\lambda}$.
\subsection{Résistances thermiques dans les trois géométries}
\paragraph{Attention} Le notion de résistance thermique nécessite un régime permanent (ou l'approximation des régimes permanents, dite ARQP) ainsi que l'absence de terme source.
\paragraph{Géométrie cartésienne} On suppose $T=T(x)$. Notre système traversé par un flux $\Phi$ et ayant une température d'entrée $T_0$ et une température de sortie $T_1$ est assimilable à un dipôle électrique où la différence de potentiel joue le rôle de la différence de température et l'intensité joue le rôle du flux thermique et on cherche à déterminer l'expression de la résistance thermique, qui joue un rôle analogue à la résistance électrique. On note $$\Phi_x=\iint_\text{section}\overrightarrow{j}(x)\cdot dS\overrightarrow{u_x}$$ On effectue un bilan de puissance au solide entre $0$ et $x$ : $$\frac{U(t+dt)-U(t)}{dt}=\Phi_0-\Phi_{x}+\Phi_\text{lat}+P_\text{source}$$ avec $P_\text{source}=0$ car il n'y a pas de terme source et $\Phi_\text{lat}=0$. Or, $\displaystyle\frac{U(t+dt)-U(t)}{dt}=0$ car on est en régime permanent. Ainsi : $\Phi_x=\Phi_{x=0}$, ce qui se traduit par $Sj_x(x)=Sj_x(0)$ car la section est constante. On en déduit que $j_x(x)=j_x(0)$ est constant. D'après la loi de Fourier, on a donc $T$ affine en fonction de $x$.  En notant $L$ la longueur de notre solide (ou plutôt son épaisseur), on a : $$T(x)=T_0+(T_1-T_0)\frac{x}{L}$$ Ainsi, on a : $$j(x)=\lambda\frac{T_0-T_1}{L}$$ et on en déduit que $$\Phi=\frac{\lambda S}{L}(T_0-T_1)$$ On a donc $$R_\text{th}=\frac{L}{\lambda S} \ \ \text{et} \ \ G_\text{th}=\frac{\lambda S}{L}$$ avec $R_\text{th}$ en K.W$\inv$ et $G_\text{th}$ en W.K$\inv$. \\ \par La résistance thermique quantifie la qualité de l'isolation par le solide. \\ \par Avec cette expression, pour augmenter $R_\text{th}$, on peut :
\begin{itemize}
	\item augmenter l'épaisseur d'isolant (augmenter $L$) ;
	\item diminuer la surface d'échange (diminuer $S$) ;
	\item choisir un matériau plus isolant (diminuer $\lambda$).
\end{itemize}

\paragraph{Géométrie cylindrique}
Pour un cylindre de rayon $R_1$ et de hauteur $L$ dans un un cylindre de rayon $R_2$ ($R_2>R_1$) de hauteur $L$ et de conductivité thermique $\lambda$, on a en régime permanent si $T=T(r)$ : $$R_\text{th}=\frac{1}{2\pi\lambda L}\ln\left(\frac{R_2}{R_1}\right)\ \ \text{et}\ \ G_\text{th}=\frac{2\pi \lambda L}{\displaystyle\ln\left(\frac{R_2}{R_1}\right)}$$ Dans la limite $R_2=R_1+\varepsilon$ avec $\varepsilon\ll R_1$, on a $$R_\text{th}\approx\frac{\varepsilon}{2\pi LR_1\lambda}$$ et on reconnaît au dénominateur la surface, ce qui approxime l'expression qu'on obtient dans le cas cartésien. \\ \par Pour la démonstration, adapter celle du cas cartésien en prenant garde au fait que la section n'étant plus constante, $j_r$ dépend de $r$

\paragraph{Géométrie sphérique}
Pour une boule de rayon $R_1$ dans une boule de rayon $R_2$ ($R_2>R_1$) et de conductivité thermique $\lambda$, on a en régime permanent si $T=T(r)$ : $$R_\text{th}=\frac{1}{4\pi\lambda}\left(\frac{1}{R1}-\frac{1}{R_2}\right)$$ Dans la limite $R_2=R_1+\varepsilon$ avec $\varepsilon\ll R_1$, on a $$R_\text{th}\approx\frac{\varepsilon}{4\pi R_1^2\lambda}$$ et on reconnaît au dénominateur la surface, ce qui approxime l'expression qu'on obtient dans le cas cartésien.

\subsection{Analogie résistance thermique / résistance électrique / capacité}

\paragraph{Rappel} On rappelle (conductimétrie de sup et chapitre EM4 de spé) qu'on a, pour une résistance électrique de longueur $L$, de section $S$ et de résistivité électrique $\rho_\text{el}$, la formule : $$R_\text{el}=\rho_\text{el}\frac{L}{S}$$
\paragraph{Analogies} On a les analogies suivantes : 
\begin{center}
	\begin{tabular}{|c|c|}
		\hline
		Capacité&Résistance thermique \\
		\hline
		$Q=Cu$&$\displaystyle\Phi=\frac{1}{R_\text{th}}(T_0-T_1)$ \\
		\hline
		Potentiel électrique&Température\\
		\hline
		$Q$ : flux de $\varepsilon_0\overrightarrow{E}$&$\Phi$ : flux de $\overrightarrow{j_\text{th}}$ \\
		\hline
		Théorème de Gauss & Bilan de puissance \\
		\hline
		Condensateur plan&Résistance thermique "plane" \\
		\hline
		$C=\displaystyle\frac{\varepsilon_0 S}{L}$&$\displaystyle R_\text{th}=\frac{L}{\lambda S}$\\
		\hline
		$\displaystyle\frac{1}{C}$&$R_\text{th}$\\
		\hline
	\end{tabular}
\end{center}
Pour retenir tout cela, il est capital de retenir la dernière ligne du tableau. Pour la retenir plus facilement, on pourra penser aux impédances complexes.
\subsection{Résistance thermique conducto-convective}
Si $\overrightarrow{n}$ est le vecteur unitaire orienté du solide vers le fluide avec une interface de surface $S$, $T_S$ est la température du solide et $T_L$ est la température loin de l'interface dans le fluide, on rappelle que la loi de Newton fournit : $$\overrightarrow{j}=h(T_S-T_L)\overrightarrow{n}$$ Ainsi, on en déduit que $$\Phi=hS(T_S-T_L)$$ c'est-à-dire que $$R_\text{th}=\frac{1}{hS}$$
\subsection{Associations de résistances thermiques}
\paragraph{Association en série} Deux résistances thermiques en série $R1$ et $R_2$ en série sont équivalentes à une résistance thermique $R_\text{eq}=R_1+R_2$ (le flux se conserve). On généralise a $R_1$, ..., $R_n$ en série : $$R_\text{eq}=\sum_{i=1}^{n}R_i$$
\paragraph{Association en parallèle} Deux résistances thermiques en série $R1$ et $R_2$ en parallèle sont équivalentes à une résistance thermique $R_\text{eq}=\displaystyle\frac{R_1R_2}{R_1+R_2}$ (le flux se sépare selon la loi des nœuds et les températures en entrée et en sortie ne changent pas). On généralise a $R_1$, ..., $R_n$ en parallèle : $$\frac{1}{R_\text{eq}}=\sum_{i=1}^{n}\frac{1}{R_i}$$
\paragraph{Application : Résistance thermique d'un igloo} on considère un sol parfaitement isolant. Un igloo est formé par une demi-sphère de rayon $R$ et une demi-sphère de rayon $R+e$. La neige qui forme la coquille a une conductivité thermique $\lambda$. On observe des transferts conducto-convectifs entre l'igloo et l'air extérieur avec un coefficient $h_\text{ext}$ et un transfert conducto-convectif entre l'igloo et l'air intérieur avec un coefficient $h_\text{int}$. Calculer la résistance thermique totale $R_\text{th}$ de l'igloo. \\ \par Il s'agit en fait de l'association en série de deux résistances conducto-convectives et d'une résistance thermique due à la neige de la coquille. Calculons celle de la coquille : $$R_{]R,R+e[}=\int_{R}^{R+e}\delta R$$ avec $\delta R$ la résistance thermique de la coquille située entre $r$ et $r+dr$. Or, $$\delta R=\frac{dr}{2\pi\lambda r^2}$$ donc on a : $$R_{]R,R+e[}=\frac{1}{2\pi\lambda}\left(\frac{1}{R}-\frac{1}{R+e}\right)$$ le résultat final de la résistance thermique totale de l'igloo est donc, puisqu'il s'agit d'une association en série : $$R_\text{th}=\frac{1}{2\pi}\left(\frac{1}{h_\text{int}R^2}+\frac{1}{\lambda}\left(\frac{1}{R}-\frac{1}{R+e}\right)+\frac{1}{h_\text{ext}(R+e)^2}\right)$$
\paragraph{Application : Intérêt du double vitrage} On place du gaz entre deux épaisseurs de vitrage. Ce gaz (du vide dans l'idéal) a une grande résistance thermique ce qui permet d'isoler de façon très performante si on néglige le rayonnement. Afin d'éviter le rayonnement, on réalise le vitrage de façon à ce que le rayonnement visible puisse traverser le vitrage mais de façon à ce que le rayonnement infrarouge (celui qui transporte le plus d'énergie thermique) soit stoppé.

\subsection{Résolution avec source d'énergie}

\paragraph{Exemple 1} On considère un cylindre isolé de l'extérieur de résistance électrique $R$, de conductivité thermique $\lambda$, de section $S=\pi b^2$ et de longueur $L$. Il est parcouru par une intensité électrique $I$ et l'extérieur est à la température $T_0$. On suppose $T=T(x)$. On veut :
\begin{enumerate}
	\item Calculer $\delta R$ entre $x$ et $x+dx$ ;
	\item Le graphe de $T(x)$ ;
	\item L'expression de $T(x)$.
\end{enumerate}
Voici la résolution.
\begin{enumerate}
	\item Puisqu'il s'agit d'une association en série, on a : $$\delta R=\frac{R}{L}dx$$
	\item On veut que cela soit une parabole dans le cylindre, qui atteint son maximum au milieu du cylindre par symétrie. On veut que la température soit constante égale à $T_0$ à l'extérieur. Il y a des points anguleux en $x=0$ et $x=L$.
	\item On effectuer un bilan de puissance à la portion de résistance entre $x$ et $x+dx$ :\begin{align*}
		 \underbrace{\frac{U(t+dt)-U(t)}{dt}}_{=0 \ \text{(régime permanent)}}&=\Phi_{\text{reçu en }x}+\Phi_{\text{reçu en }x+dx}+\underbrace{\Phi_\text{latéral}}_{=0}+P_\text{Joule} \\
		 &= \pi b^2j_x(x)-\pi b^2j_x(x+dx)+\frac{R}{L}I^2dx
	\end{align*} Par conséquent, $\displaystyle\frac{dj_x}{dx}=\frac{RI^2}{\pi b^2L}$. Or, par symétrie, on a nécessairement $j_x\displaystyle\left(\frac{L}{2}\right)=0$ donc on en déduit que $$j_x(x)=\frac{RI^2}{\pi b^2L}\left(x-\frac{L}{2}\right)$$ Or, d'après la loi de Fourier, on a : $j_x=\displaystyle -\lambda\frac{dT}{dx}$. En intégrant une fois de plus et en utilisant le fait que $T(0)=T_0$, on obtient finalement l'expression : $$T(x)=-\frac{RI^2}{2\pi b^2L}x(x-L)+T_0$$
\end{enumerate}

\paragraph{Exemple 2}
On considère un cylindre de longueur $L$, de rayon $b$, de résistance électrique $R$ et de conductivité thermique $\lambda$. On considère que ses faces gauche et droite (celles circulaires) sont isolées thermiquement et qu'il effectue un échange thermique avec l'air extérieur à la température $T_0$ avec un coefficient conducto-convectif $h$. Il est parcouru par une intensité électrique $I$. On suppose cette fois que $T=T(r)$. Exprimer $T(b^-)$ et donner l'équation vérifiée par $T$. \\ \\ On effectue un bilan de puissance à l'ensemble de la résistance : 
$$\underbrace{\frac{U(t+dt)-U(t)}{dt}}_{=0 \ \text{(régime permanent)}}=\underbrace{\Phi_{\text{gauche}}}_{=0}+\underbrace{\Phi_{\text{droite}}}_{=0}+\Phi_\text{latéral}+\underbrace{P_\text{Joule}}_{=RI^2}$$
Or, d'après la loi de Newton, $\Phi_\text{latéral}=-2\pi bLh(T-T_0)$. On en déduit que $$T(b^-)=T_0+\frac{RI^2}{2\pi bLh}$$ Ensuite, la résistance d'une portion de cylindre comprise entre les rayons $0$ et $r$ vaut : $$R_{[0,r]}=R\frac{b^2}{r^2}$$ (résistance en parallèle). La puissance dissipée par effet Joule vaut : $$P_\text{Joule}=\frac{U^2}{R}=RI^2$$ et sur une telle portion : $$P_{\text{Joule, }[0,r]}=\frac{U^2}{R_{[0,r]}}$$ En faisant le quotient pour éliminer $U^2$, on trouve que, sur une telle portion : $$P_{\text{Joule, }[0,r]}=\frac{r^2RI^2}{b^2}$$On effectue ensuite un bilan de puissance à la portion de résistance entre $0$ et $r$, qui, toujours en régime permanent, donne : $$0=-2\pi rLj_r+\frac{r^2RI^2}{b^2}$$ Par conséquent, on en déduit que $$j_r=\frac{RI^2}{2\pi Lb^2}r$$ Or, d'après la loi de Fourier : $j_r=\displaystyle -\lambda\frac{dT}{dr}$. Par conséquent, $T$ vérifie l'équation : $$\frac{dT}{dr}+\frac{RI^2}{2\pi\lambda Lb^2}r=0$$ On peut ensuite intégrer pour obtenir l'expression de $T$ en fonction de $r$ et $T(b^-)$.
\paragraph{Remarque} Remarquer comment nous avons eu des approches totalement différentes pour les deux exemples en changeant simplement quelques données dans le problème.
\section{Compléments}
\subsection{Production d'entropie par conduction thermique : entropie massique créée par unité de temps}
On la détermine à partir d'une application du deuxième principe et d'un bilan thermique. Pour simplifier, on suppose que $T=T(x,t)$. On s'intéresse à un cylindre se situant entre des abscisses $x$ et $x+dx$, de section $\Sigma$, de capacité thermique massique $c$, de masse volumique $\rho$ et de conductivité thermique $\lambda$. Le deuxième principe de la thermodynamique infinitésimal appliqué à ce système entre $t$ et $t+dt$ donne : $$dS=\delta S_\text{créée}+\delta S_{\text{échangée}}$$ Or, on sait aussi que $$dS=\Sigma dx\times \rho c\frac{dT}{T}$$ De plus, on a :
\begin{align*}
	\frac{\delta S_\text{échangée}}{dt}&=\frac{\Phi_{\text{reçu en }x}}{T_x}+\frac{\Phi_{\text{reçu en }x+dx}}{T_{x+dx}} \\
	&= \frac{j_x(x)\Sigma}{T_x}-\frac{j_x(x+dx)\Sigma}{T_{x+dx}}
\end{align*}
Par conséquent, on en déduit d'après la loi de Fourier que $$\frac{\delta S_\text{échangée}}{dt} =-\Sigma\lambda\left(\frac{1}{T(x)}\frac{\partial T}{\partial x}(x)-\frac{1}{T(x+dx)}\frac{\partial T}{\partial x}(x+dx)\right)$$ Montrons que $\delta S_\text{créée}\geqslant0$. Puisqu'on sait que $$dS=\Sigma dx\rho c\frac{1}{T}\frac{\partial T}{\partial t}dt$$ on calcule brutalement en utilisant les expressions précédentes : \begin{center}
\begin{align*}
	\delta S_\text{créée}&=dS-\delta S_\text{échangée} \\
	&= \left[\frac{\Sigma dx\rho c}{T}\frac{\partial T}{\partial t}+\Sigma\lambda\left(\frac{1}{T(x)}\frac{\partial T}{\partial x}(x)-\frac{1}{T(x+dx)}\frac{\partial T}{\partial x}(x+dx)\right)\right]dt\\
	&= \Sigma \lambda\left[\frac{dx}{T}\frac{\rho c}{\lambda}\frac{\partial T}{\partial t}+\frac{1}{T(x)}\frac{\partial T}{\partial x}(x)-\frac{1}{T(x+dx)}\frac{\partial T}{\partial x}(x+dx)\right]dt \\
	&= \Sigma \lambda\left[\frac{dx}{T}\Delta T+\frac{1}{T(x)}\frac{\partial T}{\partial x}(x)-\frac{1}{T(x+dx)}\frac{\partial T}{\partial x}(x+dx)\right]dt \\
	&= \Sigma \lambda\left[\frac{\Delta T}{T}dx+-\frac{\partial\left(\frac{1}{T}\frac{\partial T}{\partial x}\right)}{\partial x}dx\right]dt \\
	&=\Sigma\lambda \times\frac{\left(\frac{\partial T}{\partial x}\right)^2}{T^2} dtdx \\
	&\geqslant 0 \\
	&>0\ \ \text{si T n'est pas uniforme et }\lambda\ne0
\end{align*}
\end{center}
\subsection{Un exemple de résolution numérique approchée de l'équation de la chaleur}
Le but est de trouver une solution numérique approchée de l'équation de la chaleur : $$\forall t>0,\ \forall z\in]0,L[,\ \frac{\partial T}{\partial t}(z,t)=D\frac{\partial^2 T}{\partial z^2}(z,t)$$ avec la condition initiale $$\forall z\in]0,L[,\ T(x,t=0)=20^\circ\text{C}$$ et les conditions aux limites : $$\forall t>0,\ T(z=0,t)=T(z=L,t)=10^\circ\text{C}$$ Pour cela, on commence par procéder à une discrétisation du problème. On pose $$z_i=ia \ \ \text{et} \ \ t_j=j\tau$$ avec $a$ le pas spatial et $\tau$ le pas temporel. On veut calculer les $T_{i,j}$. Idéalement,, on souhaiterait que $T_{i,j}\approx T(z_i,t_j)$. On remplace les dérivées partielles par des "différences finies" : $$\frac{\partial T}{\partial t}(z_i,t_j)\longrightarrow\frac{T_{i,j+1}-T_{i,j}}{\tau}$$ $$\frac{\partial^2 T}{\partial z^2}\longrightarrow\frac{T_{i+1,j}-2T_{i,j}+T_{i-1,j}}{a^2}$$ (effectuer des DL en ces points pour vérifier la cohérence). Par conséquent, l'équation de la chaleur est approchée par l'équation : $$\frac{T_{i,j+1}-T_{i,j}}{\tau}=D\frac{T_{i+1,j}-2T_{i,j}+T_{i-1,j}}{a^2}$$ En posant $\displaystyle\alpha=\frac{D\tau}{a^2}$, on peut enfin réécrire cette équation sous la forme : $$T_{i,j+1}=\alpha T_{i-1,j}+(1-2\alpha)T_{i,j}+\alpha T_{i+1,j}$$ Cette équation appelle une résolution numérique : si on affiche la résolution au cours du temps ($j$ varie), on observera une propagation de la chaleur. Remarquons que si $0<\alpha<\displaystyle\frac{1}{2}$, alors le système est stable et $T_{i,j+1}\leqslant\max(T_{i-1,j},T_{i,j},T_{i+1,j})$ : logique, la température diminue à cause de la condition aux bords.

\part{Thermochimie}
\chapter{Enthalpie libre ; potentiel chimique}
\localtableofcontents
\section{Notion de potentiel thermodynamique ; enthalpie libre}
\subsection{Potentiels thermodynamiques}
\paragraph{Analogie mécanique} Lors d'une évolution "spontanée" (\textit{ie} sans apport extérieur), l'énergie mécanique d'un système ne peut que diminuer, c'est l'analogie à connaître.
\paragraph{Potentiel thermodynamique} Un potentiel thermodynamique est une fonction des paramètres du système et des "réservoirs" (thermostats et barostats), et qui ne peut que décroître.
\paragraph{Exemple 1 : Néguentropie} On considère un système isolé. Dans ce cas, $S$ ne peut qu'augmenter en vertu du deuxième principe. Ainsi, $-S$ est un potentiel thermodynamique, qu'on appelle \textbf{néguentropie}.
\paragraph{Exemple 2 : Énergie libre} On considère tout d'abord l'évolution isochore et monotherme (thermostat à la température $T_0$) d'un système fermé, sans variation d'énergie mécanique et sans travail utile. \\ \par Alors, en vertu du premier et du deuxième principe, on obtient : $$(U_F-T_0S_F)-(U_I-T_0S_I)=-T_0S_\text{créée}\leqslant0$$
Ainsi, $U-T_0S$ est un potentiel thermodynamique \textbf{pour une évolution monotherme et isochore}. \\ \par De la même façon, on montre que $$F=U-TS$$ est un potentiel thermodynamique \textbf{pour une évolution isochore et isotherme}. On l'appelle l'\textbf{énergie libre}. 
\paragraph{Exemple 3 : Enthalpie libre} On considère une transformation isotherme et isobare d'un système fermé, sans variation d'énergie mécanique et sans travail utile. \\ \par Alors, en vertu du premier principe pour une isobare et du deuxième principe, on obtient que $$G=H-TS=U+PV-TS$$ est un potentiel thermodynamique \textbf{pour une transformation isotherme et isobare}, appelé \textbf{enthalpie libre}. C'est une fonction d'état extensive qui s'exprime en Joule.
\paragraph{Résumé} Voici un résumé des différents potentiels thermodynamiques intéressants selon les différents types de transformations \textbf{pour un système fermé} : 
\begin{center}
	\begin{tabular}{|c|c|c|}
		\hline
		Type de transformation&Nom du potentiel&Expression \\
		\hline
		Adiabatique&Néguentropie&$-S$ \\
		\hline
		Isochore et isotherme&Énergie libre&$F=U-TS$\\
		\hline
		Isobare et isotherme&Enthalpie libre&$G=H-TS$\\
		\hline
	\end{tabular}
\end{center}
Le potentiel thermodynamique qui nous intéressera le plus est l'enthalpie libre grâce à ses applications en chimie, puisque la plupart des réactions chimiques peuvent être considérées comme isobares et isothermes.
\subsection{Identités thermodynamiques pour un système de composition constante}
\paragraph{Constituant physico-chimique} On appelle constituant physico-chimique une espèce chimique dans un état physique précis.
\paragraph{Composition constante} On dit d'un système qu'il est de composition constante lorsque pour tout constituant physico-chimique $i$, la quantité de matière de ce constituant physico-chimique est une constante. Cela implique notamment que les constituants physico-chimiques ne changent pas d'état physique au cours de la transformation du système.
\paragraph{Première identité thermodynamique} Pour un système de composition constante, une transformation réversible ou infinitésimale et de variance $2$, alors on a : $$dU=TdS-PdV$$
\paragraph{Deuxième identité thermodynamique} Pour un système de composition constante, une transformation réversible ou infinitésimale et de variance $2$, alors on a : $$dH=TdS+VdP$$
\paragraph{Troisième identité thermodynamique} Pour un système de composition constante, une transformation réversible ou infinitésimale et de variance $2$, alors on a : $$dG=VdP-SdT$$
\subsection{Cas des systèmes de composition variable}
On note $n_i$ les quantités de matière des différents constituants physico-chimiques. On a : $$dG=VdP-SdT+\sum_i \frac{\partial G}{\partial n_i}\bigg|_{T,P,n_{j\ne i}}dn_i$$ et on note $$\mu_i:=\frac{\partial G}{\partial n_i}\bigg|_{T,P,n_{j\ne i}}$$ le \textbf{potentiel chimique} du constituant physico-chimique $i$. Il s'exprime en J/mol et est une fonction d'état intensive. La relation précédente se réécrit donc $$dG=VdP-SdT+\sum_i \mu_i dn_i$$ On remarquera que l'on a alors : $$V=\frac{\partial G}{\partial P}\bigg|_{T,n_i} \ \ \text{et}\ \ -S=\frac{\partial G}{\partial T}\bigg|_{P,n_i}$$ De même, puisqu'on avait $dU=TdS-PdV$, on remarquera que l'on a les relations : $$T=\frac{\partial U}{\partial S}\bigg|_V \ \ \text{et} \ \ -P=\frac{\partial U}{\partial V}\bigg|_S$$
\subsection{Cas du corps pur monophasé}
On se place ici dans le cas d'un unique constituant physico-chimique. On a alors : $$G(T,P,n)=nG_m(T,P)$$ car $G$ est une fonction d'état extensive. Dans le cas d'un corps pur monophasé, on note $\mu\st$ son potentiel chimique, et on a alors : $$\mu\st(T,P)=G_m(T,P)$$ Et puisque $dG_m=V_mdP-S_mdT$, on a alors les relations : $$-S_m=\frac{\partial \mu\st}{\partial T}\bigg|_P \ \ \text{et}\ \ V_m=\frac{\partial \mu\st}{\partial P}\bigg|_T$$
\subsection{Application au changement d'état d'un corps pur ; formule de Clapeyron}
On se place dans le cas où on est en présence de deux constituants physico-chimiques de la même espèce sous forme liquide $l$ et sous forme gazeuse $g$. On a alors par extensivité : $$G=n_lG_{m,l}(T,P)+n_gG_{m,g}(T,P)$$ Lors d'une évolution isotherme et isobare (un changement d'état par exemple tiens !), on sait que $dG\leqslant 0$ puisque $G$ est alors un potentiel thermodynamique. Or, puisque $T$ et $P$ sont alors constantes, on a : $$dG=G_{m,l}(T,P)dn_l+G_{m,g}(T,P)dn_g$$ avec $G_{m,l}(T,P)=\mu_L\st$ et $G_{m,g}(T,P)=\mu_G\st$. Or, $n_l+n_g$ est une constante par conservation de la matière, donc $dn_l=-dn_g$. On en déduit que $$dG=dn_l(\mu_l\st-\mu_g\st)$$
\paragraph{Démonstration de la formule de Clapeyron} Si on prend $T$ et $P$ tels que $\mu_l\st>\mu_g\st$, alors on a nécessairement $dn_l\leqslant0$ si bien qu'il y a alors une vaporisation. \\ \par Précisément, sur la courbe de vaporisation, on a : $\mu_l\st(T,P)=\mu_g\st(T,P)$. En différenciant cette relation, on obtient : $$\frac{\partial \mu_l\st}{\partial T}\bigg|_PdT+\frac{\partial \mu_l\st}{\partial P}\bigg|_TdP=  \frac{\partial \mu_g\st}{\partial T}\bigg|_PdT+\frac{\partial \mu_g\st}{\partial P}\bigg|_TdP$$ ce qui se réécrit en vertu de la fin de la sous-partie précédente : $$-S_{m,l}dT+V_{m,l}dP=-S_{m,g}dT+V_{m,g}dP$$ On en déduit la \textbf{formule de Clapyeron} : 
\begin{center}
	\begin{align*}
		\frac{dP_\text{sat}}{dT}&=\frac{S_{m,g}-S_{m,l}}{V_{m,g}-V_{m,l}} \\
		&=\frac{H_{m,g}-H_{m,l}}{T(V_{m,g}-V_{m,l})}
	\end{align*}
\end{center} qu'on préfère écrire en thermodynamique sous la forme : $$\frac{dP_\text{sat}}{dT}=\frac{\Delta_\text{vap}h}{T\Delta_\text{vap}v}$$
\subsection{Identité d'Euler}
On considère une fonction de deux variables $f:(x,y)\mapsto f(x,y)$ qui vérifie $f(\alpha a,\alpha b)=\alpha f(a,b)$. En lui appliquant la règle de la chaîne en dérivant par rapport à $\alpha$, on obtient : $$a\partial_x f(\alpha a,\alpha b)+b\partial_y f(\alpha a,\alpha b)=f(a,b)$$ ce qui donne, après évaluation en $\alpha=1$ : $$a\partial_x f(a,b)+b\partial_y f(a,b)=f(a,b)$$
Dans le cas de l'enthalpie libre pour une transformation isobare et isotherme, cela donne :
\begin{center}
	\begin{align*}
		G(T,P,n_i)&=\sum_i n_i\frac{\partial G}{\partial n_i}\bigg|_{T,P,n_{j\ne i}} \\
		&=\sum_i n_i\mu_i(T,P,n_{i})
	\end{align*}
\end{center}
\section{Expressions du potentiel chimique}
Le but de cette section est de montrer que $\mu_i=\mu_i^\circ+RT\ln(a_i)$ avec $a_i$ l'activité du constituant physico-chimique $i$.
\subsection{État standard d'un constituant physico-chimique}
\paragraph{Cas d'un gaz} Pour un gaz se comportant comme un gaz parfait, l'état standard de ce gaz correspond au fait que sa pression partielle $P$ vaut $P=P^\circ$ la pression standard telle que $P^\circ=1$ bar. \\ \par Il y a donc autant d'états standards que de températures pour un gaz. 
\paragraph{Cas d'un soluté} Pour un soluté "extrapolé" des solutions diluées, l'état standard de ce soluté (état théorique jamais atteint en pratique) est lorsque sa concentration vérifie $c_i=c^\circ$ où $c^\circ$ est la concentration standard valant $c^\circ=1$ mol/L.
\paragraph{Cas d'une phase condensée} Pour une phase condensée, l'état standard de cette phase condensée correspond au fait qu'elle soit seule dans sa phase et que $P=P^\circ$.
\subsection{Potentiel chimique d'un gaz parfait}
Pour un gaz parfait, on a par définition des différentielles et d'après l'équation d'état : 
\begin{center}
	\begin{align*}
		\mu(T,P)&=\mu^\circ(T,P^\circ)+\int_{P^\circ}^{P}\frac{\partial \mu}{\partial P}\bigg|_TdP \\
		&= \mu^\circ(T,P^\circ)+\int_{P^\circ}^{P}V_m\ dP \\
		&= \mu^\circ(T,P^\circ)+\int_{P^\circ}^{P}\frac{RT}{P}\ dP \\
		&= \mu^\circ(T,P^\circ)+RT\ln\left(\frac{P}{P^\circ}\right)
	\end{align*}
\end{center}
Cela est bien en accord, dans le gaz du gaz parfait (et donc du gaz dans son état standard) avec le résultat annoncé en début de partie. \\ \par On admet la généralisation pour un mélange idéal de gaz parfaits : $$\mu_i(T,P)=\mu_i^\circ(T)+RT\ln\left(\frac{P_i}{P^\circ}\right)$$ avec $\mu_i$ le potentiel chimique du gaz parfait $i$ et $P_i$ la pression partielle du gaz parfait $i$.
\subsection{Potentiel chimique d'une phase condensée}
De même, on a $$\mu(T,P)=\mu^\circ(T,P^\circ)+\int_{P^\circ}^{P}V_mdP$$ Or, pour une phase condensée, le volume molaire est négligeable devant le volume molaire de la même espèce à l'état gazeux (on a environ un facteur 1000 entre les deux). Et on en déduit/constate par les observations expérimentales que $\displaystyle\int_{P^\circ}^{P}V_mdP$ est négligeable devant $\mu^\circ(T)$. Ainsi, pour une phase condensée : $$\mu(T,P)\approx\mu^\circ(T)$$
\paragraph{Mélange de phases condensées} Pour un mélange de phases condensées (un alliage par exemple), on admet que $$\mu_i(T,P)=\mu_i^\circ(T)+RT\ln(x_i)$$ avec $\mu_i$ le potentiel chimique du constituant $i$ de l'alliage et $x_i$ la fraction molaire en constituant $i$ de l'alliage.
\subsection{Potentiel chimique d'un soluté dans une solution diluée idéale}
On admet que $$\mu_i(T,P)=\mu_i^\circ(T)+RT\ln\left(\frac{c_i}{c^\circ}\right)$$ avec $c_i$ la concentration molaire du soluté et $c^\circ$ la concentration standard qui vaut 1 mol/L.
\subsection{Activité d'un constituant physico-chimique}
Dans tous les cas, on a : $$\mu_i(T,P)=\mu_i^\circ(T)+RT\ln(a_i)$$ où $a_i$ est \textbf{l'activité} du constituant physico-chimique $i$. C'est une grandeur sans dimension. \\ \par Pour un gaz parfait seul, $a=\displaystyle\frac{P}{P^\circ}$ \\ \par Pour un mélange idéal de gaz parfaits, $a_i=\displaystyle\frac{P_i}{P^\circ}$ avec $P_i$ la pression partielle du gaz $i$. \\ \par Pour une phase condensée seule ou pour le solvant d'une solution diluée, $a=1$. \\ \par Pour un solution dans une solution diluée $a_i=\displaystyle\frac{c_i}{c^\circ}$ avec $c_i$ sa concentration et $c^\circ$ la concentration standard.

\chapter{Évolution d'un système lors d'une réaction chimique}
\localtableofcontents
\section{Condition d'évolution spontanée lors d'une réaction chimique}
On travaille sur une \textbf{unique} réaction, caractérisée par l’avancement $\xi$ (variable de De Donder). On la note $$\sum_i \nu_iA_i=0$$ avec $A_i$ les constituants physico-chimiques et $\nu_i$ les coefficients stœchiométriques algébriques (négatifs pour les réactifs).
\subsection{Enthalpie libre de réaction ; opérateur de Lewis}
\paragraph{Enthalpie libre de réaction} On la définit par $$\Delta_rG(T,P,\xi)=\frac{\partial G}{\partial\xi}$$ Elle dépend de l’équation de réaction (coefficients stœchiométriques) et s'exprime en J/mol. On a $dG\leqslant0$ pour une isobare et isotherme donc $\Delta_rG\leqslant0$.
\paragraph{Expression pendant une réaction}  Pour l’espèce $s$ : $n_s=n_{s,0}+\nu_i\xi$. On a alors en dérivant l'identité d'Euler $\displaystyle G=\sum_i n_i\mu_i$ par rapport à $\xi$ : $$\Delta_rG=\sum_i\nu_i\mu_i$$ Pour une réaction isotherme et isobare, on a : $$dG=\Delta_rG\times d\xi$$ 
\paragraph{Opérateur de Lewis}
On généralise à toute grandeur extensive $X$ grâce à l'identité d'Euler : $$\Delta_rX=\frac{\partial X}{\partial\xi}\bigg|_{T,P}=\sum_i\nu_i\frac{\partial X}{\partial n_i}\bigg|_{T,P,n_{j\ne i}}$$ Par exemple, on peut écrire : $$\Delta_rG=\Delta_rH-T\Delta_rS$$
\subsection{Critère d'évolution d'un système lors d'une récation chimique isotherme et isobare}
\paragraph{Critère d'évolution} Pour une réaction isotherme et isobare, $dG\leqslant0$ donc, comme pour une réaction isotherme et isobare, $dG=\Delta_rGd\xi$, on en en déduit que $$\Delta_rGd\xi\leqslant0$$ Ainsi : \begin{itemize}
	\item si $\Delta_rG<0$, la réaction se produit dans le sens direct ;
	\item si $\Delta_rG>0$, la réaction se produit dans le sens indirect ;
	\item si $\Delta_rG=0$, il y a équilibre.
\end{itemize}
La bonne analogie consiste à voir $G$ du milieu comme une énergie potentielle de réaction et $\Delta_rG$ comme l'opposé d'une force.
\subsection{Arrêt de la réaction chimique : équilibre chimique, réaction totale ou blocage cinétique}
Attention aux blocages cinétiques et aux équilibres inatteignables (solides). Vérifier que $\xi_{eq}\in\left[\xi_\text{min},\xi_\text{max}\right]$. \\ \par Une réponse usuelle à la question "Pourquoi ça ne marche pas ?" en fin de sujet de chimie après des tonnes de calculs de thermodynamique bourrins est souvent "car il peut y avoir un blocage cinétique".
\subsection{Enthalpie libre standard de réaction ; quotient de réaction}
\paragraph{Enthalpie libre standard de réaction} On définit l'enthalpie libre standard de réaction par : $$\Delta_rG^\circ=\sum_i\mu_i^\circ\nu_i$$
\paragraph{Quotient de réaction} On définit le quotient de réaction par : $$Q=\prod_i a_i^{\nu_i}$$
\paragraph{Lien avec le quotient de réaction} On a : $$\Delta_rG=\Delta_rG^\circ+RT\ln(Q)$$ En effet, $\mu_i=\mu_i^\circ+RT\ln(a_i)$ donc $\Delta_rG=\displaystyle\sum_i\nu_i\mu_i=\sum_i\nu_i\mu_i^\circ+RT\ln(Q)$
\section{Calcul des grandeurs standards de réaction}
\subsection{Entropie et enthalpie standards de réaction}
On a : $$\Delta_rG^\circ=\Delta_rH^\circ-T\Delta_rS^\circ$$ avec $\Delta_rH^\circ$ l'enthalpie standard de réaction et $\Delta_rS^\circ$ l'entropie standard de réaction.
\paragraph{Enthalpie standard de formation} On définit l'état de référence comme celui du corps pur simple monophasé le plus stable à la température $T$ et à une pression de $1$ bar. Par exemple, pour l'oxygène, on considère $O_{2(g)}$ à $25^\circ$C. Pour un composant $C$, l'enthalpie standard de formation $\Delta_fH^\circ(C)$ est l'enthalpie de la réaction :
\begin{center}
	composants de $C$ dans leur état standard $\longrightarrow$ C
\end{center}
\paragraph{Entropie standard de réaction} Pas besoin de définir un état arbitraire de référence (principe de Nernst). On a donc simplement, avec les $S_m^\circ$ les entropies molaires absolues, pour $\alpha A\to\beta B$ :$$\Delta_rS^\circ =	-\alpha S_m^\circ(A)+\beta S_m^\circ(B)$$
\subsection{Variations avec la température ; approximation d'Ellingham}
\paragraph{Approximation d'Ellingham} Dans l'approximation d'Ellingham : 
\begin{itemize}
	\item $\Delta_rH^\circ$ est indépendant de $T$
	\item $\Delta_rS^\circ$ est indépendant de $T$
\end{itemize} car on considère que tant que la réaction a lieu, les réactifs et les produits n'ont pas changé d'état.
\paragraph{Variation avec la température} On a alors la \textbf{relation de Gibbs-Helmoltz} : $$\frac{d}{dt}\left(\frac{\Delta_rG^\circ}{T}\right)=-\frac{\Delta_rH^\circ}{T^2}$$ et aussi la relation : $$\frac{d\Delta_rG^\circ}{dT}=-\Delta_rS^\circ$$ En effet, dans l'approximation d'Ellingham, il suffit alors de dériver par rapport à la température la relation $$\Delta_rG^\circ=\Delta_rH^\circ-T\Delta_rS^\circ$$ et de dériver son analogue divisé par $T$ puisque l'on considère que $\Delta_rH^\circ$ et $\Delta_rS^\circ$ sont indépendantes de $T$. On admet que cette relation reste valable hors de l'approximation d'Ellingham.
\subsection{Utilisation de grandeurs tabulées ; loi de Hess}
\paragraph{Loi de Hess} Pour $\alpha A\to\beta B$, l’enthalpie standard de réaction est donné par $$\Delta_rH^\circ=-\alpha\Delta_fH^\circ(A)+\beta\Delta_fH^\circ(B)$$ (on éclate tous les composants pour les recréer). On a une variante identique pour $G$. Aussi, on a ce que l'on pourrait appeler une "loi de Hess pour $S$" : $$\Delta_rS^\circ=\sum_i\nu_iS_{m,i}^\circ$$ avec $S_{m,i}^\circ$ l'entropie molaire absolue du constituant physico-chimique $i$.
\subsection{Exemples}
On donne 
\begin{center}
	\begin{tabular}{|c|c|c|c|c|}
		\hline
		&NH$_{3(g)}$&N$_{2(g)}$&O$_{2(g)}$&H$_{2(g)}$ \\
		\hline
		$\Delta_fH^\circ$&$\alpha$&$\beta(=0)$&$\gamma(=0)$&$\delta(=0)$ \\
		\hline
		$S_m^\circ$&a&b&c&d \\
		\hline
	\end{tabular}
\end{center}
Exprimer les entropies et les enthalpies standards des réactions : 
\begin{center}
	$\text{N}_{2(g)}+3\text{H}_{2(g)}=2\text{NH}_{3(g)}$ (réaction 1)
\end{center}
\begin{center}
	$\displaystyle \frac12 \text{N}_{2(g)}+\frac32 \text{H}_{2(g)}=\text{NH}_{3(g)}$ (réaction 2)
\end{center}
D'après la loi de Hess, on a alors : $$\begin{cases}
	\Delta_rH_1^\circ= 2\alpha\\
	\Delta_rH_2^\circ= \alpha \\
	\Delta_rS_1^\circ= 2a-b-3d \\
	\displaystyle\Delta_rS_1^\circ= a-\frac12 b-\frac32 d
\end{cases}$$
\section{Interprétation et utilisation des grandeurs standards de réaction}
\subsection{Constante d'équilibre standard ; loi de Van't Hoff}
\paragraph{Constante d'équilibre standard}
On rappelle que $\Delta_rG=\Delta_rG^\circ+RT\ln(Q)$. Or, à l'équilibre, $\Delta_rG=0$ et $Q=K^\circ$, où $K^\circ$ est la constante de réaction standard. On en déduit alors que $$0=\Delta_rG^\circ+RT\ln(K^\circ)$$ ou autrement dit $$K^\circ=\exp\left(-\frac{\Delta_rG^\circ}{RT}\right)$$ Cette dernière ressemble \textbf{visuellement} à la loi d'Arrhénius, mais n'a rien à voir avec celle-ci. Cela peut cependant être un bon moyen de retenir la formule.
\paragraph{Loi de Van't Hoff} Réciproquement, puisque $\Delta_rG^\circ=-RT\ln(K^\circ)$ et d'après la relation de Gibbs-Helmoltz $\displaystyle\frac{d}{dt}\left(\frac{\Delta_rG^\circ}{T}\right)=-\frac{\Delta_rH^\circ}{T^2}$, on obtient la \textbf{loi de Van't Hoff} : $$\frac{d}{dT}\left(\ln(K^\circ)\right)=\frac{\Delta_rH^\circ}{RT^2}$$
\subsection{Évolution vers un équilibre ou réaction totale ?}
Il faut bien comprendre que $\Delta_rG\ne\Delta_rG^\circ$. \\ \par Lorsque $\xi_{eq}\in\left[\xi_\text{min},\xi_\text{max}\right]$, il y a équilibre (sauf en cas de blocage cinétique). Dans le cas contraire, la réaction est totale (sauf en cas de blocage cinétique là encore).
\subsection{Effets thermiques associés à la réaction ; interprétation du signe de $\Delta_rH^\circ$}
\paragraph{Utilisation de la loi de Van't Hoff} Lors d'une réaction isotherme et isobare, on a $$\displaystyle Q=\Delta H=\int_{\xi_0}^{\xi_\infty}\Delta_rH\ d\xi$$ Or $\Delta_rH$ varie peu donc on peut considérer que $\Delta_rH\approx\Delta_rH^\circ$ et on a alors $$Q=\left(\xi_\infty-\xi_0\right)\Delta_rH^\circ$$. On a alors différents cas : 
\begin{itemize}
	\item si $\Delta_rH^\circ>0$, la réaction dans le sens direct est endothermique ;
	\item si $\Delta_rH^\circ<0$, la réaction dans le sens direct est exothermique ;
	\item si $\Delta_rH^\circ\approx0$, la réaction dans le sens direct est athermique.
\end{itemize}
En général, $\Delta_rH^\circ$ est de l'ordre de $100$ kJ.mol$\inv$ (énergie de liaison). Pour une réaction \textbf{adiabatique et isobare}, on effectue d'abord fictivement la réaction, puis dans un second temps le transfert thermique. On a alors $\Delta H=Q=0$, et l'expression obtenue par le chemin fictif nous permet d'obtenir la température finale en fonction des données.
\subsection{Interprétation et prévision du signe de $\Delta_rS^\circ$}
\paragraph{Interprétation de $\Delta_rS$} On a $\Delta_rS^\circ=\displaystyle\sum_i\nu_iS_{m_i}$. Or, $S_{m,i}$ est négligeable si $i$ n'est pas un gaz et est de l'ordre de grandeur de $200$ J.mol$\inv$.K$\inv$ si $i$ est un gaz. On a donc : $$\Delta_rS^\circ\approx\Delta\nu_{gaz}\times\text{C}$$ avec C$\approx$200 J/K/mol.
\paragraph{Exemple} Si on considère la réaction $$3\text{H}_{2(g)}+1\text{N}_{2(g)}=2\text{NH}_{3,(g)}$$ on a alors $\Delta\nu_\text{gaz}=-2$, si bien que $$\Delta_rS^\circ=-2C<0$$
\chapter{Déplacement et rupture d'équilibre}
\localtableofcontents
\section{Caractérisation de l'état intensif d'un système en équilibre}
\subsection{Variance}
La \textbf{variance} est le nombre maximal de paramètres intensifs indépendants qu'un opérateur peut fixer sans rompre l'équilibre.
\begin{itemize}
	\item Les paramètres intensifs sont P, T et les fractions molaires de chaque espèce
	\item Ils sont relié par une équation pour chaque réaction et une pour chaque phase $\left(\displaystyle\sum_\varphi x=1\right)$
\end{itemize}
La variance est donc $$v=2+\sum_i 1-\text{nbr}_\text{réactions}-\text{nbr}_\text{phases}$$ où les $i$ sont les différents constituants physico-chimiques.
\subsection{Exemples}
\paragraph{$\text{N}_{2,(l)}=\text{N}_{2,(g)}$} Les paramètres intensifs sont $T$ et $P$. La seule relation entre ces paramètres provient de la loi d'action des masses qui donne $P_{\text{N}_{2,(g)}}=K^\circ(T)$. La variance est donc $1$ (en accord avec la formule de la première section). 
\paragraph{$\text{C}_{(s)}+\text{ZnO}_{(s)}=\text{Zn}_{(s)}+\text{CO}_{2,(s)}$}
Les paramètres intensifs sont $T$ et $P$ (tous les constituants sont seul dans leur phase). La seule relation entre les paramètres est la loi d'action des masses. La variance vaut donc $1$ (en accord avec la formule de la première section).
\paragraph{$3\text{H}_{2(g)}+\text{N}_{2(g)}=2\text{NH}_{3(g)}$} Les paramètres intensifs sont $T$, $P$, $x_{\text{N}_2}$, $x_{\text{H}_2}$ et $x_{\text{NH}_3}$. Les relations entre ces paramètres sont la loi d'action des masses et la relation $x_{\text{N}_2}+x_{\text{H}_2}+x_{\text{NH}_3}=1$. La variance est donc de $3$ (en accord avec la formule de la première section).
\section{Rupture ou déplacement d'équilibre}
\subsection{Exemple de rupture d'un équilibre}
\paragraph{Rupture d’équilibre} Si la variance est de 1, on peut faire varier le paramètre pour empêcher l’équilibre d’être atteint et ainsi obtenir une transformation totale.
\paragraph{Exemple : réaction d'oxydation du fer} On considère la réaction d'oxydation du fer : $$2\text{Fe}_{(s)}+\text{O}_{2(g)}=2\text{FeO}_{(s)}$$
\subsection{Méthode générale d'étude des déplacements d'équilibre}
On part de $Q_{\text{eq},1}=K_1^\circ(T_1)$ et on modifie soit $K_1$ (température) soit $Q_1$ (le reste).
\paragraph{Moyens de déplacer un équilibre} On perturbe un système à l’équilibre pour voir sa réaction par :
\begin{itemize}
	\item Ajout d’un réactif/produit (à $P$, $T$ et autres concentrations fixés)
	\item Modification de $T$ (à $P$ et composition constantes)
	\item Modification de $P$ (à $T$ et composition constantes)
	\item Ajout d’une espèce inerte
\end{itemize}
\paragraph{Méthode} On donne un nom aux différents états : l'état d'équilibre initiale sera ici noté (1) ; l'état hors équilibre juste après la perturbation sera noté (int) et le nouvel état d'équilibre sera noté (2). \par Comme (1) et (2) sont des états d'équilibre, on aura toujours $K^\circ(T_1)=Q_1$ et $K^\circ(T_2)=Q_2$. Comme l'état (int) est hors équilibre, on a aura toujours $K^\circ(T_\text{int})\ne Q_\text{int}$. Il faut comparer le quotient de réaction juste après la perturbation à la nouvelle constante d'équilibre.
\subsection{Influence de la température ; loi de Van't Hoff}
Modification de $T$ : regarder l’évolution de $K^\circ$ avec la loi de Van’t Hoff $$\frac{d}{dT}\left(\ln(K^\circ)\right)=\frac{\Delta_rH^\circ}{RT^2}$$
\paragraph{Exemple} On considère la réaction $$2\text{H}_{2(g)}+\text{O}_{2(g)}=2\text{H}_2\text{O}_{(g)}$$ dont on sait qu'elle est exothermique. On part d'un état d'équilibre $(T_1,P_1)$ puis on change la température de contrôle jusqu'à $T_2$ en gardant $P_2=P_1$. On a alors : \begin{center}
	\begin{align*}
		Q&=\displaystyle\frac{P_{\text{H}_2\text{O}}^2P^\circ}{P_{\text{H}_2}^2P_{\text{O}_2}} \\
		&=\displaystyle\frac{x_{\text{H}_2\text{O}}^2}{x_{\text{H}_2}^2x_{\text{O}_2}}\times\frac{P^\circ}{P_\text{tot}}
	\end{align*}
\end{center}
Or, $\displaystyle\frac{d}{dT}\left(\ln\left(K^\circ\right)\right)=\frac{\Delta_rH^\circ}{RT^2}<0$ car la réaction est exothermique et on sait que $Q_\text{int}=Q_1$ car on ne change pas la pression de contrôle ni n'ajoute de composant, et $Q_1=K^\circ(T_1)$ car il y a équilibre. \\ \par  Ainsi, si on suppose $T_2>T_1$, on a alors $K^\circ(T_2)>K^\circ(T_1)$ et $Q_\text{int}>K^\circ(T_2)$, si bien que la réaction a lieu dans le sens indirect.
\paragraph{Loi de modération liée à la température} Il faut retenir la loi de Van't Hoff qualitative : 
\begin{center}
	Élévation de T $\to$ réaction endothermique (loi de Van’t’Hoff)
\end{center} Ou encore : une réaction exothermique est favorisée par une diminution de la température.
\subsection{Influence de la pression ; loi de Le Châtelier}
Modification de $P$ : $K_2=K_1$, écrire le quotient de réaction avec les fractions molaire pour faire apparaître $P_\text{tot}$
\paragraph{Exemple} On considère la réaction : $$\text{CH}_{4(\text{g})}+\text{H}_2\text{O}_{(\text{g})}=\text{CO}_{(\text{g})}+3\text{H}_{2(\text{g})}$$ On se place dans l'état (1) à $P_1$ et $T_1$, puis on change la pression de contrôle à $P_2>P_1$ sans changer la température de contrôle ni ajouter de constituant. Juste après la perturbation, on a : $$Q=\frac{x_\text{CO}x_{\text{H}_2}^3}{x_{\text{CH}_4}x_{\text{H}_2\text{O}}}\times\frac{P_\text{tot}^2}{P^{\circ 2}}$$ Or, le facteur $\displaystyle\frac{x_\text{CO}x_{\text{H}_2}^3}{x_{\text{CH}_4}x_{\text{H}_2\text{O}}P^{\circ 2}}$ n'est pas modifié par la modification de pression, donc on a ici $$Q_\text{int}>Q_1=K^\circ(T_1)=K^\circ(T_2)$$ car la température ne change pas donc les constantes d'équilibre sont les mêmes, et $Q_1=K^\circ(T_1)$ car l'état (1) est un état d'équilibre. Ainsi, la réaction est déplacée dans le sens inverse.
\paragraph{Généralisation} On a $$Q=C\times P_\text{tot}^{\Delta_\text{gaz}\nu}$$ où $C$ est une constante indépendante de la pression. On en déduit la :
\paragraph{Loi de Le Châtelier} Si $\Delta_\text{gaz}\nu>0$, une diminution de la pression favorise la réaction dans le sens direct.
\paragraph{Principe de modération (principe de Le Châtelier)} "La réaction essaye de diminuer l'effet de la perturbation" $\to$ Ici, la réaction consomme du gaz pour essayer de diminuer la pression.
\subsection{Influence de l'introduction d'un constituant}
Modification de la composition : écrire $Q$ avec les fractions puis quantités de matières. Exprimer $dQ$
\paragraph{Cas d'un constituant actif}
On prend l'exemple de la réaction $$2\text{CO}_{(\text{g})}+O_{2(\text{g})}=2\text{CO}_{2(\text{g})}$$ avec ajout d'une petite quantité de $\text{CO}$.  L'état d'équilibre initial est à $P_1$ et $T_1$, on ajoute un peu de $\text{CO}$ et on ne change pas la pression ni la température de contrôle ($P_2=P_1$ et $T_2=T_1$). En notant $n_1$, $n_2$ et $n_3$ les quantités de matière des différents constituants à l'état initial et $dn_1$ la quantité ajoutée, on a : $$\ln(Q_\text{int})-\ln(Q_1)=\frac{dn_1}{n_\text{tot gaz}}-\frac{2dn_1}{n_1+dn_1}=\frac{n_1+dn_1-2n_\text{tot gaz}}{n_\text{tot gaz}(n_1+dn_1)}dn_1<0$$ Ainsi, $Q_\text{int}<Q_1=K^\circ(T_1)=K^\circ(T_2)$ (l'état (1) est un état d'équilibre et $T_2=T_1$), donc la réaction est déplacée dans le sens direct.
\paragraph{Cas d'un constituant inerte} On prend l'exemple de la réaction : $$3\text{H}_{2(\text{g})}+\text{N}_{2(\text{g})}=\text{NH}_{3(\text{g})}$$ On part d'un équilibre à $T_1$ et $P_1$ et on ajoute $dn$ d'hélium sans changer la température ni la pression de contrôle. On a alors juste après la perturbation : $$Q_\text{int}=\frac{n_{\text{NH}_3}^2}{n_{\text{H}_{2}}^3n_{\text{N}_2}}\times\left(n_{\text{NH}_3}+n_{\text{N}_2}+n_{\text{H}_2}+dn\right)\times\frac{P^{\circ 2}}{P_\text{tot}^2}$$ où le seul terme qui a changé par rapport à $Q_1$ est $n_{\text{NH}_3}+n_{\text{N}_2}+n_{\text{H}_2}$ qui a été remplacé par $n_{\text{NH}_3}+n_{\text{N}_2}+n_{\text{H}_2}+dn$. Ainsi, on a donc $$Q_\text{int}>Q_1=K^\circ(T_1)=K^\circ(T_2)$$ car il y a équilibre à l'état (1) et $T_2=T_1$. La réaction est donc déplacée dans le sens indirect.
\section{Optimisation d'un procédé chimique}
Lire le document sur la synthèse de l'ammoniac et répondre aux question.
\paragraph{Lois de modération} Le système tend à limiter la perturbation :
\begin{itemize}
	\item Élévation de T $\to$ réaction endothermique (loi de Van’t’Hoff) ;
	\item Élévation de P $\to$ consommation de gaz (principe de modération Le Châtelier) ;
	\item Ajout d’une espèce $\to$ consommation	de cette espèce.
\end{itemize}

\part{Optique}
\chapter{Optique géométrique}
\localtableofcontents
\section{Lois de l'optique géométrique}
\subsection{Lois fondamentales}
\paragraph{Notion de rayon lumineux} Un rayon lumineux est une ligne de champ du vecteur de Poynting de l'onde lumineuse électromagnétique : en effet, le rayon lumineux nous amène l'énergie de cette onde.
\paragraph{Négligeabilité de la diffraction} Dans le cadre de l’optique géométrique, on néglige la diffraction.
\paragraph{Principe de la propagation rectiligne} Dans un milieu homogène et isotrope, les rayons lumineux sont des portions de droites.
\paragraph{Principe d’indépendance des rayons lumineux} Les rayons lumineux sont indépendants les uns des autres. Ainsi, lors des tracés de rayons lumineux, on ne se préoccupe pas du nombre de rayons lumineux et des conséquences potentielles de leurs intersections. $\implies$ Ce principe est faux en général : en effet, une onde peut influencer la propagation d'une autre, se référer à la pâle sur le sujet X PC 2017 avec l'onde sonde et l'onde porteuse (trauma).
\paragraph{Principe du retour inverse de la lumière} Le trajet suivi par la lumière entre deux points sur le même rayon est indépendant du sens de parcours. $\implies$ Pour un exercice d'optique géométrique, \textbf{soit on ne connaît pas son cours, soit on n'a pas pensé au retour inverse de la lumière} : ce principe permet bien souvent de symétriser les problèmes pour en faciliter la résolution.
\paragraph{Lois de Snell-Descartes} Pour la réflexion sur un miroir plan
\begin{itemize}
	\item Les rayons incidents, réfléchis et la normale au miroir au point d’incidence sont dans le même plan ;
	\item L'angle du rayon incident avec la normale $i$ est égal à l'opposé de l'angle du rayon réfléchi $r$ par rapport à la normale (avec des angles orientés) : $i=-r$
\end{itemize}
Pour la réfraction entre deux milieux transparents, homogènes, isotropes et linéaires :
\begin{itemize}
	\item Les rayons incidents, réfractés et la normale au dioptre au point d’incidence sont dans le même plan ;
	\item $n_1\sin(i_1)=n_2\sin(i_2)$
\end{itemize}
Si $n_1\leqslant n_2$, il y a toujours réfraction. Sinon, il y a réflexion totale pour $\displaystyle\sin(i_1)\geqslant\frac{n_2}{n_1}$
\paragraph{Principe de Fermat} On définit un chemin optique à partir d’un chemin $[AA']$ par $$(AA') =\int_{[AA']}n\overrightarrow{r}dr$$ avec $n$ l’indice du milieu au point $r$. C’est la distance que parcourrait un photon dans le vide pendant la même durée. Le principe de Fermat s’énonce :
\begin{itemize}
	\item Historiquement : le chemin optique est minimal sur le trajet du vrai rayon comparé aux autres trajets propre
	\item Correspond à une dérivée nulle du chemin optique sur le trajet du rayon.
\end{itemize}
Ce principe contient la propagation rectiligne, le retour inverse et les lois de Snell.
\subsection{Applications}
\paragraph{Lentille} Dans une lentille avec un bord droit et un bord convexe, il n'y a pas stigmatisme car des rayons parallèles à l'axe sont déviés différemment. Cependant, dans le conditions de Gauss, il y a stigmatisme approché.
\paragraph{Le prisme} On considère un prisme d'angle au sommet $\alpha$. Le premier rayon incident arrive avec un angle $i$, est réfracté avec un angle $r$. Ce rayon réfracté arrive avec un angle $r'$ sur la dernière face et est réfracté avec un angle $i'$. Les lois de Snell-Descartes aux deux points de réfraction donnent : $$\begin{cases}
	\sin(i)=n\sin(r)\\
	\sin(i')=n\sin(r')
\end{cases}$$ Ensuite, la somme des angles d'un triangle vaut $\pi$ radians, donc $\alpha+\displaystyle\left(\frac{\pi}{2}-r\right)+\left(\frac{\pi}{2}-r'\right)=\pi$, ce qui fournit $$\alpha=r+r'$$ Enfin, la déviation $D$ est l'angle entre le prolongement du premier rayon incident et le prolongement du rayon émergent du prisme. On a donc $D=(i-r)+(i'-r')$. Cela donne avec la formule précédente : $$D=i+i'-\alpha$$
\paragraph{Résumé : les 4 formules du prisme} Pour le prisme considéré précédemment, on a donc les formules du prisme:  $$\begin{cases}
	\sin(i)=n\sin(r)\\
	\sin(i')=n\sin(r')\\
	\alpha=r+r'\\
	D=i+i'-\alpha
\end{cases}$$ Ces formules sont \textbf{à connaître par cœur} (accompagnées du schéma qui convient) : certains concours ne sont pas au courant qu'elles ne sont pas au programme...
\paragraph{Condition de non réflexion totale sur le prisme} Il n'y a pas de réflexion totale en I' si $n\sin(r)\leqslant1$. La loi de la réfraction en I donne $n\sin(r)=\sin(i)\leqslant1$. Donc $\displaystyle r\leqslant\arcsin\left(\frac{1}{n}\right)$ et $\displaystyle r'\leqslant\arcsin\left(\frac{1}{n}\right)$, si bien que $$\sin\left(\frac{\alpha}{2}\right)\leqslant\frac{1}{n}$$
\paragraph{Variation de la déviation du prisme en fonction de l'incidence} On obtient une courbe parabolique numériquement. Déterminons le minimum de déviation. Par unicité du minimum de déviation et d'après le principe du retour inverse de la lumière, on a $i=i'$ au minimum de déviation. Par conséquent, les deux premières formules du prisme donnent $r=r'$ puis la troisième fournit : $$r=r'=\frac{\alpha}{2}$$ Or, on a $\sin(i)=n\sin(r)$ donc $$\sin\left(\frac{D_m+\alpha}{2}\right)=n\sin\left(\frac{\alpha}{2}\right)$$
\paragraph{Formation d'un halo lumineux} Pour simplifier, on suppose que presque tous les rayons sont déviés de $D_m$. Pour l'arc-en-ciel, on doit tourner le dos au soleil pour le voir car les gouttes d'eau diffusent dans la direction vers la direction d'où elles reçoivent la lumière, tandis que pour voir un halo lumineux, il faut regarder vers la source de lumière (retenir que c'est opposé pour l'arc-en-ciel et qu'il est logique de regarder la source pour voir le halo). Ensuite, dispersion $n=n(\omega)$ et $D_m=f(\omega)$.
\section{Systèmes optiques centrés}
\subsection{Définitions}
\paragraph{Système centré} Un système optique est dit \textbf{centré} lorsqu'il est invariant par rotation autour d'un axe, alors appelé \textbf{axe optique}. Attention, il peut y avoir plusieurs axes optiques pour le miroir plan, mais c'est essentiellement le seul cas où cela peut arriver.
\paragraph{Exemples} Une lentille est en général un système centré. Un miroir plan est un système centré. Un dioptre plan est un système centré.
\paragraph{Foyer objet d'un système centré} Pour un système centré, si les rayons incidents (ou leurs prolongements) passent par le foyer objet $F$ du système, alors les rayons émergents (ou leurs prolongements) sont parallèles à l'axe optique.
\paragraph{Foyer image d'un système centré} Pour un système centré, si les rayons incidents (ou leurs prolongements) arrivent parallèlement à l'axe optique, alors les rayons émergents passent par le foyer image $F'$ du système.
\paragraph{Centre optique pour une lentille mince} Pour une lentille mince, le \textbf{centre optique} est l'intersection de l'axe optique et de la lentille. Pour une association de lentilles, il n'est pas évident qu'un tel point existe.
\paragraph{Plan focal objet} Le \textbf{plan focal objet} est le plan perpendiculaire à l'axe optique passant par $F$. On dit aussi que c'est le plan de front de $F$.
\paragraph{Plan focal image} Le \textbf{plan focal image} est le plan perpendiculaire à l'axe optique passant par $F'$. On dit aussi que c'est le plan de front de $F'$.
\paragraph{Foyer secondaire objet} C'est un point du plan focal objet différent du foyer objet.
\paragraph{Foyer secondaire image} C'est un point du plan focal image différent du foyer objet.
\paragraph{Stigmatisme et conjugaison}
On dit qu'il y a stigmatisme rigoureux lorsque tout rayon passant par A (ou dont le prolongement passe par A) passe après traversée du système optique par A' (ou leur prolongement). A est l’objet, A' l’image, A et A' sont conjugués. Il y a stigmatisme approché si tout rayon passant par A passe au voisinage de A' (plus petit que le grain du détecteur).
\paragraph{Aplanétisme} On dit d'un système centré qu'il est \textbf{aplanétique} lorsque l'image de tout objet dans un plan de front est contenue dans un plan de front.
\paragraph{Conditions de Gauss} Un système centré est dans les \textbf{conditions de Gauss} si on se limite aux rayons paraxiaux (proche de l’axe et à faible incidence par rapport à celui-ci). Dans ces conditions, il y a stigmatisme approché et aplanétisme approché.
\subsection{Dioptre plan ; lame à faces parallèles}
\paragraph{Dioptre plan} Pour A l’objet et H l’intersection entre la normale au dioptre passant par A et le dioptre, on a dans les conditions de Gauss : $$\frac{\overline{HA}}{n_1}=\frac{\overline{HA'}}{n_2}$$ Le grandissment transversal vaut $1$.
\paragraph{Démonstration de la relation de conjugaison du dioptre} En effet, on note $i$ l'angle d'incidence et $r$ l'angle de réfraction. Alors : $$HA\tan(i)=HA'\tan(r)$$ et d'après la loi de Descartes sur la réfraction, on a $$n_1\sin(i)=n_2\sin(r)$$ Le quotient de ces deux relations dans les conditions de Gauss (donc dans l'approximation des petits angles) fournit la formule recherchée.
\paragraph{Lame à faces parallèles} Deux lames à faces parallèles espacées d'une épaisseur $e$ délimitent un milieu d'indice $n$ de l'air. Avec le diagramme $$A\overset{D_1}{\longrightarrow}A_1\overset{D_2}{\longrightarrow}A'$$ on a, dans les conditions de Gauss, la relation de conjugaison : $$\overline{AA'}=e\left(1-\frac{1}{n}\right)$$
\paragraph{Remarques} Cette relation est homogène. Pour $e=0$ ou $n=1$, on trouve $\overline{AA'}=0$, ce qui est cohérent car c'est alors comme s'il n'y avait pas de lame.
\paragraph{Démonstration de la relation de conjugaison pour la lame à faces parallèles} On applique les relations de conjugaison aux deux dioptres plans. Pour le premier : $$\frac{\overline{H_1A}}{1}=\frac{\overline{H_1A_1}}{n}$$ Pour le deuxième : $$\frac{\overline{H_2A_1}}{n}=\frac{\overline{H_2A'}}{1}$$ On obtient alors : $$\overline{AA'}=\overline{AH_1}+\overline{H_1H_2}+\overline{H_2A'}=e\left(1-\frac{1}{n}\right)$$
\subsection{Lentilles minces}
\paragraph{Tracés par les rayons particuliers} On peut utiliser plusieurs rayons particuliers :
\begin{itemize}
	\item le rayon passant par l'axe optique qui n'est pas dévié ;
	\item le rayon incident parallèle à l'axe optique qui émerge par le foyer image ;
	\item le rayon incident passant par le foyer objet qui émerge parallèle à l'axe optique.
\end{itemize}
\paragraph{Tracés par les foyers secondaires}
Sinon, on peut utiliser les foyer secondaires : si deux rayons incidents sont parallèles entre eux, alors les rayons émergents se croisent en un foyer secondaire image (et on le contraire pour des rayons émergents parallèles entre eux et un foyer secondaire objet). 
\paragraph{Tracés par stigmatisme et aplanétisme}Enfin, on peut utiliser le stigmatisme et surtout l'aplanétisme : lorsqu'on a un point sur l'axe, placer un objet dans le plan de front, dont il est plus facile de déterminer l'image ou l'antécédent, ce qui permet de trouver la position de l'image ou de l'antécédent du point sur l'axe par aplanétisme.
\paragraph{Remarque} Si $\overline{OA}=2f$, alors $\overline{OA'}=2f'$ : c'est la méthode de Bessel, aussi appelée méthode des "2f-2f".
\paragraph{Rappel 1} Avec une unique lentille convergente, la distance minimale entre l'objet et l'écran vaut $4f'$. \textbf{Pour ne pas passer pour un guignol en TP, écarter l'objet et l'écran d'au plus 4f'}.
\paragraph{Formule de conjugaison de Newton} Elle s'énonce : $$\overline{FO}\times\overline{F'O}=\overline{FA}\times\overline{F'A'}=-f'^2$$
\paragraph{Formule de conjugaison de Descartes} Elle s'énonce : $$\frac{1}{\overline{OA'}}-\frac{1}{\overline{OA}}=\frac{1}{f'}$$
\paragraph{Grandissement} Le grandissement transversal est défini par : $$\gamma=\frac{\overline{A'B'}}{\overline{AB}}$$ Si $\lvert\gamma\rvert>1$, l'image est agrandie. SI $\gamma >0$, l'image est droite. Si $\gamma <0$, l'image est renversée.
\paragraph{Formule de grandissement de Newton} Elle s'énonce : $$\gamma =\frac{\overline{F'A'}}{\overline{F'O}}=\frac{\overline{FO}}{\overline{FA}}$$
\paragraph{Formule de grandissement de Descartes}Elle s'énonce : $$\gamma=\frac{\overline{OA'}}{\overline{OA}}$$
\section{Instruments d'optique}
\subsection{L'œil humain : performances et limites, modélisation}
\paragraph{Modélisation de l'œil} L'œil est constitué de : 
\begin{itemize}
	\item la pupille, trou qui laisse passer la lumière ;
	\item le cristallin, modélisé par une lentille ;
	\item la rétine, modélisée par un écran dont la distance à la lentille/au cristallin est fixe ;
	\item le diaphragme qui contrôle l'ouverture de l'œil et la quantité de lumière admise. 
\end{itemize}
\paragraph{Punctum proximum} C'est le point le plus proche que peut voir l'œil en accommodant. En général sans défaut de l'œil, il est situé à 15 cm.
\paragraph{Punctum remotum} (\textit{Remote} c'est la télécommande de la télé, qui est donc loin de moi quand j'en ai besoin) C'est le point le plus loin visible à l'œil sans accommoder. En général, sans défaut de l'œil
\paragraph{Œil emmétrope} Un œil emmétrope (prononcer "an-mé-trope) est un œil sans défaut.
\paragraph{Myopie} La myopie correspond à un rapprochement du punctum proximum et du punctum remotum. C'est lorsque le cristallin est trop convergent.
\paragraph{Hypermétropie} L'hypermétropie correspond à un éloignement du punctum remotum et du punctum proximum. C'est lorsque le cristallin n'est pas assez convergent.
\paragraph{Astigmatie} L'astigmatie correspond au fait que le cristallin n'est alors pas exactement un système centré, ce qui peut créer des déformations de visions des lignes droites car l'œil n'est plus aplanétique dans les conditions de Gauss.
\paragraph{Pouvoir séparateur de l'œil} Le pouvoir séparateur de l'œil $\varepsilon$ est le plus petit angle sous lequel un œil humain peut distinguer deux objets. $\varepsilon$ vaut pour un œil normal une minute d'arc, soit environ $3.10^{-5}$ radians.
\subsection{Lunettes et viseurs}
\paragraph{Lunettes, viseurs, collimateurs}
\textit{cf} TP. Il faut connaître :
\begin{itemize}
	\item la composition de ces appareils ;
	\item leur réglage ;
	\item leur utilité.
\end{itemize}
A COMPLETER AVEC LE POLY DE TP

\chapter{Introduction à l'optique physique}
\localtableofcontents
\section{Modèle scalaire des ondes lumineuses}
\subsection{Grandeur lumineuse ou grandeur scalaire de l'optique}
\paragraph{Modèle scalaire de la lumière} On définit : $$s(x,t)=\overrightarrow{E}\cdot\overrightarrow{u}$$ $s$ est une des composantes transverses du champ électrique.
\paragraph{Surfaces d'onde}Les surfaces d'onde correspondent alors au même chemin optique et à la même phase.
\paragraph{Cas particuliers des ondes planes et sphériques}On a :
\begin{itemize}
	\item pour une onde sphérique : $s(r,t)=\displaystyle \frac{s_0}{r}e^{i(\omega t-kr)}$ ;
	\item pour une onde plane : $s(z,t)=s_0\displaystyle e^{i(\omega t-kz)}$
\end{itemize}
\paragraph{Cas général}
De façon générale, pour $A$ et $B$ sur le même rayon, on a la formule : $$\bigg\lvert\frac{\varphi(A,t)-\varphi(B,t)}{2\pi}\bigg\rvert=\frac{(AB)}{\lambda}$$
\paragraph{Principe de superposition} Pour des sources $1$ et $2$, $s$ qui résulte des sources $1$ et $2$ à la fois vérifie : $$s=s_1+s_2$$ grâce à la linéarité des équations de Maxwell et aux directions de propagation "proches".
\subsection{Domaine de validité du modèle scalaire}
\paragraph{Limite} On oublie le caractère transverse et polarisé de la lumière, ce qui est valable pour des sources non polarisées. On ne peut donc pas étudier des sources polarisées avec ce modèle.
\section{Émission et réception de la lumière}
Le but est de s'intéresser à :
\begin{itemize}
	\item $T$ la période de l'onde ;
	\item $\tau_C$ le temps de cohérence de la source ;
	\item $\tau_R$ le temps de cohérence du détecteur.
\end{itemize}
On cherche à comprendre ce que représentent ces grandeurs, quelles limites elles imposent dans nos études et quelles relations elles entretiennent entre elles.
\subsection{Modèle des trains d'onde aléatoires ; durée de cohérence de la source}
\paragraph{Train d’onde} Il est émis lors de la désexcitation d’un atome : c'est une sinusoïde de période $T$ pendant une durée $\Delta T$. On considère un $\Delta T$ moyen sur tous les trains d’ondes. Les trains d'ondes expliquent le fait que la lumière est non polarisée : en effet, c'est la somme de trains polarisés aléatoirement.
\paragraph{Modèle du train d'onde aléatoire} Tous les $\tau_C$, on tira aléatoirement :
\begin{itemize}
	\item l'amplitude ;
	\item la phase à l'origine ;
	\item la polarisation
\end{itemize} mais pas la période.
\paragraph{Longueur de cohérence} La longueur de cohérence est définie par : $$l_C=\tau_C\times c$$ En TP, $l_C$ est la précision avec laquelle il faut placer les miroirs du Michelson pour obtenir des interférences.
\subsection{Détecteurs ; éclairement et intensité lumineuse}
\paragraph{Détecteurs} Les détecteurs sont sensibles à l'énergie reçue. 
\paragraph{Intensité lumineuse }On définit l'\textbf{intensité lumineuse} $I$ résultant d'une onde $s$ par : $$I=\langle s^2\rangle_{\tau_R}$$ avec $\tau_R$ le temps de réponse du détecteur.
\paragraph{Ordres de grandeur de temps de réponse} Pour l'œil humain, le temps de réponse est d'environ $1/24$ s, soit environ $4.10^{-2}$ s. Pour une photodiode, on a $\tau_R=10^{-6}$ s.
\paragraph{Condition pour observer des interférences} Pour observer (ou détecter) des interférences, il faut avoir $\tau_R \ll T$. Or, pour la lumière visible, $T\approx 2.10^{-15}$ s, donc il est très difficile d'obtenir des interférences pour les ondes lumineuses.
\paragraph{Éclairement} On définit l'\textbf{éclairement} $\xi$ par $$\xi=\langle \pi\rangle$$ En optique, on moyenne avec $\tau_R$, tandis qu'en électromagnétisme, on moyenne avec $T$. Avec les bonnes approximations, cela revient essentiellement au même.
\paragraph{Rappel} Pour l'OPPH $\underline{\overrightarrow{E}}=\underline{E_0}e^{i(\omega t-kz)}\overrightarrow{u_x}$, on rappelle qu'on a : $$\langle\pi\rangle =\frac{c\varepsilon_0}{2}\lvert\underline{E_0}\rvert^2$$
\subsection{Densité spectrale ; largeur spectrale}
\paragraph{Densité spectrale} On définit $I_\lambda$ l'intensité spectrale, ou densité spectrale, telle que $$I=\int I_\lambda\ d\lambda$$
\paragraph{Largeur spectrale} On considère $\displaystyle s(t)=\sum_\omega\underline{g}(\omega)e^{i\omega t}$. On a alors : $$\Delta\omega\Delta T\geqslant1$$ donc plus $\Delta\omega$ est petit, plus $\Delta T$ est grand. Pour une sinusoïde parfaite, $\Delta\omega=0$ et $\Delta T=+\infty$.
\paragraph{Largeur spectrale version énergétique} On a $E_\text{après}=E_\text{avant}+E_\text{photon}=E_\text{avant}+\hbar\omega_0$ et $\Delta E=h\Delta\nu$, donc : $$\Delta E\Delta T\geqslant\hbar$$
\paragraph{Cas de la lampe spectrale} Pour une lampe spectrale, la largeur spectrale en longueur d'onde vérifie : $$\Delta\lambda\ll\lambda_m$$ où $\lambda_m$ est la longueur d'onde moyenne du train d'onde. On admet alors que $$\Delta f\times\tau_C\approx 1$$ où $\Delta f$ est la largeur spectrale en fréquence.
\paragraph{Cas de la lampe au sodium et des lampes spectrales de TP} Le doublet jaune du sodium a pour longueurs d'onde $589,0$ nm et $589,6$ nm.  En général, $\Delta\lambda\approx0,1$ nm pour les lampes spectrales de TP.
\paragraph{cas de la lampe à lumière blanche} Dans ce cas, on a $\Delta\lambda\approx\lambda_m$.
\paragraph{Relation entre largeurs spectrales en fréquence et en longueur d'onde} Pour $f=\displaystyle\frac{c}{\lambda_0}$, on a la relation : $$\Delta f=\frac{c\Delta\lambda_0}{\lambda_{0,m}^2}$$ (le signe $-$ a disparu). On en déduit avec la relation admise que $$l_C\approx\frac{\lambda_{0,m}^2}{\Delta\lambda}$$ Pour une lampe spectrale de TP, l'application numérique donne $l_C\approx 4.10^{-3}$ m.
\subsection{Émission d'une lampe spectrale ; causes de l'élargissement spectral}
Voici les différentes causes de la largeur spectrale.
\paragraph{Effet Doppler} Les atomes étant agités, ils n'ont pas le même mouvement par rapport à l'observateur donc n'émettent pas la même lumière. On a : $$v\st=\sqrt{\frac{3k_BT}{m}} \ \ \text{avec} \ \ k_B=\frac{R}{N_A}$$ ce qui donne un décalage de fréquence $\Delta\nu=\displaystyle\nu_0\frac{v\st}{c}$
\paragraph{Collisions} Si le nombre de chocs par unité de temps (pression) augmente, $\Delta T$ diminue donc $\Delta\omega$ augmente et $\tau_C$ diminue.
\paragraph{Largeur spectrale intrinsèque} C'est un exercice classique où on mélange :
\begin{itemize}
	\item la formule de Larmor ;
	\item l'exercice classique du satellite qui frotte dans l'atmosphère et tombe sur la Terre (appliquer le théorème de la puissance cinétique).
\end{itemize}
\subsection{Émission de lumière par un laser (culture)}
On a une relation de phase entre les deux photons qui augmente $\tau_C$.
\section{Résumé}
\subsection{Retour sur les deux approches}
\paragraph{Approche scalaire} On a la courbe de $s$ en fonction de $t$. Dans le modèle des trains d'ondes aléatoires, on a toujours la même période, mais l'amplitude, la phase et la polarisation changent tous les $\tau_c$.
\paragraph{Approche spectrale} On a le graphe de $I_\lambda$ en fonction de $f$. $\Delta f$ est la largeur spectrale en fréquence et $f_m$ est la fréquence moyenne du train d'ondes. On a les relations : $$f_m=\frac{1}{T} \ \text{ et } \ \Delta f\times\tau_C\approx 1$$
\subsection{Les trois temps caractéristiques}
\paragraph{Trois temps caractéristiques} Les trois temps caractéristiques sont : 
\begin{itemize}
	\item la période $T$ de l'onde ;
	\item le temps de cohérence $\tau_C$ de la source, qui correspond dans le modèle des trains d'onde aléatoires à la période de réinitialisation de l'onde (en tout cas de sa phase et de son amplitude).
	\item le temps de réponde $\tau_R$ du détecteur, qui définit jusqu'à quel point le détecteur perçoit le signal comme étant continu.
\end{itemize}
\paragraph{Ordres de grandeur} On a :
\begin{itemize}
	\item $T=10^{-15}$ s pour le visible ;
	\item $\tau_C=10^{-11}$ s pour une lampe spectrale et $10^{-8}$ s pour un laser ;
	\item Détecteurs : photodiode à $\tau_R=10^{-6}$ s et œil humain à $\tau_R=0,1$ s
\end{itemize}

\chapter{Interférences lumineuses}
\localtableofcontents
\section{Différence de phase entre deux points}
\subsection{Chemin optique ; surfaces d'onde}
\paragraph{Chemin optique} Si $\Gamma$ est un chemin reliant $A$ à $B$, on définit le chemin optique : $$(AB)_\Gamma=\int_{\Gamma}n(s)\ ds$$ où $n$ est l'indice optique et $ds$ est le $dl$ de la mécanique, l'abscisse curviligne. C'est la distance que parcourrait la lumière dans le vide pendant la durée qu'elle met à aller de $A$ à $B$.
\paragraph{Chemin optique dans deux cas particuliers} Pour un milieu homogène d'indice optique constant $n$ et si $\Gamma$ et la ligne droite entre $A$ et $B$, on a : $$(AB)_\Gamma=nAB$$ Si on prend deux milieux homogènes d'indices constants $n_1$ et $n_2$ avec $A$ dans le premier milieu, $B$ dans le deuxième et $I$ à l'interface, on a : $$(AB)=n_1AI+n_2IB=(AI)+(IB)$$
\paragraph{Surface d'onde} C'est une surface définie par rapport à un point telle que tous les points de cette surface ont le même chemin optique par rapport au point de définition.
\paragraph{Différence de marche et différence de phase} On a : $$\frac{\varphi_B-\varphi_A}{2\pi}=\frac{(AB)}{\lambda_0}=\frac{p}{1}$$ où $p$ est l'ordre d'interférence.
\paragraph{Principe de Fermat (culture)} "Le chemin optique effectivement suivi par la lumière minimise le chemin optique" (version historique). En réalité, ce serait plutôt "rend stationnaire" (de dérivée nulle) à la place de minimise $\to$ cela peut servir dans l'exercice des miroirs elliptiques. Le principe de Fermat permet de retrouver :
\begin{itemize}
	\item le principe de propagation rectiligne
	\item le principe de retour inverse
	\item les 4 lois de Snell-Descartes
\end{itemize}
\subsection{Théorème de Malus ; stigmatisme}
\paragraph{Théorème de Malus} Après un nombre arbitraire de réflexions et de réfractions, les rayons lumineux sont perpendiculaires aux surfaces d’onde. Autre présentation : pour $A$ et $B$ sur le même rayon, $A'$ et $B'$ sont sur le même rayon ;  $A$ et $A'$ sont sur la même surface d’onde et $B$ et $B'$ sont sur la même surface d’onde. On a donc $(AB) = (A'B')$
\paragraph{Onde sphérique} Pour une onde sphérique : $$s(r,t)=\frac{s_0}{r}e^{i(\omega t-kr)}$$ avec $\displaystyle\frac{\omega}{k}=v_\varphi=\frac{c}{n}$ donc on a $$s(r,t)=\frac{s_0}{r}e^{i\omega\left(t-\frac{nr}{c}\right)}$$ En prenant des points $A$ et $B$ alignés correspondant à des rayons $r_1$ et $r_B$, on trouve : $$\varphi_B-\varphi_A=2\pi\left(\frac{nr_B}{c}-\frac{nr_A}{c}\right)=\frac{2\pi}{\lambda_0}n(r_B-r_A)=\frac{2\pi}{\lambda_0}$$ On retrouve bien la formule donnée précédemment.
\paragraph{Onde plane} Pour une onde plane $\underline{s}=\underline{s_0}e^{i(\omega t-k_z)}$, on retrouve aussi la formule $$\frac{\Delta\varphi}{2\pi}=\frac{(AB)}{\lambda_0}$$
\paragraph{Stigmatisme} Le théorème de Malus permet-il de déterminer les surfaces d'onde à partir de quelques rayons ? Cela est complexe. En revanche, il permet de montrer que deux chemins optiques sont égaux.
\paragraph{Cas problématique des trous d'Young} On a une source qui émet vers des trous d'Young $T_1$ et $T_2$, puis tous cela passe dans une lentille convergente qui donne une image $M$ sur un écran. But : exprimer $(SM)_{T_2}-(SM)_{T_1}$.
\subsection{Déphasage additionnels (culture)}
Une réflexion ou une transmission peut induire un déphasage supplémentaire.
\paragraph{Réflexion sur un métal} Pour une réflexion normale sur un métal, on a $\underline{r}=-1$), soit un déphasage de $\pi$.
\paragraph{Réflexion entre deux diélectriques (HP)} Pour une réflexion entre deux diélectriques, dont le premier est d'indice $n_1$ et le deuxième d'indice $n_2$, on a la formule : $$\underline{r}=\frac{n_1-n_2}{n_1+n_2}$$ Si $n_2>n_1$, la réflexion introduit un déphasage de $\pi$.
\section{Superposition de deux ondes lumineuses}
On dit que les sources $1$ et $2$ interfèrent si $I_{\text{1 et 2}}\ne I_1+I_2$. Dans la suite, on note : $$\begin{cases}
	s_1=S_1^\circ\cos(\omega_1 t-\varphi_1) \\
	s_2=S_2^\circ\cos(\omega_2 t-\varphi_2)
\end{cases}$$
\subsection{Critère de cohérence temporelle}
\paragraph{Condition de réalisation d'interférences, sources cohérentes} On a par le calcul $$\langle s_1s_2\rangle =\frac{S_1^\circ S_2^\circ}{2}\bigg\langle\cos\left((\omega_1-\omega_2)t+\varphi_1(M)-\varphi_2(M)\right)\bigg\rangle$$ Cela est nul en moyenne à condition que $\displaystyle\frac{2\pi}{\omega_1-\omega_2}\ll T_\text{détecteur}$. Or, $T_\text{détecteur}\approx 10^{-6}$s et $\displaystyle\frac{1}{\omega}\approx 10^{-15}$rad.s$\inv$, donc les interférences ne sont par réalisables avec des sources distinctes. Pour réaliser des interférences, on utilise donc la même source et on a $$I=I_1+I_2+(S_1^\circ)^2\bigg\langle\cos\left(\varphi_1(M)-\varphi_2(M)\right)\bigg\rangle$$ Mais $\varphi$ varie aléatoirement d'un train d'onde à l'autre donc il faut attendre que la différence de marche $\delta$ soit inférieure à la longueur de cohérence $c\tau_c\approx 3\mu$m. On dit alors que les sources sont \textbf{cohérentes}.
\subsection{Superposition de deux ondes cohérentes ; formule de Fresnel ; contraste}
\paragraph{Interférences à deux ondes} On considère deux sources $s_1$ et $s_2$ ainsi qu'un point $M$ : $$s_1(M,t)=S_1^\circ\cos\left(\omega t+\varphi_1+\frac{S_1M}{\lambda_1}\right) \ \ \text{et} \ \ s_2(M,t)=S_2^\circ\cos\left(\omega t+\varphi_2+\frac{S_2M}{\lambda_2}\right)$$
L'interférence de ces deux ondes résulte en une onde $s$ telle que $s=s_1+s_2$. En effet, on peut appliquer le principe de superposition car ce sont des composantes du champ $\overrightarrow{E}$.
\paragraph{Interférences constructives et destructives} On a donc $I=I_1+I_2+2\langle s_1s_2\rangle$. On dit qu'on a des interférences \textbf{constructives} lorsque $I>I_1+I_2$, et \textbf{destructives} lorsque $I<I_1+I_2$.
\paragraph{Formule de Fresnel} Pour deux sources cohérentes, on a : $$I=I_1+I_2+2\sqrt{I_1I_2}\cos\left(\varphi_2-\varphi_1\right)=I_1+I_2+2\sqrt{I_1I_2}\cos\left(\frac{2\pi\delta}{\lambda}\right)$$
\paragraph{Démonstration} La formule se démontre proprement en réutilisant les calculs faits au début de la section. Sinon, pour la retrouver rapidement, on peut passer en notation complexe. La démonstration est alors bancale car elle n'utilise pas toutes les hypothèses, mais permet d'aboutir rapidement au résultat.
\paragraph{Cas particulier de la formule de Fresnel} Si de plus $I_1=I_2=I_0$, cela se réécrit : $$I=2I_0(1+\cos(\varphi_2-\varphi_1))=2I_0\left(1+\cos\left(\frac{2\pi\delta}{\lambda}\right)\right)$$
\paragraph{Remarque} Pour vérifier la formule, on peut vérifier :
\begin{itemize}
	\item qu'elle est homogène ;
	\item que si $\varphi_1$ et $\varphi_2$ sont en opposition de phase et $I_1=I_2$, alors $I=0$ ;
	\item que si $\varphi_1$ et $\varphi_2$ sont en phase et $I_1=I_2$, alors $I=4I_0$.
\end{itemize}
\paragraph{Contraste} On définit le contraste $C$ par $$C=\frac{I_\text{max}-I_\text{min}}{I_\text{max}+I_\text{min}}$$ C'est une grandeur sans dimension. Le contraste est inchangé si on multiplie par deux les intensités. Enfin, $C=0$ pour un éclairage uniforme (attention, cela ne veut pas dire que tout est sombre !).
\paragraph{Expression du contraste avec la formule de Fresnel} La formule de Fresnel permet directement d'établir que $$C=\frac{2\sqrt{I_1I_2}}{I_1+I_2}$$ On en déduit que $C\leqslant 1$ et que $C=1$ pour $I_1=I_2$ avec l'inégalité arithmético-géométrique.
\subsection{Superposition de deux ondes incohérentes}
On a dans ce cas : $I=I_1+I_2$.
\section{Résumé}
\paragraph{Résumé}
On a :
\begin{itemize}
	\item Si $S_1$ et $S_2$ sont cohérentes : $I=I_1+I_2+2\sqrt{I_1I_2}\cos(\varphi_2-\varphi_1)$ ;
	\item Si $S_1$ et $S_2$ sont incohérentes : $I=I_1+I_2+$.
\end{itemize}


\chapter{Exemple de dispositif interférentiel par division du front d'onde : trous d'Young}
\localtableofcontents
\section{Description : les deux montages classiques}
\paragraph{Dispositif des fentes d'Young} Fentes d'Young séparées de $a\approx 1\mu$m et placées à $D\approx1$m d’un écran. On peut placer une lentille entre les fentes et l’écran telle que l’écran soit le plan focal image.
\subsection{Le rôle de la diffraction}
\paragraph{Rôle de la diffraction} La diffraction donne la zone d’interférence : c'est l'intersection des deux cônes de diffraction, $\sin(\theta)=\displaystyle\frac{\lambda}{\varepsilon}$ avec $\varepsilon$ la largeur de la fente et $\theta$ le demi-angle au sommet du cône.
\paragraph{Cas $\varepsilon$ petit} SI $\varepsilon$ est petit, on a une diffraction isotrope mais des franges peu lumineuses.
\paragraph{Cas $\varepsilon$ grand} Si $\varepsilon$ est grand, la diffraction n'est plus isotrope mais les franges sont plus lumineuses.
\subsection{Calcul de la différence de marche ; intensité lumineuse}
\paragraph{Différence de marche} On la calcule grâce au théorème de Pythagore et on trouve avec un développement limité : $$\delta=\frac{nax}{D}$$ avec $n$ l'indice optique du milieu et $x$ la distance du point de l'écran au centre de l'écran. Pour une lentille, on utilise le théorème de Malus et le principe du retour inverse de la lumière pour établir que $$\delta=\frac{nax}{f'}$$
\paragraph{Remarque : exercice du miroir de Lloyd} Il existe divers exercices, notamment celui du miroir de Lloyd, où on pourra reconnaître des trous d'Young cachés ou quelque chose qui s'y ramène, ce qui pourra permettre d'éviter à calculer à la main la différence de marche et faire les mêmes DL que dans le cours.
\paragraph{Remarque : importance de la différence marche} Il faut bien comprendre que l'expression qui donne la différence de marche contient à elle seule toute la suite du chapitre.
\paragraph{Intensité lumineuse} Les deux sources sont cohérentes. Ainsi, d'après la formule de Fresnel, on a : $$I=2I_0\left(1+\cos\left(\frac{2\pi\delta}{\lambda_0}\right)\right)=2I_0\left(1+\cos\left(\frac{2\pi ax}{\lambda_0D}\right)\right)$$
\subsection{Frange d'interférences ; interfrange}
\paragraph{Figure d'interférences} Si $\delta$ ne dépend que de $x$, on a des franges rectilignes. L'ordre d'interférence $p$ vérifie : $$p=\frac{\delta}{\lambda}=\frac{\varphi}{2\pi}$$ Si $p\in\Z$, on a des interférences \textbf{constructives}. Si $p\in\frac{1}{2}+\Z$, on a des interférences \textbf{destructives}. L'\textbf{interfrange} $i$ est la distance entre deux franges consécutives : $$i=x_{p+1}-x_p=\frac{\lambda D}{na}$$
\paragraph{Remarque : obtention de l'interfrange à partir de la différence de marche} Retenir qu'il faut différencier l'expression de la différence de marche pour obtenir l'interfrange.
\subsection{Non localisation des franges}
On conserve la figure en déplaçant l'écran : on parle d'interférences non localisées. Le contraste $C$ vaut toujours $1$, quelle que soit la position de l'écran.
\section{Effet du déplacement et de l'élargissement spatial de la source}
\subsection{Déplacement de la source}
\paragraph{Effet d’un déplacement de la source} La nouvelle différence de marche est $$\delta=\frac{nax'}{D'}+\frac{nax}{D}$$ avec $x'$ la nouvelle abscisse de la source. La figure n’est pas modifiée, mais décalée en bloc.
\paragraph{Remarque : autre montage} Pour le montage analogue avec des lentilles, on a de même $$\delta=\frac{nax_1}{f_1'}+\frac{nax_2}{f_2'}$$
\paragraph{Conséquence sur l'interfrange} L'interfrange vérifie alors toujours la relation : $$i=\frac{\lambda D}{na}$$ L'interfrange n'est donc pas modifié.
\paragraph{Conséquence sur la figure d'interférences} Ainsi, lorsqu'on décale la source, la figur d'interférences est décalée en bloc.
\subsection{Deux sources ponctuelles}
\paragraph{Situation} On considère deux sources ponctuelles $S_1$ et $S_2$ qui émettent vers les trous d'Young. Elles sont de part et d'autre de l'axe qui passe pile entre les deux trous d'Young. $S_1$ est à $x_1$ au-dessus de cet axe, et $S_2$ est à $x_2$ au-dessous. Ainsi si $S_2$ est masquée : $$\displaystyle I_1=2I_0\left[1+\cos\left(\frac{2\pi n}{\lambda_{0,1}}\left(\frac{ax}{D}+\frac{ax_1}{D'}\right)\right)\right]$$
et si $S_1$ est masquée : $$I_2=2I_0\left[1+\cos\left(\frac{2\pi n}{\lambda_{0,2}} \left(\frac{ax}{D}+\frac{ax_2}{D'}\right)\right)\right]$$
\paragraph{Conséquence} La figure d'interférence est deux fois plus lumineuse. Si $$x_1-x_2=\left(p+\frac12\right)i\frac{D'}{D}$$ avec $p\in\Z$, alors le contraste est nul.
\paragraph{Rappel} Lorsque le contraste est nul, l'éclairement est uniforme mais l'intensité lumineuse n'est pas nécessairement nulle.

\subsection{Source rectangulaire ; contraste}
\paragraph{Effet d’un élargissement de la source} Selon $y$, rien ne change. Selon $x$, il y a brouillage. En effet, pour une source rectangulaire (même puissance $\frac{I_0}{h}$ pour tout point entre $-h/2$ et $h/2$), les différents
points de la source étant incohérents on a : $$I(M)=\int_{-h/2}^{h/2}I_{x'}(M)dx'=2I_0\left(1+\text{sinc}\left(\frac{\pi a h}{\lambda D'}\right)\cos\left(\frac{2\pi ax}{\lambda D'}\right)\right)$$ Pour cela, on a calculé $I_{x'}$ avec la formule de Fresnel.
\paragraph{Expression du contraste pour une source élargie spatialement} On a donc $$C=\bigg\lvert\text{sinc}\left(\frac{\pi a h}{\lambda D'}\right)\bigg\rvert$$ On a donc un bon contraste avant la moitié de la largeur correspondant à la première annulation, soit pour $h\leqslant\displaystyle\frac{\lambda D'}{2na}$.
\paragraph{Remarque : contraste dans le cas général} De façon générale, si $I=A(1+B\cos(\Delta\varphi))$, on montre immédiatement avec les variations de la fonction cosinus que $$C=\lvert B\rvert$$
\paragraph{Ordre de grandeur} On prend $n=1$, $\lambda=6.10^{-7}$ m (milieu du spectre visible) et $a=10^{-4}$ m. Le diamètre apparent vaut alors : $$\frac{\lambda}{2a}\approx 3.10^{-3} \ \text{rad}$$
\paragraph{Remarque : diamètre apparent du Soleil} Le diamètre apparent du Soleil vaut : $$\frac{700\times 10^3\times 10^3}{150\times10^{6}\times10^{3}}\approx 5.A0^{-3}\ \text{rad}$$
\subsection{Cohérence spatiale ; critère semi-qualitatif de brouillage}
\paragraph{Critère de brouillage} On considère qu'il y a $h/2$ entre le haut et le bas de la source. Si $p_\text{haut}-p_\text{bas}=\displaystyle\frac12$, il y a brouillage, \textit{ie} $C+0$. Dit autrement, il y a brouillage lorsque $$\frac{na}{D'}\times\frac{h_1}{2}=\frac{\lambda}{2}$$ où $h_1$ correspond à la première annulation.
\subsection{Fentes d'Young}
\paragraph{Fentes d'Young} Si on décale les trous d'Young dans la direction $y$, les figures d'interférences se superposent "bien" (les franges sombres sont en face des franges sombres). On obtient alors le dispositif des fentes d'Young en perçant plein de petits trous d'Young alignés selon $y$.
\section{Influence de la largeur spectrale de la source}
On utilise ici une source ponctuelle et $x'=0$ (la source est sur l'axe). On sait qu'alors : $$\delta=\frac{nax}{D}$$
\subsection{Doublet}
\paragraph{Doublet} On considère un spectre en doublet avec des longueurs d'ondes $\lambda_1$ et $\lambda_2$. Alors $I_1$ a pour expression : $$I_1=2I_0\left(1+\cos\left(\frac{2\pi nax}{\lambda_1 D}\right)\right)$$ et $I_2$ a pour expression $$I_12=2I_0\left(1+\cos\left(\frac{2\pi nax}{\lambda_2 D}\right)\right)$$ Par conséquent, les intensités n'oscillent pas à la même pulsation et le contraste n'est alors pas uniforme.
\paragraph{Notions de coïncidence et d'anticoïncidence} Lorsque les intensités sont en même temps minimale ou maximale, on parle de \textbf{coïncidence}. Lorsque les intensités sont "opposées" (l'une nulle, l'autre maximale), on parle d'anticoïncidence.
\paragraph{Influence de la largeur spectrale} Pour des ondes incohérentes, on somme les intensités. On utilise le nombre d'onde $\sigma$. Deux exemples : 
\begin{itemize}
	\item les doubles raies du sodium : elle font la même figure pour $x$ petit, mais diffèrent et brouillent la figure plus loin ;
	\item la lampe de largeur spectrale $\Delta\sigma$ : comme on a $\tau_c\Delta\sigma\geqslant1$, le temps de cohérence diminue. Les interférences sont brouillées si $\delta$ augmente trop.
\end{itemize}
\subsection{Profil spectral rectangulaire en $\sigma$ ; contraste}
\paragraph{Nombre d'onde} On appelle nombre d'onde la quantité $\sigma$ définie par $$\sigma=\frac{1}{\lambda}$$
\paragraph{Profil rectangulaire} Pour un profil en largeur spectrale rectangulaire de largeur $\Delta\sigma$, on a : $$dI=\frac{2I_0}{\Delta\sigma}d\sigma\left(1+\cos\left(2\pi\sigma\delta\right)\right)$$
\paragraph{Conséquence pour l'intensité} Les différentes longueurs d'onde sont incohérentes entre elles donc on peut sommer (ici continûment, donc intégrer) les intensités. On obtient :
\begin{align*}
	I&=\frac{2I_0}{\Delta\sigma}\int_{\sigma_0-\frac{\Delta\sigma}{2}}^{\sigma_0+\frac{\Delta\sigma}{2}}\left(1+\cos\left(2\pi\sigma\delta\right)\right)d\sigma \\
	&=\frac{2I_0}{\Delta\sigma}\left(\Delta\sigma+\frac{1}{2\pi\delta}\left(\sin\left(2\pi\delta\left(\sigma_0+\frac{\Delta\sigma}{2}\right)\right)-\sin\left(2\pi\delta\left(\sigma_0-\frac{\Delta\sigma}{2}\right)\right)\right)\right) \\
	&=2I_0(1+\text{sinc}(\pi\Delta\sigma\delta)\cos(2\pi\sigma_0\delta))
\end{align*}
Puis tracer la courbe : oscillations dans une enveloppe qui décroît comme une fonction inverse autour d'une constante.
\paragraph{Contraste} On a donc $$C=\lvert\text{sinc}(\pi\Delta\sigma\delta)\rvert$$ On a un très bon contraste pour $\delta\ll l_c$, et un contraste nul pour $\delta >l_c$.
\paragraph{Remarque} On rappelle que $l_c=\displaystyle\frac{\lambda^2}{\Delta\lambda}$. Or, ici, $\displaystyle\sigma=\frac{1}{\lambda}$, donc en différenciant, on trouve $\displaystyle l_c=\frac{1}{\Delta\sigma}$ (surprise !).
\subsection{Lumière blanche}
\paragraph{Détecteur sensible à la couleur} Blanc au milieu, irisation bleu/vert/rouge, puis blanc d'ordre supérieur. Ce sont des oscillations sinusoïdales de l'intensité. Au-delà de la tache de diffraction, l'écran est sombre.
\paragraph{Remarque} Si on prend en compte la diffraction, l'intensité n'est plus sinusoïdale, mais est plutôt en $\text{sinc}^2$ (mais l'étude de la diffraction n'est pas au programme !).
\subsection{Cohérence temporelle ; critère semi-qualitatif de brouillage}
\paragraph{Critère de brouillage} On considère que $\sigma_1$ et $\sigma_2$ sont espacés de $\Delta\sigma/2$. Or :
\begin{itemize}
	\item pour $\sigma_1$ : $p_1=\displaystyle\frac{ax}{D}\sigma_1$ ;
	\item pour $\sigma_2$ : $p_2=\displaystyle\frac{ax}{D}\sigma_2$ ;
\end{itemize} Lorsque le contraste s'annule pour la première fois, on a $p_2-p_1=\pi\displaystyle\frac{1}{2}$. Cela se réécrit : $$\frac{ax}{D}(\sigma_2-\sigma_2)=\pm\frac{1}{2}$$ On obtient les conditions suivantes sur $x$ et $\delta$ pour la première annulation du contraste : $$\begin{cases}
x=\displaystyle\frac{D}{a\Delta\sigma} \\
\delta=\displaystyle\frac{1}{\Delta\sigma}
\end{cases}$$

\chapter{Réseaux}
\localtableofcontents
\section{Direction des pics principaux ; relation fondamentale des réseaux}
On notera dans tout le chapitre $a$ le pas du réseau, \textit{ie} sa période spatiale.
\subsection{Réseaux en transmission}
\paragraph{Dispositif} Il s'agit de $N$ répétitions d’un motif élémentaire (fente très fine) de pas $a\approx 1\mu$m. Pour établir la formule des réseaux, on considère deux fentes adjacentes, un faisceau incident parallèle (collimateur) et émergent parallèle (lunette de visée à l’infini). On établit la différence de marche à l’aide du théorème de Malus et du retour inverse de la lumière, c'est la \textbf{formule des réseaux en transmission} : $$\delta=n_0a\left(\sin(\theta)-\sin(\theta_0)\right)$$
\paragraph{Attention !} Souvent, dans les exercices, $1/a$ est noté $n$.
\paragraph{Condition d'interférences constructives} Il y a interférence constructive si $\delta=p\lambda$ pour $p\in\Z$. Ainsi, les \textit{maxima} d'intensité correspondent à : $$n_0a(\sin(\theta)-\sin(\theta_0))=p\lambda \ \ \text{avec}\ \ p\in\Z$$
\paragraph{Attention} Cette formule suppose l'utilisation :
\begin{itemize}
	\item d'un collimateur ;
	\item d'une lunette de visée à l'infini.
\end{itemize}
\subsection{Réseaux en réflexion}
\paragraph{Réseaux en réflexion} En réflexion, on raisonne de même pour obtenir la \textbf{formule des réseaux en réflexion} : $$\delta=n_0a\left(\sin(\theta)+\sin(\theta_0)\right)$$ On a bien un maximum de réflexion en $\theta=-\theta_0$, conformément aux lois de Snell-Descartes.
\paragraph{\textit{Maxima} d'intensité} En réflexion, les \textit{maxima} d'intensité correspondent à :$$n_0a(\sin(\theta)+\sin(\theta_0))=p\lambda \ \ \text{avec}\ \ p\in\Z$$
\paragraph{Note importante !} Dans toute la suite des chapitres, on ne s'intéresse qu'à des réseaux en transmission.
\section{Conséquences de la formule des réseaux}
\paragraph{Application de la formule des réseaux} Elle permet de déterminer les différents ordres de \textit{maxima} d’intensité. En monochromatique, il n’y a pas de superposition, mais en lumière blanche les ordres deux et trois se chevauchent (car $3\times 400$ nm < $2\times 800$ nm).
\subsection{Cas monochromatique : nombre de pics principaux ; minimum de déviation}
\paragraph{Nombre de pics principaux} On prend $\lambda=632$ nm (laser He-Ne en TP), $a=2$ $\mu$m et $i=0$ (incidence normale). Quelle est la direction des pics principaux ? Quel est le nombre de pic principaux ? On applique la formule des réseaux en transmission en assimilant l'air au vide ($n_0=1$) : $$a\left(\sin(\theta)-\underbrace{\sin(i)}_{=0}\right)=k\lambda_0$$ Par conséquent, tant que $\displaystyle\left\lvert\frac{k\lambda_0}{a}\right\rvert\leqslant 1$, on a : $$\theta_k=\arcsin\left(\frac{k\lambda_0}{a}\right)$$ Numériquement avec nos paramètres, $k$ est inférieur à 3 en valeur absolue, il y a donc 7 pics principaux, \textit{ie} 7 \textit{maxima} d'intensité.
\paragraph{Attention !} $i$ et $\theta$ ne sont pas toujours petits !
\paragraph{Minimum de déviation} La déviation $D$ est définie par : $$D=\theta-i$$ On cherche à minimiser $D$ en faisant varier $i$ et en travaillant à $a$, $\lambda_0$ et $k$ fixés. Or, $dD=d\theta-di=0$ au minimum de déviation. Et en différenciant la formule des réseaux, on trouve $n_0a(\cos(\theta)d\theta-\cos(i)di)=0$, ce qui fournit : $$\cos(\theta_m)=\cos(i)$$ On élimine la solution $\theta_m=i_m$ qui est sans intérêt car elle correspond à l'ordre 0. On a alors la solution $\theta_m=-i_m$ qui est \textbf{très utile en TP}, car elle permet de savoir là où il faut chercher le minimum de déviation. Enfin, en réinjectant dans la formule des réseaux, on obtient : $$\sin(i_m)=\frac{-k\lambda_0}{2na}$$
\subsection{Cas polychromatique : dispersion ; recouvrement d'ordre}
On travaille ici avec de la lumière blanche.
\paragraph{Ordre 0} L'ordre $0$ est non dispersif, \textit{ie} $\theta_0=i$ pour toutes les longueurs d'ondes.
\paragraph{Ordres supérieurs ; recouvrement} En revanche, les ordres des différentes longueurs d'onde varient selon cette dernière. Ainsi, d'après la formule des réseaux, l'ordre -2 pour le violet est le même que l'ordre -1 pour le rouge (car 2$\times$400=800). Par conséquent, ces ordres se chevauchent. De façon générale, pour une source de "lumière blanche", l'ordre 2 et l'ordre 3 se recouvrent : on dit qu'il y a recouvrement.
\section{Largeur des pics principaux}
\subsection{Fonctions réseau}
\paragraph{Cas de $N$ sources cohérentes} Pour $N$ fentes cohérentes, on pose $s_0=S_0e^{i\omega t}$ la vibration due à la fente $0$ seule en M . Donc $s_k = S_0e^{i(\omega t+k\varphi)}$. En sommant puis passant à l’intensité, on obtient $$I = I_0\frac{\sin^2\left(\frac{N\varphi}{2}\right)}{\sin^2\left(\frac{\varphi}{2}\right)}$$ On a donc annulation tous les $\displaystyle\frac{2k\pi}{N\varphi}$ sauf en $0$, $2\pi$, etc. Augmenter $N$ diminue la largeur des \textit{maxima} principaux (demi-largeur égale à $2\pi/N$) et diminue l’intensité des \textit{maxima} secondaires.
\paragraph{Fonctions réseau} On appelle fonctions réseau les fonctions : $$f_n(\varphi)=\frac{\sin^2\left(\frac{N\varphi}{2}\right)}{\left(N^2\sin\left(\frac{\varphi}{2}\right)\right)^2}$$ Ces fonctions vérifient la relation, pour $N$ sources cohérentes : $$I=N^2I_0f_N(\varphi)$$ Elles permettent de tracer plus aisément le graphe de l'intensité (\textit{cf} le paragraphe précédent).
\paragraph{Pic principal} Dans ce cas, le pic principal a pour demi-largeur à la base : $$\Delta\varphi=\frac{2\pi}{N}$$ En termes d'ordre, cela se traduit par le fait que le pic principal se situe entre $k-\displaystyle\frac{1}{N}$ et $k+\displaystyle\frac{1}{N}$ avec $k\in\Z$.
\subsection{Représentation de Fresnel}
\paragraph{Représentation de Fresnel} On représente les uns à la suite des autres les différents vecteurs complexes qui représentent notre signal.
\paragraph{Utilisation pour la première annulation} Pour retrouver la propriété de la première annulation, se dire qu'avec la représentation de Fresnel, il s'agit en fait de faire tourner d'un même angle nos $N$ vecteurs afin de revenir à l'origine. Ainsi, cela correspond à faire un $N$-gone régulier, et donc à $$\frac{2\pi}{N}$$
\subsection{Résumé} On dispose de trois niveaux d'étude : 
\begin{itemize}
	\item Niveau 0 : on considère qu'il n'y a de la lumière qu'au niveau des pics \textit{maxima} (\textit{ie} on utilise la formule des réseaux).
	\item Niveau 1 : on considère qu'il y a de la lumière au niveau des pics principaux, qui se situent entre $k-\displaystyle\frac{1}{N}$ et $k+\displaystyle\frac{1}{N}$.
	\item Niveau 2 : on considère l'ensemble de la lumière et on utilise les fonctions réseau.
\end{itemize}
\section{Utilisation du caractère dispersif}
\subsection{Pouvoir de résolution intrinsèque d'un spectromètre à réseau}
\paragraph{Remarque} Dans cette section, on n'étudie pas l'effet de la largeur de la fente source (pour cela, \textit{cf} exercice de TD).
\paragraph{Critère de Rayleigh} On considère que deux raies sont distinguables lorsque la première annulation de l’une correspond au pic d’intensité de l’autre.
\paragraph{Pouvoir de résolution} On définit le \textbf{pouvoir de résolution} $R$ par $$R=\frac{\lambda}{\Delta\lambda}$$ Avec le critère de Rayleigh, il vérifie $R=Np$ avec $p$ l’ordre et $N$ le nombre de raies éclairés : on le démontre en différenciant la formule des réseaux. Sinon, on peut le retrouver en utilisant le fait que la différence de marche est la même pour les deux : $$\delta=k(\lambda+\Delta\lambda)=\left(k+\frac{1}{N}\right)\lambda$$


\chapter{Exemple de dispositif à division d'amplitude : l'interféromètre de Michelson éclairé par une source spatialement étendue}
\localtableofcontents
\section{Introduction}
L'interféromètre de Michelson est un dispositif mis au point par Albert Michelson et son collaborateur E. Morley à la fin du XIXème siècle, dans le but de mettre en défaut l'hypothèse d'invariance de la vitesse de la lumière dans le vide par changement de référentiel galiléen. L'échec de ces tentatives a eu le mérite de valider expérimentalement avec un gran degré de précision cette hypothèse, laquelle constitue le fondement de la théorie de la relativité restreinte. \\ \par L'interféromètre initial n'a cessé d'être amélioré, son avantage est de permettre d'utiliser une source large (donc d'obtenir une figure d'interférence plus lumineuse) tout en gardant un bon contraste des franges. En 2003 a été inauguré l'interféromètre franco-italien VIRGO dont les bras font 3km de long : c'est un gigantesque interféromètre de Michelson, destiné à la détection d'ondes gravitationnelles, un phénomène prévu par la relativité générale. Le détecteur américain LIGO a en effet observé pour la première fois en 2016 le passage d'ondes gravitationnelles, ce qui constitue (s'il en faut encore...) une confirmation expérimentale spectaculaire des prévision d'Albert Einstein effectuées 100 ans auparavant. Celles-ci auraient été émises par la fusion de deux trous noirs, un phénomène cosmique assez violent pour déformer considérablement l'espace-temps et générer ces ondes gravitationnelles qui se propagent à la vitesse de la lumière dans le vide. Cette découverte ouvre la voie à une toute nouvelle physique car, jusqu'à présent, les seules ondes qui nous donnent des informations sur l'univers sont les ondes électromagnétiques, qui ne nous permettent pas l'observation des tous premiers instants de l'univers (avant 380 000 ans, le libre parcours moyen des photons était très faible car ils étaient diffusés par les électrons libres qui ne s'étaient pas encore combinés avec les protons, la température étant trop importante). Ces interféromètres nous permettraient également, par exemple, de sonde l'énergie noire, cette composante qui expliquerait l'expansion accélérée de l'univers.
\subsection{Présentation de l'interféromètre}
La lumière issue de la source pénètre dans l'interféromètre en traversant un verre anticalorique destiné à filtrer le rayonnement IR (chaleur) émis par la source et qui serait nuisible aux éléments optiques (revêtement des miroirs, lame séparatrice).
\paragraph{Division d'amplitude} La lumière incidente subit une division d'amplitude en traversant la lame séparatrice. C'est la lame de verre possédant une surface partiellement réfléchissante, qui laisse passer la moitié de l'intensité du faisceau et réfléchit l'autre moitié.
\paragraph{Premier trajet de la lumière} Le faisceau (1) qui traverse la séparatrice sans être dévié est réfléchi sur un miroir $(M_1)$ et revient sur la séparatrice. La moitié de l'intensité de ce faisceau (1) est réfléchie et se dirige vers la sortie de l'interféromètre, l'autre moitié traverse la séparatrice et repart vers l'entrée (elle est inutilisable, on ne la représente pas).
\paragraph{Deuxièem trajet de la lumière} Le faisceau (2) qui est réfléchi sur la séparatrice est ensuite réfléchi sur un miroir $(M_2)$ et revient  également sur la séparatrice. La moitie de l'intensité de ce faisceau (2) est transmise et se dirige vers la sortie de l'interféromètre, l'autre moitié est réfléchie sur la séparatrice et repart vers l'entrée (elle est inutilisable, on ne la représente pas).
\paragraph{Les miroirs} L'interféromètre comporte deux miroirs.
\begin{itemize}
	\item Le miroir $(M_1)$ est orientable par deux vis de réglage "grossier" et peut être translaté le long de l'axe entrée/séparatrice : c'est un miroir \textbf{chariotable}.
	\item Le miroir $(M_2)$ est également orientable par deux vis de réglage (plus fin), mais sa position est fixée.
\end{itemize}
\paragraph{Conclusion} Le faisceau incident est divisé en deux parties d'intensités moitié par la séparatrice. Ainsi, deux faisceaux ressortent de l'interféromètre après avoir suivi deux trajets différents. Ces trajectoires (1) et (2) sont dictées par la position et l'orientation des miroirs $(M_1)$ et $(M_2)$.
\paragraph{Intensité} Notons $4I_0$ l'intensité entrante dans l'interféromètre. Si certaines conditions que l'on va étudier sont respectées, il pourra y avoir (en sortie) interférences entre 2 ondes issues d'un même rayon incident, par division d'amplitude. L'intensité en un point $M$ du champ d'interférences s'écrira alors : $$I(M)=2I_0\left(1+\cos\left(\frac{2\pi}{\lambda_0}\delta(M)\right)\right)=2I_0\left(1+\cos\left(2\pi p(M)\right)\right)$$
\paragraph{Nécessité d'une compensatrice} L'interféromètre de Michelson est un dispositif très précis, que l'on va utiliser pour obtenir des interférences en lumière polychromatique (même en lumière blanche). Mais les conditions de cohérence temporelle imposent notamment d'avoir une différence de marche inférieure à la longueur de cohérence temporelle de la source, donc de contrôler très précisément cette différence de marche entre les deux rayons qui interféreront en un point donné. \\ \\ Or, si l'on regarde d'un peu plus près ce qui se passe au niveau de la séparatrice, on constate que le faisceau (1) travers trois fois la lame de verre tandis que le faisceau (2) ne la traverse qu'une seule fois. Cela introduit un déphasage (donc une différence de marche) entre les deux faisceaux, qui n'est pas due à une différence de la longueur des bras du Michelson (la seule contrôlable). La solution retenue est de rajouter une lame de verre de même épaisseur que la séparatrice et parallèle à celle-ci sur le trajet du faisceau (2) : on l'appelle la \textbf{compensatrice}.
\paragraph{Remarques sur la compensatrice}
\begin{itemize}
	\item Le parallélisme de la compensatrice et de la séparatrice est crucial, mais malgré cela, la compensation n'est pas idéale, car elle dépend de l'inclinaison des rayons. Elle n'est parfaite que dans le cas de l'interféromètre réglé en lame d'air avec observation à l'infini.
	\item Pour ne pas déformer les surfaces d'ondes; la planéité des miroirs, de la séparatrice et de la compensatrice est très précise (tolérance de 10 nm).
\end{itemize}
\subsection{Les deux (+$\varepsilon$) configuration de l'interféromètre}
\paragraph{Cas général} Si les deux miroirs ont des positions et des orientations quelconques, on a peu de chances d'obtenir une figure d'interférences (sauf avec un laser). En pratique, on utilisera l'une des deux configurations suivantes.
\paragraph{Réglage en lame d'air} Les miroirs $(M_1)$ et $(M_2)$ sont perpendiculaires entre eux, la séparatrice faisant un angle de $45^\circ$ avec chacun d'entre eux : on l'appelle "réglage en lame d'air". La différence de marche entre les deux rayons qui interfèrent en sortie dépend de la différence de longueur des bras, donc de la position de $(M_1)$.
\paragraph{Réglage en coin d'air} Les miroirs $(M_1)$ et $(M_2)$ ne sont pas perpendiculaires entre eux, mais les longueurs des bras sont quasiment égales : on l'appelle le "réglage en coin d'air". La différence de marche entre deux rayons qui interfèrent en sortie dépend de l'angle entre $(M_1)$ et $(M_2)$.
\paragraph{Explication des noms} On comprendra mieux les noms de ces deux types de réglages par la suite.
\section{Utilisation en configuration "lame d'air" ; anneaux d'égale inclinaison (ou anneaux de Haidinger)}
\subsection{Description ; montage équivalent}
\paragraph{Éclairage et observation : constations expérimentales}
On souhaite éclaire le Michelson par une source étendue pour avoir une meilleure luminosité. Dans ce cas, on constate, même lorsque le réglage des miroirs est idéal, que les interférences sont d'autant plus visibles sur un écran placé en sortie que l'écran est éloigné. \\ \\ \textbf{Pour un Michelson réglé en lame d'air et éclairé par une source étendue, les franges d'interférences sont localisées à l'infini.} En pratique, bien sûr, on place l'écran ou le détecteur dans le plan focal image d'une lentille convergente.
\paragraph{Remarques}
\begin{itemize}
	\item On n'est pas surpris d'obtenir un brouillage des interférences avec une source étendue (problème de cohérence spatiale), on l'est plus de parvenir à conserver une figure d'interférence très nette à l'infini, ceci quelle que soit l'étendue de la source !
	\item On pourrait éclaire le Michelson en lame d'air avec une source ponctuelle : on obtiendrait alors des interférences non localisées (observation quelle que soit la distance de l'écran), mais peu lumineuses.
\end{itemize}
\paragraph{Différence de marche entre deux rayons issues d'un point source $S$. Ordre d'interférence} On trace les deux trajets possibles d'un rayon incident issu d'un point source $S$ et faisant un angle $i$ avec l'axe du bras portant le miroir $(M_1)$. On sait que $M$ est dans le plan focal image de la lentille $(L)$ de centre $O$ donc les deux rayons (1) et (2) qui interfèrent sont parallèles à $(OM)$ avant la lentille. D'après le théorème de Malus, on a donc : $(P_1M)=(P_2M)$ (chemins optiques). Par conséquent, la différence de marche $\delta=(SM)_2-(SM)_1$ entre les deux rayons vaut : $$\delta =\left((SA)+\frac{\lambda_0}{2}+(AC)+\frac{\lambda_0}{2}+(CP_2)+(P_2M)\right)-\left((SA)+(AB)+\frac{\lambda_0}{2}+(BD)+\frac{\lambda_0}{2}+(BP_1)+(P_1M)\right)$$ où les $\displaystyle\frac{\lambda_0}{2}$ proviennent des réflexions. Par conséquent : $$\delta=\bigg((AC)+(CP_2)\bigg)-\bigg((AB)+(BD)+(DP_1)\bigg)$$
\paragraph{Remarque} Suivant le traitement de la séparatrice, il peut ne pas y avoir de différence de marche supplémentaire $\displaystyle\frac{\lambda_0}{2}$ à la dernière réflexion à l'interface verre/métal.
\paragraph{Miroir fictif} On appelle $(M_1\st)$ le symétrique du miroir $(M_1)$ par rapport à la séparatrice et on trace les symétriques rayons $(AB)$ et $(BD)$.
\paragraph{Première méthode}
En termes de chemin optique, la situation est identique au trajet réel de la lumière. En déplaçant le plan d'onde $(P_1P_2)$ juste à la sortie du miroir $(M_2)$, on a alors : $$\lvert\delta\rvert=(CP_2)-(CB')-(B'P_1)$$ Or, on a $(CB')=(B'P_1)=\displaystyle\frac{2ne}{\cos(i)}$ et $(CP_1)=\displaystyle 2ne\tan(i)$ donc $$(CP_2)=2ne\tan(i)\sin(i)=\frac{2ne\sin^2(i)}{\cos(i)}$$ On en déduit que $$\lvert\delta\rvert =2ne\frac{1-\sin^2(i)}{\cos(i)}=2ne\cos(i)$$
\paragraph{Remarques}
\begin{itemize}
	\item On garde $n$ ici même si on travaille dans l'aire car on pourra mettre en évidence la différence entre $n$ et $1$ grâce à ce dispositif.
	\item La valeur absolue permet de généraliser le résultat au cas où $(M_1\st)$ est plus proche de la séparatrice que $(M_2)$.	
\end{itemize}
\paragraph{Explication du nom} On constate que la différence de marche est équivalente à celle d'une lame d'air d'épaisseur $e$ située entre $(M_1\st)$ et $(M_2)$, d'où le nom de réglage en "lame d'air".
\paragraph{Deuxième méthode}
On peut aussi appeler $S\st$ le symétrique de $S$ par rapport à la séparatrice, et $S_1$ et $S_2$ les symétriques de $S\st$ par rapport à $(M_1\st)$ et $(M_2)$ respectivement, appelées \textbf{sources fictives secondaires}. On obtient le montage équivalent qui est le plus pratique pour déterminer la différence de marche. Du point de vue du chemin optique parcouru par les deux rayons qui interfèrent en $M$, tout se passe comme s'ils s'étaient propagés en ligne droite depuis $S_1$ et $S_2$. On se retrouve comme dans le montage des trous d'Young avec deux sources secondaires, mais on étudie l'intensité dans un plan orthogonal à l'axe $(S_1S_2)$. On constate alors immédiatement que la différence de marche s'écrit : $$\delta =(S_1H)=nS_1S_2\cos(i)=2ne\cos(i)$$
\paragraph{Conclusion} Lorsque l'interféromètre de Michelson est réglé en lame d'air, la différence de marche et l'ordre d'interférence entre deux rayons \textbf{issus du même rayon incident}, qui interfèrent à l'infini, sont donnés par : $$\lvert\delta\rvert=2ne\cos(i) \ \ \text{et} \ \ \lvert p\rvert=\frac{\lvert\delta\rvert}{\lambda_0}=\frac{2ne\cos(i)}{\lambda_0}$$
\subsection{Analyse de la figure d'interférences}
\paragraph{Étude qualitative}
On peut intuiter l'allure de la figure d'interférences à l'aide des symétries. Le montage équivalent possède une symétrie de révolution, donc la figure d'interférences également d'après le principe de Curie : celle-ci consiste en des anneaux concentriques. \\ \\ Lorsqu'on a deux sources secondaires (ici fictives), le lieu des points de l'espace tel que $S_2M-S_1M=cte$ est un hyperboloïde de révolution, de foyers $S_1$ et $S_2$. Donc le lieu des points d'intensité maximale $S_2M-S_1M=m\lambda$ avec $m\in\Z$ consiste en un réseau d'hyperboloïdes de mêmes foyers. La figure d'interférences observée sur un écran de sortie est donc l'intersection de ce réseau d'hyperboloïdes avec le plan de l'écran : ce sont bien des anneaux. \\ \\  Pour $e$ et $\lambda$ fixées, un anneau brillant (ou sombre) correspond à $i=cte$ : il est donc constitué de l'interférence de tous les rayons incidents ayant la même inclinaison $i$ avec l'axe d'entrée $O_SO_1$, c'est pourquoi on parle d'anneaux de même inclinaison, ou anneaux de Haidinger.
\paragraph{Étude quantitative : rayon des anneaux} Les rayons incidents ont des incidences faibles, donc on a : $$\cos(i)\approx 1-\frac{1}{2}\tan^2(i)=1-\frac{r^2}{2f'^2}$$ En réinjectant dans les formules précédentes, on trouve finalement : $$p(M)=\frac{2ne}{\lambda_0}\left(1-\frac{r^2(M)}{2f'^2}\right) \ \ \text{et} \ \ r(M)=f'\sqrt{2\left(1-\frac{\lambda_0}{2ne}p(M)\right)}$$
\paragraph{Commentaires}
\begin{itemize}
	\item \textbf{En lame d'air, l'ordre au centre est non nul} (sauf si on est au contact optique, \textit{ie} pour $e=0$).
	\item L'ordre est maximal au centre, en $F'$ : $p_\text{max}=\displaystyle\frac{2ne}{\lambda_0}\in\R_+$. Il correspond à l'incidence nulle $i=0$. \textbf{L'ordre diminue au fur et à mesure que le rayon augmente.} En particulier, s'il est nul au centre, il est nul partout.
	\item Nombre $N$ d'anneaux visibles : cela dépend de l'étendue angulaire du faisceau incident $i_\text{max}$ et de la taille de la zone d'observation. En général, c'est $i_\text{max}$ qui est limitant : $$N=\lfloor p_\text{max}-p_\text{min}\rfloor =\left\lfloor\frac{\delta_\text{centre}-\delta_\text{bord}}{\lambda_0}\right\rfloor =\left\lfloor \frac{ei_\text{max}^2}{\lambda_0}\right\rfloor$$ \textbf{Le nombre d'anneaux visibles est d'autant plus grande que l'épaisseur de la lame d'air est grande.}
	\item Le rayon d'un anneau d'ordre $p$ est une fonction croissante de $e$. Donc si on chariote $(M_1)$ de manière à diminuer l'épaisseur $e$ de la lame d'air, $r$ diminue et $p_\text{max}$ aussi. Ainsi, les anneaux rétrécissent et disparaissent au centre. \textbf{Pour diminuer l'épaisseur de la lame d'air, il faut charioter $(M_1)$ de manière à faire rétrécir les anneau. Il y a alors de moins en moins d'anneaux visibles sur l'écran.} Cette remarque est très utile lorsque l'on cherche à régler le Michelson de manière à avoir $e=0$. On dit alors que l'on est au \textbf{contact optique}. Lorsqu'on arrive au contact optique, la différence de marche est nulle quelle que soit l'incidence $i$, donc l'écran est uniformément éclairé : c'est la teinte plate.
	\item Étant donnée une épaisseur $e$ fixée, pour voir beaucoup d'anneaux, il faut que la lame d'air soit éclairée avec des incidences très variées. Pour cela, on fait fortement converger la lumière incidente sur les miroirs. On obtient typiquement une \textbf{demi ouverture angulaire de $10^\circ$ à $15^\circ$}.
\end{itemize}
\paragraph{Rayon de l'anneau brillant numéro $k$ en partant du centre}
Le rayon de l'anneau $k$ est alors : $$r_k=f'\sqrt{\frac{\lambda}{2ne}\left(\left\{p_\text{max}\right\}+k-1\right)}$$ où $\left\{\cdot\right\}$ désigne la fonction partie fractionnaire. \textbf{Attention, la formule est très probablement douteuse !!!} Néanmoins, le rayon des anneaux varie bien en $\sqrt{k}$ donc \textbf{les anneaux sont de plus en plus resserrés en s'éloignant du centre}.
\subsection{Source étendue ; localisation des interférences à l'infini}
On a dit qu'avec une source étendue, les lieux d'éclairement maximale dans l'espace constituent des réseaux d'hyperboloïdes. Donc, quelle que soit la position de l'écran, on observe des franges d'interférences. On dit que les interférences sont délocalisées. \\ \\ Si on étend la source $S$, les sources secondaires $S_1$ et $S_2$ s'étendent aussi, et on a une superposition des réseaux d'hyperboloïdes dus à chacun des points sources (donc décalés les uns par rapport aux autres), donc \textit{a priori} brouillage. \\ \\ Mais lorsqu'on observe à l'infini comme dans note montage, on voit que si l'on déplace $S$ dans n'importe quelle direction en $S'$, cela provoque une translation de $S_1$ et $S_2$ d'un même vecteur : l'axe $(S_1S_2)$ se retrouve donc parallèle mais non confondu avec l'axe de la lentille $(OF')$. Le rayon d'incidence $i$ donnera toujours naissance à deux rayons qui sortent de l'interféromètre avec un angle $i$, donc qui interfèrent au même point $M$ de l'écran. On constate donc que malgré un déplacement de $S$, donc de $S_1$ et $S_2$, les rayons d'incidence $i$ interféreront toujours au même point $M$ : il n'y a pas de déplacement de la figure d'interférences. \\ \\ Si l'on étend la source, les points sources incohérents entre eux donnent des figures d'interférences strictement identiques qui s'additionnent en intensité. \\ \\ Ce raisonnement  est valable uniquement en observant la figure d'interférences à l'infini, sinon il y a brouillage de la figure d'interférences. \\ \\ Conclusion : lorsqu'on étend la source, les figures d'interférences à l'infini dues à chacun des points sources se superposent exactement : on dit que les interférences sont localisées à l'infini. La figure garde donc un très bon contraste et gagne beaucoup en luminosité !
\subsection{Source polychromatique ; interférométrie par transformée de Fourier}
\paragraph{Effet qualitatif de la largeur spectrale de la source} On considère que le Michelson est réglé en lame d'air, et qu'on observe les anneaux dans le plan focal image d'une lentille convergente (donc à l'infini). On a considéré pour l'instant une source parfaitement monochromatique, ce qui n'existe pas. Même si on isole une raie spectrale à l'aide d'un filtre, la source possède tout de même une certaine largeur spectrale. On a vu que des ondes de fréquences différentes étaient incohérentes, on peut donc dire que les différentes composantes spectrales de la source donnent des figures d'interférences qui s'ajoutent en intensité (donc elles se "superposent"). \\ \\ Comme le rayon des anneaux dépend de la la longueur d'onde, les anneaux correspondant à chaque longueur d'onde se décalent les uns par rapport aux autres, et conduisent à un brouillage de la figure d'interférence si les anneaux sombres de certaines longueurs d'ondes correspondent aux anneaux brillants d'autres longueurs d'ondes.
\paragraph{Éclairement par le doublet jaune du sodium ; interférogramme}
Supposons que le dispositif interférentiel soit éclairé par une lampe spectrale à vapeur de sodium, avec un filtre ne laissant passer que le doublet jaune caractéristique : $\lambda_1=589$ nm et $\lambda_2=589,6$ nm. On néglige pour l'instant la largeur chaque raie, et on suppose que les deux raies ont la même intensité $I_0$. On note $$\lambda_m=\frac{\lambda_1+\lambda_2}{2}=589,3\ \text{nm} \ \  \text{et} \ \ \Delta\lambda=\lambda_2-\lambda_1=0,6\ \text{nm}$$ On introduit les nombres d'ondes $$\sigma_1=\frac{1}{\lambda_1} \ \ \text{et} \ \  \sigma_2=\frac{1}{\lambda_2}$$ On note également $$\sigma_m=\frac{\sigma_1+\sigma_2}{2} \ \ \text{et} \ \ \Delta\sigma=\sigma_2-\sigma_1$$ si bien que l'on a : $$\sigma_1=\sigma_m-\frac{\Delta\sigma}{2} \ \ \text{et} \ \ \sigma_2=\sigma_m+\frac{\Delta\sigma}{2}$$
\paragraph{Expression du contraste} Les sources sont incohérentes donc $$I(\delta)=I_0(1+\cos(2\pi\sigma_1\delta))+I_0(1+\cos(2\pi\sigma_2\delta))$$ En factorisant et par trigonométrie, on obtient : $$I(\delta)=2I_0\left[1+\cos(\pi\Delta\sigma\delta)\cos(2\pi\sigma_m\delta)\right]$$ On constate l'apparition d'un terme de contraste : $$C=\left\lvert\cos(\pi\Delta\sigma\delta)\right\rvert$$
\paragraph{Tracé expérimental d'un interférogramme}
On peut tracer expérimentalement la courbe $I(\delta)$ avec le Michelson réglé en lame d'air. Pour cela, on place une cellule photoélectrique au centre des anneaux (en $F'$), et on se place au contact optique ($\delta=0$). Puis on déplace le miroir $(M_1)$ à vitesse constante, grâce à un moteur, de manière à augmenter l'épaisseur de la lame d'air (donc on augmente $\delta$ aussi car $\delta=2ne$ en ce point). La cellule va alors enregistrer les variations de $I$ en fonction du temps, donc de $\delta$. \\ \\ Dans le cas du doublet, on observe un brouillage périodique : ce sont des \textbf{battements}.
\paragraph{Notion d'anticoïncidence} Lorsque $\delta=0$ (\textit{ie} $p=0$), l'intensité lumineuse est maximale quelle que soit la longueur d'onde. Mais lorsque $\delta$ et $p$ augmentent, les zones d'intensités maximales ou minimales se décalent. On dit qu'il y a anticoïncidence lorsqu'un anneau sombre de $\lambda_1$ correspond à un anneau brillant de $\lambda_2$. On a alors : $$p_1(\delta)=\delta\sigma_1=m_1+\frac{1}{2} \ \ \text{et} \ \ p_2(\delta)=\delta\sigma_2=m_2 \ \ \text{avec} \ \ (m_1,m_2)\in\Z^2$$ Les anticoïncidences ont donc lieu lorsque $\displaystyle\delta\times(\sigma_1-\sigma_2)=(m_1-m_2)+\frac{1}{2}$, donc quand : $$\delta=\frac{m+\frac{1}{2}}{\Delta\sigma} \ \ \text{avec} \ \ m\in\Z$$ On retrouve ainsi les annulations du contraste. L'écart entre deux anticoïncidences successives est donc : $$\Delta\delta=\frac{1}{\Delta\sigma}\approx\frac{\lambda_m^2}{\Delta\lambda}$$ Application numérique : pour le doublet du sodium, on trouve $\Delta\delta=578,8$ $\mu$m, soit : $$N=\frac{\Delta\delta}{\lambda_m}=\frac{\sigma_m}{\Delta\sigma}=\frac{\lambda_m}{\Delta\lambda}=982$$ franges entre deux anticoïncidences. En comptant simplement les franges entre deux anticoïncidences sur l'interférogramme, on peut donc mesure $\Delta\lambda$.
\paragraph{Éclairage en lumière blanche}
Éclairé par une source de lumière blanche, un dispositif interférentiel donne donc des franges irisées, du moins si la différence de marche n'est pas trop élevée. En effet, le brouillage intervient très rapidement dès que $\delta$ augmente, car la source est spectralement très large. Avec le Michelson réglé en lame d'air, il est indispensable d'être très près du contact optique pour observer des anneaux. \\ \\ Par exemple, si on observe la lumière reçue au centre de l'écran avec le Michelson réglé en lame d'air :
\begin{itemize}
	\item au contact optique : $\delta=0$ donc les interférences sont constructives pour toutes les longueurs d'onde, la lumière est \textbf{blanche} ;
	\item lorsque $\delta\approx0,25$ $\mu$m (environ une demi-longueur d'onde), les intensités de toutes les longueurs d'ondes sont proches de leur minimum, on obtient du \textbf{noir} ;
	\item à partir de $\delta\approx 1$ $\mu$m, il y a toujours une radiation pour laquelle l'intensité est proche de son maximum et une autre qui est proche de $0$. L'intensité totale est à peu près constante, mais on distingue des couleurs appelées \textbf{teintes de Newton}. Par exemple, pour $\delta\approx 1,7$ $\mu$m, le rouge et le bleu sont éteints, mais le vert a une intensité quasi maximale. On observe donc du vert.
	\item à partir de $\delta\approx 3$ $\mu$m, pour chaque valeur de $\delta$, plusieurs radiations sont éteintes, et d'autres sont à leur maximum d'intensité : l'œil ne discerne plus de couleur et perçoit donc du blanc. Comme ce n'est pas un blanc pur, on parle de \textbf{blanc d'ordre supérieur}. Si l'on trace le spectre de cette lumière, on discerne les longueurs d'ondes éteintes. : on dit que le spectre est \textbf{cannelé} (pschit). Plus $\delta$ est grande, plus y a de cannelures.
\end{itemize}
\paragraph{Remarque} Si le but est d'observer le spectre cannelé et non de visualiser les anneaux, on éclaire plutôt le Michelson réglé en lame d'air par un faisceau de lumière parallèle à l'axe optique (l'éclairement est alors uniforme en sortie). On peut interposer entre la sortie de l'interféromètre et l'écran un prisme afin de disperser la lumière reçue. Les cannelures sont alors bien visibles à la sortie du prisme.
\section{Utilisation en configuration "coin d'air" ; franges d'égale épaisseur (ou franges de Fizeau)}
\subsection{Description ; montage équivalent}
\paragraph{Éclairage et observation}
On souhaite toujours éclairer le Michelson par une source étendue pour avoir une meilleure luminosité. Mais à la différence du montage "lame d'air", on constate que la figure d'interférence s'observe alors à l'oeil nu "sur les miroirs". De plus, si l'étendue angulaire du faisceau incident est trop importante, il y a brouillage. On peut donc utiliser une source étendue rejetée à l'infini (\textit{ie} dans le plan focal objet d'une lentille convergente). En pratique, pour un Michelson réglé en coin d'air, on éloigne la source étendue (lampe) ou on la place dans le plan focal d'une lentille convergente, et les franges d'interférences sont localisées au voisinage des miroirs. On les regarde à l'oeil nu ou on les projette sur un écran que l'on conjugue avec le plan des miroirs grâce à une autre lentille convergente.
\paragraph{Différence de marche entre deux rayons incidents issus d'un point source S. Ordre d'interférence} Analysons les deux trajets possibles d'un rayon incident issu d'un point source S (à l'infini) et faisant un angle $i$ petit avec l'axe du bras portant le miroir $(M_1)$, les deux miroirs faisant un petit angle $\alpha$ entre eux. La différence de marche $\delta=(SM)_2-(SM)_1$ entre les deux rayons qui interfèrent en un point $M$ de l'écran vaut : $$\delta=\left[(SA)+\frac{\lambda_0}{2}+(AC)+\frac{\lambda_0}{2}+(CM)\right]-\left[(SA)+(AB)\frac{\lambda_0}{2}+(BD)+\frac{\lambda_0}{2}+(DM)\right]$$ On obtient en simplifiant l'expression : $$\delta=\left[(AC)+(CM)\right]-\left[(AB)+(BD)+(DM)\right]$$ En général, il est commode d'utiliser le miroir fictif $(M_1\st)$ avec $(M_2)$ pour visualiser le coin d'air ainsi formé, caractérisé par son arête et l'angle $\alpha$. En termes de chemin optique, la situation est identique au trajet réel de la lumière. On constate que tout se passe comme si le rayon interférait avec lui-même (en réalité, avec l'autre partie de son amplitude) en $M'$ situé très près des miroirs. La lentille de sortie sert à faire l'image de $M'$ sur l'écran, donc la condition de stigmatisme nous dit que $(M'M)_1=(M'M)_2$. Finalement, on a donc : $$\lvert\delta\rvert=(B'C)+(CM')-(B'M')$$ Il est évident que quel que soit l'angle d'incidence, la différence de marche est nulle sur l'arête, qui constitue la frange d'interférence d'ordre 0.
\paragraph{Incidence nulle} Pour une incidence $i$ nulle, le rayon arrive perpendiculairement à $(M_2)$ e est réfléchi sur lui-même. La différence de marche est donc égale à : $$\lvert\delta\rvert=2(B'C)=2ne(M')$$ On constate que la différence de marche est équivalente à celle d'un coin d'air, c'est-à-dire d'une lame d'air d'épaisseur variable $e(M')$ située entre $(M_1\st)$ et $(M_2)$, d'où le nom de réglage en "coin d'air".
\paragraph{Influence de l'angle d'incidence} Comme on utilise une source étendue pour augmenter la luminosité, on ne peut pas avoir uniquement des rayons d'incidence nulle. La surface de localisation correspondant à l'angle $i$ est le plan contenant l'arête et le point $M'$. Ce plan de localisation reste au voisinage des miroirs, mais diffère en fonction de l'angle d'incidence. De plus, l'expression de la  différence de marche varie un peu par rapport au cas de l'incidence nulle. Cependant, en pratique, si les angles $i$ et $\alpha$ sont suffisamment petits, on ne détecte pas ce décalage entre les figures d'interférences correspondant au différents angles d'incidence, et on voit des interférences "au voisinage des miroirs". \\ \par Conclusion : l'extension de la source modifie la surface de localisation des interférences, mais tant qu'on garde un angle $\alpha$ "petit" (typiquement inférieur à $1'$ d'arc, soit $3.10^{-4}$ rad) et qu'on se limite à une source "modérément étendue" (incidences inférieures à une dizaine de $^\circ$), le brouillage n'intervient que pour des ordres d'interférence très élevés. \\ \par Sous ces conditions également, l'expression de la différence de marche reste quasiment identique à celle en incidence normale : $$\lvert\delta\rvert=2ne(M')$$ Conclusion : lorsque l'interféromètre de Michelson est réglé en coin d'air, la différence de marche et l'ordre d'interférence entre deux rayons issus du \textbf{même rayon incident} qui interfèrent au voisinage des miroirs sont : $$\lvert\delta\rvert=2ne(M')\text{  et  }\lvert p\rvert=\frac{\lvert\delta\rvert}{\lambda_0}=\frac{2ne(M')}{\lambda_0}$$
\paragraph{Ordres de grandeur} On éclaire l'interféromètre avec un faisceau d'ouverture angulaire $\Delta i$ autour de l'incidence normale. Pour observer 15 franges de chaque côté de l'arête, il faut avoir $\Delta i\approx 30^\circ$, ce qui n'est pas du tout limitant. Pour observer 500 franges de chaque côté de l'arête, il faut avoir $\Delta i\approx 5^\circ$ : il faut limiter le faisceau incident par une fente assez fine, ce qui diminue la luminosité. Si on veut observer beaucoup de franges, on utilise plutôt le Michelson en lame d'air.
\subsection{Analyse de la figure d'interférences}
\paragraph{Étude qualitative} Cette fois, le montage a perdu sa symétrie de révolution, faisant apparaître une direction transverse privilégié : celle autour de laquelle le miroir $(M_1)$ a pivoté. On appelle \textbf{arête du coin d'air} l'intersection des miroirs $(M_1\st)$ et $(M_2)$. Pour comprendre l'allure de la figure d'interférences, on peut, comme pour le montage en lame d'air, représenter le dispositif équivalent avec des sources fictives secondes $S_1$ et $S_2$, en supposant que la source $S$ est ponctuelle. Cette fois, les sources $S1$ et $S_2$ sont quasiment parallèle au plan d'observation des interférences (le plan du miroir $(M_1)$). On obtient donc les mêmes hyperboles très aplaties; de la même manière qu'avec le dispositif des trous d'Young. En réalité, on observe des franges quasi-rectilignes, parallèles à l'arête du coin d'air. Une frange brillante correspond à $\delta=\text{cte}$ donc à $e(M')=\text{cte}$. On parle de \textbf{franges d'égale épaisseur}, ou encore de \textbf{franges de Fizeau}.
\paragraph{Étude quantitative} On a donc : $$I(M)=2I_0\left(1+\cos\left(\frac{2\pi}{\lambda_0}\delta\right)\right)=2I_0(1+\cos(2\pi p))$$ avec $\delta=2ne(M')$ et $\displaystyle p=\frac{\delta}{\lambda_0}=\frac{2ne(M')}{\lambda_0}$. Si on note $x$ la distance à l'arête du coin d'air, on a $e(M')=\alpha x$ car $\alpha\ll 1$ donc on trouve les expressions suivantes de $\delta(M')$ et $p(M')$ : $$\delta(M')=2N\alpha x\text{  et  }p(M')=\frac{2n\alpha x}{\lambda_0}$$ Finalement, on obtient l'expression de l'interfrange : $$i=x_{p+1}-x_p=\frac{\lambda_0}{2n\alpha}$$
\paragraph{Remarques} On peut faire les remarques suivantes.
\begin{itemize}
	\item C'est l'interfrange dans le plan de localisation, au voisinage des miroirs. Si l'on visualise les interférences en faisant l'image de ce plan sur un écran à l'aide d'une lentille, ce montage introduit un grandissement transversal $\gamma$.
	\item \textbf{La frange d'ordre 0 est située sur l'arête du coin d'air}.
	\item L'interfrange $i$ est une fonction décroissante de $\alpha$. Les franges sont d'autant plus serrées que l'angle du coin d'air est grand.
	\item Le nombre $N$ de franges rectilignes visibles dépend de l'angle du coin d'air et est limité par la taille des miroirs. Si le rayon des miroirs est $R$, on a : $$N=\left\lfloor\frac{2R}{I}\right\rfloor+1=\left\lfloor\frac{2n\alpha R}{\lambda_0}\right\rfloor+1$$ Le nombre de franges visibles est d'autant plus grand que l'angle du coin d'air est grand.
	\item Ainsi, si on tourne une vis de réglage fin de $(M_2)$ de manière à diminuer l'angle $\alpha$ du coin d'air, $i$ augmente et $N$ diminue. Les franges s'écartent et disparaissent sur les bords du miroir. Cette dernière remarque est très utile quand on cherche à régler le Michelson de manière à avoir $\alpha=0$. Pour visualiser confortablement des interférences à l'oeil nu sur des miroirs, $i$ doit être au moins de l'ordre du mm donc $\alpha$ doit être inférieur à $3.10^{-4}$ rad (1' d'arc).
\end{itemize}
\paragraph{Effet d'une translation de $(M_1)$} Si les deux bras de l'interféromètre ont exactement la même longueur, l'arête du coin d'air est juste au milieu des miroirs. Si l'on translate le miroir $(M_1)$, cela aura pour conséquence de déplacer l'arête horizontalement donc de déplacer les franges rectilignes. Les franges défilent dans le sens du glissement de l'arête.
\paragraph{Observations en lumière blanche} Étant donnée la faible longueur de cohérence, on ne peut visualiser des interférences en lumière blanche que si l'on est très près du contact optique ($\delta=0$) avec un petit angle $\alpha$ : il faut que l'arête du coin d'air soit sur les miroirs (et pas à côté). La frange centrale blanche (ou noire suivant le traitement de la séparatrice) est encadrée de deux franges noires (ou blanches) puis de franges irisées, puis de blanc d'ordre supérieur.
\subsection{Applications}
\paragraph{Visualisation d'une variation d'indice}
Si l'on vaporise un jet de gaz d'indice $n\ne n_{\text{air}}$ de diamètre $d$ devant l'un des miroirs, par exemple grâce à une bombe de dépoussiérant, on visualise une modification des franges due à l'indice du gaz différent de celui de l'air ambiant : la différence de marche supplémentaire due à la variation d'indice $\delta_\text{indice}=2(n_{\text{gaz}}-n_\text{air})d$ se traduit par un déplacement des franges. \\ En effet, en un point $m$ donné, la traversée de l'épaisseur $d$ de gaz d'indice $n_\text{gaz}>n_\text{air}$ provoque une variation de l'ordre d'interférence : $$\Delta p=\frac{\delta_\text{indice}}{\lambda_0}=\frac{2(n_\text{gaz}-n_\text{air})d}{\lambda_0}$$ Supposons que $d\approx3$ mm. On peut expérimentalement avoir $\Delta p=3$ et $\lambda_0=480$ nm, ce qui nous permet de trouver $$n_\text{gaz}-n_\text{air}\approx 2,5.10^{-4}$$ Pour effectuer une mesure quantitative plus précise, il serait préférable de travailler en lumière quasi-monochromatique, par exemple en éclairant par une lampe à vapeur de mercure munie d'un filtre vert pour isoler la raie verte de longueur d'onde $\lambda_0=546,1$ nm. On constate donc qu'on a accès à de très faibles variations de l'indice optique. \\ \\  De même, il est possible de visualiser le gradient d'indice provoqué par les variations de température au voisinage d'une flamme.
\paragraph{Mise en évidence d'une déformation d'un miroir} On a vu qu'une différence de marche supplémentaire égale à $\lambda$ provoque un décalage local d'un interfrange. On en déduit que toute déformation ou de défaut de planéité d'un miroir, de la séparatrice ou de la compensatrice de l'ordre de la longueur d'onde provoque une déformation des franges d'interférences clairement visible. Ces éléments optiques nécessitent un polissage extrêmement précis : typiquement, les défauts sont inférieurs à $\lambda/20$ d'après les constructeurs.
\section{Résumé}
\subsection{Lame d'air}
\begin{itemize}
	\item $(M_1\st)$ est parallèle à $(M_2)$, espacés de $e$ ;
	\item la différence de marche vaut $\delta=2ne\cos(i)$
	\item avec une source étendue, il faut observer à l'infini ;
	\item si on veut beaucoup d'anneaux, il faut augmenter $i_\text{max}$ $\rightarrow$ faire l'image de la source sur les miroirs.
\end{itemize}
\subsection{Coin d'air}
\begin{itemize}
	\item les miroirs $(M_1\st)$ et $(M_2)$ ne sont plus parallèles, et forment un angle $\alpha$ ;
	\item la différence de marche vaut $2n\alpha x$ ;
	\item avec une source étendue, les interférences sont localisées au voisinage des miroirs (si on met au point sur le miroir, il y a un très bon contraste) et il faut un éclairage parallèle.
\end{itemize}
\subsection{Passer d'un montage à l'autre}
\paragraph{Contact optique $\to$ coin d'air} Il faut augmenter $\alpha$, les franges se rapprochent.
\paragraph{Contact optique $\to$ lame d'air} Il faut augmenter $e$ : les anneaux émergent du centre.
\paragraph{Passage de lame d'air à coin d'air} En général, on commence par le réglage en lame d'air. Pour passer en coin d'air, on se rapproche du contact optique (en diminuant $e$, les anneaux rentrent au centre) puis en augmentant $\alpha$. 

\part{Mécanique quantique}
\chapter{Description quantique d'une particule}
\localtableofcontents
\section{Dualité onde-corpuscule}
\subsection{Onde de matière de de Broglie ; expérience de Davisson et Germer}
\paragraph{Dualité onde-corpuscule}
Mise en évidence de la nature corpusculaire :
\begin{itemize}
	\item Rayonnement du corps noir (Planck, 1900) ;
	\item Effet photoélectrique (Einstein, 1905) ;
	\item Effet Compton
\end{itemize}
Mise en évidence de la nature ondulatoire :
\begin{itemize}
	\item Diffraction, interférences ;
	\item Idée de De Broglie : diffraction de la matière. Prouvée par Davisson et Germer en observant la même figure d'interférences avec des électrons que des rayons X sur un réseau.
\end{itemize}
\paragraph{Hypothèse de de Broglie} Les électrons, photons, etc. ont aussi un caractère ondulatoire.
\paragraph{Confirmation expérimentale} L'expérience de Davisson et Germer, avec la diffraction d'électrons par un cristal de Ni, confirme cette hypothèse.
\subsection{Expérience des fentes d'Young avec des photons ou des électrons}
\paragraph{Expérience des trous d'Young photon par photon} On envoie des électrons ou des photons sur un dispositif de trous d'Young. Pour un long temps de pose, on retrouve la figure d’interférences. Mais pour un unique photon, sa position finale est aléatoire, on n’en connait que la probabilité (plus forte dans les franges claires).
\paragraph{Bilan} Ainsi, un choc unique correspond au caractère corpusculaire, et une figure de diffraction correspond au caractère ondulatoire.
\subsection{Principe de complémentarité de Bohr}
\paragraph{Extrait de Cohen-Tannoudji} Si on arrive à déterminer la variation de quantité de mouvement du plan des trous d'Young, alors la figure d'interférences est brouillée : en effet, il faut aussi considérer la plaque contenant les trous d'Young comme un objet quantique.
\paragraph{Principe de complémentarité de Bohr} Un système a soit un comportement corpusculaire (position et trajectoire connues), soit ondulatoire (diffraction, interférences) mais pas les deux. Exemple des trous d’Young avec fentes mobiles : si l’on sait par quelle fente le quanton est passé, on perd la figure d’interférence.
\subsection{Relations de Planck-Einstein ; relation de de Broglie}
\paragraph{Remarque : culture} Planck a publié en 1900 son article sur le corps noir et Einstein a publié en 1905 son article sur l'effet photoélectrique.
\paragraph{Relation de Planck-Einstein} $$E=h\nu=\hbar\omega$$ où $\hbar=\displaystyle\frac{h}{2\pi}$ est la constante de Planck réduite avec $h$ la constante de Planck ($h=6,62.10^{-34}$ J.s).
\paragraph{Remarque : effet photoélectrique} On considère un rayonnement incident sur un métal. Si $E_\text{photon}>W$ où $W$ est le travail d'extraction, alors un électron est arraché. Or, pour un photon dans le vide, $f=\displaystyle\frac{c}{\lambda}$, donc cela revient à avoir $$\lambda <\frac{hc}{W}=\lambda_\text{seuil}$$ Ce seuil ne dépend donc pas de l'intensité du rayon incident. En utilisoant la relation de de Broglie (\textit{cf} après), on trouve $E_\text{photon}=pc$.
\paragraph{Attention !} Un photon, même s'il est de masse nulle, possède une quantité de mouvement.
\paragraph{Relation de De Broglie} $$p=\frac{h}{\lambda} \ \ \text{ou} \ \ \overrightarrow{p}=\hbar\overrightarrow{k}$$
\paragraph{Remarque} Les relations de Planck-Einstein et de De Broglie lient les caractères ondulatoire et corpusculaire. Ces relations ne dépendent pas de la particule, mais la relation de dispersion liant $\overrightarrow{k}$ à $\omega$ en dépend.
\paragraph{Résumé} De façon générale, on a donc :
\begin{itemize}
	\item pour tout quanton : $E=hf$ et $\lambda=\displaystyle\frac{h}{p}$ ;
	\item pour une particule non relativiste : $p=mv$ et $E_c=\displaystyle\frac12 mv^2$ ;
	\item pour un photon dans le vide : $\displaystyle\lambda=\frac{c}{f}$.
\end{itemize}
\paragraph{Rappel} On a $h=6,62.10^{-34}$ J.s ; $1$ eV $=1,6.10^{-19}$ J et $m_e=9,1.10^{-31}$ kg.
\paragraph{Ordres de grandeur} On a les ordres de grandeur suivants :
\begin{itemize}
	\item Pour un photon visible ($\lambda=600$ nm) : $E=2,0$ eV
	\item Pour un électron avec $E_c=50$ eV, on obtient $E_c=\displaystyle\frac{p^2}{2m}$ Par conséquent : $$\lambda=\frac{h}{\sqrt{2mE_c}}$$ On trouve numériquement $\lambda\approx 17$ nm (de l'ordre d'un paramètre de maille en cristallographie).
\end{itemize}
\section{Fonction d'onde ; équation de Schrödinger}
\subsection{Interprétation probabiliste de l'expérience des fentes d'Young}
\paragraph{Interprétation probabilitse} Pour un seul photon, on a une position complètement aléatoire. Cependant, si on observe beaucoup de photons, on retrouve la figure d'interférences. On a donc perdu une partie du déterminisme classique.
\paragraph{Bonus} \textit{cf} la conférence d'Alain Aspect de l'an dernier. On a notamment :
\begin{itemize}
	\item les inégalités de Bell ;
	\item l'expérience d'Aspect qui met en évidence une violation des inégalités de Bell.
\end{itemize}
\subsection{Postulats ; fonction d'onde et équation de Schrödinger}
\paragraph{Fonction d'onde} La fonction d'onde $\Psi$ caractérise toute l'information que l'on peut avoir sur le système. C'est une fonction de $\R^4$ dans $\C$ qui prend en argument les trois coordonnées de position et le temps.
\paragraph{Postulat de Born}La probabilité qu'une particule se trouve dans un volume $d\tau$ autour de $M$ est donnée par la relation $$d\P=\lvert \Psi(M)\rvert^2d\tau$$ $\Psi$ est donc de carré intégrable.
\paragraph{Vocabulaire : amplitude et densité de probabilité} On dit que $\Psi$ est une amplitude de probabilité et que $\lvert \Psi\rvert^2$ est une densité de probabilité.
\paragraph{Remarque : VDC en roue libre} A propos de Born (qui est une personne différente de Bohr) : "Il a eu le prix Nobel un peu tard parce qu'il a fait des choses que la morale réprouve entre 33 et 45. Heisenberg aussi a eu ce problème."
\paragraph{Dimension de la fonction d'onde} Si on travaille en dimension $n\in\llbracket 1,3\rrbracket$, alors la fonction d'onde a pour dimension $L^{-n/2}$.
\paragraph{Condition de normalisation} Par conséquent, on sait que $$\int_{\R^n} \lvert\Psi(x,y,z,t)\rvert^2 \ dxdydz=1$$ et ce pour tout $t$. Cette condition permet souvent de déterminer des constantes inconnues dans des problèmes de mécanique quantique.
\paragraph{Probabilité de mesurer une particule dans un volume} Par conséquent, la probabilité de mesurer une particule dans un volume $V$ à l'instant $t$ vaut donc : $$\P(e^-\in V)=\int_{V}\lvert \Psi(x,y,z,t)\rvert^2 \ dxdydz$$
\paragraph{Équation de Schrödinger} La fonction d'onde satisfait l'\textbf{équation de Schrödinger} : 
$$i\hbar\frac{\partial\Psi}{\partial t}=-\frac{\hbar^2}{2m}\Delta\Psi+V\left(\overrightarrow{r}\right)\Psi$$ avec l'énergie potentielle, plus couramment appelée potentiel en mécanique quantique, qui est uniquement fonction de l'espace. C'est une équation déterministe et linéaire.
\paragraph{Principe de superposition} Le caractère linéaire de l'équation de Schrödinger fait qu'elle vérifie le principe de superposition : si on possède deux solutions, alors toute combinaison linéaire de ces deux solutions est solution.
\subsection{États stationnaires ; équation de Schrödinger indépendante du temps}
\paragraph{État stationnaire} On dit que $\Psi$ est un état stationnaire lorsque $\Psi$ s'écrit sous la forme :$$\Psi=f(x,y,z)g(t)$$
\paragraph{Attention !} Ne pas confondre avec une onde stationnaire : on n'exige pas ici que $f$ et $g$ soient réelles. Par exemple, $e^{i(\omega t-kx)}$ correspond à un état stationnaire mais n'est âs une onde stationnaire.
\paragraph{Recherche des états stationnaires en 1D} On cherche $\Psi$ sous la forme $f(x)g(t)$. L'équation de Schrödinger s'écrit alors : $$i\hbar \frac{g'(t)}{g(t)}=-\frac{\hbar^2}{2m}\frac{f''(x)}{f(x)}+V(x)=E$$ En effet, le tout est bien une constante puisque le premier membre ne dépend que du temps alors que le deuxième ne dépend que de l'espace. On admet qu'il s'agit de l'énergie potentielle de la particule. Ainsi $f(t)=\exp\left(-i\displaystyle\frac{E}{\hbar}t\right)$ et alors $$\Psi=f(x)e^{\displaystyle-i\frac{E}{\hbar}t}$$ Par conséquent, $\lvert\Psi\rvert^2$ ne dépend pas du temps, d'où le nom d'état stationnaire.
\paragraph{Attention !} On est allé un peu vite sur la fin : $\lvert e^{\displaystyle-i\frac{E}{\hbar}t}\rvert=1$ parce que $E\in\R$, mais on ne peut pas montrer que $E\in\R$ simplement. En fait, en exploitant la condition de normalisation à $t=0$ et pour tout $t$, on trouve directement le fait que $\lvert e^{\displaystyle-i\frac{E}{\hbar}t}\rvert=1$.
\paragraph{Équation de Schrödinger indépendante du temps} On dit alors que $f$ vérifie l'équation de Schrödinger indépendante du temps : $$Ef(x)=\frac{-\hbar^2}{2m}f''(x)+V(x)f(x)$$
\paragraph{Structure des solutions (admis)} Toute solution de l'équation de Schrödinger s'écrit comme superposition d'états stationnaires ou d'états propres (valeurs propres de l'équation).
\subsection{Interprétation en termes d'opérateurs}
\paragraph{Remarque : convention} Attention : en mécanique quantique, on utilise la convention $e^{-i\omega t}$.
\paragraph{Traduction en termes de conservation de l'énergie} L'équation de Schrödinger traduit en fait la conservation de l'énergie mécanique (avec des notations bizarres) : $E=E_c+V$.
\paragraph{Remarque} Il est explicitement écrit dans le programme : "$\displaystyle\frac{-\hbar^2}{2m}\Delta\Psi$ correspond à l'énergie cinétique".
\paragraph{Équation aux valeurs propres ; hamiltonien (HP)} L'équation de Schrödinger indépendante du temps est aussi appelée équation aux valeurs propres. En effet, elle correspond à l'équation aux valeurs propres de l'opérateur $$H:\Psi\mapsto\frac{-\hbar^2}{2m}\Delta\Psi+V\Psi$$ appelé opérateur hamiltonien.
\subsection{Remarque : choix de l'origine de l'énergie potentielle}
\paragraph{Dépendance de l'origine des potentiels} On considère $\Psi$ solution de l'équation de Schrödinger $$i\hbar\frac{\partial\Psi}{\partial t}=-\frac{\hbar^2}{2m}\Delta\Psi+V\Psi$$ On pose alors $\tilde{V}=V+C$. Comment trouver une solution de $$i\hbar\frac{\partial\tilde{\Psi}}{\partial t}=-\frac{\hbar^2}{2m}\Delta\tilde{\Psi}+\tilde{V}\tilde{\Psi}$$ Il suffit de considérer $$\tilde{\Psi}=e^{-i\frac{Ct}{\hbar}}\Psi$$ dont on vérifie qu'elle est effectivement solution de la nouvelle équation. Cela est rassurant car les mesures de position ne peuvent pas dépendre de $C$.
\paragraph{Facteur de phase} Le terme $e^{-i\frac{Ct}{\hbar}}$ est appelé facteur de phase.
\subsection{Remarque : invariance par renversement du temps}
\paragraph{Quasi invariance par renversement du temps} Puisque $V$ est réel, si $\Psi$ est solution de l'équation de Schrödinger $$i\hbar\frac{\partial\Psi}{\partial t}=-\frac{\hbar^2}{2m}\Delta\Psi+V\Psi$$ alors $\tilde{\Psi}:(x,t)\mapsto \Psi(x,-t)\st$ est solution de la même équation. Par conséquent, l'ensemble de la description (Schrödinger et Born) est donc bien invariante par renversement du temps.
\paragraph{Retour sur les équations usuelles}
Voyons quelles équations usuelles nous connaissons, quelle quantité l'on "mesure" dans ces équations, et leur comportement face au renversement du temps.
\begin{itemize}
	\item \textbf{Équation de Schrödinger} pour $V=0$ : $$i\hbar\frac{\partial\Psi}{\partial t}=-\frac{\hbar^2}{2m}\Delta\Psi$$ On "mesure" $\lvert\Psi\rvert^2$ la densité de probabilité. Cette équation est "invariante" par renversement du temps au sens précisé ci-dessus.
	\item \textbf{Équation de la diffusion thermique sans terme source} : $$\rho c\frac{\partial T}{\partial t}=\lambda\Delta T$$ On mesure $T$ et cette équation n'est pas invariante par renversement du temps (irréversibilité).
	\item \textbf{Équation de d'Alembert} : $$c^2\overrightarrow{\Delta}\overrightarrow{E}=\frac{\partial^2 \overrightarrow{E}}{\partial t^2}$$ On mesure $\overrightarrow{E}$ et cette équation est invariante par renversement du temps.
\end{itemize}
\section{Particule libre ; paquet d'onde}
\subsection{État stationnaire d'une particule libre}
\paragraph{Particule libre} On dit qu'une particule est libre lorsqu'elle n'est soumise à aucune force. Le potentiel est alors une constante, choisie arbitrairement nulle.
\paragraph{Équation de Schrödinger pour une particule libre} Pour une particule libre, l'équation de Schrödinger devient $$i\hbar\frac{\partial\Psi}{\partial t}=-\frac{\hbar^2}{2m}\Delta\Psi$$
\paragraph{Relation de dispersion} Pour une particule libre, on cherche une solution de la forme $e^{i(\overrightarrow{k}\cdot\overrightarrow{r}-\omega t)}$. On a alors la relation $$\omega=\frac{\hbar k^2}{2m}$$
\paragraph{Vitesses de phase et de groupe, caractère dispersif} Le relation de dispersion fournit les vitesses de phase et de groupe. Pour la vitesse de phase : $$v_\varphi=\frac{\omega}{k}=\frac{\hbar k}{2m}=\frac{p}{2m}$$ Pour la vitesse de groupe : $$v_g=\frac{d\omega}{dk}(k)=\frac{\hbar k}{m}=\frac{p}{m}$$ C'est donc un milieu "dispersif" puisque $v_\varphi\ne v_g$. Bien avoir en tête que $\displaystyle\frac{p}{m}$ est l'équivalent du $v$ de la mécanique classique.
\paragraph{Caractère non normalisable}
Une onde plane monochromatique n’est pas normalisée : on ne peut trouver de constante telle que $$\iiint \lvert\Psi\rvert^2\ d\tau=1$$
\subsection{Fonction d'onde d'une particule libre ; paquet d'ondes}
\paragraph{Forme du paquet d'ondes} Pour pallier ce problème de non normalisabilité, On superpose plusieurs ondes de vecteur d'onde $\overrightarrow{k}$ voisin de $\overrightarrow{k_0}$ en dimension 1 : $$\Psi(x,t)=\int_\R g(k)e^{i(kx-\omega t)}dk$$ On suppose que $\Delta k$ est faible, ce qui nous permettra de faire un DL ensuite.
\subsection{Propagation du paquet d'ondes ; lien avec la description classique}
\paragraph{Propagation du paquet d'ondes}
Au voisinage de $k_0$, on a : $$\omega(k)=\omega_0+(k-k_0)\underbrace{\frac{\partial\omega}{\partial k}(k_0)}_{=v_g}$$ donc on a : $$\Psi(x,t)=e^{-i(\omega_0 t-k_0 x)}\int_\R g(k)e^{-i(k-k_0)(v_gt-x)} \ dk$$
\paragraph{Lien avec la description classique} Par conséquent, $\lvert\Psi\rvert^2$ se propage à $v_g$. Ainsi, $v_g$ correspond à al vitesse classique.
\paragraph{Étalement du paquet d'ondes} En revanche, le paquet d'ondes s'étale.
\subsection{Inégalités de Heisenberg}
\paragraph{Remarque : autres noms} L'inégalité de Heisenberg est aussi parfois appelée principe d'incertitude, ou encore principe d'indétermination.
\paragraph{Inégalité de Heisenberg} L'inégalité de Heisenberg s'écrit : $$\Delta x\Delta p_x\geqslant \frac{\hbar}{2}$$ avec $\Delta x$ l'écart-type de la mesure de position selon $x$ et $\Delta p_x$ l'écart-type de la mesure de la quantité de mouvement selon $x$. On a de même les inégalités : $$\Delta y\Delta p_y\geqslant \frac{\hbar}{2} \ \ \text{et} \ \ \Delta z\Delta p_z\geqslant \frac{\hbar}{2}$$
\paragraph{Définition et calcul de $\Delta x$} $\Delta x$ est donné par : $$\left(\Delta x\right)^2=\langle x^2\rangle_\Psi-\langle x\rangle_\Psi^2=\int_\R\Psi\st x^2\Psi \ dx-\left(\int_\R \Psi\st x\Psi \ dx\right)^2$$ Remarquer qu'il s'agit là d'une formule de König-Huygens et d'espérances avec la formule de transfert.
\paragraph{Attention !} En revanche, on n'a aucune inégalité sur les quantités de la forme $\Delta x\Delta p_y$ !
\paragraph{Interprétation de l'inégalité de Heisenberg} On considère un dispositif de diffraction dont la largeur du trou de diffraction vaut $a$. L'inégalité de Heisenberg nous dit donc que si $a$ (et donc $\Delta x$) diminue, alors $\Delta p_x$, \textit{ie} la diffraction, augmente.
\paragraph{Interprétation plus précise par la diffraction} Avant, on a une onde plane se propageant selon $y$ donc $\Delta p_x=0$ et $\Delta x=\infty$. Après diffraction, on a une onde circulaire telle que $\Delta x=a$ et $\displaystyle\Delta p_x=p\sin(\theta)=p\frac{\lambda}{a}$ donc $$\Delta x\Delta p_x=p\lambda=h\geqslant\frac{\hbar}{2}$$
\paragraph{Détermination de $\Delta p_x$ (HP)} Si on considère $$\hat{p_x}:\Psi\mapsto\frac{\hbar}{i}\frac{\partial\Psi}{\partial x}$$ alors on a : $$\langle \hat{p_x} \rangle_\Psi=\int_\R \Psi\st(x)\frac{\hbar}{i}\frac{\partial\Psi}{\partial x}(x) \ dx=\left(\Psi \mid \hat{p_x}(\Psi)\right)$$ où $\left(\cdot\mid\cdot\right)$ correspond à un produit scalaire hermitien.
\paragraph{En pratique} Si on dispose de beaucoup de particules dans l'état $\Psi$, on mesure $x$ sur certaines et $p_x$ sur d'autres. On calcule ensuite $\Delta x\Delta px$ et on vérifie (ou on montre) l'inégalité de Heisenberg.
\paragraph{Inégalité temps-énergie} On a : $$\Delta E\Delta t\geqslant\frac{\hbar}{2}$$ avec $\Delta E$ l'écart-type en énergie et $\Delta t$ le temps de cohérence.
\subsection{Vecteur densité de courant de probabilité}
\paragraph{Vecteur densité de courant de probabilité} Le vecteur densité de courant de probabilité est défini par $$\overrightarrow{j}=-\frac{i\hbar}{2m}\left[\Psi\st\grad(\Psi)-\Psi\grad(\Psi\st)\right]$$ C'est l'analogue du vecteur de Poyting $\overrightarrow{\pi}=\displaystyle\frac{\overrightarrow{E}\wedge\overrightarrow{B}}{\mu_0}$. Cette formule générale est \textbf{hors-programme}.
\paragraph{Développement (HP) : explication de la définition de $\overrightarrow{j}$ et preuve de l'équation de continuité} On souhaite obtenir une équation locale de conservation de la probabilité (aussi appelée équation de continuité) analogue à l'équation locale de de conservation de la charge ou à l'identité de Poynting : $$\frac{\partial \rho}{\partial t}+\div\left(\overrightarrow{j}\right)=0$$ où on a posé $\rho=\lvert\Psi\rvert^2=\Psi\Psi\st$ la densité de probabilité. On va calculer la dérivée temporelle de la densité de probabilité et essayer de l'écrire comme la divergence d'un vecteur. En dérivant un produit, on a : $$\frac{\partial \rho}{\partial t}=\frac{\partial \Psi}{\partial t}\Psi\st+\Psi\frac{\partial\Psi\st}{\partial t}$$ On utilise l'équation de Schrödinger $$i\hbar\frac{\partial\Psi}{\partial t}=-\frac{\hbar^2}{2m}\Delta\Psi+V\Psi$$ pour réécrire le premier terme. Pour réécrire le deuxième, on calcule le complexe conjugué de l'équation de Schrödinger : $$-i\hbar\left(\frac{\partial \Psi}{\partial t}\right)\st=-\frac{\hbar^2}{2m}\left(\Delta\Psi\right)\st+V\st\Psi\st$$
Or, on sait que $$\left(\frac{\partial\Psi}{\partial t}\right)\st=\frac{\partial\Psi\st}{\partial t} \ \ \text{et} \ \ \left(\Delta\Psi\right)\st=\Delta\Psi\st$$ De plus, le potentiel est à valeurs réelles, donc on sait que $V\st=V$. Par conséquent, on obtient : $$-i\hbar\frac{\partial\Psi\st}{\partial t}=-\frac{\hbar^2}{2m}\Delta\Psi\st+V\Psi\st$$ En reportant l'équation de Schrödinger et son conjugué dans la dérivation temporelle de $\rho$, on obtient : $$\frac{\partial\rho}{\partial t}=\left(\frac{i\hbar}{2m}\Delta\Psi+\frac{1}{i\hbar}V\Psi\right)\Psi\st+\Psi\left(\frac{-i\hbar}{2m}\Delta\Psi\st+\frac{1}{i\hbar}V\Psi\st\right)=\frac{i\hbar}{2m}\left(\Delta\Psi \times\Psi\st-\Psi\times\Delta\Psi\st\right)$$
Or, on sait, d'après les relations d'analyse vectorielle, que $$\begin{cases}
	\div\left(\Psi\grad(\Psi\st)\right)=\Psi\div\left(\grad(\Psi\st)\right)+\grad(\Psi)\cdot\grad(\Psi\st)=\Psi\Delta\Psi\st+\grad(\Psi)\cdot\grad(\Psi\st) \\
	\div\left(\Psi\st\grad(\Psi)\right)=\Psi\st\div\left(\grad(\Psi)\right)+\grad(\Psi\st)\cdot\grad(\Psi) =\Psi\st\Delta\Psi+\grad(\Psi\st)\cdot\grad(\Psi)
\end{cases}$$ Par conséquent, on a la relation : $$\Delta\Psi \times\Psi\st-\Psi\times\Delta\Psi\st =\div\left(\Psi\st\grad(\Psi)\right)-\div\left(\Psi\grad(\Psi\st)\right)$$ D'où, en reprenant les équations de départ, l'équation de continuité : $$\frac{\partial\rho}{\partial t}+\div\left[-\frac{i\hbar}{2m}\left(\div\left(\Psi\st\grad(\Psi)\right)-\div\left(\Psi\grad(\Psi\st)\right)\right)\right]=0$$ Cela explique donc la définition de $\overrightarrow{j}$ comme $$\overrightarrow{j}=-\frac{i\hbar}{2m}\left[\Psi\st\grad(\Psi)-\Psi\grad(\Psi\st)\right]$$
\paragraph{Identité de conservation (ou équation de continuité)} De façon analogue à la conservation de la charge, on a une conservation de la probabilité : $$\frac{\partial \lvert\Psi\rvert^2}{\partial t}+\div\left(\overrightarrow{j}\right)=0$$
\paragraph{Cas particulier utile et au programme} Pour une onde de la forme $$\Psi(x,t)=A\exp(i(k_x x-\omega t)) \ \ \text{où} \ \ k_x\in\R$$ l'expression générale de $\overrightarrow{j}$ (ou bien l'équation de continuité) fournit $$j_x=-\frac{i\hbar}{2m}\left(ik_xAA\st-(-ik_x)AA\st\right)=\lvert A\rvert^2\frac{\hbar k_x}{m}=\rho\frac{p_x}{m}$$
\paragraph{Cas particulier utile et presque au programme}
Pour une onde de la forme $$\Psi(x,t)=\left[A\exp(ik_x x)+B\exp(-ik_x x)\right]\exp(-i\omega t) \ \ \text{où} \ \ k_x\in\R$$ l'expression générale de $\overrightarrow{j}$ fournit : $$j_x=\frac{\hbar k_x}{m}\left(AA\st-BB\st\right)=\frac{\hbar k_x}{m}(\lvert A\rvert^2-\lvert B\rvert^2)$$
\paragraph{Cas particulier utile mais hors programme}
Pour une onde de la forme $$\Psi(x,t)=A\exp(-i\omega t)\exp(\alpha x)\ \ \text{où} \ \ \alpha\in\R$$ l'expression générale de $\overrightarrow{j}$ fournit l'expression : $$j_x=-\frac{i\hbar}{2m}\left(\alpha AA\st-A\alpha A\st\right)=0$$
\subsection{Interprétation hydrodynamique}
\paragraph{Rappel d'électromagnétisme} En électromagnétisme, on a : $$\overrightarrow{j}=\sum_i n_i q_i \overrightarrow{v_i}$$ S'il n'y a qu'un seul type de porteur de charges, cela devient : $$\overrightarrow{j}=\underbrace{nq}_{=\rho}\overrightarrow{v}$$ (attention, cela est faux pour un câble coaxial).
\paragraph{Interprétation hydrodynamique} Il faut donc bien comprendre qu'ici, le vecteur densité de probabilité est véritablement l'analogue d'un vecteur densité de charges volumique. En effet, pour une onde plane (cas particulier utile et au programme de la section précédente), on a : $$j_x=\rho\frac{p_x}{m}$$ où $\displaystyle\frac{p_x}{m}$ est l'analogue de la vitesse classique.

\chapter{Particule quantique dans un potentiel}
\localtableofcontents
\section{Marche de potentiel}
\subsection{Exemples}
\paragraph{Interface métal/vide} Pour une interface métal/vide, on peut avoir comme marche de potentiel le travail d'extraction nécessaire pour arracher un électron au métal.
\paragraph{Hétérostructures semi-conductrices} C'est un mot stylé pour parler des semi-conducteurs. On a par exemple l'association AlAs/GaAs.
\subsection{Cas $E>V_0$ ; probabilité de réflexion et de transmission}
\paragraph{Classiquement} La particule passe toujours et sa vitesse diminue. \paragraph{Quantiquement} On a des coefficient de transmission et réflexion non nuls, qu'on cherche à calculer. L'équation de Schrödinger indépendante du temps nous donne rapidement que, pour $x<0$, on a : $$\varphi(x)=A_1e^{ik_1x}+B_1e^{-ik_1x}\ \ \text{avec} \ \ k_1=\sqrt{\frac{2mE}{\hbar^2}}$$ De même, elle donne pour $x>0$ : $$\varphi(x)=A_2e^{ik_2x}+B_2e^{-ik_2x}\ \ \text{avec} \ \ k_1=\sqrt{\frac{2m(E-V_0)}{\hbar^2}}$$ On a en réalité $B_2=0$ car il n'y a pas de source de l'autre côté de la marche de potentiel.
\paragraph{Identification des ondes} L'onde en $A_1$ se propage dans le sens des $x$ croissants, tout comme celle en $A_2$. En effet, avec la convention de la mécanique quantique : $$A_1e^{ik_1x}e^{-i\omega t}=A_1e^{-i(\omega t-k_1x)}$$ se propage dans le sens des $x$ croissants. L'onde en $B_1$ est l'onde réfléchie.
\paragraph{Coefficients de réflexion et de transmission}Le coefficient de réflexion vaut alors $$\underline{r}=\frac{B_1e^{-ik_1\times 0}}{A_1e^{ik_1\times 0}}=\frac{B_1}{A_1}$$ et le coefficient de transmission vaut : $$\underline{t}=\frac{A_2e^{ik_2\times 0}}{A_1e^{ik_1\times 0}}=\frac{A_2}{A_1}$$
\paragraph{Relations de passage} On admet les relations de passage suivantes : $$\begin{cases}
	\varphi(0^-)=\varphi(0^+)\\
	\varphi'(0^-)=\varphi'(0^-)
\end{cases}$$ La première est toujours valable, et la deuxième est valable à condition que $V_0$ soit borné au voisinage de $0$ (pour plus de précisions, voir l'avant dernière section de cette partie).
\paragraph{Calcul de $\underline{r}$ et $\underline{t}$} Les relations de passage fournissent donc $$\begin{cases}
	A_1+B_1=A_2 \\
	ik_1A_1-ik_1B_1=ik_2A_2
\end{cases}$$ On en déduit, en simplifiant les $i$ et en divisant par $A_1$ : $$\begin{cases}
1+\underline{r}=\underline{t} \\
1-\underline{r}=\displaystyle\frac{k_2}{k_1}\underline{t}
\end{cases}$$ En résolvant ce système, on trouve finalement les expressions : $$\underline{t}=\frac{2k_1}{k_1+k_2} \ \ \text{et} \ \ \underline{r}=\frac{k_1-k_2}{k_1+k_2}$$
\paragraph{Remarque : cohérence des expressions précédentes} Les expressions précédentes sont homogènes car sans dimension. Si $V_0$ (pas de marche), on a bien $\underline{t}=1$ et $\underline{r}=0$ (revenir aux expressions de $k_1$ et $k_2$)Toutefois l’onde n’est pas normalisable, elle ne décrit pas une particule, mais un flux de particules.
\paragraph{Calcul des probabilités de réflexion et de transmission $R$ et $T$}
Ces probabilités sont définies par $$R=\frac{j_\text{ref}}{j_\text{inc}}=\frac{\lvert B_1\rvert^2}{\lvert A_1\rvert^2}\frac{\frac{\hbar k_1}{m}}{\frac{\hbar k_1}{m}}=\lvert\underline{r}\rvert^2 \ \ \text{et} \ \ T=\frac{j_\text{trans}}{j_\text{inc}}=\frac{\lvert A_2\rvert^2}{\lvert A_1\rvert^2}\frac{\frac{\hbar k_2}{m}}{\frac{\hbar k_1}{m}}=\frac{k_2}{k_1}\lvert\underline{t}\rvert^2$$
Par conséquent : $$R+T=\frac{1}{(k_1+k_2)^2}\left((k_1-k_2)^2+4k_1k_2\right)=1$$ Ainsi, l'onde est soit réfléchie, soit transmise : il n'y a pas d'absorption.
\paragraph{Graphe} On trace le graphe de $\rho=\lvert\varphi\rvert^2$ en fonction de $x$. Pour $x>0$, $\rho$ est constante. Pour $x<0$, $$\rho(x)=\left\lvert A_1(e^{ik_1x}+\underline{r}e^{-ik_1x})\right\rvert^2=\lvert A_1\rvert^2\left(1+R+2\underline{r}\cos(2k_1x)\right)$$ Attention, la tangente en $0$ est horizontale !
\paragraph{Animations} Ca fait des choses jolies : l'onde s'approche, oscille près de la marche puis une partie est transmise et une autre est réfléchie. La partie réfléchie va plus vite que la partie transmise.
\subsection{Cas $E<V_0$ ; onde évanescente}
\paragraph{Classiquement} La particule rebondit sur la marche de potentiel.
\paragraph{Quantiquement} L'équation de Schrödinger indépendante du temps nous donne rapidement que, pour $x<0$, on a : $$\varphi(x)=A_1e^{ik_1x}+B_1e^{-ik_1x}\ \ \text{avec} \ \ k_1=\sqrt{\frac{2mE}{\hbar^2}}$$ De même, elle donne pour $x>0$ : $$\varphi(x)=A_2e^{k_2x}+B_2e^{k_2x}\ \ \text{avec} \ \ k_1=\sqrt{\frac{2m(V_0-E)}{\hbar^2}}$$ On a en réalité $A_2=0$ car l'onde ne peut diverger au voisinage de l'infini.
\paragraph{Coefficients de réflexion et de transmission}Le coefficient de réflexion vaut alors $$\underline{r}=\frac{B_1e^{-ik_1\times 0}}{A_1e^{ik_1\times 0}}=\frac{B_1}{A_1}$$ et le coefficient de transmission vaut : $$\underline{t}=\frac{B_2e^{-k_2\times 0}}{A_1e^{ik_1\times 0}}=\frac{A_2}{A_1}$$
\paragraph{Calcul de $\underline{r}$ et $\underline{t}$} Les relations de passage fournissent donc $$\begin{cases}
	A_1+B_1=B_2 \\
	ik_1A_1-ik_1B_1=-k_2B_2
\end{cases}$$ On en déduit, en divisant par $A_1$ : $$\begin{cases}
	1+\underline{r}=\underline{t} \\
	i(1-\underline{r})=-\displaystyle\frac{k_2}{k_1}\underline{t}
\end{cases}$$ En résolvant ce système, on trouve finalement les expressions : $$\underline{t}=\frac{2ik_1}{ik_1-k_2} \ \ \text{et} \ \ \underline{r}=\frac{k_1-ik_2}{k_1+ik_2}$$ Remarquons que $\underline{r}$ est donc de module $1$. On l'écrira sous la forme $\underline{r}=e^{i\theta}$ pour un certain $\theta\in\R$.
\paragraph{Calcul de $R$ et $T$} On en déduit que $R=1$ et donc, nécessairement, $T=0$ : l'onde est entièrement réfléchie, avec un certain déphasage.
\paragraph{Graphe} Pour $x>0$, on a une exponentielle décroissante. Pour $x<0$, après calcul de module : $$\rho(x)=\lvert A_1\rvert^2\left(1+R+2\cos(k_2x+\theta)\right)=2\lvert A_1\rvert^2\left(1+\cos(2k_2+\theta)\right)$$ Ainsi, $\rho$ s'annule périodiquement pour $x<0$. Attention : la tangente à l'origine n'est pas nulle à cause de l'exponentielle décroissante !
\paragraph{Animations} Ca fait des choses jolies : l'onde s'approche, oscille près de la marche puis est rejetée.
\subsection{Cas $E\ll V_0$}
\paragraph{Apparition d'un point anguleux} Dans le cas où $E\ll V_0$, on a l'apparition d'un point anguleux en $0$ d'après les expressions précédentes.
\paragraph{Condition de validité de la deuxième relation de passage} Par conséquent, on en déduit que la deuxième relation de passage n'est valable que si $V_0$ est borné au voisinage de $x_0$ (la première est toujours valable).
\subsection{Résumé}
Pour $E<V_0$, $R=1$ et $T=0$, puis pour $E>V_0$, $R$ diminue peu à peu et $T$ augmente peu à peu.
\section{Barrière de potentiel}
\subsection{Cas $E>V_0$ ; résonance de transmission}
\paragraph{Barrière de potentiel : cas $E>V_0$} Si $E>V_0$, on peut avoir de multiples réflexions dans la marche et donc des interférences en réflexion et en transmission.
\begin{center}
	\begin{tabular}{|c|c|c|c|}
		\hline 
		$x$&$x<0$&$0<x<L$&$L<x$ \\
		\hline
		$\varphi$&$A_1e^{ik_1x}+B_1e^{-ik_1x}$&$A_2e^{ik_2x}+B_2e^{-ik_2x}$&$A_3e^{ik_1x}+\xcancel{B_3e^{-ik_2x}}$\\
		\hline
	\end{tabular}
\end{center}
$k_1$ et $k_2$ dont donnés par l'équation de Schrödinger indépendante du temps : $$k_1=\sqrt{\frac{2mE}{\hbar^2}} \ \ \text{et} \ \ k_2=\sqrt{\frac{2m(E-V_0)}{\hbar^2}}$$
Les relations de passage donnent 4 relations entre les 5 inconnues, dont on élimine la première puisque les coefficients de réflexion et de transmission correspondent à un quotient entre ces inconnues. On admet que la résolution de ce système $4\times 4$ fournit : $$T=\frac{1}{1+\displaystyle\frac{V_0^2}{4E(E-V_0)}\sin^2(k_2L)}$$
\paragraph{Remarque : cohérence de l'expression de $T$} On remarque que l'expression de $T$ admise :
\begin{itemize}
	\item est homogène ;
	\item vérifie $T=1$ si $V_0=0$ ou si $L=0$, ce qui correspond bien à une absence de barrière de potentiel.
\end{itemize}
\paragraph{Transparences} On appelle \textbf{transparences} les points où $T$ vaut $1$. 
\paragraph{Version "réflexions multiples"} On suppose désormais qu'il y a de multiples réflexions. Si les rayons deux rayons consécutifs $a$ et $b$ émergents interfèrent constructivement, on sait que toutes les ondelettes transmises vont interférer constructivement. Or, la différence de phase entre $a$ et $b$ correspond à deux épaisseurs et deux réflexion, donc à une différence $$\underline{r}'^2e^{i2k_2L}$$ avec $\underline{r}'^2\geqslant 0$ car on sait que $\underline{r}\in\R$. Par conséquent, $a$ et $b$ sont en phase lorsque $$k_2L=2\pi n\ \ \text{pour} \ \ n\in\N\st$$
\paragraph{Remarque : réflexions multiples en optique} En optique, le système classique qu'on étudie dans ce cas s'appelle l'\textbf{interféromètre de Fabry-Pérot}. Retenir donc cette méthode de calculer le déphasage entre deux rayons consécutifs pour ce genre d'exercices.
\subsection{Cas $E<V_0$ ; effet tunnel}
\paragraph{Barrière de potentiel : cas $E<V_0$}Dans le cas où $E<V_0$, on a une onde évanescente $A_2e^{\rho'x}+B_2e^{-\rho'x}$ dans la barrière (exponentielle \textbf{réelle}) et on a alors : $$T=\frac{1}{1+\displaystyle\frac{V_0^2}{4E(V_0-E)}\sinh^2(\rho'L)}$$
\paragraph{Effet tunnel} C’est ce phénomène de transmission de la particule à travers une barrière de potentielle classiquement infranchissable. Application du microscope à effet tunnel, fait passer les électrons d’un conducteur à la point d’un microscope, comme l’effet est très sensible à l’épaisseur de la barrière, il permet de mesurer précisément de surfaces. Une autre application est la radioactivité $\alpha$ : émission de $_2^4\text{He}^{2+}$, modèle de Gamov (interaction forte (potentiel négatif) jusqu’à $r_0$ puis interaction électromagnétique (potentiel en $1/r$ )) explique la grande différence de durée de vie des particules par les variations de barrière de potentiel en fonction de l’énergie.
\subsection{Approximation de la barrière épaisse et approche documentaire}
\paragraph{Approximation de la barrière épaisse}
Dans la limite de la barrière épaisse $\displaystyle\left(a\gg \frac{1}{\rho'}\right)$, on a : $$T\approx \exp\left(-2\rho'a\frac{16E(V_0-E)}{V_0^2}\right)$$
\section{Puits de potentiel ; quantification de l'énergie}
\subsection{Puits infini 1D}
\paragraph{Puits de potentiel infini 1D} On utilise l'équation de Schrödinger et les conditions $\varphi(0)=\varphi(a)=0$ pour établir que $\varphi$ est sinusoïdale et que $k=\frac{n\pi}{a}$ est quantifié ($n\in\N\st$). Puis avec $$E=\frac{\hbar k}{2m}$$ on retrouve que l’énergie est quantifiée avec seuil plus grand que le minimum de potentiel (ce qui s’explique par l’inégalité d’Heisenberg).
\subsection{Ordre de grandeur de l'énergie de confinement ; inégalité de Heisenberg}
\subsection{Puits infini 3D}

\paragraph{Puits de potentiel infini 3D} En 3D, on cherche une solution en séparant les variables $\varphi=f(x)g(y)h(z)$. On trouve une superposition d’énergie quantifiée et des états dégénérés (plusieurs combinaisons donnent la même énergie).
\subsection{Puits fini 1D}

\paragraph{Puits de potentiel fini 1D} Principe de symétrie : si $\varphi$ est solution non-nulle de l’équation de Schrödinger alors sa partie réelle ou imaginaire l’est également, et sa partie réelle paire ou impaire aussi (l’une au moins est non-nulle puisque $\lvert \varphi^2\rvert$ est paire). On peut donc se restreindre à ce type de solutions. Avec $E<V_0$, seul un nombre fini d’états sont possibles. Augmenter la largeur du puits ou diminuer sa hauteur revient à moins confiner la particule, donc rendre plus de niveaux d’énergie possibles et diminuer l’énergie du fondamental.

\paragraph{Méthode} Pour un puits bizarre (fini ou autre), on peut écrire les relations de passage matriciellement puis dire que notre matrice n'est pas inversible car un vecteur non nul est dans son noyau. Calculer le déterminant de cette matrice permet alors d'obtenir une relation entre les différents paramètres du problème. Souvent, l'équation qui en résulte n'est pas résoluble explicitement, mais on peut en faire une résolution graphique par exemple, à l'aide des autres contraintes du problème.
\section{États non stationnaires}
\subsection{Puits infini}
\paragraph{États non stationnaires} Avec les $\varphi_i$ les fonctions propres du système, la fonction : $$\Psi(x,t=0)=\frac{1}{\sqrt{2}}\left(\varphi_0(x)+\varphi_1(x)\right)$$ est solution donc $$\Psi(x,t)=\frac{1}{\sqrt{2}}\left(\varphi_0(x)\exp\left(i\frac{E_0}{\hbar}t\right)+\varphi_1(x)\exp\left(i\frac{E_1}{\hbar}t\right)\right)$$ Pour $\lvert \Psi^2 \rvert$, on obtient une formule d'interférence semblable à celle de Fresnel. On a donc une oscillation périodique à la fréquence de Bohr : $$\nu=\frac{E_1-E_0}{2\pi\hbar}$$
\subsection{Double puits ; exemple de la molécule d'ammoniac}
\paragraph{Double puits} Si les deux puits sont collés le système est équivalent à un puits de largeur double ; s’ils sont très écartés les puits sont découplés. Or le fondamental a une énergie plus faible dans un grand puits que dans deux puits de largeur plus faible, ce qui explique la stabilité des molécules ($\text{H}_2^+$, l’énergie potentielle de proton augmente inversement à $r$, mais celle de l’électron diminue).
\subsection{Liaison chimique (culture)}
\part{Électronique numérique}
\chapter{Électronique numérique}
\localtableofcontents
\section{Échantillonage}

\subsection{Signal échantillonné}

\paragraph{Échantillonnage} Il s'agit d'une discrétisation en temps. On note $T_e$ la période d'échantillonnage. On relève une valeur tous les $T_e$.

\paragraph{Compromis} Il faut un compromis :
\begin{itemize}
	\item si $T_e$ est trop petit, il sera compliquer de stocker et transmettre le signal ;
	\item si $T_e$ est trop grand, il y aura de la perte d'information.
\end{itemize}

\paragraph{Problème majeur} Si on considère le signal sinusoïdal : $$S(t)=A\cos(2\pi ft+\varphi)$$ on peut remarquer que les signaux $$S_n(t)=A\cos(2\pi(f+nf_e)t+\varphi)$$ pour $n\in\Z$ passent par les même points donc après échantillonnage, on ne sait pas si la fréquence de départ était $f$, $f+f_e$, $f-f_e$, etc.

\subsection{Une analogie}

\paragraph{Analogie édifiante} On observe une roue qui tourne avec une période $T$. Si on prend des photos avec une période $T/4$, on obtient le même résultat que si on prenait des photos avec une période $5T/4$, donc on n'a pas accès à la période d'origine.

\subsection{Spectre du signal échantillonné}

\paragraph{Avec simplement un fondamental} On voit apparaître des pics à $f+nf_e$.

\paragraph{Avec des harmoniques ; repliement} Pour chaque harmonique, des pics $f_i+nf_e$ apparaissent. Si on diminue trop la période d'échantillonnage, le spectre se replie sur lui-même et on ne peut plus distinguer les différentes harmoniques. On dit qu'il y a repliement du spectre.

\paragraph{Remarque} En TP, il faut faire l'acquisition d'un nombre entier de périodes du signal.

\subsection{Condition et théorème de Nyquist-Shannon}

\paragraph{Critère de Nyquist-Shanon} Il faut avoir $f_e>2f_\text{max}$ pour pouvoir distinguer les fréquences. On ne fait pas la différence entre un signal à $f_e$ et $f_e+f$ ou $f_e-f$ dans l'échantillon, donc on utilise :
\begin{itemize}
	\item un filtre passe-bas anti-repliement avant l'échantillonneur (pour éviter que les fréquences $f_e-f$ hors du spectre se retrouvent à $f$)
	\item un filtre passe-bas après l'échantillonneur pour enlever les fréquences à $nf_e+f$ et $nf_e-f$
\end{itemize}

\paragraph{Exemple du CD audio} Pour un CD audio, $f_\text{max}=20$ kHz et $f_e=44,1$ kHz.

\paragraph{Remarque} Un système linéaire continu invariant ne fait pas apparaître de nouveaux pics dans le spectre. L'échantillonnage est certes linéaire, mais il n'est pas invariant donc peut faire apparaître de nouveaux pivcs.

\subsection{Filtre anti-repliement}

\paragraph{Filtre anti-repliement} C'est un filre passe-bas avant l'échantillonnage. Il permet de "fabriquer" un $f_{\text{max}}$. On choisit les caractéristiques du filtre de façon à respecter le critère de Shannon tout en essayant de conserver un maximum de fréquences.

\section{Numérisation d'un signal}

\subsection{Quantification}

\paragraph{Quantification} Il s'agit d'une discrétisation en espace. On parle de loi linéaire lorsque le pas de quantification $q$ est constant. Par exemple, $[0,q[$ est codé par $0$, $[q,2q[$ est codé par $1$, etc.

\paragraph{Pas de la quantification} $q$ est appelé le pas de la quantification.

\paragraph{Dynamique du capteur} $$\Delta s = s_\text{max} - s_\text{min}$$ est la dynamique du capteur.

\paragraph{Lien entre le nombre de bits et la dynamique du capteur} Si on stocke sur $N$ bits on a $$q= \frac{\Delta s}{2^N-1}$$

\subsection{Signal et bruit}

\paragraph{Choix du pas de la quantification pour limiter le bruit} On choisit $q$ strictement supérieur à l'amplitude du bruit et $e$ le calibre de manière à éviter la saturation.

\section{Filtrage numérique}

\paragraph{Intérêt} Le filtrage numérique est beaucoup plus simple qu'un filtrage analogique (avec $R$ et $C$).

\subsection{Principe}

\paragraph{Principe d'un filtrage numérique} On part d'un signal que l'on échantillonne et que l'on quantifie pour obtenir des signaux. On a donc effectué une première conversion analogique-numérique. Ensuite, on traite numériquement ces échantillons. Enfin, on réalise une conversion numérique-analogique pour obtenir le signal traité.

\paragraph{Exemple : passe bande du premier ordre} On prend un filtre : $$\underline{H}=\frac{\underline{s}}{\underline{e}}=\frac{H_0}{1+j\frac{\omega}{\omega_c}}$$ L'équation différentielle associée est : $$\frac{1}{\omega_c}\frac{ds}{dt}+s=e$$ On souhaite arriver à $$s_n\approx s(nT_e)$$ Or, la dérivée de $s$ par rapport au temps correspond de façon discrète à $$\frac{s_{n+1}-s_n}{T_n}$$ On obtient donc la relation de récurrence : $$\frac{1}{\omega_c T_e}s_{n+1}+\left(1-\frac{1}{\omega_c T_e}\right)s_n=e_n$$ ce qui peut encore se réécrire $$s_{n+1}=s_n(1-\omega_c T_e)+e_n\omega c T_e$$ sachant que $e_n$ est donné par le convertisseur analogique-numérique. On peut alors calculer dans les cas pratiques tous les $s_n$, puis utiliser le convertisseur numérique-analogique.

\subsection{Limites}

\paragraph{Limites} La méthode exposée précédemment est une bonne approximation pour $f_e\gg f_c$, ce qui correspond à $\omega_c T_e\ll 1$. Sinon, le taux de variation n'est pas une bonne approximation de la dérivée.

\paragraph{Amélioration : transformée de Fourier rapide} La bonne méthode est en réalité la transformée de Fourier rapide (FFT) : on applique une transformation de Fourier, on multiplie par la fonction de transfert, puis on applique la transformée de Fourier inverse.


\part{Physique statistique}
\chapter{Physique statistique}
\localtableofcontents
\paragraph{États} L'état d'un système macroscopique peut être :
\begin{itemize}
	\item son micro-état : classiquement 3 paramètres de position et de vitesse pour chaque particule, donc 6N paramètres ;
	\item son macro-état : $P$, $T$, $V$, $m$, $U$, $H$, $S$, $C_v$, etc.
\end{itemize}
\paragraph{Approche} Dans ce chapitre, on emploie une approche statistique et probabiliste.
\section{Facteur de Boltzmann}

\subsection{Pression cinétique d'un gaz parfait}

\paragraph{Calcul de la force de pression} On considère une surface $S$ de vecteur normal $\overrightarrow{u_z}$. Les atomes ayant une vitesse comprise entre $\overrightarrow{v}$ et $\overrightarrow{v}+d\overrightarrow{v}$ qui tapent cette surface entre $t$ et $t+dt$ sont au nombre de : $$n\st f(\overrightarrow{v})dv_xdv_ydv_z\times S\overrightarrow{v}\cdot\overrightarrow{u_z} dt$$Pour un tel atome, le choc est élastique (loi de Descartes) et la variation de quantité de mouvement associée est : $$-2mv_z\overrightarrow{u_z}$$ Le théorème de la résultante au gaz entre $t$ et $t+dt$ s'écrit alors : $$\delta\overrightarrow{F}_{\text{paroi}\to\text{gaz}}dt=-2mv_zn\st f(\overrightarrow{v})dv_xdv_ydv_z\times S\underbrace{\overrightarrow{v}\cdot\overrightarrow{u_z}}_{=v_z} \overrightarrow{u_z}$$ Par conséquent, on a alors : $$\overrightarrow{F}_{\text{paroi}\to\text{gaz}}=-2mn\st S\int_{-\infty}^{+\infty}dv_x\int_{-\infty}^{+\infty}dv_y\int_{-\infty}^{+\infty}v_z^2f(\overrightarrow{v})dv_z\overrightarrow{u_z}$$ Or, on a d'autre part : $$\int_{-\infty}^{+\infty}dv_x\int_{-\infty}^{+\infty}dv_y\int_{-\infty}^{+\infty}v_z^2f(\overrightarrow{v})dv_z=\langle v_z^2\rangle$$ donc on en déduit $$\overrightarrow{F}_{\text{paroi}\to\text{gaz}}=-mn\st\langle v_z^2\rangle S\overrightarrow{u_z}$$

\paragraph{Pression cinétique} On dispose donc bien d'une force de pression associée à une pression $$P=+m\langle v_z^2\rangle n\st$$ Or, par isotropie de la distribution des vitesses, on sait que $$\langle v_x^2\rangle=\langle v_y^2\rangle=\langle v_z^2\rangle=\frac{u^2}{3}$$ Par conséquent, on dispose de l'expression suivante de la pression cinétique : $$P=\frac{1}{3}mn\st u^2$$ avec : \begin{itemize}
	\item $m$ la masse d'une particule ;
	\item $n\st$ la densité particulaire (nombre de particules par unité de volume) en $m^{-3}$ ;
	\item $u=\sqrt{\langle \overrightarrow{v}^2\rangle}$ la vitesse quadratique moyenne (en $m.s\inv$).
\end{itemize}

\paragraph{Remarque} On a négligé les interactions à distance.

\subsection{Définition de la température cinétique}

\paragraph{Température cinétique} Si on part de la loi des gaz parfaits $PV=nRT$, on a $PV=\displaystyle\frac{n\st V}{N_A}RT$ donc avec $k_B=\displaystyle\frac{R}{N_A}$ et l'expression de la pression cinétique, on obtient l'expression de la température cinétique avec : $$\frac{1}{2}mu^2=\frac{3}{2}k_B T$$

\paragraph{Attention} $n$ est en mol et $n\st$ est en $m^{-3}$.

\paragraph{Remarque} On a ici affaire à un cas particulier du théorème d'équipartition (\textit{cf} la suite du chapitre).

\subsection{Modèle de l'atmosphère isotherme}

\paragraph{Modèle de l'atmosphère comme un gaz parfait isotherme} Cela permet d'obtenir une équation différentielle sur la pression $P$ et qui permet d'aboutir à l'expression : $$P(z)=P(0)\exp\left(-\frac{mgz}{k_BT}\right)$$ Cela se généralise à $n\st$.
\subsection{Facteur de Boltzmann et interprétation probabiliste}

\paragraph{Apparition d'un facteur de Boltzmann}
On a donc : $$n\st=n_0\st\exp\left(-\frac{Mgz}{RT}\right)=n_0\st\exp\left(-\frac{mgz}{k_BT}\right)$$ où $m$ désigne la masse d'une particule de gaz.

\paragraph{Remarque} $k_B T$ est donc une énergie.

\paragraph{Interprétation probabiliste} Si on choisit une particule au hasard, la probabilité qu'elle soit entre $z$ et $z+dz$ est alors : $$d\P=\frac{1}{Z}\exp\left(-\frac{mgz}{k_BT}\right)$$ où $Z$ est la fonction de partition qui dépend de la condition de normalisation.

\section{Système à spectre discret en équilibre avec un thermostat}

\subsection{Loi de Boltzmann ; fonction de partition}

\paragraph{Contexte} On considère un système fermé en contact avec un thermostat à la température $T$.

\paragraph{Probabilité d'être dans un état d'énergie} La probabilité que le système soit dans le microétat $\alpha$ d'énergie $E_\alpha$ est : $$\P_\alpha=\frac{1}{Z}\exp\left(-\frac{E_\alpha}{k_B T}\right)$$ où $Z$ est la fonction de partition qui ne dépend pas de $\alpha$ mais peut dépendre de $T$.

\paragraph{Expression de la fonction de partition} La condition de normalisation impose que la somme des probabilités fasse $1$. Par conséquent, en sommant l'expression précédente, on en déduit : $$Z=\sum_\alpha\exp\left(-\frac{E_\alpha}{k_B T}\right)$$

\subsection{Particules identiques et sans interaction}

\paragraph{Exemple de particules indépendantes} Deux masses attachées aux bouts de deux ressorts eux-mêmes reliés à des murs différents. L'énergie totale du système vaut : $$E=\frac{1}{2}m_1v_1^2+\frac{1}{2}k_1(x_1-l_1)^2+\frac{1}{2}m_2v_2^2+\frac{1}{2}k_2(x_2-l_2)^2$$

\paragraph{Exemple de particules non indépendantes} Deux masses reliées par un ressort entre elles deux. L'énergie ne peut pas s'écrire comme une somme $f(x_1,v_1)+g(x_2,v_2)$.

\paragraph{Conséquence de l'indépendance} Si les particules sont indépendantes, on a : $$\P(1\in \alpha,\ 2\in\beta)=\P(1\in\alpha)\P(2\in\beta)$$

\paragraph{Système de $N$ particules identiques et indépendantes} Pour $N$ particules identiques et indépendantes, on a : $$Z=\sum_{\alpha_1}\cdots\sum_{\alpha_N}\exp\left(-\frac{E_{1\in\alpha_1}+\cdots+E_{N\in\alpha_N}}{k_BT}\right)$$ soit encore $$Z=\left(\sum_{\alpha_1}\exp\left(-\frac{E_{1\in\alpha_1}}{k_BT}\right)\right)\times\cdots\times\left(\sum_{\alpha_N}\exp\left(-\frac{E_{N\in\alpha_N}}{k_BT}\right)\right)$$ On en déduit que dans ce cas : $$Z=Z_{\text{particule}}^N$$

\paragraph{Loi de Boltzmann} Pour un système constitué de $N$ particules identiques et indépendantes en contact avec un thermostat à la température $T$, la probabilité qu'une particule soit dans l'état d'énergie $E_\alpha$ est : $$\P_\alpha=\frac{1}{Z_{\text{particule}}}\exp\left(-\frac{E_\alpha}{k_BT}\right)$$ On appelle facteur de Boltzmann la quantité $\displaystyle\exp\left(-\frac{E_\alpha}{k_BT}\right)$.

\subsection{Énergie interne ; capacité thermique}

\paragraph{Énergie interne} On appelle énergie interne du système la quantité : $$U=\langle E_{\text{système}}\rangle$$

\paragraph{Cas du système de $N$ particules identiques et indépendantes} Dans le cas de $N$ particules identiques et indépendantes, on : $$U=\sum_{\text{particules}}\langle E_{\text{particule}}\rangle=N\langle E_{\text{particule}}\rangle$$ où on a $$\langle E_{\text{particule}}\rangle=\sum_{\alpha}E_\alpha\P_\alpha$$

\paragraph{Expression de $\langle E_{\text{particule}}\rangle$ en fonction de $\beta$} En posant $\beta=\displaystyle\frac{1}{k_BT}$, les expressions de $Z_{\text{particule}}$ et $\langle E_{\text{particule}}\rangle$ donnent : $$\langle E_{\text{particule}}\rangle=-\frac{d}{d\beta}(\ln(Z_{\text{particule}}))$$

\paragraph{Capacité thermique à volume constant} On définit la capacité thermique à volume constant par : $$C_V=\frac{\partial \langle E_{\text{système}}\rangle}{\partial T}\bigg\rvert_V$$ Ainsi, lorsque $V$ est constant (c'est-à-dire lorsqu'on ne change pas les valeurs de $E_\alpha$), on a : $$C_V=\frac{d \langle E_{\text{système}}\rangle}{d T}$$

\subsection{Fluctuations relatives de l'énergie}

\paragraph{Fluctuation pour le système de $N$ particules identiques et indépendantes} Pour un système constitué de $N$ particules identiques et indépendantes, la fluctuation relative de l'énergie vaut : \begin{align*}
	\frac{\sigma(E_\text{système})}{\langle E_\text{système}\rangle}&=\frac{\sqrt{\mathbb{V}(E_\text{système})}}{\langle E_\text{système}\rangle} \\
	&=\frac{1}{\sqrt{N}}\times\frac{\sqrt{\langle E_\text{particule}^2\rangle -\langle E_\text{particule}\rangle^2}}{\langle E_\text{particule}\rangle}
\end{align*}donc les fluctuations relatives de l'énergie sont inversement proportionnelles à $\sqrt{N}$.

\paragraph{Généralisation} De façon générale, les fluctuations de l'énergie d'un système macroscopique sont presque invisibles.

\subsection{Remarque : choix de l'origine de l'énergie potentielle}

On considère $C$ une constante indépendante de $T$ et de $\alpha$ et on change l'origine des énergies potentielles : $$E_\alpha'=E_\alpha+C$$

\paragraph{Influence sur la probabilité} En factorisant, on trouve rapidement : $$Z'=\exp\left(-\frac{C}{k_BT}\right)Z$$ puis on en déduit : $$\P_\alpha'=\P_\alpha$$ Ouf ! La probabilité n'est pas modifiée.

\paragraph{Influence sur la capacité thermique} Dans ce cas, on a : $$\langle E_\text{système}'\rangle=C+\langle E_\text{système}\rangle$$ donc on a bien $C_V'=C_V$, ce qui est tout aussi rassurant !

\section{Système à deux niveaux d'énergie non dégénérés en équilibre avec un thermostat}

\subsection{Exemples}

\paragraph{Moment magnétique} Un moment magnétique (un électron par exemple) dans un champ magnétique constitue un tel système.

\paragraph{Rappels sur le magnétisme} On dit qu'un matériau est : 
\begin{itemize}
	\item paramagnétique lorsque le moment magnétique $\overrightarrow{M}$ s'oriente selon $\overrightarrow{B}_{\text{ext}}$. On définit alors l'aimantation $\displaystyle\overrightarrow{\mathcal{M}}=\frac{\delta\overrightarrow{M}}{\delta V}$ en $A/m$.
	\item Ferromagnétique lorsque le matériau conserve son orientation alors qu'on a cessé d'appliquer $\overrightarrow{B}_{\text{ext}}$.
	\item Diamagnétique lorsque le moment magnétique $\overrightarrow{M}$ s'oriente selon $-\overrightarrow{B}_{\text{ext}}$.
\end{itemize}

\paragraph{Niveaux hyperfins du fondamental du césium} Deux niveaux hyperfins du fondamental du césium conviennent (\textit{cf} la définition d'une seconde).

\paragraph{Spin} Pour un spin $1/2$ dans le champ $B_0\overrightarrow{u_z}$ : $$E_\text{particule}=-\overrightarrow{M}\cdot\overrightarrow{B}=-M_zB_0=-(\pm\hbar B_0\gamma)$$ avec $\gamma=\displaystyle\frac{-e}{2m_e}$ le facteur gyromagnétique de l'électron.

\paragraph{Exemple gardé dans la suite} $N$ particules indépendantes ayant pour niveaux d'énergie :
\begin{itemize}
	\item $-\varepsilon$ (niveau 1) ;
	\item $\varepsilon$ (niveau 2)
\end{itemize} avec $\varepsilon>0$.

\subsection{Populations ; énergie interne}

\paragraph{Calcul de la fonction de partition pour une particule} On a : $$Z_\text{particule}=\exp\left(-\frac{-\varepsilon}{k_BT}\right)+\exp\left(-\frac{\varepsilon}{k_BT}\right)=2\cosh\left(\frac{\varepsilon}{k_BT}\right)$$

\paragraph{Probabilités} On a donc : $$\P_1=\frac{\exp\left(\frac{\varepsilon}{k_BT}\right)}{2\cosh\left(\frac{\varepsilon}{k_BT}\right)}\ \ \text{et}\ \ \P_2=\P_1=\frac{\exp\left(\frac{-\varepsilon}{k_BT}\right)}{2\cosh\left(\frac{\varepsilon}{k_BT}\right)}$$

\paragraph{Populations} On a ensuite : $$\langle N_1\rangle=N\P_1\ \ \text{et} \ \ \langle N_2\rangle=N\P_2$$

\paragraph{Énergie moyenne d'une particule} Ainsi : $$\langle E_\text{particule}\rangle=E_1\P_1+E_2\P_2=-\varepsilon\tanh\left(\frac{\varepsilon}{k_BT}\right)$$

\paragraph{Énergie moyenne du système} On en déduit finalement : $$\langle E_\text{système}\rangle =-N\varepsilon\tanh\left(\frac{\varepsilon}{k_BT}\right)$$

\subsection{Capacité thermique ; fluctuations de l'énergie}

\paragraph{Calcul de la capacité thermique} En dérivant par rapport à $T$, on obtient alors : $$C_V=\frac{N\varepsilon^2}{k_BT^2\cosh^2\left(\frac{\varepsilon}{k_BT}\right)}$$

\paragraph{Fluctuations} On a : $$\mathbb{V}(E_\text{particule})=\langle E_\text{particule}^2\rangle-\langle E_\text{particule}\rangle^2=\varepsilon^2-\left(\varepsilon\tanh\left(\frac{\varepsilon}{k_BT}\right)\right)^2=\frac{\varepsilon^2}{\cosh^2\left(\frac{\varepsilon}{k_BT}\right)}$$ donc on observe qu'on dispose de la relation : $$C_V=\frac{\mathbb{V}(E_\text{système})}{k_BT^2}$$ C'est un cas particulier du théorème de fluctuation-dissipation.

\paragraph{Théorème de fluctuation-dissipation}Vérifions que la relation précédente est vraie dans le cas général. On a : \begin{align*}
	C_V&=\frac{d\langle E\rangle}{dT} \\
	&= \frac{1}{Z}\sum_{i=1}^{N}\frac{E_i^2}{k_BT^2}e^\frac{-E_i}{k_BT}+\frac{-1}{Z^2}\left(\sum_{i=1}^{N}\frac{E_i^2}{k_BT^2}e^\frac{-E_i}{k_BT}\right)\left(\sum_{i=1}^{N}E_ie^\frac{-E_i}{k_BT}\right) \\
	&= \frac{1}{k_BT^2}\left[\sum_{i=1}^{N}E_i^2\P_i-\left(\sum_{i=1}^{N}E_i\P_i\right)^2\right] \\
	&=\frac{1}{k_BT^2}\left[\langle E^2\angle-\langle E\rangle^2\right]
\end{align*} On en déduit la relation $$\mathbb{V}(E_\text{système})=C_Vk_BT^2$$ La variance de l'énergie, qui caractérise les fluctuations d'un système à l'équilibre à la température $T$, est donc proportionnelle à la capacité thermique, qui caractérise les variations (\textit{ie} la réponse) de l'énergie moyenne lors d'une modification de la température. On parle de théorème de fluctuation-réponse ou fluctuation-dissipation.

\section{Particules identiques, indépendantes et discernables dans un puits de potentiel}
\subsection{États d'une particule dans un puits infini (rappel)}
\paragraph{Partie spatiale de la fonction d'onde pour l'état $n\in\N\st$} Elle vaut : $$\varphi_n(x)=\sqrt{\frac{2}{L}}\sin\left(\frac{n\pi}{L}x\right)$$

\paragraph{Énergie associée à l'état $n\in\N\st$} Elle vaut : $$E_n=\frac{\hbar^2k_n^2}{2m}=\frac{h^2}{8mL^2}n^2$$

\subsection{Approximation continue}

\paragraph{Approximation continue} Lorsque $k_BT\gg $ l'écart d'énergie entre deux niveaux "en haut", on peut se placer dans l'approximation continue. La description usuelle est donc valable lorsque $T$ est faible.

\paragraph{Fonction de partition dans l'approximation continue pour le puits de potentiel $1$D} Elle vaut alors, avec $k_n=\displaystyle\frac{n\pi}{L}$ : $$Z=\sum_{n\in\N\st}\exp\left(-\frac{\hbar^2k_n^2}{2mk_BT}\right)=\frac{L}{\pi}\int_{0}^{+\infty}\exp\left(-\frac{\hbar^2k^2}{2mk_BT}\right)dk$$
\subsection{Énergie moyenne}

\paragraph{Calcul de l'énergie moyenne dans l'approximation continue} On a donc : $$\langle E_\text{particule}\rangle=\displaystyle\frac{\frac{L}{\pi}\int_{0}^{+\infty}\frac{\hbar^2k^2}{2m}\exp\left(-\frac{\hbar^2k^2}{2mk_BT}\right)dk}{\frac{L}{\pi}\int_{0}^{+\infty}\exp\left(-\frac{\hbar^2k^2}{2mk_BT}\right)dk}$$ En réalisant une IPP, on peut éliminer l'intégrale du bas pour finalement obtenir, dans la limite continue : $$\langle E_\text{particule}\rangle=\frac{1}{2}k_BT$$

\paragraph{Remarque} Pour un gaz parfait monoatomique, on a $C_{V,m}=\displaystyle\frac{R}{2}$.

\paragraph{Remarque} Classiquement, on a plutôt : $$Z=\frac{1}{h}\int\int \exp\left(-\frac{E(x,p)}{k_BT}\right)dx\ dp$$

\section{Capacité thermique des gaz et des solides}
\subsection{Limite classique ; approximation continue ; théorème d'équipartition}

\paragraph{Attention} Le théorème d'équipartition est faux en mécanique quantique.

\paragraph{Limite classique} C'est quand on peut négliger la quantification de l'énergie, c'est-à-dire lorsque $k_BT$ est grand devant les écarts d'énergie. $Z$ s'écrit alors : $$Z=\int e^{-\beta(x_1,\ldots,p_{3N})}dx_1\ldots dx_{3N}dp_1\ldots dp_{3N}$$

\paragraph{Remarque} Il n'existe pas de limite classique intéressante pour le système à deux niveaux.

\paragraph{Théorème d'équipartition} Dans la limite classique, si l'énergie s'écrit $aX^2+b$ où $X$ vaut $x_i$ ou $p_i$, et $a$ et $b$ sont indépendants de $X$, alors : $$\langle aX^2\rangle=\frac{k_BT}{2}$$ Pour le démontrer, intégrer par parties  en faisant sortir tout sauf $X$ de la dernière intégrale.
\subsection{Distribution de Maxwell-Boltzmann des vitesses (complément)}
\paragraph{Loi de Maxwell-Boltzmann} La loi de Maxwell-Boltzmann décrit un gaz parfait avec $N$ molécules sans interactions. Elle donne la probabilité de mesure une particule avec une vitesse entre $v_x$ et $v_x+dv_x$, $v_y$ et $v_y+dv_y$ et $v_z+dv_z$ par : $$dp_{v_x,v_y,v_z}=N\exp\left(-\frac{mv^2}{2k_BT}\right)dv_xdv_ydv_z$$ Par conséquent, la probabilité de mesurer une particule avec une norme entre $v$ et $v+dv$  est $$dp_v=N\exp\left(-\frac{mv^2}{2k_BT}\right)4\pi v^2dv$$ (coordonnées sphériques pour la vitesse). Cela permet d'obtenir un ordre de grandeur pour $u$ : $u\approx 500$ m.s$\inv$. On a donc $v_\text{son}<u<v_\text{libération}$.
\subsection{Capacité thermiques des gaz parfaits monoatomiques}

\paragraph{Capacité thermique des gaz parfaits monoatomiques} Pour un gaz parfait monoatomique, on a 3 degrés de liberté, donc $$E=\frac{p_x^2}{2m}+\frac{p_y^2}{2m}+\frac{p_z^2}{2m}$$ Or, d'après le théorème d'équipartition, on a : $$\overline{\left(\frac{p_x^2}{2m}\right)}=\frac{k_BT}{2}$$ donc $$\overline{E}=\frac{3}{2}k_BT$$ puis pour un gaz parfait à $N$ particules : $$\overline{E_\Sigma}=\frac{3}{2}Nk_BT=\frac{3}{2}nRT$$ On en déduit que $\displaystyle C_{v,m}=\frac{3}{2}R$ si $1\ll T\ll 10^{4}$.

\subsection{Capacité thermique des gaz parfaits diatomiques}

\paragraph{Capacité thermique des gaz parfaits diatomiques}De même, pour un gaz parfait diatomique, on a suivant les paliers :
\begin{itemize}
	\item 3 degrés de translation et alors $C_{v,m}=\frac{3}{2}R$ ;
	\item 3 degrés de translation et 2 degrés de rotation et alors $C_{v,m}=\frac{5}{2}R$ ;
	\item 3 de degrés de rotation, 2 degrés de rotation propre et 2 degrés de translation, et alors $C_{v,m}=\frac{7}{2}R$
\end{itemize}

\subsection{Modèle classique des solides ; loi de Dulong et Petit}
\paragraph{Modèle classique des solides} Pour un solide, on considère que chaque atome est relié à une position d'équilibre stable par un ressort. Ainsi, on a : $$E_\text{atome}=\frac{1}{2}m(v_x^2+v_y^2+v_z^2)+\frac{1}{2}k(x^2+y^2+z^2)$$ donc d'après le théorème d'équipartition (on est dans la limite classique), on a : $$\langle E_\text{système}\rangle=6N\times\frac{1}{2}k_BT=3Nk_BT$$

\paragraph{Loi de Dulong et Petit} On en déduit la loi de Dulong et Petit : dans la limite classique pour un solide isolant et on non magnétique, on a : $$C_{V,m}=3R$$

\paragraph{Limites} Expérimentalement, la capacité thermique molaire des solides s’effondre à basse température, d’où la nécessité d’un modèle quantique. Cependant, pour la plupart des cristaux (sauf le diamant), être au-dessus de $T=300$ K suffit pour être dans la limite continue.
\subsection{Modèle quantique d'Einstein (1907)}

\paragraph{Résultat admis sur l'oscillateur harmonique quantique} On admet que l'équation $$E\varphi=-\frac{\hbar^2}{2m}\varphi''+\frac{1}{2}kx^2\varphi$$ a pour spectre $$E_n=\left(n+\frac{1}{2}\right)\hbar\omega$$ avec $n\in\N$ et $\omega=\displaystyle\sqrt{\frac{k}{m}}$

\paragraph{Fonction de partition pour une particule} On a : $$Z_{OH1D}=\sum_{n\in\N}\exp\left(-\left(n+\frac{1}{2}\right)\hbar\omega\beta\right)=\frac{1}{2\sinh\left(\frac{\hbar\omega\beta}{2}\right)}$$

\paragraph{Énergie moyenne d'une particule} On en déduit, en dérivant par rapport à $\beta$ : $$\langle E_{OH1D}\rangle=\frac{\hbar\omega}{2}\frac{1}{\tanh\left(\frac{\hbar\omega\beta}{2}\right)}$$ Cela est homogène.

\paragraph{Énergie moyenne d'un solide} Pour $N$ particules, comme on a $3$ degrés de liberté, on obtient : $$\langle E_\text{solide}\rangle=\frac{3}{2}\frac{N\hbar\omega}{\tanh\left(\frac{\hbar\omega\beta}{2}\right)}$$

\paragraph{Capacité thermique d'un solide} On en déduit : $$C_V=3Nk_B\times\left(\displaystyle\frac{\displaystyle\frac{\hbar\omega}{2k_BT}}{\sinh\left(\frac{\hbar\omega}{2k_BT}\right)}\right)^2$$ Cela est homogène, et lorsque $T$ tend vers $+\infty$, on retrouve la loi de Dulong et Petit.

\paragraph{Température d'Einstein} On peut définir la température d'Einstein : $$\theta_E=\frac{\hbar\omega}{k_B}$$ pour laquelle se produit une inversion nette de la courbure de la courbe de $C_V$. 

\paragraph{Retour sur les gaz parfaits diatomiques} Dans ce cas, $\hbar\omega$ est compris entre $0,1$ eV et $1$ eV (c'est-à-dire l'énergie d'un photon dans l'infrarouge). On a alors $\theta_E=10^3$ K en ordre de grandeur, avec $\theta_E=\theta_{vib}$.

\subsection{Résumé : activation ou gel des degrés de liberté}

\paragraph{Cas $T\ll \theta_E$} Si $T\ll \theta_E$, le système reste dans le fondamental. Dans ce cas, on a : $$\langle E_{OH1D}\rangle\approx\frac{\hbar\omega}{2}\ \ \text{et}\ \  C_V\to 0$$

\paragraph{Cas $T\gg\theta_E$} Si $T\gg \theta_E$, on peut utiliser le théorème d'équipartition.

\paragraph{Remarque} On a un degré de liberté "électronique" et "nucléaire" à $T=300$ K.


\part{Électrochimie}
\chapter{Courbes intensité-potentiel}
\localtableofcontents
\section{Cinétique des réactions d'oxydo-réduction}
\subsection{Vitesse de réaction et intensité électrique ; conventions}

\paragraph{Demi-pile} On appelle demi-pile l'ensemble formé :
\begin{itemize}
	\item d'un conducteur solide qu'on appelle l'électrode ;
	\item d'une solution comportant un couple oxydant-réducteur ;
	\item d'un pont salin.
\end{itemize}

\paragraph{Exemples} Par exemple, on a :
\begin{itemize}
	\item une électrode de platine avec une solution d'ions Fe$^{2+}$ et Fe$^{3+}$ ;
	\item une électrode de cuivre avec une solution contenant des ions Cu$^{2+}$ et SO$_4^{2-}$.
\end{itemize}

\paragraph{Convention d'orientation du courant pour une demi-pile} Le courant $i$ est orienté de l'électrode vers la solution. Ainsi, si : \begin{itemize}
	\item $i>0$, l'électrode est une anode, le lieu d'une oxydation ;
	\item $i<0$, l'électrode est une cathode, le lieu d'une réduction.
\end{itemize}

\paragraph{Vitesse de réaction et intensité} Pour une réaction d’oxydoréduction, à coefficients stœchiométriques fixés, on note $\xi$ l’avancement et $p$ le nombre d’électrons échangés par réaction. On a alors $$\lvert I\rvert=pF\frac{d\xi}{dt}$$ où $F=N_Ae$ est la constante de Faraday. Mesurer l’intensité donne donc accès à la vitesse de réaction.

\subsection{Facteurs influençant la cinétique aux électrodes}

\paragraph{Facteurs cinétiques} Au-delà de ceux habituels (température, surface de contact, agitation, concentrations, catalyseurs, lumière), on considère le temps de transfert de charge entre l’électrode et le solvant (dépend de la différence de potentiel entre les deux) et le transfert de matière (le solvant amène les ions à l’électrode). Ce mécanisme se fait par diffusion selon la loi de Fick $$\overrightarrow{j}=-D\grad([X])$$ et par migration (mise en mouvement des charges car création d’un champ électrique).

\subsection{Courbes intensité-potentiel ; lien avec la formule de Nernst}

\paragraph{Courbes intensité-potentiel} On trace le courant $i$ traversant l’électrode en fonction du potentiel $E$, qui est la tension que l’on mesurerait par rapport à l’ESH électrode standard hydrogène ($E^\circ(H^+/H_2)=0$ à toute température).

\paragraph{Allure des courbes } En oxydation on a une courbe croissante d’intensité positive pour $E>E_\text{th}$. En réduction on a une courbe croissante d'intensité négative pour $E<E_{\text{th}}$ (nulle en $E_\text{th}$, électrode à l’équilibre, $E_\text{th}$ est donné par la formule de Nernst).

\subsection{Relevé expérimental de courbes intensité-potentiel}

\paragraph{Montage à trois électrodes}
On le mesure expérimentalement avec trois électrodes, celle de travail mesuré, relié d’une part par un voltmètre à une de référence, et d’autre part à une contre-électrode via un ampèremètre.

\section{Allure des courbes intensité-potentiel}
\subsection{Deux mécanismes}
\paragraph{Transfert de charge} Migration des électrons par le pont salin.

\paragraph{Transfert de matière} Consommation ou création de l'électrode.

\subsection{Systèmes rapides et systèmes lents ; cas des électrodes de référence}

\paragraph{Système électrochimique} On appelle système électrochimique l'ensemble :
\begin{itemize}
	\item couple oxydant-réducteur ;
	\item une électrode ;
	\item un solvant.
\end{itemize}

\paragraph{Système rapide} Un système électrochimique est dit rapide lorsque sa courbe intensité-potentiel présente une franche coupure avec l'axe des abscisses.

\paragraph{Exemple de système rapide} L'électrode de platine avec une solution d'ions Fe$^{2+}$ et Fe$^{3+}$ (et des anions) est un système électrochimique rapide.

\paragraph{Système lent} A l'inverse, pour un système électrochimique lent, la courbe intensité-potentiel présente un plateau pour $i\approx 0$. On ne sait donc pas précisément où se situe le potentiel de Nernst.

\paragraph{Surtension (ou surpotentiel)} Pour un système lent, on définit la surtension par : $$\eta=E-E_\text{Nernst}$$ où le $E$ correspond à une démarcation franche de $i$ avec l'axe des abscisses. Cette grandeur est algébrique car elle permet de définit :\begin{itemize}
	\item $\eta_a$ la surtension seuil en oxydation ;
	\item $\eta_c$ la surtension seuil en réduction, aussi dite surtension cathodique. Cette dernière est négative.
\end{itemize}
En général, on a $\eta_a\approx 0,1 V$.

\paragraph{Cas des électrodes de référence} Pour nos électrodes de référence, il faut donc choisir un système électrochimique rapide pour avoir $E$ très proche de $R_\text{ref}=E_\text{Nernst}$ lorsque $i$ est faible.

\subsection{Existence éventuelle de paliers de diffusion}

\paragraph{Effet du transfert de matière} Palier de transmission au-delà et en deçà de $E_\text{th}$, car il faut pouvoir amener l’ion suffisamment rapidement à l’électrode. Ces paliers peuvent être à $i=0$ en l’absence d’un des membres du couple, ou non présent si l’élément en question est l’électrode ou le solvant (pas besoin de l’amener). Dans le deuxième cas on parle de mur du solvant, on ne peut pas
sonder au-delà. \par On peut modéliser ces plateaux par un brassage parfait/concentration homogène loin de l’électrode et un gradient de concentration dans une couche limite au voisinage de celle-ci (modèle conducto-convectif). On trouve ainsi $$i_\text{max}=\text{stœchio}\times S\times D\times \frac{[X]_\infty}{e}\mathcal{F}$$ avec $S$ la surface et $D$ un coefficient de diffusion.

\subsection{Possibilité de vagues successives}

\paragraph{Possibilité de vagues successives} Dans une solution où l’on a plusieurs couples oxydant réducteurs sur la même électrode, la courbe intensité-potentiel est la somme des courbes intensité-potentiel des différents couples.

\subsection{Limitations par le solvant}

\section{Utilisation des courbes intensité-potentiel}
\subsection{Potentiel mixte}

\paragraph{Utilisation des courbes} Si l’on a deux couples sur une électrode, son potentiel est celle qui vérifie $i_a=-i_c$ car les électrons ne sont pas stables en solution aqueuse. Sur deux électrodes, les intensités de l’une et de l’autre sont opposées et les potentiels sont libres. On peut tracer les courbes caractéristiques $u-i$ en regardant comment l’écart entre les point de fonctionnement sur la courbe
anode à $+I$ et cathode à $-I$ évoluent avec $I$. Cela permet de déterminer le sens spontané d’une réaction. \par Réciproquement, par électrolyse on impose l’intensité et donc les points de fonctionnements

\subsection{Étude des piles : cinétique et capacité}

\paragraph{Capacité d'une pile} C'est la charge que peut débiter une pile sur l'ensemble de sa durée de vie. Elle s'exprime souvent en mA.h. On rappelle que $1$ mA.h$=3,6$ C.

\subsection{Étude cinétique des électrolyseurs}

\paragraph{Principe d'une électrolyse} Le sens du courant est imposé par un générateur extérieur. La circulation des électrons se fait donc de manière forcée et non plus naturelle comme dans le cas des piles.

\subsection{Charge et décharge d'un accumulateur}

\paragraph{Vocabulaire} Un accumulateur est ce que l'on appelle communément une pile rechargeable.



\chapter{Corrosion humide du fer}
\localtableofcontents

\section{Aspect thermodynamique}

\subsection{Diagramme E-pH}

\paragraph{Rappel important} Pour utiliser la formule de Nernst, on écrit avant la demi-équation, que l'on équilibre avec des $\text{H}^+$.

\paragraph{Rappel : diagrammes potentiel-pH} On place les espèces par nombre d’oxydation croissant avec $E$ et basicité croissante avec le $pH$. Pour les solides, il s’agit de diagramme d’existence et pour les espèces solubles des diagrammes de prédominance. On distingue le domaine d’immunité, celui du solide, le domaine de passivation (oxydes ou hydroxydes solides qui forment une couche protectrice ralentissant la réaction) et le domaine de corrosion (oxyde et hydroxydes solubles).

\paragraph{Rappel : tracé du diagramme} Avant tout, on liste les différents nombres d'oxydation. Ensuite, on déterminer les frontières verticales avec l'apparition des précipités. Puis on déterminer les autres frontières en utilisant la formule de Nernst et la continuité du potentiel.

\paragraph{Rappel : lecture du diagramme} On utilise la formule de Nernst et les valeurs sur la courbe.

\paragraph{Rappel : conventions de tracé} On dispose de plusieurs conventions pour les espaces en solutions : 
\begin{enumerate}
	\item pour $X$ soluble, on prend $[X]=c_{\text{tracé}}$ à la frontière ;
	\item on prend une somme des concentrations en espèces égale à la concentration de tracé ;
	\item on prend une somme de concentrations en atomes égale à la concentration de tracé.
\end{enumerate}

\subsection{Enthalpie libre électrochimique $\Delta_{1/2}G^\circ$}

\paragraph{$\Delta_{1/2}G^\circ$} Pour une demi-équation d'oxydoréduction Ox$+ne^-=$ Red, on définit l'enthalpie libre électrochimique : $$\Delta_{1/2}G^\circ=-nFE^\circ$$ comme dans la formule de Nernst.

\paragraph{Exemple} Par exemple, pour Cu$^{2+}+2e^-=$ Cu, on a donc : $$\Delta_{1/2}G^\circ=-2FE_{\text{Cu}^{2+}/\text{Cu}}^\circ$$

\subsection{$\Delta_r G^\circ$, $\Delta_rS^\circ$, $\Delta_rH^\circ$, $e^\circ$ et $E^\circ$}

\paragraph{Rappel : enthalpie libre} On rappelle que $G=H-TS$ est l'enthalpie libre : c'est une fonction d'état extensive, homogène à une énergie.

\paragraph{Rappel : opérateur de Lewis} Une fois l'équation de la réaction (\textit{ie} les coefficients stœchiométriques) choisie, on définit, pour toute grandeur extensive du système : $$\Delta_rX=\frac{\partial X}{\partial\xi}\bigg|_{T,P}$$ On a alors en particulier la relation : $$\Delta_rG^\circ=\Delta_rH^\circ-T\Delta_rS^\circ$$

\paragraph{Rappel : obtention de $\Delta_rH^\circ$ et $\Delta_rS^\circ$ à partir de $\Delta_rG^\circ$} On rappelle que l'on dispose des formules : $$\frac{d}{dT}\left(\Delta_rG^\circ\right)=-\Delta_rS^\circ\ \ \text{et}\ \  \frac{d}{dT}\left(\frac{\Delta_rG^\circ}{T}\right)=-\frac{\Delta_rH^\circ}{T^2}$$

\paragraph{Rappel : activités, quotient de réaction, constante de réaction standard} On rappelle la relation fondamentale : $$\mu_i=\mu_i^\circ+RT\ln(a_i)$$ qui nous donne : $$\Delta_rG=\sum_i\mu_i\nu_i=\sum_i\mu_i^\circ\nu_i+RT\ln(Q)=\Delta_rG^\circ+RT\ln(Q)$$ Or, à l'équilibre, $Q=K^\circ$ et $\Delta_rG=0$ donc on en déduit : $$\Delta_rG^\circ=-RT\ln(K^\circ)$$ ou bien encore $$K^\circ=\exp\left(-\frac{\Delta_rG^\circ}{RT}\right)$$

\paragraph{Attention !} Dans la formule de Nernst, il y a un $\log$ (décimal) alors que dans les relations de thermochimie, ce sont des $\ln$ (naturels).

\paragraph{Exemple 1}On considère la réaction $$2\text{Ag}^{+}+\text{Cu}=2\text{Ag}+\text{Cu}^{2+}$$ Cette réaction provient des deux demi-équations suivantes :\begin{itemize}
	\item $\text{Cu}^{2+}+2e^-=\text{Cu}$ (1) ;
	\item $\text{Ag}^++e^-=\text{Ag}$ (2).
\end{itemize} La réaction est considérée s'écrit $-1\times(1)+2\times(2)$. Ainsi, l'enthalpie standard de réaction vaut : $$\Delta_rG^\circ=2\Delta_{1/2}G^\circ(2)-\Delta_{1/2}G^\circ(1)=2(-FE_{\text{Ag}^+/\text{Ag}}^\circ)-(-2FE_{\text{Cu}^{2+}/\text{Cu}}^\circ)=-2F(E_{\text{Ag}^+/\text{Ag}}^\circ-E_{\text{Cu}^{2+}/\text{Cu}}^\circ)$$ On peut alors exprimer $K^\circ$ en fonction des $E^\circ$ : $$K^\circ=\exp\left(-\frac{\Delta_rG^\circ}{RT}\right)=\exp\left(\frac{2F}{RT}(E_{\text{Ag}^+/\text{Ag}}^\circ-E_{\text{Cu}^{2+}/\text{Cu}}^\circ)\right)$$ Or, on a $E_{\text{Cu}^{2+}/\text{Cu}}^\circ=0,34$ V et $E_{\text{Ag}^+/\text{Ag}}^\circ=0,80$ V donc la réaction est très avancée.

\paragraph{Remarque} De façon générale, si les potentiels standards sont distants d'au moins $0,25$ V, la réaction est très avancée.

\paragraph{Variante "sup" de l'exemple 1} On aurait pu calculer cette constante de réaction avec les outils de sup en écrivant la formule de Nernst pour les deux couples, en utilisant l'égalité des potentiels de Nernst à l'équilibre et en utilisant la loi d'action des masses. On pourra vérifier que'on obtient exactement le même résultat.

\paragraph{Exemple 2} On note : \begin{itemize}
	\item $E_a^\circ$ le potentiel standard du couple $\text{Fe}^{3+}/\text{Fe}$ ;
	\item $E_b^\circ$ le potentiel standard du couple $\text{Fe}^{3+}/\text{Fe}^{2+}$ ;
	\item $E_c^\circ$ le potentiel standard du couple $\text{Fe}^{2+}/\text{Fe}$.
\end{itemize} On cherche une relation entre ces trois potentiels standards. Or : \begin{itemize}
\item pour la demi-équation $a$ : $\text{Fe}^{3+}+3e^-=\text{Fe}$, $\Delta_{1/2}G^\circ(a)=-3FE_a^\circ$ ;
\item pour la demi-équation $b$ : $\text{Fe}^{3+}+e^-=\text{Fe}^{2+}$, $\Delta_{1/2}G^\circ(b)=-FE_b^\circ$ ;
\item pour la demi-équation $a$ : $\text{Fe}^{2+}+2e^-=\text{Fe}$, $\Delta_{1/2}G^\circ(c)=-2FE_a^\circ$.
\end{itemize} Comme on a $c=a-b$ (en termes de demi-équations), on en déduit : $$2E_c^\circ=-E_b^\circ+3E_a^\circ$$

\paragraph{Variante "sup" de l'exemple 2} Avec les outils de sup, on écrit les 3 formules de Nernst et on utilise l'égalité des potentiels de Nernst à l'équilibre dans une solution où barbotent les 3 couples.

\paragraph{Exemple 3} On donne $E^\circ$ le potentiel standard pour le couple $\text{Fe}^{3+}/\text{Fe}$ et on note $E'^\circ$ le potentiel standard pour le couple $\text{Fe}(\text{OH})_3/\text{Fe}$. On note $K_s$ le produit de solubilité de la réaction R : $$\text{Fe}(\text{OH})_3=\text{Fe}^{3+}+3\text{HO}^-$$ On souhaite exprimer $E'^\circ$. On dispose des équations : \begin{itemize}
	\item a : $\text{Fe}^{3+}+3e^-=\text{Fe}$ ;
	\item b : $3\text{H}^++\text{Fe}(\text{OH})_3+3e^-=\text{Fe}+3\text{H}_2\text{O}$ ;
	\item $c$ : $\text{H}_2\text{O}=\text{H}^++\text{HO}^-$.
\end{itemize} Or, la réaction R s'écrit R=b-a+3c. En passant à l'enthalpie libre de réaction standard, on en déduit : $$-RT\ln(K_s)=-3FE^\circ-(-3FE'^\circ)-3RT\ln(K_e)$$

\paragraph{Variante "sup" de l'exemple 3} Appliquer la formule de Nernst, la loi d'action des masses et l'égalité des potentiels de Nernst à l'équilibre.

\paragraph{$e$ et $e^\circ$} Si on choisit de mettre $e$ dirigé d'une électrode de cuivre baignant dans une solution de ses ions et vers une électrode d'argent baignant dans une solution de ses ions, on a alors : $$e=E_{\text{Ag}^+/\text{Ag}}-E_{\text{Cu}^{2+}/\text{Cu}}\ \ \text{et}\ \ e^\circ=E_{\text{Ag}^+/\text{Ag}}^\circ-E_{\text{Cu}^{2+}/\text{Cu}}^\circ$$

\subsection{Travail électrique}

\paragraph{Remarque} C'est un exemple de "travail utile".

\paragraph{Travail électrique et thermodynamique} On travaille sur une pile Argent/Cuivre. On considère le système \{pile\} auquel on applique le premier principe. Les variations d'énergie mécanique sont négligeables donc : $$U_F-U_I=W_{\text{pression}}+Q+W_{\text{électrique}}$$ Si on suppose la réaction isotherme à la température $T_0$ et isobare à la pression $P_0$, on a : $$W_\text{pression}=-P_0(V_F-V_I)$$ puis le deuxième principe s'écrit : $$S_F-S_I=\frac{Q}{T_0}+S_\text{créée}$$ Alors : $$-W_\text{électrique}=-P_0\Delta V+T_0\Delta V-T_0S_\text{créée}=-\Delta G-T_0S_\text{créée}$$ donc on en déduit $-W\text{électrique}<-\Delta G$

\paragraph{Remarque : critère d'évolution} On retrouve le critère d'évolution : pour un système fermé, lors d'une transformation isotherme et isobare sans travail utile, $G$ diminue.

\paragraph{Travail électrique} On a aussi en convention générateur, un travail électrique reçu qui vaut : $$\delta W_\text{électrique}=-uidt$$

\section{Aspect cinétique : cas de la corrosion uniforme}

\subsection{Potentiel et intensité de corrosion}

\paragraph{Intensité de corrosion} Si on abandonne un morceau de fer dans de l'eau, on a $i_\text{tot}=0$ car l'électrode n'est reliée à rien.

\subsection{Une eau aérée accélère la corrosion}

\paragraph{Effet de l'aération de l'eau} Ajouter $O_{2(\text{aq})}$ permet :\begin{itemize}
	\item d'augmenter le potentiel de corrosion ;
	\item d'augmenter le courant de corrosion
\end{itemize} donc accélère la corrosion.

\section{Aspect cinétique : cas de la corrosion différentielle}

Dans ce cas, certaines zones du métal sont des anodes et d'autres sont des cathodes.

\subsection{Pile de concentration}

\paragraph{Exemple 1} On considère deux électrodes de cuivre. Celle de gauche trempe dans une solution d'ions $\text{Cu}^{2+}$ de concentration $[\text{Cu}^{2+}]_g=0,01$ mol/L et celle de droite trempe dans une solution d'ions $\text{Cu}^{2+}$ de concentration $[\text{Cu}^{2+}]_d=0,1$ mol/L. Dans les deux cas, les demi-équations sont : $$\text{Cu}^{2+}+2e^-=2\text{Cu}$$ En appliquant la formule de Nernst aux deux systèmes électrochimiques, on obtient successivement : $$E_g=E_{\text{Cu}^{2+}/\text{Cu}}^\circ+\frac{0,06}{2}\log\left([\text{Cu}^{2+}]_g\right)$$ et $$E_d=E_{\text{Cu}^{2+}/\text{Cu}}^\circ+\frac{0,06}{2}\log\left([\text{Cu}^{2+}]_d\right)$$ donc dans ce cas on a : $$e=E_g-E_d=\frac{0,06}{2}\log(10\inv)=-0,03\ V$$

\paragraph{Exemple 2} On considère deux électrodes d'argent qui ne sont pas reliées par un fil. Celle de gauche trempe dans une solution de sulfate d'argent $(2\text{Ag}^+,\ \text{SO}_4^{2-})$ à la concentration $c_0$. L'électrode de droite baigne dans une solution de chlorure de potassium $(\text{K}^+,\ \text{Cl}^-)$ à la concentration $c_0$. On verse dans le bécher de droite une goutte de nitrate d'argent qui contribue à la formation d'un précipité d'AgCl. On donne le produit de solubilité de la réaction de dissociation du précipité $K_s$. On souhaite déterminer $e$ la force électromotrice de la pile, prise de l'électrode de droite vers celle de gauche \textit{a priori}. Pour la réaction $$\text{Ag}^++e^-=\text{Ag}$$ qui se déroule dans le bécher de droite, la formule de Nernst donne : $$E_g=E^\circ+\frac{0,06}{1}\log(2c_0)$$ Le fait qu'une précipité apparaisse dans le bécher de droite fournit : $$K_s=[\text{Ag}^+]_d[\text{Cl}^-]$$ On en déduit : $$e=E_g-E_d=0,06\log(2c_0)-0,06\log\left(\frac{K_s}{c_0}\right)=0,06\log\left(\frac{2c_0^2}{K_s}\right)$$

\subsection{Pile d'Evans ; aération différentielle}

\paragraph{Aération différentielle} Si l’on place deux électrodes de fer dans un bécher oxygéné et un désoxygéné reliés par un pont salin, il y a oxydation de celle dans le bêcher désoxygéné. Donc une plaque de fer dans de l’eau de mer rouillera majoritairement sur les endroits moins oxygénés, c'est-à-dire dans les trous, ce qui la mène à se percer rapidement.

\subsection{Piles avec des électrodes de natures différentes}

\paragraph{Corrosion} Électrodes de fer et de cuivre reliées, oxydation du fer, tandis qu’avec fer-zinc, c’est le zinc qui s’oxyde.

\section{Protection contre la corrosion}

\subsection{Acier inoxydable}
\paragraph{Acier} L'acier est un alliage de fer et de carbone, comportant entre $0,15\%$ et $0,85\%$ de carbone.  Le carbone vient se glisser entre les différents noyaux de fer.

\paragraph{Acier inoxydable} L'acier inoxydable est de l'acier recouvert d'une couche de $\text{Cr}_2\text{O}_3$. Cette dernière couche est imperméable et protège donc l'acier.

\paragraph{Passivation de l'acier inoxydable} En cas de rayure qui exposerait l'acier, une couche de $\text{Cr}_2\text{O}_3$ se reforme au-dessus de l'acier. Cette passivation permet à l'acier inoxydable de s'auto-protéger de la corrosion.

\subsection{Traitement de surface}

\paragraph{Anode sacrificielle} Pour protéger une coque en acier (pour un bateau par exemple) de la corrosion due à l'eau ou à l'air, on peut apposer un morceau de zinc sur la coque, dit "anode sacrificielle", car le contact avec le zinc diminue le potentiel de corrosion du fer et ralentit donc son oxydation et sa corrosion. Il faut alors mettre beaucoup d'anodes sacrificielles et les remplacer régulièrement, ce qui peut coûter cher.

\paragraph{Peinture} On peut aussi peindre la coque avec une peinture qui protège de la corrosion. L'inconvénient est qu'une telle peinture peut provoquer la formation d'algues ou bien libérer des substances toxiques dans l'eau.

\subsection{Protection électrochimique}

\paragraph{Protection anodique par courant imposé} On relie la conduite à protéger à une autre électrode, les deux étant reliées par un générateur de courant. Le tout baigne dans un électrolyte. Alors, le fait d'imposer un courant avec un générateur permet de limiter la corrosion sur la conduite à protéger et consomme plutôt l'autre électrode. Cette technique peut s'avérer coûteuse en énergie.

\subsection{Passivation}

\paragraph{Exemple de l'aluminium} On fait tremper une électrode d'aluminium dans de l'eau. La courbe intensité-potentiel présente une bosse avec les réactions : $$\text{Al}\to\text{Al}^{3+}\ \ \text{et}\ \ \text{Al}^{3+}\to\text{Al}_2\text{O}_3$$ La bosse se situe vers $\text{Al}^{3+}$. Alors, $\text{Al}_2\text{O}_3$ est imperméable et bloque l'oxydation de $\text{Al}$ : le solide se protège en créant une couche imperméable à sa surface, c'est la passivation.

\part{Chimie de sup}

\chapter{Cinétique chimique}
A COMPLETER
\localtableofcontents
\section{}
\subsection{}
\subsection{}
\section{}
\subsection{}
\subsection{}

\chapter{Solutions aqueuses}
A COMPLETER
\localtableofcontents
\section{}
\subsection{}
\subsection{}
\section{}
\subsection{}
\subsection{}

\chapter{Architecture de la matière}
A COMPLETER
\localtableofcontents
\section{}
\subsection{}
\subsection{}
\section{}
\subsection{}
\subsection{}

\chapter{Cristallographie}
\localtableofcontents
\section{Définitions générales}
\paragraph{Cristal parfait} Une cristal parfait est un ensemble de particules empilées régulièrement dans l'espace. il forme une structure périodique sans défaut de taille infinie.
\paragraph{Réseau} Entité géométrique définie par un ensemble infini de points obtenus à partir d'un point choisi arbitrairement comme origine et par translations d'un vecteur
\paragraph{Nœud} Un nœud est un point du réseau obtenu par translation de ce vecteur.
\paragraph{Motif} Un motif est la plus petite entité discernable qui se répète périodiquement. A chaque nœud est associé un motif. Un cristal est la donné d'un réseau et d'un motif.
\paragraph{Maille élémentaire} Une maille élémentaire correspond à la portion de l'espace de volume minimal qui assure le pavage de l'espace et redonne le réseau par translation.
\paragraph{Nombre de nœuds par maille} On partage les nœuds entre les différentes mailles auxquelles ils appartiennent.
\paragraph{Population} La population est le nombre de motifs par maille. Elle correspond donc au nombre de nœuds par maille.
\paragraph{Masse volumique d'un cristal} Elle vérifie : $$\rho=\frac{ZM}{N_1V}$$ où $Z$ est le nombre de motifs par maille ; $M$ est la masse molaire du motif ; $V$ est le volume de la maille et $N_A$ est le nombre d'Avogadro.
\paragraph{Coordinence} La coordinence d'une entité $A$ par rapport à une autre entité $B$ est le nombre de plus proches voisins de $A$ qui sont des entités de type $B$. On peut la noter $\displaystyle\frac{A}{B}$
\paragraph{Compacité} La compacité $C$ d'un solide cristallin est la fraction de volume occupée par la matière dans la maille n considérant les atomes et les ions comme des sphères dures.
\section{Structures compactes}
\subsection{Cubique à faces centrées}
\paragraph{Description} Un atome à chaque sommet et un atome au centre de chaque face. Il n'y en a pas au milieu du cube. Elle correspond à l'empilement ABC ABC.
\paragraph{Population} La population vaut $$8\times\frac{1}{8}+6\times\frac{1}{2}=4$$
\paragraph{Coordinence} Elle vaut 12 (trouvée pour n'importe quel atome).
\paragraph{Relation entre $a$ et $R$} Les atomes se touchent sur la diagonale d'une face donc $$4R=a\sqrt{2}$$
\paragraph{Compacité} En utilisant la relation précédente et la population de $4$, on trouve $$C=\frac{\pi}{3\sqrt{2}}\approx 0,74$$ C'est une structure compacte : c'est la meilleure compacité que l'on puisse obtenir.
\paragraph{Masse volumique} La masse volumique vaut alors : $$\rho=\frac{4M}{N_Aa^3}$$
\paragraph{Sites octaédriques} Il y en a un sur chaque arête et un au centre soit $$12\times\frac{1}{4}+1=4$$ par maille. Il y en a donc autant que d'atomes constituant le réseau. Ils occupent les milieux des arêtes et le centre du cube. Au centre du cube, on a la relation : $$2R_O+2R=a$$ Par conséquent, on a : $$R_O=(\sqrt{2}-1)R\approx 0,414 R$$
\paragraph{Sites tétraédriques} Il y en a un dans chacun des sous-cubes de côté $a/2$ de la structure cubique à faces centrées, donc $8$. Il y en a donc deux fois plus que de sites $O$, et donc deux fois plus que de nœuds dans le réseau hôte. Sur la demi-diagonale d'un petit cube, on obtient : $$R_T+R=\frac{a}{4\sqrt{3}}$$ Avec la relation générale entre $a$ et $R$, on en déduit que $$R_T=\left(\sqrt{\frac{3}{2}}-1\right)R\approx 0,225R$$ Les sites $T$ sont donc plus petits que les sites $O$.
\paragraph{Structure hexagonale compacte} \textbf{A COMPLETER}
\paragraph{Compacité} On obtient aussi $0,74$.
\section{Structures usuelles}
\subsection{Diamant}
\paragraph{Description} Les atomes de carbone occupent une structure cubique à faces centrées avec occupation de la moitié des sites $T$ en alternance.
\paragraph{Population} On a donc une population de $$8\times\frac{1}{8}+6\times\frac12+4\times 1=8$$
\paragraph{Coordinence} Elle vaut $4$, obtenue sur un atome d'un site $T$.
\paragraph{Relation entre $a$ et $R$} Sur la demi-diagonale d'un petit cube, on a : $$2R=\frac{a}{4\sqrt{3}}$$
\paragraph{Compacité} La compacité vaut alors : $$C=\frac{\pi\sqrt{3}}{16}\approx 0,34$$ C'est une structure non compacte.
\subsection{Glace type diamant}
\paragraph{Description} Les atomes d'oxygène occupent une structure identique à celle du diamant.
\subsection{Structure type CsCl}
\paragraph{Description} Les ions Cl$^-$ occupent un réseau cubique simple et les ions Cs$^+$ occupent le centre du cube.
\paragraph{Coordinences} Les coordinences de Cs$^+$ par rapport à Cl$^-$ et de Cl$^-$ par rapport à Cs$^+$ valent toutes les deux $8$. Les coordinences de Cs$^+$ par rapport à Cs$^+$ et de Cl$^-$ par rapport à Cl$^-$ valent toutes les deux 6.
\paragraph{Populations} Dans chaque maille, il y a un ion Cl$^-$ et un ion Cs$^+$.
\paragraph{Masse volumique} Elle vaut $$\rho=\frac{M_{\text{Cs}}+M_\text{Cl}}{N_Aa^3}$$
\paragraph{Relation entre les paramètres} Le long de la diagonale du cube : $$2r_++2r_-=a\sqrt{3}$$
\paragraph{Condition de non contact anion/anion} Le long d'une arête : $$a>2r_-$$
\paragraph{Condition d'existence de la structure} Nécessairement, on doit avoir : $$x>\sqrt{3}-1\approx 0,732$$ où $x=\displaystyle\frac{r_+}{r_-}$
\paragraph{Compacité} Elle vaut alors : $$C=\frac{\pi\sqrt{3}}{2}\left(\left(1+x\right)\inv+\left(1+\frac{1}{x}\right)\inv\right)$$

\subsection{Structure type NaCl}
\paragraph{Description} Les ions Cl$^-$ occupent un réseau cubique à faces centrées et les ions Na$^+$ occupent un réseau cubique à faces centrées décalé (ou bien d'un autre point de vue tous les sites octaédriques).
\paragraph{Coordinences} Les coordinences de Na$^+$ par rapport à Cl$^-$ et de Cl$^-$ par rapport à Na$^+$ valent toutes les deux $6$. Les coordinences de Na$^+$ par rapport à Na$^+$ et de Cl$^-$ par rapport à Cl$^-$ valent toutes les deux 12.
\paragraph{Populations} Dans chaque maille, il y a 4 ions Cl$^-$ et 4 ions Na$^+$.
\paragraph{Masse volumique} Elle vaut $$\rho=4\frac{M_{\text{Na}}+M_\text{Cl}}{N_Aa^3}$$
\paragraph{Relation entre les paramètres} Le long de la l'arête du cube : $$2r_++2r_-=a$$
\paragraph{Condition de non contact anion/anion} Le long d'une diagonale d'une face : $$a\sqrt{2}>4r_-$$
\paragraph{Condition d'existence de la structure} Nécessairement, on doit avoir : $$x>\sqrt{2}-1\approx 0,414$$ De plus, on doit aussi avoir $x<\sqrt{3}-1$, sinon la structure type CsCl serait plus optimisée.
\paragraph{Compacité} Elle vaut alors : $$C=\frac{2\pi}{3}\left(\left(1+x\right)^{-3}+\left(1+\frac{1}{x}\right)^{-3}\right)$$

\subsection{Structure type ZnS, dite blende}
\paragraph{Description} Les ions S$^-$ occupent un réseau cubique à faces centrées et les ions Zn$^+$ occupent la moitié des sites tétraédriques en alternance.
\paragraph{Coordinences} Les coordinences de Zn$^+$ par rapport à S$^-$ et de S$^-$ par rapport à Zn$^+$ valent toutes les deux $4$. Les coordinences de Zn$^+$ par rapport à Zn$^+$ et de S$^-$ par rapport à S$^-$ valent toutes les deux 12.
\paragraph{Populations} Dans chaque maille, il y a 4 ions S$^-$ et 4 ions Zn$^+$.
\paragraph{Masse volumique} Elle vaut $$\rho=4\frac{M_{\text{Zn}}+M_\text{S}}{N_Aa^3}$$
\paragraph{Relation entre les paramètres} Le long de la demi-diagonale d'un petit cube : $$r_++r_-=\frac{a\sqrt{3}}{4}$$
\paragraph{Condition de non contact anion/anion} Le long de la diagonale d'une face : $$a\sqrt{2}>4r_-$$
\paragraph{Condition d'existence de la structure} Nécessairement, on doit avoir : $$x>\sqrt{\frac{3}{2}}-1\approx 0,225$$ De plus, on doit aussi avoir $x<\sqrt{3}-1$, sinon la structure type NaCl serait plus optimisée.
\paragraph{Compacité} Elle vaut alors : $$C=\frac{\sqrt{3}\pi}{4}\left(\left(1+x\right)^{-3}+\left(1+\frac{1}{x}\right)^{-3}\right)$$

\subsection{Structure type CaF$_2$ : type AB$_2$}
\paragraph{Description} Association de 8 petits cubes simples d'anions. Occupation alternée de la moitié des sites cubiques. Ce n'est pas une structure type ZnS car il n'y a pas d'atomes au milieu du cube ni au milieu des arêtes dans ZnS.
\paragraph{Deuxième description} On peut considérer que le réseau hôte est le réseau cationique de type CFC ont les ions F$^-$ occupent tous les sites tétraédriques.
\paragraph{Coordinences} La coordinence de Ca$^{2+}$ par rapport à F$^-$ vaut $8$. La coordinence de F$^-$ par rapport à Ca$^{2+}$ vaut 4. La coordinence de Ca$^{2+}$ par rapport à Ca$^{2+}$ vaut 12. La coordinence de F$^-$ par rapport à $F^-$ vaut 6.
\paragraph{populations} Il y a en propre dans une maille 8 ions F$^-$ et 4 ions Ca$^{2+}$.
\paragraph{Masse volumique} Elle vaut : $$\rho=4\frac{2M_{\text{F}}+M_\text{Cl}}{N_Aa^3}$$

\end{document}